\pdfoutput=1
%% For double-blind review submission, w/o CCS and ACM Reference (max submission space)
\documentclass[acmsmall,screen]{acmart}\settopmatter{printfolios=true,printccs=true,printacmref=true}
%% For double-blind review submission, w/ CCS and ACM Reference
%\documentclass[acmsmall,review,anonymous]{acmart}\settopmatter{printfolios=true}
%% For single-blind review submission, w/o CCS and ACM Reference (max submission space)
%\documentclass[acmsmall,review]{acmart}\settopmatter{printfolios=true,printccs=false,printacmref=false}
%% For single-blind review submission, w/ CCS and ACM Reference
%\documentclass[acmsmall,review]{acmart}\settopmatter{printfolios=true}
%% For final camera-ready submission, w/ required CCS and ACM Reference
%\documentclass[acmsmall]{acmart}\settopmatter{}


%% Journal information
%% Supplied to authors by publisher for camera-ready submission;
%% use defaults for review submission.
\setcopyright{rightsretained}
\acmPrice{}
\acmDOI{10.1145/3591291}
\acmYear{2023}
\copyrightyear{2023}
\acmSubmissionID{pldi23main-p578-p}
\acmJournal{PACMPL}
\acmVolume{7}
\acmNumber{PLDI}
\acmArticle{177}
\acmMonth{6}
\received{2022-11-10}
\received[accepted]{2023-03-31}

%\acmJournal{PACMPL}
%\acmVolume{1}
%\acmNumber{CONF} % CONF = POPL or ICFP or OOPSLA
%\acmArticle{1}
%\acmYear{2018}
%\acmMonth{1}
%\acmDOI{} % \acmDOI{10.1145/nnnnnnn.nnnnnnn}
\startPage{1}

%% Copyright information
%% Supplied to authors (based on authors' rights management selection;
%% see authors.acm.org) by publisher for camera-ready submission;
%% use 'none' for review submission.
%\setcopyright{none}
%\setcopyright{acmcopyright}
%\setcopyright{acmlicensed}
%\setcopyright{rightsretained}
%\copyrightyear{2018}           %% If different from \acmYear

%% Bibliography style
\bibliographystyle{ACM-Reference-Format}
%% Citation style
%% Note: author/year citations are required for papers published as an
%% issue of PACMPL.
\citestyle{acmauthoryear}   %% For author/year citations


%%%%%%%%%%%%%%%%%%%%%%%%%%%%%%%%%%%%%%%%%%%%%%%%%%%%%%%%%%%%%%%%%%%%%%
%% Note: Authors migrating a paper from PACMPL format to traditional
%% SIGPLAN proceedings format must update the '\documentclass' and
%% topmatter commands above; see 'acmart-sigplanproc-template.tex'.
%%%%%%%%%%%%%%%%%%%%%%%%%%%%%%%%%%%%%%%%%%%%%%%%%%%%%%%%%%%%%%%%%%%%%%


%% Some recommended packages.
\usepackage{booktabs}   %% For formal tables:
                        %% http://ctan.org/pkg/booktabs
\usepackage{subcaption} %% For complex figures with subfigures/subcaptions
                        %% http://ctan.org/pkg/subcaption

% \usepackage{algorithmic}
% \usepackage{algorithm}
% \usepackage{array}
% \usepackage[caption=false,font=normalsize,labelfont=sf,textfont=sf]{subfig}
\usepackage{textcomp}
% \usepackage{stfloats}
\usepackage{url}
\usepackage{verbatim}
% \usepackage{graphicx}
% \usepackage{cite}

\usepackage{graphicx}
% \usepackage{caption}
% \captionsetup{size=footnotesize, skip=5pt, position = bottom, labelsep=period}
% \usepackage[font={small}]{caption}
% \captionsetup{size=footnotesize,skip=5pt, position = bottom}
% \captionsetup{skip=5pt, position = bottom}
% \usepackage[font={small}]{caption}
% \usepackage[english]{babel}
\usepackage{amsmath,amssymb,enumerate}
\usepackage{amsfonts}
\usepackage{mathtools}
\usepackage[mathscr]{euscript}
\usepackage{mathrsfs}
\usepackage{times}
\usepackage[usenames,dvipsnames,svgnames,table, xcdraw]{xcolor} %http://ctan.org/pkg/xcolor
% \usepackage{url}\urlstyle{tt}
% \usepackage{subcaption}
\usepackage[caption=false,font=normalsize,labelfont=sf,textfont=sf]{subfig}
\usepackage{booktabs}
% \usepackage[font={small}]{caption}
% \captionsetup{size=footnotesize,
%     skip=5pt, position = bottom}
% \usepackage[caption=false,font=footnotesize]{subfig}
\usepackage{multirow}
\usepackage[ruled,vlined,linesnumbered]{algorithm2e}
\SetKw{Continue}{continue}
\SetKw{Break}{break}
% \usepackage{xcolor} %http://ctan.org/pkg/xcolor
\let\labelindent\relax
\usepackage{enumitem}
% \usepackage[inline]{enumitem}

% \usepackage[margin=1in]{geometry}
% \usepackage[parfill]{parskip}


% \let\proof\relax 
% \let\endproof\relax
% \usepackage{amsthm}
% \usepackage{breqn}

% \usepackage{blindtext}
% \usepackage{gensymb}
% \usepackage{xparse}
% \usepackage{lipsum}
% \usepackage{mathrsfs}
% \usepackage[mathscr]{euscript}
\usepackage{cite} % to group references
% \usepackage{multicol}
% \usepackage{cleveref}
\usepackage[utf8]{inputenc}
% \usepackage[T1]{fontenc}
% \usepackage{textcomp}
\usepackage{balance}
% \usepackage{soul}
% \usepackage{multirow}
\usepackage{float}
\usepackage{listings}
\usepackage{tcolorbox}
\usepackage[justification=centering]{caption}
% \usepackage{float}
% \usepackage{lscape}

%\usepackage{epsfig} % for postscript graphics files
% \usepackage{mathptmx} % assumes new font selection scheme installed
\usepackage[normalem]{ulem}
% \usepackage{soul}
\usepackage{pifont}


%% the rest of your preamble here
\newcommand\blfootnote[1]{%
  \begingroup
  \renewcommand\thefootnote{}\footnote{#1}%
  \addtocounter{footnote}{-1}%
  \endgroup
}







\newcommand\eva[1]{\textcolor{pink}{Eva: #1}}
\newcommand\toptile[1]{top_{#1}}
\newcommand\bottomtile[1]{bot_{#1}}
\newcommand\lefttile[1]{lef_{#1}}
\newcommand\righttile[1]{rig_{#1}}
\newcommand{\coms}{\xspace\scalebox{0.5}{\begin{tikzpicture}\draw[dash pattern=on 3pt off 3pt, line width = 1mm ] (0,0) -- (0.8,0); \node at (0,-0.08) {}; \end{tikzpicture}}\xspace}
\newcommand{\moves}{\rightarrow}
\newcommand{\colorellipse}{orange!20!white}
\newcommand{\addCommunicationB}[1]
{
	\node[above right = 0.5mm and 4mm of #1, inner sep=0mm] (nb2) {B};
	\draw[tobehere,communication] (#1) -- (nb2);
}

\newcommand\setnodes{V}
\newcommand\basenode{B}
\newcommand{\sourcenode}{s}
\newcommand{\targetnode}{t}
\newcommand\USTCONN{USTCONN}
\newcommand\UCONN{UCONN}
\newcommand{\myundirectedgraphbounded}{\mathfrak G}
\newcommand{\myundirectedgraph}{\mathfrak G'}

\newcommand\pb[2]{${#1}_{#2}$~}

\newcommand\pbReach[1]{\pb{Reachability}{#1}}
\newcommand\pbBReach[1]{\pb{bReachability}{#1}}
\newcommand\pbCoverage[1]{\pb{Coverage}{#1}}
\newcommand\pbBCoverage[1]{\pb{bCoverage}{#1}}

\newcommand{\idDir}{dir}	% identifier for directed topological graph
\newcommand{\idNC}{nc}		% identifier for neighbor-communicable topological graph
\newcommand{\idSM}{sm}		% identifier for sight-moveable topological graph
\newcommand{\idCC}{cc}		% identifier for complete-communication topological graph
\newcommand{\idUnd}{und}		% identifier for complete-communication topological graph

\newcommand{\pbReachDir}{\pbReach{\idDir}}
\newcommand{\pbBReachDir}{\pbBReach{\idDir}}
\newcommand{\pbCoverageDir}{\pbCoverage{\idDir}}
\newcommand{\pbBCoverageDir}{\pbBCoverage{\idDir}}

\newcommand{\pbReachNC}{\pbReach{\idNC}}
\newcommand{\pbBReachNC}{\pbBReach{\idNC}}
\newcommand{\pbCoverageNC}{\pbCoverage{\idNC}}
\newcommand{\pbBCoverageNC}{\pbBCoverage{\idNC}}

\newcommand{\pbReachSM}{\pbReach{\idSM}}
\newcommand{\pbBReachSM}{\pbBReach{\idSM}}
\newcommand{\pbCoverageSM}{\pbCoverage{\idSM}}
\newcommand{\pbBCoverageSM}{\pbBCoverage{\idSM}}

\newcommand{\pbReachCC}{\pbReach{\idCC}}
\newcommand{\pbBReachCC}{\pbBReach{\idCC}}
\newcommand{\pbCoverageCC}{\pbCoverage{\idCC}}
\newcommand{\pbBCoverageCC}{\pbBCoverage{\idCC}}

\newcommand{\pbReachUnd}{\pbReach{\idUnd}}
\newcommand{\pbBReachUnd}{\pbBReach{\idUnd}}
\newcommand{\pbCoverageUnd}{\pbCoverage{\idUnd}}
\newcommand{\pbBCoverageUnd}{\pbBCoverage{\idUnd}}

\newtheorem{theorem}{Theorem}
\newtheorem{fact}[theorem]{Fact}
\newtheorem{definition}[theorem]{Definition}
\newtheorem{proposition}[theorem]{Proposition}
%\newtheorem{proof}{Proof}
\newtheorem{lemma}[theorem]{Lemma}

\tikzstyle{copyofG} = [fill=gray!20!white, draw=none, rounded corners=8pt]
\tikzstyle{communication} = [dash pattern=on 1pt off 1pt] %[line cap=rect, dot diameter=1pt, dot spacing=2pt, dots] %[dash pattern=on 2pt off 2pt on \the\pgflinewidth off 2pt]
\tikzstyle{tobehere} = [color=blue!50!white]
\newcommand\set[1]{\{#1\}}
\newcommand\suchthat{\mid}
\newcommand{\nbdrones}{n}

\newcommand{\booleanvariablegeneric}{{x}}
\newcommand{\clausegeneric}{\mathsf{c}}

\newcommand{\ok}{\includegraphics[height=2mm]{check.png}}
\newcommand{\imgDrone}{\includegraphics[height=2mm]{drone_black.png}}

\usetikzlibrary{decorations.pathmorphing}
\usetikzlibrary{decorations.pathreplacing}
\usetikzlibrary{arrows}
\usetikzlibrary{positioning}

\makeatletter
\tikzset{
	dot diameter/.store in=\dot@diameter,
	dot diameter=3pt,
	dot spacing/.store in=\dot@spacing,
	dot spacing=10pt,
	dots/.style={
		line width=\dot@diameter,
		line cap=round,
		dash pattern=on 0pt off \dot@spacing
	}
}
\makeatother

%%% Local Variables:
%%% mode: latex
%%% TeX-master: "main"
%%% End:




\renewcommand{\smallskip}{}
\setlength{\textfloatsep}{3pt plus 1pt minus 2pt}
%\setlength{\intextsep}{8pt plus 1.0pt minus 2.0pt}
\setlength{\abovecaptionskip}{4pt plus 1pt minus 1pt}
\setlength{\parskip}{0.4em}
%\setlength{\abovedisplayskip}{2pt}
%\setlength{\belowdisplayskip}{2pt}
%\setlength{\abovedisplayshortskip}{1pt}
%\setlength{\belowdisplayshortskip}{1pt}
%\renewcommand{\bottomfraction}{0.8}

%\setlength{\textfloatsep}{8pt plus 1.0pt minus 2.0pt}
%\setlength{\intextsep}{8pt plus 1.0pt minus 2.0pt}
%\setlength{\abovecaptionskip}{0pt plus 1pt minus 1pt}

%\excludeversion{arxiv}
%\includeversion{pldi}

\excludeversion{pldi}
\includeversion{arxiv}

\sloppy
\begin{document}

%% Title information
\title{Sound Dynamic Deadlock Prediction in Linear Time}
%\title[Short Title]{Full Title}         %% [Short Title] is optional;
                                        %% when present, will be used in
                                        %% header instead of Full Title.
%\titlenote{with title note}             %% \titlenote is optional;
                                        %% can be repeated if necessary;
                                        %% contents suppressed with 'anonymous'
%\subtitle{Subtitle}                     %% \subtitle is optional
%\subtitlenote{with subtitle note}       %% \subtitlenote is optional;
                                        %% can be repeated if necessary;
                                        %% contents suppressed with 'anonymous'


%% Author information
%% Contents and number of authors suppressed with 'anonymous'.
%% Each author should be introduced by \author, followed by
%% \authornote (optional), \orcid (optional), \affiliation, and
%% \email.
%% An author may have multiple affiliations and/or emails; repeat the
%% appropriate command.
%% Many elements are not rendered, but should be provided for metadata
%% extraction tools.

\author{H\"{u}nkar Can Tun\c{c}}
\orcid{0000-0001-9125-8506}
\affiliation{
	\institution{Aarhus University}            %% \institution is required
	\country{Denmark}                    %% \country is recommended
}
\email{tunc@cs.au.dk}          %% \email is recommended


\author{Umang Mathur}
\orcid{0000-0002-7610-0660}
\affiliation{
  \institution{National University of Singapore}            %% \institution is required
  \country{Singapore}                    %% \country is recommended
}
\email{umathur@comp.nus.edu.sg}          %% \email is recommended

\author{Andreas Pavlogiannis}
\orcid{0000-0002-8943-0722}
\affiliation{
	\institution{Aarhus University}            %% \institution is required
	\country{Denmark}                    %% \country is recommended
}
\email{pavlogiannis@cs.au.dk}          %% \email is recommended


\author{Mahesh Viswanathan}
\orcid{0000-0001-7977-0080}
\affiliation{
	\institution{University of Illinois, Urbana Champaign}            %% \institution is required
	\country{USA}                    %% \country is recommended
}
\email{vmahesh@illinois.edu}

%\author{First1 Last1}
%\authornote{with author1 note}          %% \authornote is optional;
                                        %% can be repeated if necessary
%\orcid{nnnn-nnnn-nnnn-nnnn}             %% \orcid is optional
%\affiliation{
%  \position{Position1}
%  \department{Department1}              %% \department is recommended
%  \institution{Institution1}            %% \institution is required
%  \streetaddress{Street1 Address1}
%  \city{City1}
%  \state{State1}
%  \postcode{Post-Code1}
%  \country{Country1}                    %% \country is recommended
%}
%\email{first1.last1@inst1.edu}          %% \email is recommended

%% Author with two affiliations and emails.
%\author{First2 Last2}
%\authornote{with author2 note}          %% \authornote is optional;
                                        %% can be repeated if necessary
%\orcid{nnnn-nnnn-nnnn-nnnn}             %% \orcid is optional
%\affiliation{
%  \position{Position2a}
%  \department{Department2a}             %% \department is recommended
%  \institution{Institution2a}           %% \institution is required
%  \streetaddress{Street2a Address2a}
%  \city{City2a}
%  \state{State2a}
%  \postcode{Post-Code2a}
%  \country{Country2a}                   %% \country is recommended
%}
%\email{first2.last2@inst2a.com}         %% \email is recommended
%\affiliation{
%  \position{Position2b}
%  \department{Department2b}             %% \department is recommended
%  \institution{Institution2b}           %% \institution is required
%  \streetaddress{Street3b Address2b}
%  \city{City2b}
%  \state{State2b}
%  \postcode{Post-Code2b}
%  \country{Country2b}                   %% \country is recommended
%}
%\email{first2.last2@inst2b.org}         %% \email is recommended


%% Abstract
%% Note: \begin{abstract}...\end{abstract} environment must come
%% before \maketitle command

\begin{abstract}

Constraint Programming (CP) and Machine Learning (ML) face challenges in text generation due to CP's struggle with implementing ``meaning'' and ML's difficulty with structural constraints. This paper proposes a solution by combining both approaches and embedding a Large Language Model (LLM) in CP. The LLM handles word generation and meaning, while CP manages structural constraints. This approach builds on 
%
GenCP, an improved version of On-the-fly Constraint Programming Search (OTFS) using LLM-generated domains.
Compared to Beam Search (BS), a standard NLP method, this combined approach
%
(GenCP with LLM)
is faster and produces better results, ensuring all constraints are satisfied. This fusion of CP and ML presents new possibilities for enhancing text generation under constraints.


\end{abstract}


%% 2012 ACM Computing Classification System (CSS) concepts
%% Generate at 'http://dl.acm.org/ccs/ccs.cfm'.
\begin{CCSXML}
<ccs2012>
<concept>
<concept_id>10011007.10011074.10011099</concept_id>
<concept_desc>Software and its engineering~Software verification and validation</concept_desc>
<concept_significance>500</concept_significance>
</concept>
<concept>
<concept_id>10003752.10010070</concept_id>
<concept_desc>Theory of computation~Theory and algorithms for application domains</concept_desc>
<concept_significance>300</concept_significance>
</concept>
<concept>
<concept_id>10003752.10010124.10010138.10010143</concept_id>
<concept_desc>Theory of computation~Program analysis</concept_desc>
<concept_significance>300</concept_significance>
</concept>
</ccs2012>
\end{CCSXML}

\ccsdesc[500]{Software and its engineering~Software verification and validation}
\ccsdesc[300]{Theory of computation~Theory and algorithms for application domains}
\ccsdesc[300]{Theory of computation~Program analysis}


%% Keywords
%% comma separated list
\keywords{concurrency, runtime analyses, predictive analyses}  %% \keywords are mandatory in final camera-ready submission

% \ucomment{\begin{compactenum}
% \item We need some additional examples/figures in the technical parts of the paper
% \item Maybe mention in the contributions that we have one deadlock that is not sync-preserving (but we use correct-reorderings as the only possible witness)
% \item HCT: I think we never outline the differences between syncp and syncpd
% \end{compactenum}
% }

%% \maketitle
%% Note: \maketitle command must come after title commands, author
%% commands, abstract environment, Computing Classification System
%% environment and commands, and keywords command.
\maketitle


\section{Introduction}

One of the most fundamental problems in combinatorial optimization is the traveling salesperson problem (TSP), formalized as early as 1832 (c.f. \cite[Ch 1]{ABCC07}).
In an instance of  TSP we are given a set of $n$ cities $V$ along with their pairwise symmetric distances, $c:V\times V \to\R_{\geq 0}$. The goal is to find a Hamiltonian cycle of minimum cost. In the metric TSP problem, which we study here, the distances satisfy the triangle inequality. Therefore, the problem is equivalent to finding a closed Eulerian connected walk of minimum cost.%\footnote{Given such an Eulerian cycle, we can use the triangle inequality to shortcut vertices visited more than once to get a Hamiltonian cycle.}

It is NP-hard to approximate TSP within a factor of $\frac{123}{122}$ \cite{KLS15}.  An algorithm of Christofides-Serdyukov~\cite{Chr76,Ser78} from four decades ago gives a $\frac32$-approximation for TSP.
Over the years there have been numerous attempts to improve the Christofides-Serdyukov algorithm and exciting progress has been made for various special cases of metric TSP, e.g., \cite{OSS11,MS11,Muc12,SV12,HNR21, KKO20, HN19, GLLM21}.
 Recently, ~\cite{KKO21} gave the first improvement for the general case by demonstrating that the so-called ``max entropy" algorithm of \cite{OSS11} gives a randomized $\frac{3}{2}-\epsilon$ approximation for some $\epsilon > 10^{-36}$.% (see \cite{VS20} for a historical note about TSP)

%After a long line of work %~\cite{Wol80,SW90,BP91,Goe95,CV00,GLS05,BM10,BC11,SWV12, HNR17,HN19, KKO20a} 
	%the best known approximation algorithm for the general case of the problem is $\frac{3}{2}-\epsilon$ for some $\epsilon > 10^{-36}$ due to ~\cite{KKO21}, a result that built upon the work of the third author, Saberi, and Singh ~\cite{OSS11}. 
	The method introduced in \cite{KKO21} exploits the optimum solution to the following linear programming relaxation of metric TSP studied by \cite{DFJ59,HK70,BG93}, also known as the subtour elimination LP:
\begin{equation}\label{eq:tsplp}
\begin{aligned}
	\min \quad& \sum_{u,v} x_{\{u,v\}} c(u,v)& \\
	\text{s.t.,} \quad &  \sum_{u} x_{\{u,v\}} = 2&\forall v\in V,\\
	& \sum_{u\in S, v\notin S} x_{\{u,v\}}\geq 2,&\forall S \subsetneq V, S\not= \emptyset\\
	& x_{\{u,v\}}\geq 0 &\forall u,v\in V.
\end{aligned}	
\end{equation} 
	
	 However, ~\cite{KKO21} did not show that the integrality gap of the subtour elimination polytope is bounded below $\frac{3}{2}$, and therefore did not make progress towards the ``4/3 conjecture" which posits that the integrality gap of LP \eqref{eq:tsplp} is $\frac{4}{3}$. In this work we remedy this discrepancy by proving the following theorem, improving upon the bound of $\frac{3}{2}$ from Wolsey~\cite{Wol80} in 1980:

\begin{theorem}\label{thm:main}
	Let $x$ be a solution to LP \eqref{eq:tsplp} for a TSP instance. For some absolute constant $\epsilon > 10^{-36}$, the \hyperlink{tar:alg}{max entropy algorithm} outputs a TSP tour with expected cost at most $\frac{3}{2}-\epsilon$ times the cost of $x$. Therefore the integrality gap of the subtour elimination LP is at most $\frac{3}{2} - \epsilon$. 
\end{theorem} 

To prove \cref{thm:main}, we amend Section 4 of \cite{KKO21} but keep the remainder of the analysis essentially the same. Unlike \cite{KKO21}, this argument now preserves the integrality gap by avoiding the use of the optimum solution in bounding the cost of the matching. See \cref{sec:overview} for a discussion of our new approach.
%We note that the analysis in this paper is not specialized to the max entropy algorithm (although we rely on many results from \cite{KKO21} to obtain \cref{thm:main} itself); instead, it is valid for any algorithm which samples a spanning tree from the support of a solution to LP \eqref{eq:tsplp} and then adds the minimum cost matching on the odd degree vertices of the tree.  
%Instead, we use the polygon representation of near minimum cuts \cite{Ben95,BG08} to bound  the cost of the matching (see the following section for an overview of our new findings). %An added benefit of avoiding the use of OPT in the analysis is  %We remark this makes the analysis constructive 
%We remark that this allows future analyses to explicitly compute and possibly utilize the relevant laminar family of near minimum cuts (whereas previously one needed to know OPT to find the laminar family used in the analysis in \cite{KKO21}).
%In particular, we show that to get a bound better than $\frac{3}{2}$ for this class of algorithm it is (essentially) sufficient to handle the case in which the near minimum cuts of $x$ are a laminar family.

\subsection{Other Consequences}
\paragraph{Path TSP} In recent exciting work, Traub, Vygen, Zenklusen \cite{TVZ20} showed that an $\alpha$-approximation algorithm for metric TSP can be used as a black box to get a $\alpha(1+\eps)$ approximation algorithm for Path TSP. This together with \cite{KKO21} implies that there is a $3/2-\eps$ approximation algorithm for Path TSP (for $\eps>10^{-36}$). On the other hand, it is a folklore result that the integrality gap of the natural LP relaxation of Path TSP is at least $3/2$.  Therefore, a consequence of the above theorem is that although the best possible approximation factors of the two problem are the same (up to polynomial reductions), the natural LP relaxation of metric TSP has a strictly smaller integrality gap.


\paragraph{2-ECSM} In the 2-edge-connected multi-subgraph problem, or 2-ECSM for short, we are given a weighted graph $G$ and we want to find a minimum cost 2-edge-connected spanning subgraph, where an edge can be chosen multiple times.
The classical Christofides-Serdyukov algorithm gives a 3/2-approximation for 2-ECSM and despite significant attempts \cite{CR98,BFS16,SV14,BCCGISW20} improved algorithms were designed only for special cases of the problem.
Since in \cite{BG93} it is shown that LP \eqref{eq:tsplp} is a valid relaxation for 2-ECSM, we obtain:

\begin{corollary}	
For some absolute constant $\epsilon > 10^{-36}$ the \hyperlink{tar:alg}{max entropy algorithm} is a randomized $\frac{3}{2}-\epsilon$ approximation for the 2-edge-connected multi-subgraph problem.
\end{corollary}
%Beyond these theorems, we believe the analysis in this paper will open new avenues to improve the arguments in ~\cite{KKO21}. The analysis in that work is by nature non-constructive because it uses information about the optimal solution. Here we remove this weakness and could in principle construct the proposed fractional matching in polynomial time. Although of course this has no practical benefit since the algorithm always finds the minimum cost matching, this may allow future works to manipulate the algorithm to better serve the analysis.

%We analyze the max-entropy rounding algorithm introduced in \cite{OSS11} and slightly modified in \cite{KKO20, KKO21}. 

%In other words, we design a feasible vector for the $O$-join polytope to ``satisfy'' all near min cuts ``crossed on both  sides'' 


%Whereas Section 4 of ~\cite{KKO21} only deals with the near minimum cuts of $x$ (where $x$ is a solution to LP \eqref{eq:tsplp}) which lie along the optimal Hamiltonian cycle, we deal with all near minimum cuts of $x$ using the so-called polygon representation of near minimum cuts ~\cite{Ben97,BG08}. %The results give new intuition for the structure of cuts that are within $\frac{6}{5}$ or less of the edge connectivity of the graph.

 %: we show that we can incur a cost of $O(\eta^2) \cdot c(x)$ to ensure that the set of cuts with $x(\delta(S)) \le 2+\eta$ is a laminar family.


\subsection{New techniques and contributions}\label{sub:newtechniques}

This paper can be seen as a case study on how to reason about and deal with {\em near} minimum cuts. One can deduce from the classical cactus representation of a graph $G$ \cite{DKL76} (i) the structure of {\em all} min cuts of $G$ and (ii) the structure of the edges of $G$ in the sense that every edge $\{u,v\}$ maps to a unique {\em path} in the cactus between the images of $u$ and $v$. Furthermore, such a path intersects every cycle of the cactus on at most one cactus edge. The theory has found many application from designing fast algorithms
\cite{Kar00,KP09} to the analysis of approximation algorithms for TSP \cite{KKO20} and connectivity augmentation \cite{BGJ20,CTZ21}.

Two decades later, the theory of min cuts was extended to near min cuts in works of Bencz\'ur and Goemans \cite{Ben95, BG08} where they introduced the polygon representation which represents all cuts of a graph with at most $\frac{6}{5}k$ edges, where $k$ is its edge connectivity. Although these works completely classify the structure of all near min cuts of a given graph $G$, they do not characterize the structure of the \textit{edges} of $G$ with respect to these cuts, which can be important in applications (for example, in many of the recent applications of min cuts,
 one also needs to exploit the structure of the edges in relation to the cactus).
The structure on the edges turns out to be highly relevant in this work as well, and as a byproduct of our analysis we make progress towards classifying the way in which the edges of $G$ relate to the structure of the polygon representation.
 
 % and (to some extent) a classification of the set of edges of $G$ with respect to the polygon representation of Bencz\'ur and Goemans.
 
  %i
 %s to give a better understanding of the structure of edges of $G$ with respect to its near min cuts.

  %One can partition the edges of $G$ into sets $F_1\dots,F_m$ such that the set of edges in every min cut $(S,\overline{S})$ of $G$ is the union of edges in a pair $F_i,F_j$ for $i\ neq j$.
%\Nathan{Shayan can add something} For example...

For motivation, consider a generic family of network design problems in which we want to construct a network such that every pair $u,v$ of vertices has connectivity at least $c_{u,v}$. A natural approach is to write an LP relaxation to find a (minimum cost) vector $x: E \to \R_{\ge 0}$ such that for every cut $S$ separating $u$ and $v$, $x(\delta(S))\geq c_{u,v}$. We can round this LP using independent rounding or a dependent rounding scheme such as sampling from max entropy distributions. Using classical concentration bounds one can show that if $x(\delta(S))\gg c_{u,v}$ then with high probability the rounded solution has at least $c_{u,v}$ edges across this cut. So the main challenge is to ``fix'' near tight cuts, i.e., cuts where $x(\delta(S))\approx c_{u,v}$.  For an explicit instantiation of this scheme see \cite{KKOZ22}. A better understanding of the global structure of the family of near tight cuts has the potential to significantly simplify or even improve the approximation factor of such rounding algorithms. A classical technique to design algorithms for such network design problems is to apply uncrossing to extreme point solutions of the LP. One can view our contribution as an approximate uncrossing technique that deals with all near tight cuts (instead of just tight cuts) as we explain next.
%Next, we explain how our main theorem can be used to give global structure for near tight cuts in the case that $c_{u,v}=2$ for all $u,v$ and we contrast it with the classical uncrossing technique which only deals with tight/min cuts. 


\paragraph{An Approximate Uncrossing Technique.} A fundamental technique in the field of approximation algorithms is the uncrossing technique\footnote{See e.g. \cite{LRS11} for a number of applications of this technique.} of Jain \cite{Jai01}. Given a graph $G=(V,E)$,  a weight vector $x:E\to\R_{\geq 0}$, and a  function $f:V\to\R$, suppose that $x(\delta(S))\geq f(S)$ for all $S\subseteq V$. Let $\cN$ be the family of sets $S$ such that $x(\delta(S)) = f(S)$, i.e., the family of {\em tight} sets with respect to $f$. The uncrossing technique says that if $f$ is (weakly) supermodular then we can refine $\cN$ to a laminar family of sets, $\cH$, such that if all sets of $\cH$ are tight, then all sets of $\cN$ are tight as well. For a concrete example, suppose $f$ is a constant function, say $f(S)=2$ for all $\emptyset\subsetneq S\subsetneq V$. Then, sets of $\cH$ can be constructed using the cactus representation \cite{DKL76} of cuts in $\cN$. The significance of this method is that if $x$ is a basic feasible solution to a LP with constraints $x(\delta(S))\geq f(S)$ for all $S$, one can use this machinery to argue that the support of $x$ has size $O(|V|)$.

Informally, we prove the following, which 
can be seen as  an {\em approximate uncrossing technique}: 
\begin{theorem}[Informal]\label{thm:uncrossing}Suppose we have a vector $x:E\to\R_{\geq 0}$ such that $x(\delta(S))\geq f(S)$ for all $S$; define $\cN$ to be sets $S$ where $x(\delta(S))\leq f(S)(1+\eps)$ for some fixed $\eps>0$. If $f(.)$ is constant, say $f(S)=2$ for all $S$, then there is a set $\cN^*\subseteq \cN$ and a collection of edge sets $F_1,\dots,F_m\subseteq E$ such that the following hold:
\begin{itemize}
	\item $|\cN^*|= O(|V|)$, $m= O(|V|)$.
	\item $x(F_i)\geq 1-\eps/2$ for all $1\leq i\leq m$.
	\item Every edge $e$ is in at most $O(1)$ of the $F_i$'s.
	\item For every set $S\in \cN\smallsetminus \cN^*$ there exists $1\leq i<j\leq m$ such that $F_i\cap F_j=\emptyset$ and $F_i\cup F_j\subseteq \delta(S)$ and for every $S\in \cN^*$, there exists $1\leq i\leq m$ such that $F_i\subseteq \delta(S)$. 
\end{itemize}
\end{theorem}
In words, although we cannot simply refine $\cN$ to a linear number of sets, we can refine the edges in cuts of $\cN$ to a linear number of sets $F_1,\dots, F_m$ such  that we can essentially capture the edges of $\delta(S)$ for any $S\in \cN\smallsetminus \cN^*$ by a pair of disjoint $F_i$'s. We give a slightly weaker condition for cuts in $\cN^*$; namely we only capture half of their edges by $F_i$'s.

\begin{example}For a simple example of the above theorem, suppose $\eps=0$, i.e. $\cN$ is the set of min cuts of a graph $G$. Furthermore, suppose that every proper  cut in $\cN$ is \hyperlink{tar:crossing}{crossed} (recall that $S$ is proper if $1<|S|<|V|-1$) and that $\cN$ has at least one proper cut. 
Then, one can use an uncrossing technique, namely that if $A,B\in \cN$ then $A\cap B\in \cN$, to prove that $G$ must be cycle, namely we can order vertices of $G$, $v_0,\dots,v_{n-1}$ such that $x_{\{v_i,v_{i+1\text{ mod n}}\}}=1$.
In such a case we let $\cN^*=\emptyset$ and $F_i=E(v_i,v_{i+1\text{ mod }n})$.
%partition $V$ into sets $a_0,\dots,a_{m-1}$ such that 
%Let $\C$ be a connected component of crossing cuts of $\cN$, namely, for any pair of sets $A,B\in \C$ there is a path of crossing cuts all from $\C$ that goes from $A$ to $B$.
% and further suppose that $\cN$ can be represented by a cycle $C$ in the sense every min cut of $\cN$ corresponds to a min cut of $C$ and vice versa. Here we assume $a_0,\dots,a_{m-1}$ are the nodes of $C$ where each $a_i$ is identified with a disjoint set of vertices where $V=\uplus_{i=1}^m a_i$. In such a case, we can simply let $\cN^*=\emptyset$ and $F_i=E(a_i,a_{i+1\text{ mod }m})$. 
\label{eg:cycle}\end{example}

\begin{example}\label{eg:laminar}
For a second example, suppose again $\eps=0$ and $\cN$ is the set of mincuts of a graph $G$ where $\cN$ forms a laminar family (no two cuts cross). It turns out that we cannot decompose edges of cuts of $\cN$ into a linear sized collection of sets where every edge appears only a constant number of times. The main reason is that some edges may appear in an unbounded number of cuts. In this case we let $\cN^*=\cN$ and for every $A\in \cN$ (with immediate parent $B\in \cN$ in the laminar family) we add a set $F_A=\delta(A)\smallsetminus \delta(B)$  to our collection.  It is straightforward to show, using the structure of min cuts, that $x(F_A)\geq 1$; furthermore, since the size of a laminar family is linear in $V$, this gives a valid decomposition in the sense of above theorem.
\end{example}
Lastly, if $\eps=0$ and $\cN$ is the set of min cuts of an arbitrary graph, one can represent all min cuts of $\cN$ by a cactus \cite{DKL76} which can be seen as a tree of cycles. In such a case, one can use a construction similar to \cref{eg:cycle} for each cycle where instead of a vertex $v_i$ we have a set $a_i \subseteq V$ and one similar to \cref{eg:laminar} for the tree part of the cactus. For a concrete application of such a decomposition of min cuts see \cite{KKO20}.
%More generally, if $\cN$ corresponds to the set of min cuts of an arbitrary graph, the cuts of $\cN$ can be represented by a {\em cactus graph}. In such a case we add one $F_i$ for every edge of a cycle of the cactus. 


%and further for simplicity assume that there is a single connected component of crossing cuts in $\cN$, namely we can go from any $A$ to $B$ for $A,B\in\cN$ simply following crossing cuts of $\cN$. Then, one can represent cuts in $\cN$ by the set of min cuts of a cycle, namely we can contract vertices of $G$ 

%For a concrete application , suppose we need at least two edges in every set in $\cN^*$, say in a network optimization problem. Then, if we make sure that we have at least one edge in each $F_i$, all typical cuts constraints, $\cN\smallsetminus \cN^*$,  are satisfied, so we  reduce the problem to cuts in $\cN^*$. 


One of the main challenges in dealing with near min cuts relative to min cuts is that if $x(\delta(A)),x(\delta(B))\leq 2+\eps$ then $x(\delta(A\cap B))\leq 2+2\eps$. Therefore, if $\eps=0$, then min cuts are closed under intersection, set difference and union, but this is no longer true when $\eps>0$. So, to employ the classical uncrossing machinery one should be very careful to "uncross" only a constant number of times (independent of $\eps$) to make sure that every cut remains within $2+O(\eps)$. This is the main reason that the polygon representation of near min cuts (see below) is more sophisticated, e.g., we can no longer argue $x(E(a_i, a_{i+1}))\approx 1$, see \cref{fig:nearmincutbadexample}.

Although we don't study it here, we believe it may be worthwhile to find generalizations of \cref{thm:uncrossing} which hold for any (weakly) supermodular function.% That could be helpful in many questions based on the network optimization framework of Jain \cite{Jai01}.

\begin{remark} 
 We do not explicitly prove \cref{thm:uncrossing} in this extended abstract, as it is not used to prove \cref{thm:main}. However it can be deduced from arguments in \cref{sec:twoside} and \cref{app:oneside}. 
%In \cref{sec:overview} we discuss the main ideas of the proof of \cref{thm:uncrossing}. Here, let us explain the main challenge: In principal one might try to simply extend the above decomposition for the case $\eps=0$. The main challenge is that if $x(\delta(A)),x(\delta(B))\leq 2+\eps$ then $x(\delta(A\cap B))\leq 2+2\eps$. Therefore, if $\eps=0$, then min cuts are closed under intersection, set difference and union, but this is no longer true when $\eps>0$. So, to employ the classical uncrossing machinery one should be very careful to "uncross" only a constant number of times (independent of $\eps$) to make sure that every cut remains within $2+O(\eps)$. This is the main reason that the polygon representation of near min cuts (see below) is more sophisticated, e.g., we can no longer argue $x(E(a_i, a_{i+1}))\approx 1$, see \cref{fig:nearmincutbadexample}.
\end{remark}





\paragraph{Extensions to the Polygon Representation} To obtain our uncrossing framework we prove new properties of the polygon representation.
Given a graph $G=(V,E)$, let $k$ be the edge-connectivity of $G$, i.e. the number of edges in a minimum cut of $G$. For $\eps>0$, consider the set of $(1+\eps)$-near minimum cuts of $G$: cuts $(S,\overline{S})$ where $|E(S,\overline{S})| < (1+\eps)k$.
Bencz\'ur \cite{Ben95} and Bencz\'ur, Goemans \cite{BG08} proved that if $\eps \le 1/5$ then the near minimum cuts of $G$ admit a {\em polygon representation}. Namely, every connected component $\cC$ of \hyperlink{tar:crossing}{crossing} $(1+\eps)$ near min cuts can be represented by the diagonals of a convex polygon. In this polygon, the vertices of $G$ are partitioned into sets called \textit{atoms}, and every atom is mapped to a cell of this polygon defined by the diagonals and the boundary of the polygon itself (see \cref{sec:polyrep} for more details). 

The polygon representation can be seen as a generalization of the well-known cactus representation \cite{DKL76} of minimum cuts where a cycle of the cactus is replaced by a convex polygon. Unlike a cycle, some vertices/atoms map to the interior of the polygon, which are called ``inside'' atoms. The inside atoms at first look like a mystery and one can ask many questions about them such as how many can exist and what structures they can exhibit.



 Here, we explain two lemmas we proved which might find further applications beyond TSP in the future. 
%Our results give new intuition and understanding about the structure of polygon representations. These guide our analysis of the integrality gap of the subtour LP.
 %For example, one of our new observations is a 
 First, we give a necessary condition for a cell of a polygon to contain an inside atom:
\begin{lemma}[Informal, see \cref{thm:halfplanes}]
	Consider a polygon $P$ for a connected component $\C$ of a family of $1+\eps$ near min cuts for $\eps \le 1/5$ (where representing diagonals correspond to cuts in $\C$). Any cell of $P$ that has an inside atom must have at least $\Omega(1/\eps)$ many sides. 
\end{lemma}
This can be seen as a generalization of \cite[Lem 22]{BG08} to the case in which the cell is allowed to be adjacent to vertices of the polygon $P$.

Now, we explain our second extension: it follows from the cactus representation of minimum cuts that for a graph $G$ and a min cut $S$ one can partition the set of all min cuts that cross $S$ into two groups ${\cal A}=\{A_1,\dots,A_k\}$ and ${\cal B}=\{B_1,\dots,B_l\}$ for some $k,l\geq 0$ such that $S\cap A_1\subseteq S\cap A_2 \subseteq \dots S\cap A_k$ and, similarly, $S\cap B_1\subseteq \dots\subseteq S\cap B_l$. We prove a generalization of this fact for near min cuts:
\begin{lemma}[Informal, see \cref{lem:crosschain}]
Consider the set of $1+\eps$ near min cuts of a graph $G$ for $\eps\leq 1/10$; for any such near min cut $S$, one can partition the $1+\eps$ near min cuts crossing $S$ into two groups ${\cal A}=\{A_1,\dots,A_k\}$ and ${\cal B}=\{B_1,\dots,B_l\}$ such that $S\cap A_1 \subseteq S\cap A_2\subseteq \dots \subseteq S\cap A_k$ and similarly for cuts in ${\cal B}$.
\end{lemma}

\subsection{Outline of rest of paper} After reviewing preliminaries in \cref{sec:prelims}, we give a high-level overview of our proof technique in \cref{sec:overview}. The main new technical contributions of this paper are in \cref{sec:polyrep} and  \cref{sec:twoside}. The remaining content of the paper essentially follows from ~\cite{KKO21}. %Therefore, the reader may want to skip \cref{sec:proof-of-main}. 



%%  LaTeX support: latex@mdpi.com 
%  For support, please attach all files needed for compiling as well as the log file, and specify your operating system, LaTeX version, and LaTeX editor.

%=================================================================
\documentclass[drones,article,submit,moreauthors,pdftex]{mdpi}
% For posting an early version of this manuscript as a preprint, you may use "preprints" as the journal and change "submit" to "accept". The document class line would be, e.g., \documentclass[preprints,article,accept,moreauthors,pdftex]{mdpi}. This is especially recommended for submission to arXiv, where line numbers should be removed before posting. For preprints.org, the editorial staff will make this change immediately prior to posting.

%--------------------
% Class Options:
%--------------------
%----------
% journal
%----------
% Choose between the following MDPI journals:
% acoustics, actuators, addictions, admsci, adolescents, aerospace, agriculture, agriengineering, agronomy, ai, algorithms, allergies, alloys, analytica, animals, antibiotics, antibodies, antioxidants, applbiosci, appliedchem, appliedmath, applmech, applmicrobiol, applnano, applsci, aquacj, architecture, arts, asc, asi, astronomy, atmosphere, atoms, audiolres, automation, axioms, bacteria, batteries, bdcc, behavsci, beverages, biochem, bioengineering, biologics, biology, biomass, biomechanics, biomed, biomedicines, biomedinformatics, biomimetics, biomolecules, biophysica, biosensors, biotech, birds, bloods, blsf, brainsci, breath, buildings, businesses, cancers, carbon, cardiogenetics, catalysts, cells, ceramics, challenges, chemengineering, chemistry, chemosensors, chemproc, children, chips, cimb, civileng, cleantechnol, climate, clinpract, clockssleep, cmd, coasts, coatings, colloids, colorants, commodities, compounds, computation, computers, condensedmatter, conservation, constrmater, cosmetics, covid, crops, cryptography, crystals, csmf, ctn, curroncol, currophthalmol, cyber, dairy, data, dentistry, dermato, dermatopathology, designs, diabetology, diagnostics, dietetics, digital, disabilities, diseases, diversity, dna, drones, dynamics, earth, ebj, ecologies, econometrics, economies, education, ejihpe, electricity, electrochem, electronicmat, electronics, encyclopedia, endocrines, energies, eng, engproc, ent, entomology, entropy, environments, environsciproc, epidemiologia, epigenomes, est, fermentation, fibers, fintech, fire, fishes, fluids, foods, forecasting, forensicsci, forests, foundations, fractalfract, fuels, futureinternet, futureparasites, futurepharmacol, futurephys, futuretransp, galaxies, games, gases, gastroent, gastrointestdisord, gels, genealogy, genes, geographies, geohazards, geomatics, geosciences, geotechnics, geriatrics, hazardousmatters, healthcare, hearts, hemato, heritage, highthroughput, histories, horticulturae, humanities, humans, hydrobiology, hydrogen, hydrology, hygiene, idr, ijerph, ijfs, ijgi, ijms, ijns, ijtm, ijtpp, immuno, informatics, information, infrastructures, inorganics, insects, instruments, inventions, iot, j, jal, jcdd, jcm, jcp, jcs, jdb, jeta, jfb, jfmk, jimaging, jintelligence, jlpea, jmmp, jmp, jmse, jne, jnt, jof, joitmc, jor, journalmedia, jox, jpm, jrfm, jsan, jtaer, jzbg, kidney, kidneydial, knowledge, land, languages, laws, life, liquids, literature, livers, logics, logistics, lubricants, lymphatics, machines, macromol, magnetism, magnetochemistry, make, marinedrugs, materials, materproc, mathematics, mca, measurements, medicina, medicines, medsci, membranes, merits, metabolites, metals, meteorology, methane, metrology, micro, microarrays, microbiolres, micromachines, microorganisms, microplastics, minerals, mining, modelling, molbank, molecules, mps, msf, mti, muscles, nanoenergyadv, nanomanufacturing, nanomaterials, ncrna, network, neuroglia, neurolint, neurosci, nitrogen, notspecified, nri, nursrep, nutraceuticals, nutrients, obesities, oceans, ohbm, onco, oncopathology, optics, oral, organics, organoids, osteology, oxygen, parasites, parasitologia, particles, pathogens, pathophysiology, pediatrrep, pharmaceuticals, pharmaceutics, pharmacoepidemiology, pharmacy, philosophies, photochem, photonics, phycology, physchem, physics, physiologia, plants, plasma, pollutants, polymers, polysaccharides, poultry, powders, preprints, proceedings, processes, prosthesis, proteomes, psf, psych, psychiatryint, psychoactives, publications, quantumrep, quaternary, qubs, radiation, reactions, recycling, regeneration, religions, remotesensing, reports, reprodmed, resources, rheumato, risks, robotics, ruminants, safety, sci, scipharm, seeds, sensors, separations, sexes, signals, sinusitis, skins, smartcities, sna, societies, socsci, software, soilsystems, solar, solids, sports, standards, stats, stresses, surfaces, surgeries, suschem, sustainability, symmetry, synbio, systems, taxonomy, technologies, telecom, test, textiles, thalassrep, thermo, tomography, tourismhosp, toxics, toxins, transplantology, transportation, traumacare, traumas, tropicalmed, universe, urbansci, uro, vaccines, vehicles, venereology, vetsci, vibration, viruses, vision, waste, water, wem, wevj, wind, women, world, youth, zoonoticdis 

%---------
% article
%---------
% The default type of manuscript is "article", but can be replaced by: 
% abstract, addendum, article, book, bookreview, briefreport, casereport, comment, commentary, communication, conferenceproceedings, correction, conferencereport, entry, expressionofconcern, extendedabstract, datadescriptor, editorial, essay, erratum, hypothesis, interestingimage, obituary, opinion, projectreport, reply, retraction, review, perspective, protocol, shortnote, studyprotocol, systematicreview, supfile, technicalnote, viewpoint, guidelines, registeredreport, tutorial
% supfile = supplementary materials

%----------
% submit
%----------
% The class option "submit" will be changed to "accept" by the Editorial Office when the paper is accepted. This will only make changes to the frontpage (e.g., the logo of the journal will get visible), the headings, and the copyright information. Also, line numbering will be removed. Journal info and pagination for accepted papers will also be assigned by the Editorial Office.

%------------------
% moreauthors
%------------------
% If there is only one author the class option oneauthor should be used. Otherwise use the class option moreauthors.

%---------
% pdftex
%---------
% The option pdftex is for use with pdfLaTeX. If eps figures are used, remove the option pdftex and use LaTeX and dvi2pdf.

%=================================================================
% MDPI internal commands
\firstpage{1} 
\makeatletter 
\setcounter{page}{\@firstpage} 
\makeatother
\pubvolume{1}
\issuenum{1}
\articlenumber{0}
\pubyear{2022}
\copyrightyear{2022}
%\externaleditor{Academic Editor: Firstname Lastname}
\datereceived{} 
%\daterevised{} % Only for the journal Acoustics
\dateaccepted{} 
\datepublished{} 
%\datecorrected{} % Corrected papers include a "Corrected: XXX" date in the original paper.
%\dateretracted{} % Corrected papers include a "Retracted: XXX" date in the original paper.
\hreflink{https://doi.org/} % If needed use \linebreak
%\doinum{}
%------------------------------------------------------------------
% The following line should be uncommented if the LaTeX file is uploaded to arXiv.org
%\pdfoutput=1

%=================================================================
% Add packages and commands here. The following packages are loaded in our class file: fontenc, inputenc, calc, indentfirst, fancyhdr, graphicx, epstopdf, lastpage, ifthen, lineno, float, amsmath, setspace, enumitem, mathpazo, booktabs, titlesec, etoolbox, tabto, xcolor, soul, multirow, microtype, tikz, totcount, changepage, attrib, upgreek, cleveref, amsthm, hyphenat, natbib, hyperref, footmisc, url, geometry, newfloat, caption

\usepackage{xcolor}
\usepackage{subcaption}

%=================================================================
%% Please use the following mathematics environments: Theorem, Lemma, Corollary, Proposition, Characterization, Property, Problem, Example, ExamplesandDefinitions, Hypothesis, Remark, Definition, Notation, Assumption
%% For proofs, please use the proof environment (the amsthm package is loaded by the MDPI class).

%=================================================================
% Full title of the paper (Capitalized)
\Title{Drone Detection and Tracking in Real-Time by Fusion of Different Sensing Modalities}

% MDPI internal command: Title for citation in the left column
\TitleCitation{Title}

% Author Orchid ID: enter ID or remove command
\newcommand{\orcidauthorA}{0000-0000-0000-000X} % Add \orcidA{} behind the author's name
%\newcommand{\orcidauthorB}{0000-0000-0000-000X} % Add \orcidB{} behind the author's name

% Authors, for the paper (add full first names)
\Author{Fredrik Svanström $^{1,\dagger}$, Fernando Alonso-Fernandez $^{2,}$*\orcidA{0000-0002-1400-346X} and Cristofer Englund $^{2,3,\ddagger}$}

%\longauthorlist{yes}

% MDPI internal command: Authors, for metadata in PDF
\AuthorNames{Fredrik Svanström, Fernando Alonso-Fernandez and Cristofer Englund}

% MDPI internal command: Authors, for citation in the left column
\AuthorCitation{Lastname, F.; Lastname, F.; Lastname, F.}
% If this is a Chicago style journal: Lastname, Firstname, Firstname Lastname, and Firstname Lastname.

% Affiliations / Addresses (Add [1] after \address if there is only one affiliation.)
\address{%
$^{1}$ \quad Air Defence Regiment. Swedish Armed Forces\\
$^{2}$ \quad School of Information Technology, Halmstad University, SE 301 18 Halmstad, Sweden\\
$^{3}$ \quad RISE, Lindholmspiren 3A, SE 417 56 Gothenburg, Sweden}

% Contact information of the corresponding author
\corres{Correspondence: feralo@hh.se}

% Current address and/or shared authorship
\firstnote{Email: DroneDetectionThesis@gmail.com} 
\secondnote{Email: cristofer.englund@hh.se}
% The commands \thirdnote{} till \eighthnote{} are available for further notes

%\simplesumm{} % Simple summary

%\conference{} % An extended version of a conference paper

% Abstract (Do not insert blank lines, i.e. \\) 
\abstract{
Automatic detection of flying drones is a key issue where its presence, especially if unauthorized, can create risky situations or compromise security.
%
Here, we design and evaluate a multi-sensor drone detection system.
%
In conjunction with standard video cameras and microphone sensors, we explore the use of thermal infrared cameras, pointed out as a feasible and promising solution that is scarcely addressed in the related literature.
%
Our solution integrates a fish-eye camera as well to monitor a wider part of the sky and steer the other cameras towards objects of interest.
%
The sensing solutions are complemented with an ADS-B receiver, a GPS receiver, and a radar module. However, our final deployment has not included the latter due to its limited detection range.
%
The thermal camera is shown to be a feasible solution as good as the video camera, even if the camera employed here has a lower resolution. 
%
Two other novelties of our work are the creation of a new public dataset of multi-sensor annotated data that expands the number of classes compared to existing ones, as well as the study of the detector performance as a function of the sensor-to-target distance.
%
Sensor fusion is also explored, showing that the system can be made more robust in this way, mitigating false detections of the individual sensors.
%
}

% Keywords
\keyword{Drone detection, UAV detection, Anti-drone systems} 

% The fields PACS, MSC, and JEL may be left empty or commented out if not applicable
%\PACS{J0101}
%\MSC{}
%\JEL{}

%%%%%%%%%%%%%%%%%%%%%%%%%%%%%%%%%%%%%%%%%%
% Only for the journal Diversity
%\LSID{\url{http://}}

%%%%%%%%%%%%%%%%%%%%%%%%%%%%%%%%%%%%%%%%%%
% Only for the journal Applied Sciences
%\featuredapplication{Authors are encouraged to provide a concise description of the specific application or a potential application of the work. This section is not mandatory.}
%%%%%%%%%%%%%%%%%%%%%%%%%%%%%%%%%%%%%%%%%%

%%%%%%%%%%%%%%%%%%%%%%%%%%%%%%%%%%%%%%%%%%
% Only for the journal Data
%\dataset{DOI number or link to the deposited data set if the data set is published separately. If the data set shall be published as a supplement to this paper, this field will be filled by the journal editors. In this case, please submit the data set as a supplement.}
%\datasetlicense{License under which the data set is made available (CC0, CC-BY, CC-BY-SA, CC-BY-NC, etc.)}

%%%%%%%%%%%%%%%%%%%%%%%%%%%%%%%%%%%%%%%%%%
% Only for the journal Toxins
%\keycontribution{The breakthroughs or highlights of the manuscript. Authors can write one or two sentences to describe the most important part of the paper.}

%%%%%%%%%%%%%%%%%%%%%%%%%%%%%%%%%%%%%%%%%%
% Only for the journal Encyclopedia
%\encyclopediadef{For entry manuscripts only: please provide a brief overview of the entry title instead of an abstract.}

%%%%%%%%%%%%%%%%%%%%%%%%%%%%%%%%%%%%%%%%%%
\begin{document}

%%%%%%%%%%%%%%%%%%%%%%%%%%%%%%%%%%%%%%%%%%

% The order of the section titles is: Introduction, Materials and Methods, Results, Discussion, Conclusions for these journals: aerospace,algorithms,antibodies,antioxidants,atmosphere,axioms,biomedicines,carbon,crystals,designs,diagnostics,environments,fermentation,fluids,forests,fractalfract,informatics,information,inventions,jfmk,jrfm,lubricants,neonatalscreening,neuroglia,particles,pharmaceutics,polymers,processes,technologies,viruses,vision


\section{Introduction}

Drones, also referred to small and remotely controlled unmanned aerial vehicles (UAVs), can fulfil valuable societal roles such as law enforcement, medical, construction, search and rescue, parcel delivery, remove area exploration, topographic mapping, forest/water management, or inspection of big infrastructures like power grids \cite{drones6060147}. 
%
%Examples include delivery of defibrillators \cite{[Sanfridsson19]}, fire fighting \cite{INNOCENTE201980}, law enforcement, construction, search and rescue, parcel delivery, etc. 
%
Their low cost and ease of operation have caused drones to find their way into consumer use just for recreation and entertainment as well \cite{[FAI]}.
%
Unfortunately, they can also be intentionally or unintentionally misused, threatening the safety of others. %In the worst case,
For example, an aircraft can be severely damaged if it collides with a consumer-sized drone, even at moderate speeds %, as shown by researchers at the University of Dayton
\cite{[Dayton]}, and an ingested drone can rapidly disable an aircraft engine.
%
An increasingly common risk is the report of drone sightings in restricted airport areas, which ultimately has led to total closure of the airport and cancellation of hundreds of flights \cite{UAVincidents}. 
%
Several near-misses and verified collisions with UAVs have involved hobbyist drone operators violating aviation safety regulations, sometimes without knowledge.
%
This rapid development in the use and potential misuse of UAVs has consequently produced an increase in research on drone detection \cite{Taha19,[GoogleTrends]} to counteract potential risks due to intrusion in restricted areas, either intentional or unintentional. 
%

In this work, we address the design and evaluation of an automatic multi-sensor drone detection and tracking system.
%
Our developments are built on state-of-the-art machine learning techniques, extending methods from conclusions and related literature recommendations \cite{Guvenc18,Taha19}.
%
In addition to effective detection, classification and tracking methods, the existing literature also points out sensor fusion as a critical open area to achieve more accuracy and robustness compared to a single sensor.
%
Despite this, research in sensor fusion for drone detection is scarce \cite{Samaras_2019,Guvenc18,Diamantidou19,Shi18}.
%
Our work also encompasses collecting and annotating a public dataset to accomplish training and evaluation of the system.
%
Another fundamental challenge is the lack of public reference databases that serve as a benchmark for researchers \cite{Taha19}.
%
Three different consumer-grade drones are included in the dataset together with birds, airplanes and helicopters, in which constitutes the published dataset with the biggest number of target classes (drone, bird, airplane and helicopter, in comparison to others which only contain three, two or one of these classes only).
%
To achieve effective detection of drones, in building the classes, we have considered including other flying objects that are likely to be mistaken for a drone \cite{Saqib17,Aker17}.
%
Another missing piece in previous studies that we address here is the system's classification performance as a function of the distance to the target, with annotations of the database including such information as well.

A preliminary version of this article appeared at a conference~\cite{drone20icpr}.
%
In the present contribution, we substantially increase the amount of reported results, e.g. the previous paper only reported the precision, recall and F1-score of the individual sensors provided in Tables~\ref{tab:results-IRcam}, \ref{tab:results-Vcam} and \ref{tab:results-audio}, as well as the fusion results of Figure~\ref{fig:fusion-results1}.
%
Here, we extensively analyze the effect of internal parameters of the different detectors on their performance for the various sensors.
%
We also report results with a radar module and provide comments about the fish-eye camera motion detector, all of them missing in the previous publication.
%
Additional results on the fusion of sensors are also provided, including an Appendix with complementary observations and visual examples.
%
New detailed information about the system architecture, hardware and software employed is also provided, including details about implementation and design choices not included in the previous publication.
%
The related work is also described in more detail.

The rest of the paper is organized as follows.
%
Section~\ref{sect:soa} describes related work.
%
Section ~\ref{sect:materialsmethods} extensively describes the proposed system, including the architecture, the hardware components, the involved software, the Graphical User Interface, and the dataset. 
%
The experimental results are presented and discussed in Section~\ref{sect:results}. 
%
Finally, the conclusions are given in Section~\ref{sect:conclusions}.


%%%%%%%%%%%%%%%%%%%%%%%%%%%%%%%%%%%%%%%%%%
\section{Related Work}
\label{sect:soa}

%To automatically detect flying drones, some form of sensor system is needed. As shown in \cite{Samaras_2019}, the fusion of data from multiple sensors %, i.e. using several sensors in combination 
%to achieve more accurate results than derived from single sensors, while compensating for their individual weaknesses, is well-founded when it comes to the drone detection task.

Fusing data from multiple sensors allows for more accurate results than a single sensor, while compensating for their individual weaknesses \cite{Samaras_2019}.
%
The sensors used for drone detection include:
%that may be considered, %for drone detection and classification tasks, 
%and hence can be found in the related scientific literature, are:
$i$) radar (on several different frequency bands, both active and passive), $ii$) cameras in the visible spectrum, $iii$) cameras detecting thermal infrared emission (IR), $iv$) microphones to detect acoustic vibrations, i.e. sound, $v$) sensors to detect radio frequency signals to and from the drone and the controller (RF), and $vi$) scanning lasers (Lidar). As mentioned in \cite{Samaras_2019} and explored further in \cite{Boddhu13}, even humans are employed for the task. It has also been successfully demonstrated that animals can be trained for this role \cite{[GuardFromAboveBV]}.
%
Systems for drone detection utilizing one or more of the sensors mentioned above may also be combined with some effector to try to bring the flying drone down or take control of it in some other way. An effector component, however, is not part of this work.

An introduction to the subject and a comparison of drone detection and tracking techniques is given in the 2018 paper \cite{Guvenc18}. It highlights as open research the use of fusion techniques to exploit data from multiple sensors, and the development of effective machine learning techniques for detection, classification and tracking in various scenarios of interest. 
%
%An introduction to the subject of drone detection can be found in \cite{Guvenc18}, published in 2018. It contains a comparison of advantages and disadvantages of different drone detection and tracking techniques. Regarding sensor fusion, the authors state that ``For accurate and quick detection/tracking of UAVs, data fusion techniques that can simultaneously use information from multiple types of sensors carry critical importance, and this constitutes an open research area''.
%
%One of the conclusions of \cite{Guvenc18} is also that ``Future research directions include developing effective machine learning methods for detecting, classifying, and tracking drones in various scenarios of interest''. 
The paper also briefly provides an overview of ways to interdict unauthorized drones.

A comprehensive review of about 50 references is given in the 2019 paper \cite{Taha19}, which comments the different machine learning techniques based on the type of sensor, including its limitations.
%
The lack of public reference datasets is identified as an essential issue. Furthermore, no study analyzes the classification performance in relation to the distance to the drone. %, or the of the regression models. 
In addition, the sensing device, drone type, detection range, or dataset used is usually not specified, all key aspects to make works reproducible. 
%
%The paper \cite{Taha19} presents a comprehensive review of the literature on machine learning for drone detection and classification. Published in 2019, it summarizes the results of %the related work found in 
%about 50 references %into several tables, 
%based on the type of sensor, describing the specifications and limitations of the different techniques.
%
%The authors stress on more than one occassion the 
%
%%One of the take-home messages is that the authors, on more than one occasion, stresses the current 
%lack of publicly available reference datasets. Furthermore, they also write that ``No study was presented which investigated the classification performance as a function of drone distance not to speak of determining the range using regression models'', a direction which is pinpointed as ``very interesting''.
%
%They also highlight that most of the research in drone detection using visual sensors do not specify the sensing device, drone type, detection range, or dataset used, all being key aspects to make works reproducible. 
%
%%This can be a very interesting research area in the future'', and that ``Most of  the research in visual drone detection fails to specify the type of the acquisition device, the drone type, the detection range, and the dataset used in their research. These details are key to validate the work and make it comparable with related literature''.
%
%%A weakness of \cite{Taha19} is that it does not include any papers that make use of thermal infrared cameras and, as we will see next, this is something that will be of importance to this work.
%

Also from 2019, the paper \cite{Samaras_2019} has an exhaustive 178 references, not specific to drone detection and classification but also regarding the general techniques employed. 
%
%Within these, there are not only articles specific to drone detection and classification, but also regarding the general techniques employed, e.g. useful aspects of sensor fusion and machine learning.
%
It %directly influences the present work, 
emphasizes as well the lack of works that use %tackling UAV detection and classification using 
thermal cameras, despite the successful use of such sensors together with deep learning-based methods for general object detection. %As such, the authors motivate researchers to ``turn their attention to this novel subject''.
%
%%such as "For the case of thermal cameras, we present deep learning based methods for general object detection and clas- sification because, to the best of authors’ knowledge, there is currently no work tackling the UAV detection and classification problem", and "Deep learning based object detection and classification architectures have been suc- cessfully utilized on thermal imagery for generic targets yet not for UAVs. This could be the motivation for researchers to turn their attention to this novel subject".

\subsection{Sensors Detecting Thermal Infrared Emission}

The work \cite{ANDRASI2017183}, from 2017, does not utilize machine learning, but %. Instead, detection and classification is done by 
a human looking at the output video stream.
%
The sensor is a low-cost FLIR Lepton 80$\times$60 pixels thermal camera. Connected to a Raspberry Pi, the authors are able to detect three different drone types up to %a distance of 
100m. One conclusion of the paper is that the drone battery, and not the motors (as one may presume), is the most significant source of heat radiation.
%
With the background from this paper and the ones above, the present work will try to extend these findings using a higher resolution sensor (FLIR Boson with 320$\times$256 pixels) in combination with machine learning methods. % to build a system that automatically detects, classifies and tracks drones. 
The IR camera will also be combined with at least one additional sensor.

Thermal cameras combined with deep-learning detection and tracking are explored in the 2019 paper \cite{wang_chen_choi_kuo_2019}. The IR videos are of 1920$\times$1080, but the sensor is not specified.
%
Detection is done with a Faster-RCNN. Given the difficulties of acquiring enough data, it uses a modified Cycle-GAN (General Adversarial Network) to produce synthetic thermal training data.
%
Via precision-recall curves, the thermal detector is shown to achieve better performance than a visible sensor used for comparison.
%
Sadly, no details about the sensor-to-target distance are given.
%
The database is said to be public as ``USC drone detection and tracking dataset'', but without a link. The dataset is also mentioned in \cite{Wu18}, but the link in that paper is not working.
%
Compared to \cite{wang_chen_choi_kuo_2019}, the present paper uses three different drone types instead of one. We also expand the number of target classes to four and, additionally, we explore the detection performance as a function of sensor-to-target distance.

A thermal infrared camera is also used as one of the sensors in \cite{Diamantidou19}, but %just as pointed out in \cite{Taha19}, 
the paper fails to specify the type, field of view or even the resolution of the sensors used, so even if there are useful tables of the results, any comparison is unfortunately hard to achieve.

\subsection{Sensors in the Visible Range}

A widespread method to detect drones %, as described in recent scientific papers, 
is to combine a video camera with a detector based on a convolutional neural network (CNN).
%
The paper \cite{Park17}, from 2017, studies six different CNN models, providing metrics for training time, speed performance (frames per second) and precision-recall curves. The comparison shows that considering the speed and accuracy trade-off, YOLOv2 \cite{Redmon17} seems to be the most appropriate detection model.

As it turns out from other studies on drone detection in the visible range, the YOLOv2 architecture is prevalent \cite{Wu18}, \cite{Liu_2018}, \cite{Saqib17} \cite{Aker17}. A lightweight version of the more recent YOLOv3 is utilized in \cite{Unlu19}.
%
The success of the YOLOv2 motivates its use in the present paper. This choice will also enable comparison to the papers mentioned above. %To be able to match the results from the thermal infrared sensor, the YOLOv2 is also used in that part.


The use of pan/tilt platforms to steer cameras in the direction of suspicious objects has also led to the use of wide-angle sensors. %, as in the present paper.
%
%The fact that the two primary sensors (IR and visible) of the drone detector system in the present paper are placed on a pan/tilt platform also arises the need for a wide-angle sensor to steer the sensors in directions where suspicious objects appear. 
%
In \cite{Unlu19}, a static camera with 110$^{\circ}$ field of view (FoV) is used together with a YOLOv3 detector to align a rotating narrow-field camera. To find objects of interest with the wide-angle camera, the paper \cite{Unlu19} employs a Gaussian Mixture Model (GMM) foreground detector \cite{Stauffer99}, a strategy also followed in the present paper. However, our wide-angle sensor has an even wider FoV (180$^{\circ}$). As pointed out in \cite{Unlu19}, this setup is prone to produce false alarms in some situations, but as described later, %in Section 3.1.3, 
this can be mitigated by tuning the associated detection and tracking algorithms.

Of all papers found using versions of the YOLO architecture for detection, \cite{Liu_2018} has the most output target classes with three (drone, airplane and helicopter), followed by \cite{Aker17} with two classes (drone and bird). However, none of the YOLO papers reports the detection accuracy as a function of the sensor-to-target distance.

\subsection{Acoustic Sensors}

Numerous papers have also explored the use of acoustic sensors. Some like \cite{Kim17}, \cite{Siriphun18}, \cite{Park15} utilize the Fast Fourier Transform (FFT) to extract features from the audio signals.
%
However, the Mel Frequency Cepstrum Coefficients (MFCC) seems to be the most popular technique, as used in \cite{Anwar19}, \cite{Liu17}, \cite{Jeon17}, \cite{Bernardini17}. The MFCC consists of a non-linear mapping of the original frequency according to the auditory mechanism of the human ear, and it is the most commonly used audio feature in current drone recognition tasks \cite{Liu_2018}.
%

When comparing %three different 
classification models, %for the drone detection task, 
the authors of \cite{Jeon17} conclude that Long Short-Term Memory (LSTM) \cite{Hochreiter97} achieves the best performance and F1-score. In that paper, the classification is binary (drone or background). The present work expands the range of output classes of \cite{Jeon17} by adding a helicopter class.
%
Also, the maximum acoustic detection range in the reviewed papers is 290 m, using a 120-element microphone array and a DJI Phantom 2 drone \cite{Busset15}. It is worth noting that the DJI Flame Wheel F450, one of the drones used in this work, is detected at a distance of 160 m by the microphone array.

\subsection{Radar}

Since radar is the most common technology to detect flying objects, it is not far-fetched to apply it to drones. 
%
However, a system designed to detect aircrafts often has features to reduce unwanted echoes from small, slow and low-flying objects, which is precisely what characterises UAVs. 
%
%Doing this with a radar system designed to detect aircrafts is however not as straight forward as it might seem, since those systems often have features to reduce unwanted echoes from small, slow and low-flying objects, which is precisely what characterise the drones in question.
%
The small Radar Cross Sections (RCS) of medium-sized consumer drones are described in \cite{Patel18}, and from \cite{Herschfelt17} we have that the RCS of the DJI Flame Wheel F450 is -17 dBsm (0.02 m$^2$). The paper \cite{Gong19} points out that flying birds have similar RCS, which can lead to false targets. % when detecting drones.
%
The F450 drone is also used in \cite{Fuhrmann17}, where the micro-doppler characteristics of drones are investigated. These are typically echoes from the moving blades of the propellers, and they can be detected on top of the bulk motion doppler signal of the drone. Since the propellers are generally made from plastic, the RCS of these parts are even smaller, and in \cite{Patel18}, it is stated that the echoes from the propellers are 20 to 25 dB weaker than the drone body itself. Nevertheless, papers like \cite{Bjorklund18}, \cite{Drozdowicz16} and \cite{Rahman18} accompany \cite{Fuhrmann17} in exploring the possibility of classifying drones using the micro-doppler signature.

\subsection{Other Drone Detection Techniques}

Very few drones are autonomous in the flight phase. Generally, they are controlled by ground equipment, and often send information on some radio frequency (RF), which can be used to detect them as well. The three drones used here are all controlled in real-time. The information sent out range from just simple telemetry such as battery level (DJI Flame wheel F450), a live video stream (Hubsan H107D+), to both a video stream and extensive position and status information (DJI Phantom 4 Pro).
%
Utilizing the RF fingerprint is described in \cite{Birnbach17}, and in \cite{Shorten18}, a CNN is used with data from an antenna array so that the direction to the drone controller can be calculated within a few degrees. In \cite{Ezuma20}, signals from 15 different drone controllers are classified with an accuracy of 98.13\% using only three RF features (shape factor, kurtosis and variance) with a K-Nearest Neighbour (KNN) classifier.
%
The use of LiDAR (Light Detection And Ranging) and LADAR (LAser Detection And Ranging) has also been explored \cite{Kim_2018}, 
%in combination with background subtraction and a nearest neighbour classifier is shown to 
successfully detecting drones up to 2 km. % in \cite{Kim_2018}.

\subsection{Sensor Fusion}

The paper \cite{Samaras_2019} mentioned data fusion from multiple sensors as a way to improve accuracy, since sensor combination helps to compensate for individual weaknesses. 
%
As mentioned above, the benefits of sensor fusion are also pointed out in \cite{Guvenc18}.
%
Considering that missed detection of an intruding drone will bring more security threats than false alarms, the authors of \cite{Shi18} conduct audio, video and RF detection in parallel, using a logical OR operation to fuse the results of the three detectors.
%
They also highlight that combining such heterogeneous data sources is an open research question since the simple combination of the results separately obtained by audio, video, and RF surveillance can induce significant information loss. In contrast, it would be of greater significance to develop reliable fusion techniques at the feature level as well. 
%
%Furthermore under the section Challenges and Open Research Is- sues they conclude that "Although we have realized basic functions for drone surveillance in our anti-drone system, new challenges have been raised, and there are still some open research issues to be addressed: Heterogeneous Information Fusion: The result of drone detection should be more than a sim- ple combination of the results separately obtained by audio, video, and RF surveillance, which will induce great information loss. It is of great signif- icance to develop reliable techniques for the fusion of audio, video, and RF information from the aspect of feature fusion and decision fusion".
%
Without any further specifications of the sensors used besides that they are visible, thermal and 2D-radar, the work \cite{Diamantidou19} presents promising results from experiments using a multilayer perceptron (MLP) to perform sensor fusion in a drone detector/classifier with just one output class.
%
Just as in \cite{Samaras_2019}, the present paper also considers early and late sensor fusion and differentiates these two principles based on whether the sensor data is fused before or after the detection element.

\subsection{Drone Detection Datasets}

As mentioned, the work \cite{Taha19} points out the lack of publicly available datasets. This is also highlighted in \cite{Samaras_2019}, especially with thermal infrared cameras.
%
%They conclude that creating a dataset for UAV detection with thermal images might be budget-wise out of reach of many universities and research centers.
%
The latter paper also states the strong need for real-world UAV audio datasets. % that could serve as research benchmarks.

In the references of \cite{Samaras_2019}, there are two useful links to datasets for visible video detectors. One of these is \cite{Reiser20}, where 500 annotated drone images can be found. 
%
This is far away from the 203k images of our database (plus audio clips).
%
The other link leads to the dataset \cite{SafeShore} of the drone-vs-bird challenge held by the Horizon2020 SafeShore project consortium. However, the dataset is only available upon request and with restrictions to the usage and sharing of the data. The drone-vs- bird challenge is also mentioned in \cite{Saqib17}, \cite{Aker17} and by the winning team of the 2017 challenge \cite{Schumann17}.

The dataset used in \cite{Unlu19} is not available due to confidentiality. Since the French Ministry of Defence funded the work, one can presume that the dataset, in one way or another, is a property of the French Government or the French Armed Forces.



%%%%%%%%%%%%%%%%%%%%%%%%%%%%%%%%%%%%%%%%%%
\section{Materials and Methods}
\label{sect:materialsmethods}

This section describes the proposed methodology and outlines the automatic drone detection system, first on a system level and then in deeper detail. We detail the hardware components, how they are connected to the main computational resource, and the involved software running in the drone detection system, including the graphical user interface.
%
The methods used for the composition of our dataset are also described, including a description of the dataset, its division per sensor type, target class and sensor-to-target distance.
%
Our descriptive approach is motivated by \cite{Taha19}, which highlights that most works on visual drone detection do not specify the acquisition device, drone type, detection range, or dataset, all being key details that allow replication and comparison with other works.


\begin{figure} [htb]
\centering
\includegraphics[width=0.7\textwidth]{system.png}
\caption{System architecture, hardware and software parts. 
%Top: system architecture with hardware and software parts. Bottom left: main hardware parts of the detection system. Bottom right: 
%
The deployed system is shown just north of the runway at Halmstad airport (IATA/ICAO code: HAD/ESMT). Pictures were originally appearing in \cite{drone20thesis} and published in \cite{drone20icpr}. Reprinted with permission.
\label{fig:system}}
\end{figure}


\subsection{Architecture of the System}

A drone detection system must be able to both cover a large volume of airspace and have sufficient resolution to distinguish small objects like drones and tell them apart from other types of objects.
%
Combining wide and narrow field of view (FoV) cameras is one way to accomplish this \cite{Unlu19}. Another way, shown in \cite{Liu_2018}, is to use an array of high-resolution cameras. 
%
Here, we follow the first approach since our solution has been designed with portability in mind (Figure~\ref{fig:system}c).
%
Since the present paper uses only one infrared sensor with a fixed FoV, there is no possibility of having neither a wide-angle infrared sensor nor an array of such sensors. The proposed way to achieve the desired volume coverage with the IR-sensor is to have it on a moving platform, as shown in Figure~\ref{fig:system}b. This platform can either have objects assigned to it or search by itself at moments when the sensors are not busy detecting and classifying objects.
%
The overall system architecture, detailing the hardware and software components employed, is shown in Figure~\ref{fig:system}a. The layout of the Graphical User Interface (GUI) elements is shown in the bottom right part (boxes in blue). 


To be able to react to moving objects and also to have the ability to track those, the combined time constraints of the detection cycle and the control loop of the moving platform means that the system must work in close to real-time. Hence, all the detection and classification processes must be done efficiently and with as little delay as possible. The feedback loop of the moving platform must run at a sub-second speed.
%
In putting together such a system involving several sensors and mechanical parts, choosing the proper methods is critical. 
%
All these constraints, in turn, impose demands on the efficiency of the software as well.
%
Another aspect is that to detect the drones with high efficiency, the system must also recognize and track other flying objects that are likely to be mistaken for drones. For some of these drone-like objects, this is indigenous hard, e.g. birds. For others, it is technically possible since some of them announce their presence and location via radio, e.g. ADS-B, over which most aircrafts regularly transmit messages with varied content.
%
Combining the data from several sensors under the time constraints described above must be kept simple and streamlined too. This, together with the fact that very few papers have explored sensor fusion techniques, is the motivation to have a system where the inclusion and weights of the sensors can be altered at runtime to find a feasible setting.
%


\subsection{Hardware}

As primary electro-optical sensors, we use a thermal infrared camera (denoted as IRcam in Figure~\ref{fig:system}a) and a video camera (Vcam). 
%
Our system can keep track of cooperative aircrafts via ADS-B information that is made available with an antenna, which collects the aircraft’s position, velocity vectors and identification information broadcasted by aircrafts equipped with such a system.
%
We also include audio information through a microphone, which is employed to distinguish drones from other objects in the vicinity, such as helicopters.
%
All computations are made in a standard laptop, also used to present the results to the user via the designed GUI.
%

Since the primary cameras have a limited FoV, a fish-eye lens camera (Fcam) covering 180$^{\circ}$ horizontally and 90$^{\circ}$ vertically is also used. The fish-eye camera is used to detect moving objects in its FoV and then steer the IRcam and Vcam towards the detected objects using the pan/tilt platform. 
%
If the Fcam detects nothing, the platform can be set to move in two different search patterns to scan the sky.
%
As additional sensors, our system also includes a GPS receiver.
%
The hardware components are mounted on a standard surveyor tripod to provide stability to the system. This solution also facilitates the deployment of the system outdoors, as shown in Figure~\ref{fig:system}c. 
%
Due to the nature of the system, it must also quickly be transported to and from any deployment. Hence, % a transport solution is available, where 
the system can be disassembled into a few large parts and placed in a transport box.



\subsubsection{Thermal Infrared Camera (IRcam)}

We employ a FLIR Breach PTQ-136 using the Boson 320$\times$256 pixels detector (Y16 with 16-bit grey-scale). The FoV of the IRcam s 24$^{\circ}$ horizontally and 19$^{\circ}$ vertically.
%
%An example image taken from the IRcam video stream can be seen in Figure~\ref{fig:system_GUI}, bottom right.
%
It is worth noting that this sensor has a higher resolution than the FLIR Lepton sensor with 80$\times$60 pixels used in \cite{ANDRASI2017183}. 
%
In that paper, the authors detected three drone types up to a distance of 100m. However, it was done manually by a person looking at the live video stream. 
%
In contrast, the present paper employs an automatic detection solution.
%
The signal of the IRcam is sent to the laptop at 60 frames per second (FPS) using the USB-C port, which also powers the IRcam. 
%

\subsubsection{Video Camera (Vcam)}

To capture video in the visible range, we employ a Sony HDR-CX405 video camera.
%
The feed is taken from the HDMI port, which is captured with an Elgato Cam Link 4K frame grabber that provides a 1280$\times$720 video stream in YUY2-format (16 bits per pixel) at 50 FPS.
%
The FoV of the Vcam can be made wider or narrower using its adjustable zoom lens.
%
In this work, it is adjusted to have about the same FoV as the IRcam. 
%

\subsubsection{Fish-eye Lens Camera (Fcam)}

To counteract the limited FoV of the IRcam and Vcam, a fish-eye lens camera is used to monitor a wider area of the sky and then steer and focus these two towards the detected objects to ascertain if they are drones or something else. 
%
The fish-eye lens camera employed is an ELP 8 Megapixel with a FoV of 180$^{\circ}$ degrees, which provides a 30FPS 1024$\times$768 video stream in Mjpg-format at 30 FPS via USB.

\subsubsection{Microphone}

The microphone is used to distinguish drones from other objects emitting sounds, such as, for example, helicopters. Here, we use a Boya BY-MM1 mini cardioid directional microphone connected directly to the laptop.
%
Data is stored in .wav format, with a sampling frequency of 44100 Hz.


\subsubsection{ADS-B Receiver}

To track aircraft equipped with transponders, an ADS-B receiver is also used. This consists of an antenna and a NooElec Nano 2+ Software Defined Radio receiver (SDR). This is tuned to 1090 MHz so that the identification and positional data sent out as a part of the 1 Hz squitter message can be decoded and displayed. The Nano 2+ SDR receiver is connected to the laptop via USB.

\subsubsection{GPS receiver}

To correctly present the decoded ADS-B data, the system is equipped with a G-STAR IV BU-353S4 GPS receiver connected via USB. The receiver outputs messages following the National Marine Electronics Association (NMEA) standard.


\begin{figure} [htb]
\centering
\includegraphics[width=0.6\textwidth]{data_collection.jpg}
\caption{The data collection setup using a lighter version of the system. Picture originally appearing in \cite{drone20thesis}. \label{fig:data_collection}}
\end{figure}


\subsubsection{Pan/tilt platform and servo controller}

The pan/tilt platform is a Servocity DDT-560H direct drive tilt platform together with the DDP-125 Pan assembly, also from Servocity. To achieve the pan/tilt motion, two Hitec HS-7955TG servos are used.
%
A Pololu Mini Maestro 12-Channel USB servo controller is included so that the respective position of the servos can be controlled from the laptop. Since the servos have shown a tendency to vibrate when holding the platform in specific directions, a third channel of the servo controller is also used to give the possibility to switch on and off the power to the servos using a small optoisolated relay board.
%

To supply the servos with the necessary voltage and power, both a net adapter and a DC-DC converter are available. The DC-DC solution is used when the system is deployed outdoors, and, for simplicity, it uses the same battery type as one of the available drones.
%
Some other parts from Actobotics are also used in mounting the system, and the following have been designed and 3D-printed: adapters for the IR, video and fish-eye lens cameras, a radar module mounting plate and a case for the servo controller and power relay boards.

A lighter version of the IRcam and Vcam mounting without the pan/tilt platform has also been prepared. This is used on a lightweight camera tripod when collecting the dataset, simplifying transportation and giving the possibility to set its direction manually. The data collection setup is shown in Figure~\ref{fig:data_collection}.

An unforeseen problem occurring when designing the system was actually of mechanical nature. Even though the system uses a pan/tilt platform with ball-bearings and very high-end titanium gear digital servos, the platform was observed to oscillate in some situations. This phenomenon was mitigated by carefully balancing the tilt platform and introducing some friction in the pivot point of the pan segment. It might also be the case that such problems could be overcome using a servo programmer. Changing the internal settings of the servos could also increase their maximum ranges from 90$^{\circ}$ to 180$^{\circ}$. This would extend the volume covered by the thermal infrared and video cameras so that all targets tracked by the fish-eye camera could be investigated, not just a portion of them, as now.

\subsubsection{Computer}

A Dell Latitude 5401 laptop handles the computational part of the system. It is equipped with an Intel i7-9850H CPU and an Nvidia MX150 GPU. The computer is connected to all the sensors mentioned above and the servo controller using the built-in ports and an additional USB hub, as shown in Figure~\ref{fig:system}c.

It is observed, regarding the sensors, that there is a lower limit of around 5 FPS, where the system becomes so slow that the ability to track flying objects is lost. All actions taken by the system must be well balanced, and just such a simple thing as plotting the ADS-B panel with a higher frequency than necessary can cause a drop in the FPS rate. Such problems can be overcome by using more than one computational resource.

\subsubsection{Radar module}

As indicated earlier, our solution does not include a radar module in its final deployment. 
%
However, since one was available, it was included in our preliminary tests. 
%
It is a radar module from K-MD2, whose specifications are shown in Figure~\ref{fig:radar} \cite{RFbeam}.
%
Its exclusion was motivated by its short practical detection range.
%
Interestingly, the K-MD2 radar module is also used in another research project connected to drones \cite{Mostafa18}, not to detect drones, but instead mounted on board one as part of the navigation aids in GNSS2 denied environments.


\begin{figure} [htb]
\centering
\includegraphics[width=0.6\textwidth]{radar.jpg}
\caption{The radar module K-MD2. Picture taken from\\ https://www.rfbeam.ch/files/products/21/downloads/Datasheet\_K-MD2.pdf \label{fig:radar}}
\end{figure}


\subsection{Software}

Matlab is the primary development environment in which the drone detection system has been developed.
%
The software for drone detection consists of the main script and five separate `workers', one per sensor, as shown in Figure~\ref{fig:system}a.
%
The main script and the workers can run asynchronously and in parallel, enabled by the Matlab parallel computing toolbox. 
%
This allows each detector to run independently of the others. This also allows the different sensors to run asynchronously, handling as many frames per second as possible without inter-sensor delays and waiting time.
%
The main script communicates with the workers using pollable data queues.
%
The \textit{Fcam worker} utilizes a foreground/background detector via GMM background subtraction \cite{Stauffer99,gmmmatlab} and a multi-object Kalman filter tracker \cite{kalmanmatlab}. After calculating the position of the best-tracked target (defined as the one with the longest track history), it sends the azimuth and elevation angles to the main script, which then controls the pan/tilt platform, so that the moving object can be analysed further by the IR and video cameras.
%
The \textit{IRcam} and \textit{Vcam} workers are similar in their basic structure, and both import and run a trained YOLOv2 detector, fine-tuned with annotated ground truth to work with data from each camera.
%
The information sent to the main script is the region of the image where the object has been detected, the class of the detected target found in the ground truth of the training dataset, the confidence, and the horizontal and vertical offsets in degrees from the centre point of the image. The latter is used to calculate servo commands and track the object.
%
%
The \textit{Audio} worker sends information about the class and confidence to the main script. 
%
It uses a classification function built on an LSTM architecture, which is applied to MFCC features extracted from audio signals captured with the microphone. 
%
Unlike the others, the \textit{ADS-B} worker has two output queues, one consisting of current tracks and the other of the history tracks. 
%
The ``current'' queue contains the current positions and additional information (Id, position, altitude, calculated distance from the system, azimuth and elevation in degrees as calculated relative to the system, time since the message was received and the target category).
%
The ``history'' tracks queue is just a list of the old positions and altitudes.
%
This partition in two queues saves computational resources by reducing the amount of data from the worker drastically, so that only the information needed/used by the main process is sent. 
%
It also makes easier to control the data flow, since the length of the history track queue is easily set if it is separated.
%
%This is done so that the presentation clearly shows the heading and altitude changes of the targets.
%
All of the above workers also send a confirmation of the command from the main script to run the detector/classifier or to be idle. The number of frames per second currently processed is also sent to the main script.

Table~\ref{tab:output-classes}
shows the different classes that each worker can provide to the main script.
%
Note that not all sensors can output all the target classes. 
%
The audio worker has an additional ``background'' class, and the ADS-B will output a ``no data'' class if the vehicle category field of the received message is empty (since it is not a mandatory field of ADS-B messages).



\begin{table}[htb]
\caption{Output classes of the sensors and their corresponding class colours. Table originally appearing in \cite{drone20thesis}. \label{tab:output-classes}}

%\small
%\footnotesize
%\scriptsize
\begin{center}
\begin{tabular}{c}

\includegraphics[width=0.85\textwidth]{output_classes.png}  \\ 

\end{tabular}

\end{center}

\end{table}
\normalsize



%We use three machine learning techniques, two of them supervised, and one unsupervised.
%
%
%The second supervised machine learning technique is a classification function built on an LSTM-architecture \cite{Hochreiter97}, which is applied to features extracted from audio signals captured with a microphone. 
%
%Finally, an unsupervised technique is used to find moving objects in the FoV of the wide-angle camera (Fcam) via GMM background subtraction \cite{Stauffer99}.


\subsubsection{Main Script}

This is the core of the system. Besides starting the five workers (threads) and setting up the queues to communicate with these, it also sends commands to and reads data from the servo controller and the GPS receiver.
%
After the start-up sequence, the script goes into a loop that runs until the program is stopped by the user via the graphical user interface (GUI).

Updating the GUI and reading user inputs are the most frequent tasks on every loop iteration. The main script interacts with the workers and the servo controller at regular intervals. Servo positions are read, and queues are polled ten times a second. The system results, i.e. the system output label and confidence, are also calculated using the most recent results from the workers. Furthermore, at a rate of 5 Hz, new commands are sent to the servo controller for execution. Every two seconds the ADS-B plot is updated. Different intervals for various tasks make the script more efficient since, for example, an aircraft sends out its position via ADS-B every second. Hence, updating the ADS-B plots too often would only be a waste of computational resources. 

\subsubsection{IRcam worker}

Raw images are processed with Matlab function \texttt{imflatfield}\footnote{2-D image flat-field correction using Gaussian smoothing with a standard deviation of $\sigma$ to approximate the shading component of the input image. It cancels the artefacts caused by variations in the pixel-to-pixel sensitivity of the detector and by distortions in the optical path. As a result, the corrected image has more uniform brightness.} for 2-D image flat-field correction (with $\sigma$=30) followed by \texttt{imadjust}\footnote{Increases the contrast of the image by mapping intensity values so that 1\% of the data is saturated at low and high intensities} to increase image contrast.
%
Flat-field correction uses Gaussian smoothing to approximate the shading component of the input image.
%
Next, the input image is processed by the YOLOv2-detector, with a given detection threshold and the execution environment set to GPU. The output from the detector consists of an array of class labels, confidence scores and bounding boxes for all objects detected and classified. The detector output may be no data at all, or just as well, data for several detected objects. In this implementation, only the detection with the highest confidence score is sent to the main script.
%
Images from the IR camera are 320$\times$256 pixels.
%
To present the result in the GUI at the same size as the Vcam output, the image is resized to 640$\times$512.
%
Then, the bounding box, class label, and confidence score are inserted into the image. To clearly indicate the detected class, the inserted annotation uses the same colour scheme as in Table~\ref{tab:output-classes}.
%
The current state of the detector and its performance (frames per second) is also inserted in the top left corner of the image. Such information is also sent to the main script with the detection results.

The YOLOv2 detector is formed by modifying a pretrained MobileNetv2 following \cite{yolo2matlab} so that the first 12 layers out of 53 are used for feature extraction. 
%
The input layer is changed to 256$\times$256$\times$3. 
%
Six detection layers and three final layers are also added to the network. Besides setting the number of output classes of the final layers, the anchor boxes used are also specified.
%
To estimate the numbers of anchor boxes to use and the size of these, the training data is processed using the \texttt{estimateAnchorBoxes} Matlab function. This function uses a $k$-means clustering algorithm to find suitable anchor boxes sizes given the number of boxes to be used, returning as well the mean intersection-over-union (IoU) value of the anchor boxes in each cluster. We test the number of anchor boxes from one to nine to provide a plot over the IoU as a function of the number of boxes, as shown in Figure~\ref{fig:IRcam-anchors-yolo}.
%

\begin{figure} [htb]
\centering
\includegraphics[width=0.7\textwidth]{IRcam-anchors-yolo.png}
\caption{Plot used to assess the number of anchor boxes to be implemented in the IRcam worker YOLOv2 detector. Picture originally appearing in \cite{drone20thesis}.}
\label{fig:IRcam-anchors-yolo}
\end{figure}

When choosing the number of anchor boxes to use, the trade-off to consider is that a high IoU ensures that the anchor boxes overlap well with the bounding boxes of the training data, but, on the other hand, using more anchor boxes will also increase the computational cost and may lead to over-fitting. After assessing the plot, the number of anchor boxes is chosen to be three, and the sizes of these (with the scaling factor of 0.8 in width to match the downsize of the input layer from 320 to 256 pixels) are taken from the output of the \texttt{estimateAnchorBoxes} function.

The detector is trained with \texttt{trainYOLOv2ObjectDetector}\footnote{Returns an object detector trained using YOLO v2 with the specified architecture, in our case a pretrained MobileNetv2 \cite{yolo2matlab} with the modifications mentioned in the main text.} using data picked from the available dataset (Section~\ref{sect:results}).
%
%after the evaluation data have been selected and put aside (Section~\ref{sect:database}). The training data for the IRcam YOLOv2 detector consists of 120 video clips, each one just over 10 seconds and evenly distributed among all classes and all distance bins, making the total number of annotated images in the training set 37428. 
The detector is trained for five epochs using the stochastic gradient descent with momentum (SGDM) optimizer and an initial learning rate of 0.001.
%
Using a computer with an Nvidia GeForce RTX2070 8GB GPU, the time for one epoch is 39 min. 
%
The training function includes pre-processing augmentation consisting of horizontal flipping (50\% probability), scaling (zooming) by a factor randomly picked from a continuous uniform distribution in the range [1, 1.1], and random colour jittering for brightness, hue, saturation, and contrast.

\subsubsection{Vcam worker}

To be able to compare their results, the methods and settings of the Vcam worker are very similar to the IRcam worker above, with some exceptions. 
%
Images are of 1280$\times$720 pixels, which are then resized to 640$\times$512 without further pre-processing. 
%
Given the bigger image size, the input layer of the YOLOv2 detector here has a size of 416$\times$416$\times$3. 
%
%The training set here consists of 37519 images, with the detector trained for five epochs as well.
%
Due to the increased image size, the training time is also extended compared to the IR case. When using a computer with an Nvidia GeForce RTX2070 8GB GPU, the time for one epoch is 2 h 25 min. 
%
%The training set consists of 37519 images, and 
The detector is trained for five epochs, just like the IRcam detector.

\subsubsection{Fcam worker}

Initially, the fish-eye lens camera was mounted upwards, but this caused the image distortion to be significant in the area just above the horizon where the interesting targets usually appear. After turning the camera so that it faces forward (as seen in Figure \ref{fig:system}b), the motion detector is less affected by image distortion.
%

The images from the camera are of 1280$\times$768 pixels, but the lower half (showing the area below the sky horizon) is cropped so that 1024$\times$384 pixels remain to be processed. 
%
Images are then analysed using the Matlab \texttt{ForegroundDetector} function \cite{gmmmatlab}, which compares an input video frame to a background model to determine whether individual pixels are part of the background or the foreground. 
%
The function uses a background subtraction algorithm based on Gaussian Mixture Models (GMM) \cite{Stauffer99}, producing a binary mask with pixels of foreground objects set to one.
%
The mask is next processed with the \texttt{imopen} function, which performs a morphological opening (erosion followed by dilation) to eliminate noise. The structural element is set to 3$\times$3, so that very small objects are deleted.
%
Then, the \texttt{BlobAnalysis} function is applied, which outputs the centroids and bounding boxes of all objects in the binary mask provided by the \texttt{ForegroundDetector} function. 
%
All centroids and bounding boxes are sent finally to a multi-object tracker based on Kalman filters \cite{kalmanmatlab}, created with the \texttt{configureKalmanFilter} function, which tracks objects across multiple frames. 
%
A Kalman filter is used to predict the centroid and bounding box of each track in a new frame based on their previous motion history. Then, the predicted tracks are assigned to the detections given by the foreground/background detector by minimizing a cost function that takes into account the Euclidean distance between predictions and detections.
%
%
The object with the longest track history is picked and retained as the best. 
%

In the Fcam presentation window, all tracks (both updated and predicted) are visualised, and the track considered to be the best is marked with red (Figure~\ref{fig:Fcam_red}).
%
The output from the Fcam worker is the FPS status, together with the elevation and azimuth angles of the best track, if such track exists at the moment. 
%
Of all the workers, the Fcam is the one with the most tuning parameters, as will be seen in Section~\ref{sect:results}. This involves choosing and tuning the image processing operations, foreground detector and blob analysis settings, and the multi-object Kalman filter tracker parameters.


\begin{figure} [htb]
\centering
\includegraphics[width=0.85\textwidth]{Fcam_red.png}
\caption{Fish-eye camera image with a detected track in red. Picture originally appearing in \cite{drone20thesis}.}
\label{fig:Fcam_red}
\end{figure}

\subsubsection{Audio worker}

The audio worker collects acoustic data in a one-second-long buffer (44100 samples), set to be updated 20 times per second. 
%
To classify the sound source in the buffer, it is first processed with the Matlab \texttt{mfcc} function, which extracts MFCC features. 
%
We employ the default parameters (a Hamming window with a length of 3\%  and an overlap of 2\% of the sampling frequency, and the number of coefficients per window equal to 13). 
%
Based on empirical trails, the parameter \texttt{LogEnergy} is set to \texttt{Ignore}, meaning that the log-energy is not calculated. Then the extracted features are sent to the classifier.
%
The extracted features are then sent to an LSTM classifier consisting of an input layer, two bidirectional LSTM layers with a dropout layer in-between, a fully connected layer, a softmax layer and a classification layer. The classifier builds on \cite{[LSTM]} but increasing the number of classes from two to three and an additional dropout layer between the bidirectional LSTM layers.

%The audio database has 30 ten-second clips of each of the output classes. Five clips from each class are set aside to be used for evaluation. The remaining audio clips are used for the training process. Out of each ten-second clip, the last second is used for validation and the rest for training. 
%
The classifier is trained from scratch for 120 epochs, after which it shows signs of over-fitting. 
%
%Since the audio worker is essentially only a classifier and not a detector and classifier (like YOLOv2), it 
%
It includes a class of general background sounds (Table~\ref{tab:output-classes}) recorded outdoors in the typical deployment environment of the system. Also, it has some clips of the sounds from the servos moving the pan/tilt platform. Like the other workers, the output is a class label with a confidence score.
%
A graphical presentation of the audio input and extracted features is also available, as shown in Figure~\ref{fig:audio-features}, but this is not included in the final GUI layout.


\begin{figure} [htb]
\centering
\includegraphics[width=0.85\textwidth]{audio-features.png}
\caption{Two examples of the audio worker plots with the audio  amplitudes and the extracted MFCC features. Pictures originally appearing in \cite{drone20thesis}.}
\label{fig:audio-features}
\end{figure}


\subsubsection{ADS-B worker}

We implement the ADS-B decoding in Matlab using functions from the Communications toolbox ~\cite{adsbmatlab}.
%
All vehicle categories that can be seen as subclasses to the airplane target label are clustered together. These are all the classes ``Light'', ``Medium'', ``Heavy'', ``HighVortex'', ``VeryHeavy'' and ``HighPerformanceHighSpeed''. The class ``Rotorcraft'' is translated into helicopter. Interestingly, there is also a ``UAV'' category label. This is also implemented in the ADS-B worker, translated into drone.

One might wonder if there are any such aircraft sending out that belong to the UAV vehicle category. Examples are found by looking at the Flightradar24 service (Figure~\ref{fig:flightradar}). Here we can find one such drone flying at Gothenburg City Airport, one of the locations used when collecting the dataset. The drone is operated by the company Everdrone AB, involved in the automated external defibrillators delivery trails of \cite{[Sanfridsson19]}.
%
Another example is the UK Coastguard/Border Force surveillance drone that is regularly flying at night over the straight of Dover since December of 2019. This is naturally a large drone with a wingspan of 7.3 m.

\begin{figure} [htb]
\centering
\includegraphics[width=0.95\textwidth]{flightradar.png}
\caption{Example of drones sending ADS-B information.
%Left: A drone flying at Gothenburg City Airport. Right: Surveillance drone over the straight of Dover. 
Images from www.flightradar24.com}
\label{fig:flightradar}
\end{figure}

As mentioned above, not all aircrafts will send out their vehicle category as part of the ADS-B squitter message. 
%
However, in our experience, about half of the aircrafts send out their category. This justifies its inclusion in this work, as one of our principles is to detect and keep track of other flying objects that are likely to be mistaken by a drone.
%
In the output message of the ADS-B worker, the confidence of the classification is set to 1 if the vehicle category message has been received. If not, the label is set to airplane (the most common aircraft type) with the confidence to 0.75 so that there is a possibility for any of the other sensors to influence the final classification.



\subsection{Graphical User Interface (GUI)}
\label{sect:GUI}

The Graphical User Interface of the system is shown in Figure~\ref{fig:system_GUI}.
%
The GUI is part of the main script but is described separately. 
%
It shows the output of the different workers, including video streams of the cameras, providing various options to control the system configuration.
%
The Matlab command window is made visible in the bottom centre so that messages (e.g. exceptions) can be monitored during the development and use of the system.


\begin{figure} [htb]
\centering
\includegraphics[width=0.95\textwidth]{system_GUI1_labelled.jpg}
\caption{Graphical User Interface of the system with labels of the different elements. 
%See Section~\ref{sect:GUI} for further details. 
Picture originally appearing in \cite{drone20thesis} and published in \cite{drone20icpr}. Reprinted with permission.}
\label{fig:system_GUI}
\end{figure}

The ADS-B presentation area (left) consists of a PPI-type (Plan Position Indicator) display and an altitude display. Besides presenting the ADS-B targets, the PPI display also shows system information. The solid green line is the main direction of the system relative to the north. The other green lines present the field of motion of the pan/tilt platform (dashed) and the field of view of the fish-eye lens camera (dotted).
%
The actual direction of the pan/tilt platform is presented with a solid red line, and the field of view of the thermal infrared and video cameras are represented using dashed red lines. If the fish-eye lens camera worker tracks any object, its direction is indicated by a solid cyan line.
%
ADS-B targets are presented using the class label colours of Table~\ref{tab:output-classes}, together with the track history plots. The altitude information is presented in a logarithmic plot to make the lower altitude portion more prominent.

The area directly below the ADS-B presentation is the control panel (seen in better detail in Figure~\ref{fig:gui-fig18-fig19}a).
%
Starting from the top left corner, we have radio buttons for the range settings of the ADS-B PPI and altitude presentations. Next is the number of ADS-B targets currently received, and below that, the set orientation angle relative to the north. The ``Close GUI'' button is used to shut down the main script and the workers.
%
The GPS-position presentation changes colour to green when the GPS receiver has received a correct position after pressing the ``Get GPS-pos'' button. Pressing the ``Set ori-angle'' button opens an input dialogue box so that the orientation angle of the system can be entered. Below that two buttons for switching the detectors between running and idle mode and a choice to display the raw Fcam image or the moving object mask only (not shown).
%
The servo settings can be controlled with the buttons of the middle column.
%
To aid the Fcam in finding targets, the pan/tilt can be set to move in two different search patterns. One where the search is done from side to side using a static elevation of 10$^{\circ}$, so that the area from the horizon up to 20$^{\circ}$ is covered, and another one where the search is done with two elevation angles to increase the coverage.
%
The pan/tilt platform can also be controlled by the elevation and azimuth angles from one of the workers. This is set by the ``assign'' buttons of the fourth column, placed in priority order from top to down. The ``IR\&V assign'' setting means that a target has to be detected by both the IRcam and Vcam workers, and if so, the pan/tilt platform is controlled by the angular values from the IRcam worker.
%
The rightmost column of the control panel shows the status information regarding the performance in FPS of the workers and the elevation and azimuth angles of the Fcam worker target (if such a target exists). The status displays are red if the worker is not connected, yellow if the detector is idle, and green if it is running.

The results panel (seen in better detail in Figure~\ref{fig:gui-fig18-fig19}b) features settings for the sensor fusion and presents the workers and system results.
%
The first column (servo controller) indicates the source of information currently controlling the servos of the pan/tilt platform.
%
In the bottom left corner, the angles of the servos are presented. 
%
The settings for the sensor fusion and the detection results presentation are found in the middle of the panel. 
%
The information in the right part of the panel is the current time and the position of the system. The system elevation and azimuth relative to the north are also presented here. Note the difference in azimuth angle compared to the bottom left corner where the system internal angle of the pan/tilt platform is presented.
%
The last part of the results panel (bottom right corner) presents offset angles for the ADS-B target if one is in the field of view of the thermal infrared and video cameras. These values are used to detect systematic errors in the orientation of the system. The sloping distance to the ADS-B target is also presented. 


\begin{figure} [htb]
\centering
\includegraphics[width=0.95\textwidth]{gui-fig18-fig19.jpg}
\caption{Detail of the control panel and results panel of the Graphical User Interface. In the control panel, note a bird detected and tracked by the IRcam worker. Pictures originally appearing in \cite{drone20thesis}.}
\label{fig:gui-fig18-fig19}
\end{figure}

\subsection{Database}
\label{sect:database}

A dataset has also been collected to accomplish the necessary training of the detectors and classifiers. The annotation has been done so that others can inspect, edit, and use the dataset. The fact that the datasets for the thermal infrared and the visible video sensors are collected under the same conditions and using the same targets ensures that a comparison between the two sensor types is well-founded.
%
%During this research, we have also captured a dataset to train and evaluate our system. 
%
The dataset is fully described in \cite{DiBdataset} and available at \cite{svanstrom_fredrik_2020_5500576}.
%
The videos and audios are recorded at locations in and around Halmstad Airport (IATA/ICAO code: HAD/ESMT), Gothenburg City Airport (GSE/ESGP) and Malmö Airport (MMX/ESMS).
%
Three different drones are used to collect and compose the dataset. These are of the following types: Hubsan H107D+, a small-sized first-person-view (FPV) drone; the high-performance DJI Phantom 4 Pro; and finally, the medium-sized kit drone DJI Flame Wheel. This can be built as a quadcopter (F450) or a hexacopter configuration (F550). The version used in this work is an F450 quadcopter. All three types can be seen in Figure~\ref{fig:db_drones}.
%
These drones differ a bit in size, with Hubsan H107D+ being the smallest, having a side length from motor to motor of 0.1 m. The Phantom 4 Pro and the DJI Flame Wheel F450 are larger with 0.3 and 0.4	m motor-to-motor side lengths, respectively.
%
To comply with regulations (drones must be flown within visual range), the dataset is recorded in daylight, even if the thermal infrared or acoustic sensors could be effective at night. 


\begin{figure} [htb]
\centering
\includegraphics[width=0.85\textwidth]{db_drones.png}
\caption{The three drones of our dataset. %Left: Hubsan H107D+. Middle: DJI Phantom 4 Pro. Right: DJI Flame Wheel F450. 
Pictures were originally appearing in \cite{drone20thesis} and published in \cite{drone20icpr} and \cite{DiBdataset}. Reprinted with permission.}
\label{fig:db_drones}
\end{figure}


The videos and audio files are cut into ten-second segments to facilitate annotation. 
%
The acquisition was performed during the drastic reduction in air traffic due to the COVID19 pandemic. Therefore, to get a more comprehensive dataset, both in terms of aircraft types and sensor-to-target distances, the data is complemented with non-copyrighted material of \cite{[VIRTUALAIRFIELD]}. %, in particular 11 plus 38 video clips in the airplane and helicopter categories, respectively.
%
Overall, the dataset contains 90 audio clips and 650 videos (365 IR and 285 visible, of ten seconds each), with a total of 203328 annotated images (publicly available).
%
The thermal infrared videos have a resolution of 320$\times$256 pixels, and the visible videos have 640$\times$512.
%
%
The greatest sensor-to-target distance for a drone in the dataset is 200 m.
%
Audio files consist of 30 ten-second clips of each of the three output audio classes (Table~\ref{tab:output-classes}), and the amount of videos among the four output video classes is shown in Table~\ref{tab:db-stats}.
%
The background sound class contains general background sounds recorded outdoor in the typical deployment environment of the system, and also includes some clips of the sounds from the servos moving the pan/tilt platform.


\begin{table}[htb]
\caption{Distribution of the thermal infrared and visible videos. Each video has approximately 10 seconds. Reprinted with permission from \cite{drone20icpr}. \label{tab:db-stats}}

%\small
%\footnotesize
%\scriptsize
\begin{center}
\begin{tabular}{cccccccccc}

%\multicolumn{5}{c}{} \\

\multicolumn{1}{c}{} & \multicolumn{4}{c}{thermal infrared (365)}  & & \multicolumn{4}{c}{visible (285)} \\  \cline{2-5} \cline{7-10}

bin & airpl. & bird & drone & helicop.  & &  airpl. & bird & drone & helicop. \\ \hline
Close & 9 & 10 & 24  & 15  & &  17 & 10 & 21  & 27 \\
Medium & 25 & 23 & 94 & 20  & & 17 & 21 & 68 & 24 \\
Distant & 40 & 46 & 39 & 20  & & 25 & 20 & 25 & 10 \\ \hline

\end{tabular}

\end{center}

\end{table}
\normalsize

The video part of the database is divided into three category bins: Close, Medium and Distant.
%
This is because one of our aims is to explore performance as a function of the sensor-to-target distance.
%
The borders between these bins are chosen to follow the industry-standard Detect, Recognize and Identify (DRI) requirements \cite{[DRI]}, building on the Johnson criteria \cite{[Chevalier16]}, as shown in Figure~\ref{fig:DRI}.
%
The Close bin is from 0 m to a distance where the target is 15 pixels wide (the requirement for `identification', e.g. explicitly telling the specific drone model, aircraft, helicopter, bird, etc.). The Medium bin is from where the target is from 15 down to 5 pixels (`recognition' of the target, e.g. a drone, an aircraft, a helicopter... albeit without the possibility of identifying the model), and the Distant bin is beyond that (`detection', e.g. there is something).
%


\begin{figure} [htb]
\centering
\includegraphics[width=0.7\textwidth]{DRI-fig24.jpg}
\caption{The DRI requirements (from \cite{[DRI]}).}
\label{fig:DRI}
\end{figure}








\begin{table}[htb]
\caption{Results with the thermal infrared sensor (confidence threshold and IoU requirement of 0.5). The average of the three F1-scores is 0.7601. Reprinted with permission from \cite{drone20icpr}. \label{tab:results-IRcam}}

%\small
%\footnotesize
%\scriptsize
\begin{center}
\begin{tabular}{cccccc}

%\multicolumn{6}{c}{} \\

\multicolumn{1}{c}{} & \multicolumn{5}{c}{distance bin: \textbf{CLOSE}}   \\ \cline{2-6}

 & airplane & bird & drone  & helicopter  &  average \\ \hline

\textbf{Precision} & 0.9197 & 0.7591 & 0.9159 & 0.9993& 0.8985 \\
\textbf{Recall} & 0.87367 & 0.85087 & 0.87907 & 0.87927 & 0.8706   \\ \hline
\textbf{F1-score} & & & & & 0.88447   \\
\hline


\multicolumn{6}{c}{} \\

\multicolumn{1}{c}{} & \multicolumn{5}{c}{distance bin: \textbf{MEDIUM}}   \\ \cline{2-6}

 & airplane & bird & drone  & helicopter  &  average \\ \hline

\textbf{Precision} &  0.82817 & 0.50637 & 0.89517 & 0.95547 & 0.7962  \\
\textbf{Recall} &  0.70397 & 0.70337 & 0.80347 & 0.83557 & 0.7615 \\ \hline
\textbf{F1-score} & & & & & 0.77857   \\ \hline


\multicolumn{6}{c}{} \\

\multicolumn{1}{c}{} & \multicolumn{5}{c}{distance bin: \textbf{DISTANT}}   \\ \cline{2-6}

 & airplane & bird & drone  & helicopter  &  average \\ \hline

\textbf{Precision} &  0.78227 & 0.61617 & 0.82787 & 0.79827 & 0.7561  \\
\textbf{Recall} &  0.40437 & 0.74317 & 0.48367 & 0.45647 & 0.5218 \\ \hline
\textbf{F1-score} & & & & & 0.61757   \\ \hline

\end{tabular}

\end{center}

\end{table}
\normalsize




\begin{table}[htb]
\caption{Results with the thermal infrared sensor (confidence threshold 0.8, IoU requirement 0.5). Table originally appearing in \cite{drone20thesis}. \label{tab:results-IRcam-th0.8}}

%\small
%\footnotesize
%\scriptsize
\begin{center}
\begin{tabular}{ccccc}

\multicolumn{1}{c}{} & \multicolumn{3}{c}{distance bin} & \multicolumn{1}{c}{Average}  \\ \cline{2-4}

 & close & medium & distant &  \\ \hline

\textbf{Precision} & 0.9985 & 0.9981 & 1.0000 & 0.9987\\
\textbf{Recall} & 0.2233 & 0.1120 & 0.0019 & 0.1124 \\ \hline

\end{tabular}

\end{center}

\end{table}
\normalsize









%%%%%%%%%%%%%%%%%%%%%%%%%%%%%%%%%%%%%%%%%%
\section{Results}
\label{sect:results}

We provide evaluation results of the individual sensors in terms of precision, recall and F1-score.
%
Precision is the fraction of the total number of true positive detections. In other words, it measures how many of the detected objects are relevant.
%
Recall is the fraction of the total number of labelled samples in the positive class that are true positive. In other words, how many of the relevant objects are detected.
%
With the video detectors, these metrics can be obtained using the Matlab function \texttt{bboxPrecisionRecall}, which measures the accuracy of bounding box overlap between detected and ground truth boxes.
%
Since we also have confidence scores, the \texttt{evaluateDetectionPrecision} function can be used to plot precision curves as a function of the recall value.
%
%
After that, we perform sensor fusion experiments, a direction barely followed in the related literature \cite{Diamantidou19,Shi18,Unlu19}.
%
To carry out experiments, two disjoint sets of videos were created for training and testing (i.e. a video selected for the training set is not used in the testing set). 
%
The training set comprises 120 infrared and 120 visible clips (10 for each class and target bin per spectrum), resulting in 37428 infrared and 37519 visible frames.
%
The evaluation set comprises 60 infrared and 60 visible clips (5 for each class and target bin per spectrum), resulting in 18691 infrared and 18773 visible frames.
%
Since the duration of each clip is roughly the same ($\sim$10 seconds), the amount of samples per class and target bin is approximately balanced.
%
%a part of the video dataset is set aside for evaluation (120 infrared and 120 visible clips, 5 for each class and target bin per spectrum). 
%
%Since the videos do not have all exactly the same duration (10.4 seconds on average), the total number of images per class varies slightly, as shown in Table XX.
%
%Among the remaining videos not set aside for evaluation, 240 are selected as evenly distributed as possible to create the training set.
%
For the audio classifier, five 10-second clips from each output category are selected for evaluation, and the remaining clips for training. 
%
Since the audio classifier processes a one-second input buffer, the clips cut into that length, with an overlap of 0.5 seconds, resulting in 20 samples per clip.
%
This results in 297 clips in the evaluation set, 99 from each class.

In evaluating the YOLOv2 detector, it provides an array of class labels, detection confidence, and bounding boxes of detected objects.
%
Here, not only the classification label but also the placement of the bounding box must be taken under consideration.
%
When it produces multiple bounding boxes for the same object, only the strongest one is used, chosen as the box with the highest IoU (Intersection Over Union) with the annotations in the training data.
%
To assign a box to a class, it must also have a minimum IoU with the training object it is supposed to detect.
%
In related works, an IoU of 0.5 is usually employed \cite{Park17,Saqib17,Aker17}, although a lower IoU of 0.2 is used in \cite{Schumann17}.
%
In this work, we will stick to 0.5.
%
A threshold to the detection confidence can also be imposed, so bounding boxes with small confidence can be rejected, even if their IoU is above 0.5.




\begin{figure} [htb]
\centering
\includegraphics[width=0.6\textwidth]{F1score-IRcam.png}
\caption{F1-score with the thermal infrared sensor as a function of detection threshold, using all the 18691 images in the evaluation dataset. Picture originally appearing in \cite{drone20thesis}.}
\label{fig:results-IRcam-F1}
\end{figure}


\begin{figure} [htb]
\centering
\includegraphics[width=0.85\textwidth]{precision-recall-IRcam.png}
\caption{Precision and recall curves with the thermal infrared sensor. The achieved values with a detection threshold of 0.5 are marked by stars. Pictures originally appearing in \cite{drone20thesis}.}
\label{fig:results-IRcam-PR}
\end{figure}


\subsection{Thermal Infrared Camera (IRcam)}

The precision, recall and F1-score of the IRcam worker detector with a confidence threshold set to 0.5 and an IoU requirement of 0.5 are shown in Table~\ref{tab:results-IRcam}.
%
We can observe that the precision and recall values are well balanced.
%
Altering the setting to a higher decision threshold of 0.8 leads to higher precision at the cost of a lower recall value, as shown in Table~\ref{tab:results-IRcam-th0.8}. The drop in recall with increasing sensor-to-target distance is also prominent.




To further explore the detection threshold setting, we run the evaluation with values from 0.1 up to 1.0 in steps of 0.1. The results in the form of an F1-score are shown in Figure~\ref{fig:results-IRcam-F1}.
%
Using not only the bounding boxes and class labels but also the confidence scores, the detector can be evaluated in the form of precision vs recall curves as well (Figure~\ref{fig:results-IRcam-PR}).
%
Note that the average precision results output from the Matlab \texttt{evaluateDetectionPrecision} function is defined as the area under the precision vs recall curve. Hence it is not the same as the actual average precision values of Table~\ref{tab:results-IRcam}.
%
To distinguish that we mean the area under the curve, we denote this as AP in Figure~\ref{fig:results-IRcam-PR}, just as in the original YOLO paper \cite{Redmon17}.
%
This is also  the definition used in \cite{Saqib17}, and to further adopt the notation of these papers, we denote the mean AP taken over all classes as the mAP, which is given in Table~\ref{tab:results-IRcam-mAP}.
%



\begin{table}[htb]
\caption{Mean values over all classes of the area under the precision vs recall curve (mAP) of the thermal infrared sensor. Table originally appearing in \cite{drone20thesis}. \label{tab:results-IRcam-mAP}}

%\small
%\footnotesize
%\scriptsize
\begin{center}
\begin{tabular}{ccccc}

\multicolumn{1}{c}{} & \multicolumn{3}{c}{distance bin} & \multicolumn{1}{c}{Average}  \\ \cline{2-4}

 & close & medium & distant &  \\ \hline

\textbf{mAP} & 0.8704 & 0.7150 & 0.5086 & 0.7097 \\ \hline

\end{tabular}

\end{center}

\end{table}
\normalsize


The choice of the detection threshold will affect the achieved precision and recall values. The stars in Figure~\ref{fig:results-IRcam-PR} indicate the results with a detection threshold of 0.5 (as reported in Table~\ref{tab:results-IRcam}). We can conclude that such a threshold results in a balanced precision-recall combination near the top right edge of the respective curves. 
%
Compare this balanced behaviour to the  precision and recall values obtained with a detection threshold of 0.8 (Table~\ref{tab:results-IRcam-th0.8}).
%
From observations of the behaviour when running the drone detection system, we can also point out that a common source of false alarms of the thermal infrared sensor is small clouds and edges of large clouds lit up by the sun. An example of this can be seen in Figure~\ref{fig:results-IRcam-false-target}.

\begin{figure} [htb]
\centering
\includegraphics[width=0.85\textwidth]{IRcam-false-fig29.jpg}
\caption{A false target of the thermal infrared sensor caused by a small cloud lit by the sun. Pictures originally appearing in \cite{drone20thesis}.}
\label{fig:results-IRcam-false-target}
\end{figure}

In comparing our work with other research using thermal infrared sensors, results can be found in \cite{ANDRASI2017183}. 
%
However, the sensor used in that article to detect drones (up to a distance of 100 m) has a resolution of just 80$\times$60 pixels, and the system does not involve any form of machine learning feature (the task is done by a human looking at the output video stream). 
%
%In that paper, three different drones were used, and a DJI Phantom 4 was detected by the setup up to a distance of 51 m.
%
In \cite{wang_chen_choi_kuo_2019}, curves for the precision and recall of a machine learning-based thermal detector are presented. It is stated that the video clips used for training and evaluation have a frame resolution of 1920$\times$1080. Unfortunately, the paper fails to mention if this is also the resolution of the sensor. Neither is the input size of the detection network specified in detail, other than the images are re-scaled so that the shorter side has 600 pixels.
%
The most substantial results to relate to in \cite{wang_chen_choi_kuo_2019} is that since the system also contains a video camera with the same image size as the thermal one, the authors can conclude that the thermal drone detector performs over 8\% better than the video detector.


\subsection{Video Camera (Vcam)}

The results of the previous sub-section will be compared to the results of the video detector in what follows.
%
To enable comparison, the same methods and settings as the infrared camera are used here, leading to the precision, recall and F1-score results of Table~\ref{tab:results-Vcam} with a confidence threshold and IoU requirement of 0.5.
%
These results differ no more than 3\% from the results of the thermal infrared detector. Recall that the input layers of the YOLOv2 detectors are different, and hence, the resolution of the visible sensor is 1.625 higher than the thermal infrared detector (416$\times$416 vs 256$\times$256).
%
So even with a lower resolution and using images in grey-scale and not in colour, the thermal infrared sensor performs as well as the visible one. 
%
This is in line with \cite{ANDRASI2017183}, where the infrared detector outperforms the visible one when the image size is the same, although in our case, this happens even with a smaller image size.



\begin{table}[htb]

\caption{Results with the visible camera (confidence threshold and IoU requirement of 0.5). The average of the three F1-scores is 0.7849. Reprinted with permission from \cite{drone20icpr}. \label{tab:results-Vcam}}
%\small
%\footnotesize
%\scriptsize
\begin{center}
\begin{tabular}{cccccc}

%\multicolumn{6}{c}{} \\

\multicolumn{1}{c}{} & \multicolumn{5}{c}{distance bin: CLOSE}   \\ \cline{2-6}

 & airplane & bird & drone  & helicopter  &  average \\ \hline

\textbf{Precision} & 0.8989 & 0.8284 & 0.8283 & 0.9225 & 0.8695 \\
\textbf{Recall} & 0.7355 & 0.7949 & 0.9536 & 0.9832 & 0.8668 \\ \hline
\textbf{F1-score} & & & & & 0.8682   \\
\hline


\multicolumn{6}{c}{} \\

\multicolumn{1}{c}{} & \multicolumn{5}{c}{distance bin: MEDIUM}   \\ \cline{2-6}

 & airplane & bird & drone  & helicopter  &  average \\ \hline

\textbf{Precision} & 0.8391 & 0.7186 & 0.7710 & 0.9680 & 0.8242 \\
\textbf{Recall} & 0.7306 & 0.7830 & 0.7987 & 0.7526 & 0.7662 \\ \hline
\textbf{F1-score} & & & & & 0.7942   \\ \hline


\multicolumn{6}{c}{} \\

\multicolumn{1}{c}{} & \multicolumn{5}{c}{distance bin: DISTANT}   \\ \cline{2-6}

 & airplane & bird & drone  & helicopter  &  average \\ \hline

\textbf{Precision} & 0.7726 & 0.6479 & 0.8378 & 0.6631 & 0.7303 \\
\textbf{Recall} & 0.7785 & 0.7841 & 0.5519 & 0.5171 & 0.6579 \\ \hline
\textbf{F1-score} & & & & & 0.6922   \\ \hline

\end{tabular}

\end{center}

\end{table}
\normalsize


Just as for the thermal infrared camera, we can also explore the effects of the detection threshold setting. This can be seen in Figure~\ref{fig:results-Vcam-F1}.
%
The precision vs recall curves of the video camera images for the different target classes and distance bins are shown in Figure~\ref{fig:results-Vcam-PR}.
%
In a similar vein, a detection threshold of 0.5 results in a balanced precision-recall combination near the top right edge of the respective curves.
%
Notably, when inspecting the precision-recall curves, the video camera detector performs outstandingly when it comes to distant airplanes. This has its explanation in that such targets often presents a very large signature consisting not only of the airplane itself but also contrails behind it.
%
Also, calculating the mAP from these results, we obtain Table~\ref{tab:results-Vcam-mAP}, with an average of 0.7261.
%
Once again, it is not far from the 0.7097 mAP of the thermal infrared detector of the previous sub-section. 



\begin{figure} [htb]
\centering
\includegraphics[width=0.6\textwidth]{F1score-Vcam.png}
\caption{F1-score with the visible camera as a function of detection threshold, using all the 18773 images in the evaluation dataset. Picture originally appearing in \cite{drone20thesis}.}
\label{fig:results-Vcam-F1}
\end{figure}


\begin{figure} [htb]
\centering
\includegraphics[width=0.85\textwidth]{precision-recall-Vcam.png}
\caption{Precision and recall curves with the visible camera. The achieved values with a detection threshold of 0.5 are marked by stars. Pictures originally appearing in \cite{drone20thesis}.}
\label{fig:results-Vcam-PR}
\end{figure}


\begin{table}[htb]
\caption{Mean values over all classes of the area under the precision vs recall curve (mAP) of the visible camera. Table originally appearing in \cite{drone20thesis}. \label{tab:results-Vcam-mAP}}

%\small
%\footnotesize
%\scriptsize
\begin{center}
\begin{tabular}{ccccc}

\multicolumn{1}{c}{} & \multicolumn{3}{c}{distance bin} & \multicolumn{1}{c}{Average}  \\ \cline{2-4}

 & close & medium & distant &  \\ \hline

\textbf{mAP} & 0.8474 & 0.7477 & 0.5883 & 0.7261 \\ \hline

\end{tabular}

\end{center}

\end{table}
\normalsize




The most frequent problem with the video detector is the auto-focus feature of the video camera. For this type of sensor, clear skies are not the ideal weather, but rather scenery with objects that can help the camera to set the focus correctly. However, note that this fact is not heavily affecting the evaluation results of the visible camera detector performance, as presented above, since only videos where the objects are seen clearly, and hence are possible to annotate, are used.
%
Figure~\ref{fig:results-Vcam-autofocus-fail} shows an example of when the focus is set wrongly so that only the thermal infrared worker detects the drone.
%
This issue justifies the multi-sensor approach followed here. Also, cameras not affected by this issue (such as the thermal infrared or the fish-eye) could be used to aid the focus of the video camera. % towards the area where the object is appearing.


%The result is also ahead to what is presented in \cite{Saqib17}, where a mAP of 0.66 achieved, albeit using a detector with drones as the only output class and giving no information about the sensor-to-target distances.
%
Comparing the results to other papers, our results are ahead of what is presented in \cite{Saqib17}, where a mAP of 0.66 was achieved, albeit using a detector with drones as the only output class and giving no information about the sensor-to-target distances.
%
Also, we see that the YOLOv2 detector in \cite{Park17} achieves an F1-score of 0.728 with the same detection threshold and IoU-requirement. This F1-score is just below the results of the thermal infrared and video camera workers.
%
However, one notable difference lies in that the detector in \cite{Park17} has only one output class. This fact could confirm the doctrine of this work, i.e. that the detectors should also be trained to recognize objects easily confused for being drones. Unfortunately, there is no notation of the sensor-to-target distance other than that ``75\% of the drones have widths smaller than 100 pixels''. Since the authors implement an original YOLOv2 model from darknet, it is assumed that the input size of the detector is 416$\times$416 pixels.
%
A YOLOv2 architecture with an input size of 480$\times$480 pixels is implemented in \cite{Aker17}. The detector has two output classes, birds and drones. Based on the presented precision-recall curve, a precision and recall of 0.9 can be achieved simultaneously.
%
%
To summarize this comparison, we provide the performance of the thermal infrared and video camera detectors together with the reported comparable results in Table~\ref{tab:results-IRcam-Vcam-SOA}. The table also shows the output classes used.




\begin{table}[htb]
\caption{Results from the related works and the thermal infrared and video camera detectors. A = Airplane, B = Bird, D = Drone and H = Helicopter. Table originally appearing in \cite{drone20thesis}. \label{tab:results-IRcam-Vcam-SOA}}

%\small
%\footnotesize
%\scriptsize
\begin{center}
\begin{tabular}{ccccccc}

\multicolumn{1}{c}{Work} & \multicolumn{1}{c}{Sensor} & \multicolumn{4}{c}{Results} & \multicolumn{1}{c}{Classes}   \\ \cline{3-6}

 & & Precision & Recall & F1-score & mAP &  \\ \hline

%Wang et al. \cite{wang_chen_choi_kuo_2019} & Thermal &  &  &  &  &   \\ 

%Wang et al. \cite{wang_chen_choi_kuo_2019} & Visible &  &  &  &  &   \\ 
 
Park et al. \cite{Park17} & Visible &  &  & 0.73 & & D  \\ 

Liu  et al. \cite{Liu_2018} & Visible & 0.99 & 0.80  &  & & A, D, H   \\ 

Saqib et al. \cite{Saqib17} & Visible &  &  &  & 0.66 & D   \\ 

Aker and Kalkan. \cite{Aker17} & Visible & $\sim$0.9 & $\sim$0.9 &  &  & B, D  \\ 

\textbf{thermal infrared} & Thermal & 0.82 & 0.72 & 0.76 & 0.71 & A, B, D, H  \\ 

\textbf{video camera} & Visible & 0.81 & 0.76 & 0.78 & 0.73 & A, B, D, H   \\  

\hline

\end{tabular}

\end{center}

\end{table}
\normalsize



\begin{figure} [htb]
\centering
\includegraphics[width=0.85\textwidth]{Vcam-autofocus-fail-fig33.jpg}
\caption{A drone is detected only by the thermal infrared sensor since the auto-focus of the video camera is wrongly set. Pictures originally appearing in \cite{drone20thesis}.}
\label{fig:results-Vcam-autofocus-fail}
\end{figure}


%\begin{figure} [htb]
%\centering
%\includegraphics[width=0.7\textwidth]{Vcam-standalone-fig34.jpg}
%\caption{Airplanes, a helicopter and birds detected when running the visible camera detector as a stand-alone 5application on an aircraft video from \cite{[VIRTUALAIRFIELDhttps://www.overleaf.com/project]}. Note the FPS-performance.}
%\label{fig:results-Vcam-standalone}
%\end{figure}


\subsection{Fish-Eye Camera Motion Detector (Fcam)}

The fish-eye lens camera is included to be able to cover a larger airspace volume than covered by just the field of view of the thermal infrared and video camera. However, the drone detection system does not totally rely on the fish-eye worker to detect objects of interest since the system also includes the search programs moving the pan/tilt platform when no other targets are detected. The search programs can easily be turned on or off using the control panel of the GUI.

It was initially observed that the foreground/background-based motion detector of the fish-eye worker was sensitive to objects moving with the wind, as shown in Figure~\ref{fig:results-Fcam-false-targets}.
%
The false targets have been mitigated by extensive tuning of the image processing operations, the GMM foreground detector \cite{gmmmatlab}, the blob analysis settings and the parameters of the multi-object Kalman filter tracker from \cite{kalmanmatlab} as follows. All other parameters or computation methods not mentioned here are left at the default values described in \cite{gmmmatlab},\cite{kalmanmatlab}.

\begin{itemize}


\item Some of the most critical points in this were first to remove the \texttt{imclose}\footnote{Morphological closing.} and \texttt{imfill}\footnote{Flood-fill of holes, i.e. background pixels that cannot be reached by filling in the background from the edge of the image.} functions that were initially implemented after the \texttt{imopen}\footnote{Morphological opening.} function in the image processing operations.


\item To only detect small targets, in the blob analysis settings of the \texttt{BlobAnalysis} function, a maximum blob area of 1000 pixels was implemented.

%In the blob analysis settings, the minimum blob area was reduced, so that smaller objects are detectable, and a maximum blob area of 1000 pixels was also implemented.

%To only detect small targets the parameter 'MaximumBlobArea' of the function 'BlobAnalysis' was set to 1000.


\item The parameters of the \texttt{ForegroundDetector} function were changed so that the model would adapt faster to changing conditions, and hence react quicker. 
%
The number of initial video frames for training the background model was reduced from 150 (default) to 10, and the learning rate for parameter updates increased from 0.005 (default) to 0.05. 
%
Furthermore, to reduce false alarms, the threshold to determine the background model (i.e. minimum possibility for pixels to be considered background values) was increased from 0.7 (default) to 0.85.
%
The number of Gaussian modes in the mixture model is left to the default value of 5.

%So the parameters 'NumTrainingFrames' and 'LearningRate'  were set to 10 and 0.05 respectively (default is 150 and 0.005). To reduce false alarms the parameter 'MinimumBackgroundRatio' was set to 0.85 instead of the default value of 0.7. 
%
%Furthermore, the minimum background ratio setting of the foreground detector was increased together with the learning rate, so that the GMM more quickly adapts to changing conditions.


\item Tuning the parameters of the Kalman filter multi-object tracker was also important. These were altered to make the tracker slower to start new tracks and quicker to terminate them if no moving objects were present at the predicted positions.
%
The motion model in the \texttt{configureKalmanFilter} function was chosen to be of the type ``constant velocity''.
%
The initial location of unassigned objects is given by the centroid computed by the foreground detector.
%
The initial estimate error was changed from [200, 50] (default) to [200, 200]. This specifies the variance of the initial estimates of location and velocity of the tracked object, with the initial state estimation error covariance matrix built as 2$\times$2 diagonal matrix with these two values on the main diagonal.
%
Larger values here help the filter to adapt to the detection results faster. %
This affect the first few detections, after which the estimate error is obtained from the noise and input data. The function assumes a zero initial velocity at the position of the initial location.
%
The motion noise was changed from [100, 25] (default) to [50, 50], with the process noise covariance matrix built as a 2$\times$2 diagonal matrix with these two values on the main diagonal. These values specify the tolerance (variance) for the deviation on location and velocity, compensating for the difference between the actual motion and that of the constant velocity model.
%
Here the changes were made since we must assume that the initial estimate error and the motion noise are uniform in x- and y-dimensions. 
%
The measurement noise covariance was increased from 10 (default) to 100, specifying the variance inaccuracy of the detected location.
%
Increasing this value enables the Kalman filter to remove more noise from the detections. This parameter was changed by trial and error.
%
Both the values of the motion noise and measurement noise stay constant.


%Regarding the parameters of the Kalman-filter: The motion model was chosen to be of the type 'ConstantVelocity' and the parameters 'InitialEstimateError, MotionNoise, MeasurementNoise' were set to [200, 200], [50, 50], 10. In the example described in the matlab help-page this was [200, 50], [100, 25], 100. Here the changes were made since we must assume that the initial estimate error and the motion noise are uniform in x- and x'-dimensions. The measurement noise parameter was changed by trial and error.

% Create a Kalman filter object.
%kalmanFilter = configureKalmanFilter('ConstantVelocity', ...
 %    centroid,       [200, 50],           [100, 25],  100);
 %    InitialLocation,InitialEstimateError,MotionNoise,MeasurementNoise
 %                    [200, 200],          [50, 50],   10

\end{itemize}


%Some of the most critical points in this was first to remove the \texttt{imclose}\footnote{Morphological closing.} and \texttt{imfill}\footnote{Flood-fill of holes, i.e. background pixels that cannot be reached by filling in the background from the edge of the image.} functions that were initially implemented after the \texttt{imopen}\footnote{Morphological opening.} function in the image processing operations.
%
%Furthermore, the minimum background ratio setting of the foreground detector was increased together with the learning rate, so that the GMM more quickly adapts to changing conditions. In the blob analysis settings, the minimum blob area was reduced, so that smaller objects are detectable and a maximum blob area was also implemented.
%
%Tuning the parameters of the Kalman filter multi-object tracker was also important. These were altered to make the tracker slower to start new tracks and quicker to terminate them if no moving objects were present at the predicted positions.


%
%\textbf{We employ the default parameters of the function (Gaussians in the mixture model=5, 
%threshold to determine background model=0.7, 
%initial video frames for training the background model=150, 
%learning rate for parameter updates across time=0.005).
%}


\begin{figure} [htb]
\centering
\includegraphics[width=0.85\textwidth]{Fcam-false-targets-fig35.jpg}
\caption{False targets in a fish-eye camera image. Picture originally appearing in \cite{drone20thesis}.}
\label{fig:results-Fcam-false-targets}
\end{figure}


With an image size of 1024$\times$384, the fish-eye worker moving object detector and tracker has been found, during the evaluation sessions, to be an effective way to assign objects to the pan/tilt platform up to a distance of 50 m against drones. 
%
Beyond this distance, the drone appears 
%falls into the ``distant'' category \cite{DiBdataset}, being theoretically 
so small (in pixels) that it is deleted by the \texttt{imopen} function.
%
The maximum resolution of the fish-eye camera employed is 3264$\times$2448, so a greater detection range should theoretically be achievable. However, using a higher resolution is also more demanding in computational resources, leading to a reduction in the FPS performance for the fish-eye worker and the other workers. Since the fish-eye lens camera is also complemented by the search program, where the pan/tilt platform can be controlled by the output of the other workers, the choice has been made to settle with this resolution and the limited detection range that follows from this.
%
Figure~\ref{fig:fig36-37}a shows a drone tracked by the fish-eye worker. At this moment, the pan/tilt platform is controlled by the output of the fish-eye worker. Just a moment later, as seen in Figure~\ref{fig:fig36-37}b, the drone is detected by the thermal infrared and video camera workers, and the thermal infrared worker output therefore controls the pan/tilt platform.

\begin{figure} [htb]
\centering
\includegraphics[width=0.95\textwidth]{fig36-37.png}
\caption{Tracking of a drone with different workers. Pictures originally appearing in \cite{drone20thesis}.
%Left: a drone tracked by the fish-eye worker. Right: the drone tracked by the thermal infrared worker just a second later (note the timestamp).
}
\label{fig:fig36-37}
\end{figure}

\subsection{Microphone}

The precision and recall results of the different classes with this detector are shown in Table~\ref{tab:results-audio}. 
%
Against a drone, the practical range of the acoustic sensor is 35-45 m, depending on how the drone is flying. This is in parity with the 50 m of \cite{Park15}, but far from the 160 m against a F450 drone reported in
\cite{Busset15} with its much more complex microphone configuration (a 120 elements microphone array).


\begin{table}[htb]
\caption{Results with the audio detector. Reprinted with permission from \cite{drone20icpr}. \label{tab:results-audio}}

%\small
%\footnotesize
%\scriptsize
\begin{center}
\begin{tabular}{ccccc}

%\multicolumn{6}{c}{} \\

\multicolumn{5}{c}{}    \\ \cline{2-5}

 & drone & helicopter & background  & average \\ \hline

\textbf{Precision} & 0.9694 & 0.8482 & 0.9885 & 0.9354 \\
\textbf{Recall} & 0.9596 & 0.9596 & 0.8687 & 0.9293 \\ \hline
\textbf{F1-score} & & & & 0.9323   \\
\hline

\end{tabular}

\end{center}

\end{table}
\normalsize



Our average F1-score is 0.9323, which is higher compared to \cite{Jeon17}, which also uses MFCC features. Out of the three classifier types tested in \cite{Jeon17}, the one comprising a LSTM-RNN performs the best, with a F1-score of 0.6984. The classification problem in that paper is binary (drone or background).
%
Another paper applying MFCC-features and Support Vector Machines (SVM) as classifier is \cite{Bernardini17}, with a precision of 0.983. Five output classes are used (drone, nature daytime, crowd, train passing and street with traffic), and the classification is based on a one-against-one strategy. Hence ten binary SVM classifiers are implemented. 
%
The final output label is then computed using the max-wins voting principle. 
%
The paper \cite{Kim17}, on the other hand, applies the Fast Fourier Transform (FFT) to extract features from the audio signals.
%
The results of the audio detector, together with results reported in other studies, are summarized in Table~\ref{tab:results-audio-SOA}.


\begin{table}[htb]
\caption{Results from the related works and the audio detector. D = Drone, H = Helicopter, BG = Background. The work \cite{Bernardini17} defines four background classes: nature daytime, crowd, train passing and street with traffic. Table originally appearing in \cite{drone20thesis}. \label{tab:results-audio-SOA}}

%\small
%\footnotesize
%\scriptsize
\begin{center}
\begin{tabular}{ccccc}

\multicolumn{1}{c}{Work} & \multicolumn{3}{c}{Results} & \multicolumn{1}{c}{Classes}   \\ \cline{2-4}

 & Precision & Recall & F1-score &   \\ \hline

Kim et al. \cite{Kim17} & 0.88  & 0.83 & 0.85  &  D, BG  \\ 

Jeon et al. \cite{Jeon17} & 0.55 & 0.96 & 0.70  &  D, BG  \\ 

Bernardi et al. \cite{Bernardini17} & 0.98  &  &   &  D, BG  \\ 

\textbf{audio worker} & 0.94 & 0.93 & 0.93 &  D, H, BG  \\ 

\hline

\end{tabular}

\end{center}

\end{table}
\normalsize



\subsection{Radar Module}

From the datasheet of the K-MD2 \cite{RFbeam}, we have that it can detect a person with a Radar Cross Section (RCS) of 1 m$^2$ up to a distance of 100 m. Since we have from \cite{Patel18} that the RCS of the F450 drone is 0.02 m$^2$, it is straightforward to calculate that, theoretically, the F450 should be possible to detect up to a distance of $100 \times \sqrt[4]{{0.02/1}} = 37.6$ m. 
%
Furthermore, given that the micro-doppler echoes from the rotors are 20 dB below that of the drone body, these should be detectable up to a distance of $100 \times \sqrt[4]{{0.02/1 \times 100}} = 11.9$ m.
%

We have observed that the F450 drone is detected and tracked by the K-MD2 up to a maximum distance of 24 m. This is, however, the maximum recorded distance, and it is observed that the drone is generally detected up to a distance of 18 m.
%
Due to the short practical detection range of the radar module, it is not included in the drone detection system of this work.
%
The micro-doppler signature can also be detected at short distances, as shown in Figure~\ref{fig:fig40-41}a. Corresponding range-doppler plots showing the micro-doppler of flying drones can be found in \cite{Drozdowicz16} and \cite{Rahman18} as well.
%
Compared to the echo of a person walking in front of the radar, as we can see in Figure~\ref{fig:fig40-41}b, no such micro-doppler signature is present in that case.


\begin{figure} [htb]
\centering
\includegraphics[width=0.6\textwidth]{fig40-41.jpg}
\caption{Micro-doppler signature of different elements in front of the radar module. Pictures originally appearing in \cite{drone20thesis}.
%Top: micro-doppler signature of the F450 drone. Bottom: the echo of a person walking in front of the radar module.
}
\label{fig:fig40-41}
\end{figure}


\subsection{Sensor Fusion}

We investigated the early or late sensor fusion choice based on whether the sensor data is fused before or after the detection element.
%
By early sensor fusion, we mean to fuse the images from the thermal infrared and video cameras before classification, either the raw images (sensor level fusion), or features extracted from them (feature level). 
%
%before processing them in a detector and classifier. 
%
Late sensor fusion will, in this case, be to combine the output decision from the separate detectors running on each camera stream, weighted by the confidence score of each sensor (decision level fusion). Other existing late fusion approaches entail combining the confidence scores (score level fusion) or the ranked class output of each system (rank level) \cite{FIERREZ201857}.
%

When investigating early sensor fusion at the raw image level in this research, the pixel-to-pixel matching of the images was the biggest issue. Even if this is possible in a static scenario, it turned out to be an unfeasible solution against moving objects with the available equipment due to the small but still noticeable difference in latency between the cameras. 
%
An interesting early sensor fusion method is also found in \cite{Unlu19} where the image from the narrow-angle camera is inserted into the image from the wide-angle camera and then processed by a single YOLO detector. It is unclear how they avoid situations when the inserted image obscures the object found in the wide-angle image.
%
In \cite{Diamantidou19}, they implemented early sensor fusion by concatenating feature vectors extracted from three different sensors (visible, thermal and 2D-radar), which are fed to a trained multilayer perceptron (MLP) classifier.
%
To do so here would likely require much more training data, not only on an individual sensor level but especially on a system level. Such amounts of data have not been possible to achieve within the scope of this work. 

The mentioned issues are the motives for implementing non-trained late sensor fusion in this work.
%
The sensor fusion implemented consists of utilizing the class outputs and the confidence scores of the available sensors in a weighed manner, smoothing their result over time.
%
A weighted fusion approach is shown to be more robust compared to other non-trained techniques such as voting, majority rule, the arithmetic combination of confidences (via, e.g. mean, median or product), or taking the most confident classifier (max rule) \cite{FIERREZ201857}.
%
To carry out the fusion, every time the main script polls the worker’s queues, it puts the results in a 4$\times$4 matrix, organized so that each class is a column and each sensor is a row. 
%
The matrix values depend not only on the class label and the confidence but also on the setting of which sensors to include and the weight of the specific sensor, i.e. how much we trust it at the moment.
%
A new 1$\times$4 matrix is then formed by column-wise sum.
%
This array is in turn placed as a new row in a 10$\times$4 first-in-first-out time-smoothing matrix. Since we have close to ten FPS from the workers, this will be the results of approximately the last second. Once again, the 10 columns are summarized into a 1$\times$4 matrix, and the column with the highest value will be the output system class.
%
The system output confidence is calculated by normalizing the value found to be highest over the number of sensors included at the moment. An additional condition before presenting the system result is that the minimum number of sensors that detects any object must fulfil the GUI setting, as shown in Figure~\ref{fig:fusion-options}.
%
The figure also shows an example of how sensor fusion is enabled in the GUI, including the weight choice for each sensor and the minimum number required.
%
With such a dynamical setting, it is possible to use not only the OR-function, as in \cite{Shi18}, but more sophisticated variants by varying the number of sensors included and required for detection, including their weights.


\begin{figure} [htb]
\centering
\includegraphics[width=0.45\textwidth]{fusion1.png}
\includegraphics[width=0.45\textwidth]{fusion2.png}
\caption{Two examples of sensor fusion results. Pictures were originally appearing in \cite{drone20thesis} and published in \cite{drone20icpr}. Reprinted with permission.}
\label{fig:fusion-options}
\end{figure}


Evaluating the fusion in operational conditions was more challenging than expected since this research was mostly done during the COVID19 pandemic. 
%
Regular flights at the airports considered in this work were cancelled, hence the possibility for a thorough system evaluation against airplanes.
%
Using the screen recorder, it was possible to do a frame-by-frame analysis of a typical drone detection, as shown in Figure~\ref{fig:fusion-results1}a.
%
The ``servo'' column indicates the current servo controlling sensor and the ``Fcam'' column specifies if the fish-eye camera motion detector is tracking the drone. 
%
The respective class output labels of each worker are shown in the rest of the columns. Note that the system output is more stable and lasts for more frames than the thermal infrared and video camera individually, indicating the benefit of the sensor fusion. 
%
Since there is no information from the ADS-B receiver in this case, that column has been omitted from the table.
%
Figure~\ref{fig:system_GUI} above is the third frame from 14:46:18. As it can be seen, the thermal infrared, video camera, and audio workers detect and classify the drone correctly. The fish-eye camera worker also tracks the drone, and the thermal infrared worker controls the pan/tilt platform.


\begin{figure} [htb]
\centering
%\includegraphics[width=0.12\textwidth]{blank.png}
%\includegraphics[width=0.12\textwidth]{blank.png}
%\includegraphics[width=0.12\textwidth]{blank.png}
%\includegraphics[width=0.45\textwidth]{fusion_results1.png}
%\includegraphics[width=0.45\textwidth]{fusion_results2.png}
\includegraphics[width=0.95\textwidth]{fusion_results1-2.png}
\caption{Evaluating the efficiency the sensor fusion.
%Left: frame-by-frame analysis of drone detection during one evaluation session. Right: false detections appearing in a ten minutes long section of screen recording from an evaluation session, including the type of object causing the false detection. 
Pictures were originally appearing in \cite{drone20thesis} and published in \cite{drone20icpr}. Reprinted with permission.
}
\label{fig:fusion-results1}
\end{figure}


To measure the efficiency of the sensor fusion, we can consider occasions such as the one described in Figure~\ref{fig:fusion-results1}a as a \textit{detection opportunity}. If we define this to be when the drone is continuously observable in the field of view of the thermal infrared and video cameras, and hence possible for the system (including the audio classifier) to analyse and track, we have carried out 73 such opportunities in the screen recordings from the evaluation sessions.
%
The duration of the individual detection opportunities is from just fractions of a second up to 28 seconds. This is highly dependent on how the drone is flying and whether the system can track the drone. We can see that Figure~\ref{fig:fusion-results1}a describes the frame-by-frame analysis of a detection opportunity lasting for three seconds.
%
Comparing the system results after the sensor fusion with the output from the respective sensors, we can observe that the system outputs a drone classification at some time in 78\% of the detection opportunities. Closest to this is the performance of the video camera detector that outputs a drone classification in 67\% of the opportunities.


We have also looked at the system behaviour without a drone flying in front of it to analyse false detections.
%
To do this, a ten-minute section of videos from the evaluation sessions was reviewed frame-by-frame.
%
Figure~\ref{fig:fusion-results1}b shows the timestamps, sensor types, causes of false detection, and resulting output labels. 
%
Setting the minimum number of sensors option to two prevents all the false detections in the figure from becoming false detections on a system level.
%
The false detections caused by insects flying just in front of the sensors are very short-lived, while the ones caused by clouds can last longer, up to several seconds.
%
Figure~\ref{fig:fig43} shows the false detections of the thermal infrared at 02:34 and the video camera at 02:58.
%
As described earlier, the individual weaknesses observed for the primary sensors are  sensitivity to clouds (thermal infrared) and autofocus problems (video camera). However, the fusion detection pipeline shows that such individual shortcomings can be overcome with a multi-sensor solution. 

Some other screenshots from the system evaluation sessions and interesting complementary observations are also pointed out in Appendix~\ref{sect:appendix}.


\begin{figure} [htb]
\centering
\includegraphics[width=0.7\textwidth]{fig43.jpg}
\caption{False detections of the thermal infrared at 02:34 and the video camera at 02:58 indicated in Figure~\ref{fig:fusion-results1}b. Pictures originally appearing in \cite{drone20thesis}.}
\label{fig:fig43}
\end{figure}



%%%%%%%%%%%%%%%%%%%%%%%%%%%%%%%%%%%%%%%%%%
%\section{Discussion}
%
%Authors should discuss the results and how they can be interpreted from the perspective of previous studies and of the working hypotheses. The findings and their implications should be discussed in the broadest context possible. Future research directions may also be highlighted.
%


%%%%%%%%%%%%%%%%%%%%%%%%%%%%%%%%%%%%%%%%%%
\section{Conclusions}
\label{sect:conclusions}

The increased use of drones for many recreational and professional purposes is creating concerns about the safety and security of premises where such vehicles may create risky situations or compromise security. 
%
Here, we explore the design of a multi-sensor drone detection system that employs state-of-the-art feature extraction and machine learning techniques, e.g. YOLOv2 detector \cite{Redmon17}, GMM background subtraction \cite{Stauffer99,gmmmatlab}, Kalman filters \cite{kalmanmatlab}, MFCC audio features or LSTM classifiers \cite{Hochreiter97}.
%
We employ a standard video camera and audio microphone complemented with a thermal infrared camera.
%
These two cameras are steered towards potential objects of interest with a fish-eye lens camera with a wider field of view that is used to detect moving objects at once over a bigger portion of the sky. 
%
We also include an ADS-B receiver, which allows for tracking cooperative aircrafts that broadcast their vehicle type, and a GPS receiver.
%
We have also evaluated a radar module, which has been finally discarded due to its very short practical range.
%
The sensors are all placed on a pan/tilt platform mounted on a standard surveyor's tripod.
%
This allows easy deployment of the solution outdoors and portability since the system can be disassembled into a few large parts and placed in a transport box.


The performance of the individual sensors is evaluated in terms of precision, recall, F1-score and mean average precision (mAP).
%
We observe that thermal infrared sensors are suitable for the drone detection task with machine learning techniques. 
%
The infrared detector achieves an F1-score of 0.7601, showing similar performance as the visible video detector with an F1-score of 0.7849. The audio classifier achieves an F1-score of 0.9323.
%
By sensor fusion, we also make the detection and classification more robust than any of the sensors individually, showing the efficiency of sensor fusion as a means to mitigate false detections. 
%
Another novelty is the investigation of detection performance as a function of the sensor-to-target distance.
%
We also contribute with a multi-sensor dataset, which overcomes the lack of publicly available data for this task. %
Apart from the feasibility of the sensors employed, especially the thermal infrared one, our dataset also uses an expanded number of target classes compared to related papers.
%
Our dataset is also especially suited for the comparison of infrared and visible video detectors due to the similarities in conditions and target types in the set.
%
To the best of our knowledge, we are also the first to explore the benefits of including ADS-B data to better separate targets prone to be mistaken for drones.

Future research could implement a distance estimation function based on the output from the detectors. Such research could also include investigating distributed detection and tracking and further using the temporal dimension to separate drones from other targets based on their behaviour over time as the system tracks them.
%
In a distributed scenario \cite{Yin20ojsp_cooperative}, several agents (detection stations in our case) could cooperate, each having access to only a few sensing modalities or varying computing capabilities or battery power.
%
Replicating one collection station with many different sensors may be cost-prohibitive or difficult if transport and deployment need to be done quickly (e.g. we used a lighter version with only two sensors and manual steering for dataset collection, Figure~\ref{fig:data_collection}). 
%
However, cooperative operation entails challenges, such as handling orchestration and task distribution between units, connectivity via 4G/5G networks, or cloud/edge computing if the deployed hardware lacks such capability. 
%
Another challenge is the optimal fusion of information from sources that may have different data quality (e.g. sensors with different features), different modalities (e.g. visible, NIR, audio...), or different availability at a given time (e.g. a specific sensor is not deployed or the object is not visible from that location). In the case of visual sensors, there is also a need to match targets observed by sensors in different physical locations since the object is seen from a different point of view.  On the other hand, the first station that detects a target can notify it and provide helpful information to aid in the detection by the others.  
%

%
Another direction is to implement the YOLO detector in the fish-eye camera. However, this would demand a training dataset or skewing video camera images to match the fish-eye lens distortion.
%
It would also be interesting to use YOLO v3 instead since it is more efficient in detecting small objects according to \cite{Unlu19}, or even more recent versions, which currently go up to YOLO v7.
%
In this paper, we kept using YOLO v2 since it is the predominant choice in the literature (as seen in recent surveys \cite{9765451}), enabling a fairer comparison of our results with previous works.
%
The performance of the audio classifier with respect to the sensor-to-target distance could also be explored in the same way as the video sensors. 
%
Furthermore, the influence of covariates given by different weather conditions could also be investigated.

A radar with a good range would have contributed significantly to the system results since it is the only available sensor that can measure the distance to the target efficiently. 
%
%Another way to increase the efficiency of the system could also be to exploit the temporal dimension better, i.e. to use the flight paths and the behaviour of the object to better classify them.
%
Likewise, it would be interesting to test the system against a drone equipped with a transponder to see the performance of the ADS-B worker at a short distance. Such transponders weigh as little as 20 grams \cite{uAvioni}. However, the price of such equipment ($\sim$3k€) is still an issue for non-professional drone users.
%
Since ADS-B information is included in the system, this could also be implemented as ground truth for an online learning feature. Images of ADS-B targets could be saved, and the detectors would be retrained at regular intervals using these additional images. To further enlarge the training data set, all images of detected objects could be saved and annotated manually at a later stage.

This work can also be helpful in other areas.
%
One example is road traffic surveillance, since most parts and scripts (except the ADS-B receiver) are applicable after appropriate retraining to detect and track vulnerable road users (pedestrians), vehicles, etc. and even specific vehicle types such as light ones (bikes, motorcycles) or heavy ones (trucks).
%
Another application involving the detection of people is surveillance of large areas where access is restricted, and operator monitoring can be very resource-consuming if not automated, such as outdoors.
%



%%%%%%%%%%%%%%%%%%%%%%%%%%%%%%%%%%%%%%%%%%
%\section{Patents}
%
%This section is not mandatory, but may be added if there are patents resulting from the work reported in this manuscript.

%%%%%%%%%%%%%%%%%%%%%%%%%%%%%%%%%%%%%%%%%%
\vspace{6pt} 

%%%%%%%%%%%%%%%%%%%%%%%%%%%%%%%%%%%%%%%%%%
%% optional
%\supplementary{The following supporting information can be downloaded at:  \linksupplementary{s1}, Figure S1: title; Table S1: title; Video S1: title.}

% Only for the journal Methods and Protocols:
% If you wish to submit a video article, please do so with any other supplementary material.
% \supplementary{The following supporting information can be downloaded at: \linksupplementary{s1}, Figure S1: title; Table S1: title; Video S1: title. A supporting video article is available at doi: link.}

%%%%%%%%%%%%%%%%%%%%%%%%%%%%%%%%%%%%%%%%%%
\authorcontributions{
%For research articles with several authors, a short paragraph specifying their individual contributions must be provided. The following statements should be used ``Conceptualization, X.X. and Y.Y.; methodology, X.X.; software, X.X.; validation, X.X., Y.Y. and Z.Z.; formal analysis, X.X.; investigation, X.X.; resources, X.X.; data curation, X.X.; writing---original draft preparation, X.X.; writing---review and editing, X.X.; visualization, X.X.; supervision, X.X.; project administration, X.X.; funding acquisition, Y.Y. All authors have read and agreed to the published version of the manuscript.'', please turn to the  \href{http://img.mdpi.org/data/contributor-role-instruction.pdf}{CRediT taxonomy} for the term explanation. Authorship must be limited to those who have contributed substantially to the work~reported.
%
This work has been carried out by Fredrik Svanström under the supervision of Fernando Alonso-Fernandez and Cristofer Englund in the context of his Master's Thesis at Halmstad University (Master's Programme in Embedded and Intelligent Systems). The thesis is available at \cite{drone20thesis}.  
%
Fredrik Svanström: Conceptualization, Methodology, Investigation, Data curation, Writing –review \& editing; Fernando Alonso-Fernandez: Conceptualization, Supervision, Funding acquisition, Writing – original draft; Cristofer Englund: Conceptualization, Supervision, Writing – review \& editing.
}

\funding{
%Please add: ``This research received no external funding'' or ``This research was funded by NAME OF FUNDER grant number XXX.'' and  and ``The APC was funded by XXX''. Check carefully that the details given are accurate and use the standard spelling of funding agency names at \url{https://search.crossref.org/funding}, any errors may affect your future funding.
%
Author F. A.-F. thanks the Swedish Research Council (VR) for funding his research.
%
Authors F. A.-F. and C. E. thank the Swedish Innovation Agency (VINNOVA) for funding their research.
}

%\institutionalreview{In this section, you should add the Institutional Review Board Statement and approval number, if relevant to your study. You might choose to exclude this statement if the study did not require ethical approval. Please note that the Editorial Office might ask you for further information. Please add “The study was conducted in accordance with the Declaration of Helsinki, and approved by the Institutional Review Board (or Ethics Committee) of NAME OF INSTITUTE (protocol code XXX and date of approval).” for studies involving humans. OR “The animal study protocol was approved by the Institutional Review Board (or Ethics Committee) of NAME OF INSTITUTE (protocol code XXX and date of approval).” for studies involving animals. OR “Ethical review and approval were waived for this study due to REASON (please provide a detailed justification).” OR “Not applicable” for studies not involving humans or animals.}

%\informedconsent{Any research article describing a study involving humans should contain this statement. Please add ``Informed consent was obtained from all subjects involved in the study.'' OR ``Patient consent was waived due to REASON (please provide a detailed justification).'' OR ``Not applicable'' for studies not involving humans. You might also choose to exclude this statement if the study did not involve humans.
%
%Written informed consent for publication must be obtained from participating patients who can be identified (including by the patients themselves). Please state ``Written informed consent has been obtained from the patient(s) to publish this paper'' if applicable.}

\dataavailability{
%In this section, please provide details regarding where data supporting reported results can be found, including links to publicly archived datasets analyzed or generated during the study. Please refer to suggested Data Availability Statements in section ``MDPI Research Data Policies'' at \url{https://www.mdpi.com/ethics}. If the study did not report any data, you might add ``Not applicable'' here.
%
The data used in this paper is fully described in \cite{DiBdataset} and publicly available at \cite{svanstrom_fredrik_2020_5500576}.
} 

%\acknowledgments{
%In this section you can acknowledge any support given which is not covered by the author contribution or funding sections. This may include administrative and technical support, or donations in kind (e.g., materials used for experiments).
%}

\conflictsofinterest{The authors declare no conflict of interest.}
%{Declare conflicts of interest or state ``The authors declare no conflict of interest.'' Authors must identify and declare any personal circumstances or interest that may be perceived as inappropriately influencing the representation or interpretation of reported research results. Any role of the funders in the design of the study; in the collection, analyses or interpretation of data; in the writing of the manuscript; or in the decision to publish the results must be declared in this section. If there is no role, please state ``The funders had no role in the design of the study; in the collection, analyses, or interpretation of data; in the writing of the manuscript; or in the decision to publish the~results''.} 

%%%%%%%%%%%%%%%%%%%%%%%%%%%%%%%%%%%%%%%%%%
%% Optional
%\sampleavailability{Samples of the compounds ... are available from the authors.}

%% Only for journal Encyclopedia
%\entrylink{The Link to this entry published on the encyclopedia platform.}

\abbreviations{Abbreviations}{
The following abbreviations are used in this manuscript:\\

\noindent 
\begin{tabular}{@{}ll}
MDPI & Multidisciplinary Digital Publishing Institute\\
DOAJ & Directory of open access journals\\
TLA & Three letters acronym\\
LD & Linear dichroism
\end{tabular}
}

%%%%%%%%%%%%%%%%%%%%%%%%%%%%%%%%%%%%%%%%%%
%% Optional
\appendixtitles{yes} % Leave argument "no" if all appendix headings stay EMPTY (then no dot is printed after "Appendix A"). If the appendix sections contain a heading then change the argument to "yes".
\appendixstart
\appendix
\section[\appendixname~\thesection]{Complementary screenshots and observations to the sensor fusion evaluation sessions}
\label{sect:appendix}

%\subsection[\appendixname~\thesubsection]{}
%The appendix is an optional section that can contain details and data supplemental to the main text---for example, explanations of experimental details that would disrupt the flow of the main text but nonetheless remain crucial to understanding and reproducing the research shown; figures of replicates for experiments of which representative data are shown in the main text can be added here if brief, or as Supplementary Data. Mathematical proofs of results not central to the paper can be added as an appendix.

%\section[\appendixname~\thesection]{}
%All appendix sections must be cited in the main text. In the appendices, Figures, Tables, etc. should be labeled, starting with ``A''---e.g., Figure A1, Figure A2, etc.

In this section, some screenshots from the system evaluation sessions and some interesting observations are pointed out.
%
The images also indicate the FPS performance. 
%
The system can process 6 FPS or more from all cameras and over 10 processing cycles per second for the input audio stream. %Note that the ADS-B solution chosen is the more computationally demanding out of the two alternatives evaluated. 
%
To be able to evaluate the system performance, a screen recording software has also been running on the computer at the same time as the system software. 
%
%This setup was a necessity since the drone flying took all the computational power of the thesis author during the evaluation sessions.
%

The ideal situation is that all the sensors output the correct classification of the detected target and that the fish-eye camera tracks the object. This is, however, far from the case at all times.
%
Nevertheless, after the sensor fusion, the system output class is observed to be robust, as shown in Figure~\ref{fig:fig44} and \ref{fig:fig45}, where the video and thermal infrared cameras classify the drone incorrectly, but still with a correct system output.

\begin{figure} [htb]
\centering
\includegraphics[width=0.9\textwidth]{fig44.jpg}
\caption{A correct system output even if the video camera classifies the drone as a bird. Picture originally appearing in \cite{drone20thesis}.}
\label{fig:fig44}
\end{figure}

\begin{figure} [htb]
\centering
\includegraphics[width=0.9\textwidth]{fig45.jpg}
\caption{A correct system output even if the thermal infrared camera classifies the drone as a helicopter. Picture originally appearing in \cite{drone20thesis}.}
\label{fig:fig45}
\end{figure}

Since the output class depends on the confidence score, the result is sometimes the opposite, as shown in Figure~\ref{fig:fig46}, so a very confident sensor (audio) causes the system output to be wrong. 
%
If this turns out to be frequent, the weight of the sensor can easily be adjusted, or the sensor can be excluded entirely from the system result. 
%
The time smoothing procedure of the sensor fusion will also reduce the effect of an occasional misclassification so that the system output stays correct, as seen in Figure~\ref{fig:fig47}. Naturally, there are also times when all sensors are wrong, as evident in Figure~\ref{fig:fig48}.

\begin{figure} [htb]
\centering
\includegraphics[width=0.9\textwidth]{fig46.jpg}
\caption{The high confidence score of the audio classifier causes the system output to be incorrect, just like the audio output. Picture originally appearing in \cite{drone20thesis}.}
\label{fig:fig46}
\end{figure}

\begin{figure} [htb]
\centering
\includegraphics[width=0.9\textwidth]{fig47.jpg}
\caption{The time smoothing part of the sensor fusion reduces the effect of an occasional misclassification, even if that has a high confidence score. Picture originally appearing in \cite{drone20thesis}.}
\label{fig:fig47}
\end{figure}

\begin{figure} [htb]
\centering
\includegraphics[width=0.9\textwidth]{fig48.jpg}
\caption{Several sensors misclassify the drone at the same time. Picture originally appearing in \cite{drone20thesis}.}
\label{fig:fig48}
\end{figure}


To the best of our knowledge, the inclusion of an ADS-B receiver in a drone detection system has not yet been described in the scientific literature, so it is interesting to see how this information is utilized.
%
In Figure~\ref{fig:fig32}, an airplane is detected and classified at a sloping distance of more than 35000 m. The ADS-B information guides the pan/tilt platform in the direction of the airplane so that the video camera can detect it.
%
At this distance, even a large commercial airplane is only about 1.4 pixels, hence well below the detection limit of the DRI criteria.
%
The reason behind the fact that the ADS-B FoV-target distance display is empty is that the main script will not present that information until the target is within a 30000 m horizontal distance. This limit is set based on the assumption that no target beyond 30000 m should be detectable, which turned out to be wrong.
%

\begin{figure} [htb]
\centering
\includegraphics[width=0.9\textwidth]{fig32.jpg}
\caption{An airplane detected and classified correctly by the video camera worker at a sloping distance of more than 35000 m. See the text for details. Picture originally appearing in \cite{drone20thesis}.}
\label{fig:fig32}
\end{figure}


Looking at Figure~\ref{fig:fig49}, we can see that the ADS-B information will appear in the results panel when the airplane comes within 30000 m horizontal distance from the system. At this moment, the sloping distance is 32000 m, and the offset between the camera direction and the calculated one is zero. Moreover, since the system has not yet received the vehicle category information, the target is marked with a square in the ADS-B presentation area, and the confidence score of the ADS-B result is 0.75.
%
The next interesting event, shown in Figure~\ref{fig:fig50}, is when the system receives the vehicle category message. To indicate this, the symbol in the ADS-B presentation is changed to a circle, and the confidence is set to one since we are sure that it is an airplane. At a distance of 20800 m, it is also detected and classified correctly by the thermal infrared worker, as shown in Figure~\ref{fig:fig51}.

\begin{figure} [htb]
\centering
\includegraphics[width=0.9\textwidth]{fig49.jpg}
\caption{When the airplane is within 30000 m horizontal distance the ADS-B information is presented in the results panel. Picture originally appearing in \cite{drone20thesis}.}
\label{fig:fig49}
\end{figure}

\begin{figure} [htb]
\centering
\includegraphics[width=0.9\textwidth]{fig50.jpg}
\caption{The system has received the vehicle category message, so the confidence for the airplane classification is set to 1. The airplane is now at a distance of 23900m. Picture originally appearing in \cite{drone20thesis}.}
\label{fig:fig50}
\end{figure}


\begin{figure} [htb]
\centering
\includegraphics[width=0.9\textwidth]{fig51.jpg}
\caption{The thermal infrared worker detects the airplane at a distance of 20800 m. Picture originally appearing in \cite{drone20thesis}.}
\label{fig:fig51}
\end{figure}

\clearpage

%%%%%%%%%%%%%%%%%%%%%%%%%%%%%%%%%%%%%%%%%%
\begin{adjustwidth}{-\extralength}{0cm}
%\printendnotes[custom] % Un-comment to print a list of endnotes

\reftitle{References}

% Please provide either the correct journal abbreviation (e.g. according to the “List of Title Word Abbreviations” http://www.issn.org/services/online-services/access-to-the-ltwa/) or the full name of the journal.
% Citations and References in Supplementary files are permitted provided that they also appear in the reference list here. 

%=====================================
% References, variant A: external bibliography
%=====================================
%\bibliography{your_external_BibTeX_file}

\bibliography{bibliography}

%=====================================
% References, variant B: internal bibliography
%=====================================
%\begin{thebibliography}{999}
%% Reference 1
%\bibitem[Author1(year)]{ref-journal}
%Author~1, T. The title of the cited article. {\em Journal Abbreviation} {\bf 2008}, {\em 10}, 142--149.
%% Reference 2
%\bibitem[Author2(year)]{ref-book1}
%Author~2, L. The title of the cited contribution. In {\em The Book Title}; Editor 1, F., Editor 2, A., Eds.; Publishing House: City, Country, 2007; pp. 32--58.
%% Reference 3
%\bibitem[Author3(year)]{ref-book2}
%Author 1, A.; Author 2, B. \textit{Book Title}, 3rd ed.; Publisher: Publisher Location, Country, 2008; pp. 154--196.
%% Reference 4
%\bibitem[Author4(year)]{ref-unpublish}
%Author 1, A.B.; Author 2, C. Title of Unpublished Work. \textit{Abbreviated Journal Name} year, \textit{phrase indicating stage of publication (submitted; accepted; in press)}.
%% Reference 5
%\bibitem[Author5(year)]{ref-communication}
%Author 1, A.B. (University, City, State, Country); Author 2, C. (Institute, City, State, Country). Personal communication, 2012.
%% Reference 6
%\bibitem[Author6(year)]{ref-proceeding}
%Author 1, A.B.; Author 2, C.D.; Author 3, E.F. Title of presentation. In Proceedings of the Name of the Conference, Location of Conference, Country, Date of Conference (Day Month Year); Abstract Number (optional), Pagination (optional).
%% Reference 7
%\bibitem[Author7(year)]{ref-thesis}
%Author 1, A.B. Title of Thesis. Level of Thesis, Degree-Granting University, Location of University, Date of Completion.
%% Reference 8
%\bibitem[Author8(year)]{ref-url}
%Title of Site. Available online: URL (accessed on Day Month Year).
%\end{thebibliography} 

% If authors have biography, please use the format below
%\section*{Short Biography of Authors}
%\bio
%{\raisebox{-0.35cm}{\includegraphics[width=3.5cm,height=5.3cm,clip,keepaspectratio]{Definitions/author1.pdf}}}
%{\textbf{Firstname Lastname} Biography of first author}
%
%\bio
%{\raisebox{-0.35cm}{\includegraphics[width=3.5cm,height=5.3cm,clip,keepaspectratio]{Definitions/author2.jpg}}}
%{\textbf{Firstname Lastname} Biography of second author}

% For the MDPI journals use author-date citation, please follow the formatting guidelines on http://www.mdpi.com/authors/references
% To cite two works by the same author: \citeauthor{ref-journal-1a} (\citeyear{ref-journal-1a}, \citeyear{ref-journal-1b}). This produces: Whittaker (1967, 1975)
% To cite two works by the same author with specific pages: \citeauthor{ref-journal-3a} (\citeyear{ref-journal-3a}, p. 328; \citeyear{ref-journal-3b}, p.475). This produces: Wong (1999, p. 328; 2000, p. 475)

%%%%%%%%%%%%%%%%%%%%%%%%%%%%%%%%%%%%%%%%%%
%% for journal Sci
%\reviewreports{\\
%Reviewer 1 comments and authors’ response\\
%Reviewer 2 comments and authors’ response\\
%Reviewer 3 comments and authors’ response
%}
%%%%%%%%%%%%%%%%%%%%%%%%%%%%%%%%%%%%%%%%%%
\end{adjustwidth}
\end{document}



%\newcommand{\E}{\mathbf{E}}

%\newcommand{\Dist}{\mathcal{D}}
%\newcommand{\Inter}[1]{\widetilde{#1}}
%\newcommand{\Poly}[1]{poly(#1)}
%\newcommand{\CC}[1]{\widehat{#1}}
%\newcommand{\Src}{\overline{src}}
%\newcommand{\Snk}{\overline{snk}}
%\newcommand{\Ol}[1]{\overline{#1}}
%\newcommand{\Ul}[1]{\underline{#1}}



\section{Preliminaries}
\label{sec:prelim}

\paragraph{Single-agent Mechanisms}
%%
%% services
%%

%A service provider faces $n$ agents and may offer service to any
%subset of size at most $k$ of them.

We consider the provisioning of an abstract service.  This service may
be parameterized by an {\em attribute}, e.g., quality of service, and
may be accompanied by a required payment.  We denote the outcome
obtained by an agent as $\outcome \in \outcomespace$.  We view this
outcome as giving an indicator for whether or not an agent is served and
as describing attributes of the service such as quality of service and
monetary payments.  Let $\ALLOC(\outcome) \in \{0,1\}$ be an indicator
for whether the agent is served or not; let $\PAYMENT(\outcome)
\in \reals$ denote any payment the agent is required to make.  In a
randomized environment (e.g., randomness from a randomized mechanism or Bayesian
environment) the outcome an agent receives is a random variable from a
distribution over $\outcomespace$.  The space of all such
distributions is denoted $\outcomedistspace$.

%%
%% types, utilities, prior
%%
The agent has a type $\type$ from a finite type space $\typespace$.
This type is drawn from distribution $\dens \in \distover{\typespace}$
and we equivalently denote by $\dens$ the probability mass function.
I.e., for every $\type \in \typespace$, $\dens(\type)$ is the
probability that the type is~$\type$.  The utility function $\util
\,:\, \typespace \times \outcomespace \to \reals$ maps the agent's
type and the outcome to real valued utility.  The agent is a von
Neumann--Morgenstern expected utility maximizer and we extend $\util$
to $\distoutcomespace$ linearly, i.e., for $\outcome
\in \distoutcomespace$, $\util(\type, \outcome)$ is the expectation of $\util$ where the outcome is drawn according
to~$\outcome$.
%$\util(\type,\outcome) = \expect[v \sim \outcome]{\util(\type,v)}$.
We do not require the usual assumption of quasi-linearity.



%%
%% incentives, service rule
%%
%By the revelation principle, we will restrict attention to Bayesian
%incentive compatible mechanisms, i.e., ones in which truthtelling is a
%Bayes-Nash equilibrium.  For an implicit mechanism we will denote the
%{\em outcome rule} by $\outcomes(\types)$ where
%$\outcome\agind(\types)$ is the service to agent $\agent$.  Similarly
%we will denote by $\allocs(\types)$ and $\prices(\types)$ the {\em
%  allocation} and {\em payment rule}.  These may be random variables
%and we denote their expectations by $\expallocs(\types) =
%\expect{\allocs(\type)}$ and $\expprices(\types) =
%\expect{\prices(\types)}$.
%When the agents' types are drawn from the
%product distribution $\dists = (\dist\agind[1] \times \cdots \times
%\dist\agind[n])$, agent $\agent$'s {\em interim outcome rule} is
%$\outcome\agind(\type\agind) =
%\evalat{\outcome\agind(\types)}{\type\agind}$ for $\types \sim \dists$ denoting the
%distribution over the service that agent $\agent$ receives when other
%types are drawn from the distribution.

%%
%% outcome rule
%%
A single-agent mechanism, without loss of generality by the revelation
principle, is just an {\em outcome rule}, a mapping from
the agent's type to a distribution over outcomes.  We denote an {\em
  outcome rule} by $\toutcome \,:\, \typespace
\to \outcomedistspace$.  We say that an outcome rule $\toutcome$ is
{\em incentive compatible} (IC) and {\em individually rational} (IR) if for
all $\type,\type' \in \typespace$, respectively,
\begin{align}
\tag{IC}\label{eq:IC}
\util(\type,\toutcome(\type)) &\geq \util(\type,\toutcome(\type')),\\
\tag{IR}\label{eq:IR}
\util(\type,\toutcome(\type)) &\geq 0.
\end{align}


%%
%% allocation rule
%%
%As the mechanism designer faces a constraint only on the allocation
%portion of the service it will be useful to separate out that part.
We refer to restriction of the outcome rule to the indicator for
service as the {\em allocation rule}.  As the allocation to
each agent is a binary random variable, distributions over allocations
are fully described by their expected value.  Therefore the allocation
rule $\talloc \,:\, \typespace \to [0,1]$ for a given outcome rule
$\outcome$ is $\talloc(\type) = \expect{\ALLOC(\toutcome(\type))}$.


%%
%% Tiger: added the following Examples
%%
We give two examples to illustrate the abstract model described above.
The first example is the standard quasi-linear risk-neutral preference
which is prevalent in auction theory.  Here the agent's type space is
$\typespace \subset \posreals$ where $\type \in \typespace$ represents the
agent's valuation for the item.  The outcome space is $\outcomespace =
\{0, 1\} \times \posreals$ where an outcome $\outcome$ in this space
indicates whether or not the item is sold to the agent, by
$\ALLOC(\outcome)$, and at what price, by $\PAYMENT(\outcome)$.  The
agent's quasi-linear utility function is $\util(\type, \outcome) =
\type \cdot \ALLOC(\outcome) - \PAYMENT(\outcome)$.
%
% example 2
%
The second example is that of an $\numservice$-item unit-demand (also
quasi-linear and risk-neutral) preference.  Here the type space is
$\typespace \subset \posreals^\numservice$ and a type $\type
\in \typespace$ indicates the agent's valuation for each of the items
when the agent's value for no service is normalized to zero.  An
outcome space is $\outcomespace = \{0,\ldots,\numservice\} \times
\posreals$.  The first coordinate of $\outcome$ specifies which item
the agent receives or none and $\ALLOC(\outcome) = 1$ if it is
non-zero; the second coordinate of $\outcome$ specifies the required
payment $\PAYMENT(\outcome)$.  The agent's utility for $\outcome$ is
the value the agent attains for the item received less her payment.
Beyond these two examples, our framework can easily incorporate more
general agent preferences exhibiting, e.g., risk aversion or a budget
limit.


\Xcomment{
%%
%% Tiger: added the following Examples
%%
We give two examples to illustrate the set of notations above.  The
first example is the standard quasi-linear risk-neutral preference
which is prevalent in auction theory.  Here the type of the agent is
simply her valuation for the item, an outcome $\outcome$ is a tuple
$(\sdalloc, \sdprice) \in \{0, 1\} \times \posreals$ that indicates
(by $\sdalloc$) whether or not the item is sold to the agent and (by
$\sdprice$) at what price.  $\ALLOC(\sdalloc,\sdprice) = \sdalloc$ and
$\PAYMENT(\sdalloc,\sdprice)$ is $\sdprice$.  The agent's quasi-linear
utility function is $\util(\type, (\sdalloc,\sdprice)) = \type
\sdalloc - \sdprice$.
%
% example 2
%
The second example is that of a $\numservice$-item unit-demand (also
quasi-linear and risk-neutral) preference.  Here the type $\type :
\{0,\ldots, \numservice\} \to \posreals$ indicates the agent's
valuation for each of the items or $\mdtype{0} = 0$ for no service.
An outcome $\outcome = (\mdalloc, \mdprice) \in
\{0,\ldots,\numservice\} \times \posreals$.  Here, $\mdalloc$
specifies which item the agent receives or none if $\mdalloc = 0$,
therefore, $\ALLOC(\mdalloc,\mdprice) = 1$ if $\mdalloc > 0$ and zero
otherwise, $\PAYMENT(\mdalloc,\mdprice) = \mdprice$.  The agent's
utility is $\util(\type,(\mdalloc,\mdprice)) = \mdtype{\mdalloc} -
\mdprice$.  Of course, our framework can easily incorporate more
general settings such as those involving a risk averse agent with a
budget limit.
}


%%
%% quantile
%%
% An allocation rule $\talloc$ induces a partial ordering on the type
% space.  Given two types $\type$ and $\type'$ we can order them by
% which is more likely to be allocated and refer to the one that is more
% likely to be allocated as the stronger one.  Furthermore, given the
% distribution $\dist$ over type space, we can associate with each type
% a quantile which refers to the strength of a type relative to the
% distribution.  The quantile $\quant \in [0,1]$ of a type is the
% probability that a random type from the distribution is stronger.  Let
% $\quantl$ be the probability that a random type is strictly stronger
% and let $\quante$ be the probability that a type is equally strong,
% i.e., $\quantl = \prob[\type'\sim\dist]{\talloc(\type') >
%   \talloc(\type)}$ and $\quante = \prob[\type'\sim\dist]{\talloc(\type')
%   = \talloc(\type)}$.  The quantile for a type $\type$ and type
% allocation rule $\alloc$, denoted $\QUANT(\type,\talloc)$, is a random
% draw from the distribution $\quantl + U[0,\quante]$.  Importantly,
% the distribution over quantiles of a random type $\type \sim \dist$ is
% $U[0,1]$.  (The randomization of $\quante$ mass is simply a random
% tie-breaking rule; we encourage the reader to assume no ties, i.e.,
% that $\quante = 0$, and that $\QUANT(\type,\talloc)$ is deterministic.)
%

%%
%% single-agent quantile-allocation-constrained mechanism design
%%
% The transformation from type to quantile and a type allocation rule
% induces a {\em (quantile) allocation rule} $\qalloc \,:\, [0,1] \to
% [0,1]$ that maps quantiles to allocation probabilities; for a given
% type-allocation rule $\alloc$ and the mapping of type to quantile that
% it induces, the quantile allocation rule $\qalloc$ is defined as
% $\qalloc(\QUANT(\type,\talloc)) = \talloc(\type)$.  Notice that there
% may be multiple type allocation rules with the same quantile
% allocation rule.  A single-agent mechanism design problem is the
% following: {\em Given a desired quantile allocation rule, design the
%   optimal type outcome rule that induces this quantile allocation
%   rule.}
%
%%
%% single-agent downward-closed quantile-allocation-constrained mechanism design
%%
% A more permissive requirement which is consistent with a standard
% downward-closure assumption is to require the outcome rule to induce a
% quantile allocation rule that dominated by the desired quantile
% allocation rule.  The {\em cumulative allocation rule} is the integral
% of the quantile allocation rule, i.e., $\cumalloc(\quant) =
% \int_0^\quant \qalloc(\dquant)\, d \dquant$ for quantile allocation
% rule $\qalloc$.  A cumulative allocation rule $\cumalloc$ {\em
%   dominates} $\cumalloc'$ if for all $\quant$, $\cumalloc(\quant) \geq
% \cumalloc'(\quant)$.  An outome or allocation rule dominates another
% outcome or allocation rule of their induced cumulative allocation
% rules satisfy the dominance condition.  Our single-agent mechanism
% design problem is the following: {\em Given a desired quantile
%   allocation rule, design the optimal type outcome rule that is
%   dominated by the desired quantile allocation rule.}
%
%%
%% single-agent ex ante constrained mechanism design
%%
% A special case of quantile allocation rule is the one that places an
% ex ante constraint on the probability that the agent is allocated.
% For a given constraint $\aquant$ this allocation rule is a step
% function, denoted $\aqalloc[\aquant]$, which steps from 0 to 1 at $\quant =
% \aquant$.  Such a rule says to allocate to the agent with ex ante
% probability $\aquant$.
%

%%
%% single-agent optimal mechanisms
%%
% Given a quantile allocation rule $\qalloc$ for a type space
% $\typespace$ and a distribution over types $\dist$, a ``solution'' to
% the single-agent problem constructs an optimal type outcome rule
% $\RULE(\qalloc)$ with expected performance which we denote
% $\REV(\qalloc)$.  For the special case above of an ex ante allocation
% probability constraint $\aquant$, i.e., the quantile allocation rule
% $\aqalloc[\aquant]$, we overload notation and write $\REV(\aquant)$
% and $\RULE(\aquant)$ for $\REV(\aqalloc[\aquant])$ and
% $\RULE({\aqalloc[\aquant]})$, respectively.
%

%%
%% single-agent problem
%%

Consider the following single-agent mechanism design problem.  A
feasibility constraint is given by an upper bound $\talloc(\type)$ on
the probability that the agent is served as a function of her type
$\type$; the distribution on types in $\typespace$ is given by
$\dens$.  The single-agent problem is to find the outcome rule
$\toutcome^*$ that satisfies the allocation constraint of $\talloc$
and maximizes the performance, e.g., revenue.  This problem is
described by the following program:
\begin{align*}
\max_{\toutcome}:  & \quad \expect[\type \sim \dens, \toutcome(\type)]{\PAYMENT(\toutcome(\type))}
\tag{SP}\label{eq:rule-rev-def} \\
\hbox{s.t. } & \quad \expect[\toutcome(\type)]{\ALLOC(\toutcome(\type))} \leq \talloc(\type), \qquad \forall \type \in \typespace \\
 & \quad \toutcome \text{ is IC and IR}.
\end{align*}
We denote the outcome rule $\toutcome^*$ that optimizes this program
by $\RULE(\talloc)$ and its revenue by $\REV(\talloc) = \expect[\type
  \sim \dens, \toutcome^*(\type)]{\PAYMENT(\toutcome^*(\type))}$.  We note
that, although this paper focuses on revenue maximization, the same
techniques presented can be applied to maximize (or minimize) general
separable objectives such as social welfare.

\paragraph{Multi-agent Mechanisms}


%%
%% feasibility
%%
There are $n$ independent agents.  Agents need not be identical,
i.e., agent $\agent$'s type space is~$\typespace\agind$, the
probability mass function for her type is~$\dens\agind$,
her outcome space is $\outcomespace\agind$, and her utility function
is $\util\agind$. The profile of agent types is denoted by $\types =
(\type\agind[1],\ldots,\type\agind[n]) \in \typespace\agind[1] \times
\cdots \times \typespace\agind[n] = \typespaces$, the joint
distribution on types is $\denss \in \distover{\typespace_1}\times \cdots \times \distover{\typespace_n}$, a vector of outcomes is
%% Tiger:  removed "\outcomes =" and "\allocs = ".  we never use them in the paper for deterministic outcomes, and the
%%     current notation coincides with the randomized outcomes, which may cause confusion.  Also, replaced all
%%     semi-colons by commas in this sentence
%$\outcomes =
$(\outcome\agind[1], \cdots, \outcome\agind[n]) \in \outcomespaces$, and an allocation is
%$\allocs =
$(\alloc\agind[1],\ldots,\alloc\agind[n]) \in \{0,1\}^n$.  The
mechanism has an inter-agent feasibility constraint that permits
serving at most $k$~agents, i.e., $\sum_\agent \alloc\agind \leq
k$.\footnote{Furthremore, in Section~\ref{sec:border_matroid}, we
  review the theory of {\em matroids} and extend our basic results
  environments with feasibility constraint derived from a matroid set
  system.}  A mechanism that obeys this constraint is {\em
  feasible}. The mechanism has no inter-agent constraint on attributes
or payments.


%%
%% induced allocation rules
%%
A mechanism maps type profiles to a (distribution over) outcome
vectors via an {\em ex post outcome rule}, denoted $\epoutcomes \,:\,
\typespaces \to \outcomedistspaces$ where $\epoutcome\agind(\types)$
is the outcome obtained by agent $\agent$. We will similarly define
$\epallocs \,:\, \typespaces \to [0,1]^n$ as the {\em ex post
  allocation rule} (where $[0,1] \equiv \distover{\{0,1\}}$).  The ex
post allocation rule~$\epallocs$ and the probability mass
function~$\denss$ on types induce {\em interim} outcome and allocation
rules.  For agent~$\agent$ with type~$\type\agind$ and $\types \sim
\distribution[\types]{\types \given \type\agind}$ the interim outcome
and allocation rules are $\intoutcome\agind(\type\agind) =
\distribution[\types]{\epoutcome\agind(\types) \given \type\agind }$
and $\intalloc\agind(\type\agind) =
\distribution[\types]{\epalloc\agind(\types) \given \type\agind}
\equiv \expect[\types]{\epalloc\agind(\types) \given
  \type\agind}$.\footnote{We use notation $\distribution{X \given E}$
  to denote the distribution of random variable $X$ conditioned on the
  event $E$.}  A profile of interim allocation rules is feasible if it
is derived from an ex post allocation rule as described above; the set
of all feasible interim allocation rules is denoted by $\InAllocSpace$.
A mechanism is Bayesian incentive compatible and interim individually
rational if equations~\eqref{eq:IC} and~\eqref{eq:IR}, respectively,
hold for all $\agent$ and all~$\type\agind$.



%%
%% examples revisited
%%
Consider again the examples described previously of quasi-linear
single-dimensional and unit-demand preferences.  For the
single-dimensional example, the multi-agent mechanism design problem
is the standard single-item $k$-unit auction problem.  For the
unit-demand example, the multi-agent mechanism design problem is an
{\em attribute auction}.  In this problem there are $k$-units
available and each unit can be configured in one of $\numservice$
ways.  Importantly, the designer's feasibility constraint restricts
the number of units sold to be $k$ but places no restrictions on how
the units can be configured.  E.g., a restaurant has $k$ tables but
each diner can order any of the $m$ entrees on the menu.

% The interim type allocation rule of an agent, as described above, induces a mapping from type space to quantile space, and this mapping induces an interim quantile allocation rule.  We call the resulting profile $\qallocs = (\qalloc\agind[1],\ldots,\qalloc\agind[n])$ of quantile allocation rules {\em feasible} as there exists a ex post feasible outcome rule that, for the given distribution, induces them.

%% Tiger:  added \agind for all variables (since it makes more sense for talking about interim allocations).  added a pointer to Section single-agent
A reduction from multi-agent mechanism design to single-agent
mechanism design as we have described above would assume that for any
types pace $\typespace\agind$, any probability mass
function~$\dens\agind$, and interim allocation rule $\intalloc\agind$,
the optimal outcome rule $\RULE(\intalloc\agind)$ and its performance
$\REV(\intalloc\agind)$ can be found efficiently (see
Section~\ref{sec:single-agent} for examples).  The goal then is to
construct an optimal multi-agent auction from these single-agent
mechanisms.  Our approach to such a reduction is
%% Tiger:  replaced "the following" by "as follows".
%the following.
as follows.
\begin{enumerate}
\item \label{step:optimize}
Optimize, over all feasible profiles of interim allocation rules
$\intallocs = (\intalloc\agind[1],\ldots,\intalloc\agind[n]) \in \InAllocSpace$, the
%% Tiger:  replaced "total performance of the optimal outcome rules" by "sum of performances of the allocation rules"
%total performance of the optimal outcome
sum of performances of the allocation
rules $\sum_\agent \REV(\intalloc\agind)$.
\item \label{step:expostize} Implement the profile of interim outcome
  rules $\intoutcomes$ given by $\intoutcome\agind =
  \RULE(\intalloc\agind)$ with a feasible ex post outcome rule $\epoutcomes$.
\end{enumerate}

Two issues should be noted.  First, Step~\ref{step:expostize} requires
an argument that the existence of a feasible ex post outcome rule for
a given profile of interim allocation rules implies the existence of
one that combines the optimal interim outcome rules
from~$\RULE(\cdot)$.  We address this issue in Section~\ref{sec:alg}.
Second, Step~\ref{step:optimize} requires that we optimize over
jointly feasible interim allocation rules, and after solving for
$\intallocs$, its implementation by an ex post allocation rule is
needed to guide Step~\ref{step:expostize}.  We address this issue in
Section~\ref{sec:ssa}.  For single-unit (i.e., $k=1$) auctions a
characterization of the necessary and sufficient condition for interim
feasibility was provided by Kim Border.

%% Jason: i am not sure why the discussion below is in the
%% preliminaries section.

%% When $k$, the
%% number of agents we are allowed to serve, is~$1$, we give a general
%% structural theorem in Section~\ref{sec:int2expost} showing, by an
%% efficient algorithm, that every set of feasible interim allocation
%% rules can be implemented by a \emph{stochastic sequential allocation}
%% procedure. %azarakhsh changed the name
%% This allows us to describe all
%% feasibility constraints on interim allocation rules by polynomially
%% many linear inequalities, with some additional variables.  In
%% geometric terms, we identify the (generally exponentially-faceted)
%% polytope of all jointly feasible interim allocation rules as a
%% projection of a higher-dimensional polytope that has only polynomially
%% many facets. This gives, as a corollary, a generalization of Border's
%% inequality, which was known only for single-item auctions before this
%% work:

% For instance, single-item auction environments, i.e., $k=1$ and $\outcomespace = \{0,1\} \times \reals$,\footnote{In a single-item auction the outcome for an agent is the allocation and payment.} \citet{MR84} and \citet{M84} independently provided a condition necessary condition for feasibility.  The following theorem, due to \citet{B91} states that in these environments the same condition is also sufficient.

%\begin{reltheorem}[\citealp{B91}]
\begin{theorem}[\citealp{B91}]
\label{thm:border}%
In a single-item auction environment, interim allocation rules
$\intallocs$ are feasible (i.e., $\intallocs \in \InAllocSpace$) if
and only if the following holds:
\begin{align}
%% Tiger:  removed \type\agind\in
%\type\agind \in
\forall S\agind[1] \subseteq \typespace\agind[1],
\cdots, \forall S\agind[n] \subseteq \typespace\agind[n]: &\qquad
\sum_{\agent=1}^n \expect{\intalloc\agind(t_i) \given \type\agind \in
    S\agind}\cdot \prob{\type\agind \in S\agind}
        \le \prob[\types \sim \denss]{\exists i \in[n]: \type\agind \in S\agind
    } \tag{MRMB} \label{eq:MRM}
\end{align}
%\end{reltheorem}
\end{theorem}


%% \begin{example}[single-item, common budget]
%% Supply is $k = 1$, budget is $B \in \posreals$, type space is
%% $\typespace = \posreals$, service space is $\outcomespace =
%% \{0,1\}\times \emptyset \times \posreals$, and agent utility
%% $\util(\type,(\alloc,\emptyset,\price)$ is $\type \alloc - \price$ if $\price
%% \leq B$ and $-\infty$, otherwise.  The designer's objective is welfare
%% \citep{M00} or revenue \citep{LR96}.
%% \end{example}

%% \begin{example}[single-item, concave utility] Supply is $k = 1$, concave
%% utility function is $U : \posreals \to \posreals$, type space is
%% $\typespace = \posreals$, service space is $\outcomespace =
%% \{0,1\}\times \emptyset \times \posreals$, and agent utility
%% $\util(\type,(\alloc,\emptyset,\price)) = U(\type \alloc - \price)$.  The
%% designer's objective is revenue \citep{M84,MR84}.
%% \end{example}

%% \begin{example}[multi-unit, attribute auctions]
%% Supply is $k$, number of attributes is $m$, service space is
%% $\outcomespace = \{0,1\} \times \{1,\ldots,m\} \times \posreals$, type
%% space is $\typespace = \posreals^m$, and agent utility is quasi-linear
%% $\util(\type,(\alloc,\attrib,\price)) = \alloc (\type \cdot \attrib) -
%% \price$ (where $\type \cdot \attrib$ is the vector inner product). The
%% designer's objective is revenue.
%% \end{example}



\Xcomment{
\paragraph{Old Preliminary Section}



Consider a simple auction with $n$ participating agents for getting
%%  Tiger Nov 25: rephrased
%some type
one of $m$~services. In this
paper, we assume that we can serve at most $k$ of these agents. Each
agent $i \in [n]$ has a private type that defines his
%%  Tiger Nov 25: changed to plural forms to emphasize multi-dimensionality
values for the services.  The type of agent $i$ is denoted by $t_i \in T_i$,
%%  Tiger Nov 25: added the following explanation for valuation
and we use $v_i(t_i) = (v_i(t_i))_{j=1}^m \in \mathbb R^m$ to represent the vector of bidder~$i$'s valuations for the
$m$ services when his type is~$t_i$.
$T_i$ is
called the types pace of agent $i$.  Suppose that the agents' types
are drawn from independent but not
%%  Tiger Nov 25: added "necessarily"
necessarily identical distributions,
%%  Tiger Nov 25: added the following
publicly known as $\Dist = \Dist_1 \times \Dist_2 \times \cdots \times \Dist_n$.
$f_i(t_i)$ denotes the probability that agent $i$ has type $t_i$.
Furthermore, for every $S_i \subset T_i$, we define $f_i(S_i) =
\sum_{t_i \in S_i} f_i(t_i)$.
%%  Tiger Nov 25: deleted the following sentence on BNE
%A Bayes-Nash equilibrium (BNE) is a profile of strategies for mapping
%%  Tiger Nov 25: changed "values" to "types"
%values
%types to bid to obtain best response assuming all other agents are using the best response strategy and their values are coming from %%  Tiger Nov 25: added "the known"
%the known distribution.
A mechanism solicits bids from the agents and determines the outcome which consists
of an allocation $x = (x_1(t_1),\ldots, x_n(t_n))$ and payments
$p = (p_1(t_1), p_2(t_2), \ldots, p_n(t_n))$,
%%  Tiger Nov 25: added the following explanation
where $x_i(t_i) = (x_i(t_i))_{j=1}^m \in \mathbb [0, 1]^m$ represents the probabilities with which bidder~$i$ receives service~$j$.
%%  Tiger Nov 25: deleted the following sentence
%When agents are bidding in the mechanism, they know their own type and other agents' distributions.


%%  Tiger Nov 25: deleted the following sentences
%Interim allocation is the expected allocation assigned to an agent with a given type where the expectation is over all other agents' types.
%Suppose the profile of types $t=(t_1,\cdots, t_n)$ is distributed according to a publicly known distribution $\Dist$.
Given any allocation rule $x$
%%  Tiger Nov 25: added "and payment rule p"
and payment rule $p$, we can define a corresponding
\emph{interim allocation rule}
%%  Tiger Nov 25: changed [0,1] to [0, 1]^m; otherwise I don't see how v_i(t_i) \cdot x_i(t_i) is defined.
$\Inter{x}_i : T_i \to [0, 1]^m$
%%  Tiger Nov 25: added interim payment rule
\emph{interim allocation rule} $\Inter{p}_i: T_i \to \mathbb R_+$
for each agent $i$ that specifies the expected
%%  Tiger Nov 25: replaced the phrase
%probability of winning
allocation and payment
for each type of agent $i$, i.e., $\Inter{x}_i(t_i) = \E_{t_{-i}
\sim \Dist_{-i}}[x_i(t_i, t_{-i})]$,
%%  Tiger Nov 25: added definition of interim payment
and $\Inter{p}_i(t_i) = \E_{t_{-i} \sim \Dist_{-i}} [p_i(t_i, t_{-i})]$.
The allocation $x$ that is
assigned to a profile of types by the mechanism is also called
ex-post allocation.
%%  Tiger Nov 25: rewrote the definition of Bayesian truthfulness
% To design BNE mechanisms, the mechanism should satisfy the following constraint:
If for each agent~$i$ and all types $t_i, t_j \in T_i$, a mechanism satisfies
\[
v_i(t_i)\cdot \Inter{x}_i(t_i) - \Inter{p}_i(t_i) \ge v_i(t_i)\cdot
\Inter{x}_i(t_j) - \Inter{p}_i(t_j)
\]
then the mechanism is said to be \emph{Bayesian truthful}.
%%  Tiger Nov 25: changed BNE to Bayesian truthful
Note that the constraint for Bayesian truthfulness is in terms only of interim allocation
and payments.
%%  Tiger Nov 25: rewrote the following before the enumeration
%As a result, to design a BNE mechanism we can formulate by one constraint per type of an agent.  However, we should be able to map the interim allocation to feasible ex post allocations to maintain the interim allocation probability for all agents and also to satisfy the feasibility constraint for the available service(items).
%To be able to build our mechanism based on computed interim allocations, we need to find the answer to the following questions:
As a result, when designing a Bayesian truthful mechanism, one may check all incentive constraints by inspecting only
the interim allocation and payment rules for each bidder, instead of going through the exponentially large description
of ex-post rules.  However, not every interim allocation rule $(\Inter{x}_1, \Inter{x}_2, \ldots, \Inter{x}_n)$ can be implemented by a feasible ex-post allocation; even if it
can, it is still challenging to find the right ex-post allocation implementing it.  In this paper, we give efficient algorithmic
solutions to the following two problems, which enable us to design optimal mechanisms in terms of reduced forms
described by interim allocations, for domains which would otherwise seem intractable:

\begin{enumerate}
\item
Are $\Inter{x}_1, \cdots, \Inter{x}_n$ feasible?  In other words, is
there a corresponding feasible ex-post allocation rule~$x$ such that
%%  Tiger Nov 25: rid "the following holds?" and added "for all i, for all t_i$
%the following holds?
$\forall i \in [n], \forall t_i \in T_i$, $\Inter{x}_i(t_i) = \E_{t_{-i} \sim \Dist_{-i}}[x_i(t_i, t_{-i})]$?

\item
%%  Tiger Nov 25: added the following, and shorten the question
If the answer to~\label{prob:A} is affirmative,
how can such an $x$ be computed?
\end{enumerate}

%%  Tiger Nov 25: refined the description
%For the case of $k=1$,
For the domain where one service ($m=1$) is to be provided to at most one agent $(k=1)$,
Maskin and Riley~\cite{MR84} and
Matthews~\cite{M84} independently
%%  Tiger Nov 25: rewrote the sentence.  "came up with" sounds too informal
provided a condition for the first problem, which they proved to be necessary.
%came up with necessary conditions for \eqref{prob:A}.
The next theorem
%which is
due to Border~\cite{B91} states that the same condition is in fact both
necessary and sufficient for that domain.
%$k=1$.


\begin{theorem}[\citealp{B91}]
\label{thm:border}%
In an $n$-agent
%% Tiger:  replaced "single-item" by "single-service"
%single-item
single-service
auction environment, interim allocation
rules $(\Inter{x}_1, \cdots, \Inter{x}_n)$
%%  Tiger Nov 25: restructured the sentence
%, there exists a feasible ex-post allocation rule $x$ that implements the interim allocation rules
can be implemented by a feasible ex-post allocation rule~$\epallocs$
if and only if the following holds:
\begin{align}
    \forall S_1 \subseteq T_1, \cdots, \forall S_n \subseteq T_n: &\qquad \sum_{i=1}^n \sum_{t_i\in S_i} \Inter{x}_i(t_i) f_i(t_i)
        \le \Pr_{t \sim \denss}\left[\exists i \in[n]: t_i \in S_i
    \right] \tag{MRM} \label{eq:MRM}
\end{align}
\end{theorem}
}

%!TEX root = main.tex

\section{The Complexity of Dynamic Deadlock Prediction}
\seclabel{lower-bounds}

Detecting deadlock patterns and predictable deadlocks is clearly a problem in $\NP$,
as any witness for either problem can be verified in polynomial time.
However, little has been known about the hardness of the problem in terms of rigorous lower bounds.
Here we settle these questions, by proving strong intractability results.
Due to space constraints, we state and explain the main results here, and refer to\begin{pldi}~our technical report \cite{arxiv}\end{pldi}\begin{arxiv}~\cref{sec:sec:app_proofs_lower_bounds}\end{arxiv} for the full proofs.

\myparagraph{Parametrized hardness for detecting deadlock patterns}{
We show that the basic problem of checking the existence of a deadlock \emph{pattern} is itself
hard parameterized by the size $k$
of the pattern.
% \ucomment{Can be rewritten as: We first consider the problem of checking the existence of deadlock \emph{patterns}. Intuitively, a solution to the problem of deadlock prediction must resolve this basic problem first. Our first result shows that this problem is hard, parametrized by the size of the pattern we are looking for. In particular, the following result implies}

\begin{restatable}{theorem}{patternwonehardness}
\thmlabel{pattern-w1-hardness-pattern}
Checking if a trace $\tr$ contains a deadlock pattern of size $k$
is $\W{1}$-hard in the parameter $k$.
Moreover, the problem remains $\NP$-hard even when the lock-nesting depth of $\tr$ is constant.
\end{restatable}
}
%!TEX root = main.tex

\begin{proof}
	We show that there is a polynomial-time fixed parameter tractable reduction from
	INDEPENDENT-SET(c) to the problem of checking the existence of deadlock-patterns
	of size $c$.
	Our reduction takes as input an undirected graph $G$ and outputs a trace $\tr$
	such that $G$ has an independent set of size $c$ iff $\tr$
	has a deadlock pattern of size $c$.
	%Here present the construction and refer to \appref{proofs} for the proof of correctness.
	
	\begin{figure}[]
		\input{figures/construction-w1-hardness}
		\hfill ~~ %\unskip\ \vrule\ \hfill
		%!TEX root = ../main.tex
%\begin{figure}
\begin{subfigure}[b]{0.475\textwidth}
	\newcommand{\xdisposition}{0}
	\newcommand{\ydisposition}{0}
	\newcommand{\xtstep}{0.75}
	\newcommand{\ytstep}{1}
	\newcommand{\ybias}{-0.3 }
	\newcommand{\xstep}{2.5}
	\newcommand{\ystep}{-0.475}
	\newcommand{\xtscale}{0.8}
	\def \numevents{9.5}
	\newcommand{\eventA}[4]{
		\node[event, draw=black, fill=white] (A#1) at (#1*\xstep, #2*\ystep) {\footnotesize $#2(x_{#3})$};
		%\node[] at (0*\xstep-\xtstep, {#1*\ystep}) {\small $#2(x_{#3})$};
	}
	\scalebox{0.9}{
		\begin{tikzpicture}[thick,
			pre/.style={<-,shorten >= 2pt, shorten <=2pt, very thick},
			post/.style={->,shorten >= 3pt, shorten <=3pt,   thick},
			seqtrace/.style={line width=2},
			und/.style={very thick, draw=gray},
			event/.style={rectangle, minimum height=0.8mm, minimum width=15mm,  line width=1pt, inner sep=0.5,},
			virt/.style={circle,draw=black!50,fill=black!20, opacity=0}]
			\footnotesize
			
			
%			\node[] at (0.5*\xstep, -1.5*\ystep){
%				\small
%				$\cs(\lk_i,\lk_j)=\acq(\lk_i) \cdot\acq(\lk_j) \cdot\rel(\lk_j) \cdot \rel(\lk_i)$
%			};
			

				
			\begin{scope}[shift={(0,-3*\ystep)}]
				
				\draw[dashed] (-0.9*\xstep,6.5*\ystep) rectangle (-0.5*\xstep,8.5*\ystep);
				\node (A1) at (-0.7*\xstep,7*\ystep) {\normalsize [1, 1]};
				\node (A2) at (-0.7*\xstep,8*\ystep) {\normalsize [1, 0]};
				\node (A) at (-0.7*\xstep,9*\ystep) {\large  $A$};
				
				\draw[dashed] (1.5*\xstep,6.5*\ystep) rectangle (1.9*\xstep,8.5*\ystep);
				\node (B1) at (1.7*\xstep,7*\ystep) {\normalsize [1, 0]};
				\node (B2) at (1.7*\xstep,8*\ystep) {\normalsize [0, 1]};
				\node (B) at (1.7*\xstep,9*\ystep) {\large  $B$};
				
				
			\end{scope}
			

			
			\node[] (S11) at (0*\xstep,0.15) {\normalsize $t_A$};
			\node[] (S12) at (0*\xstep,\numevents * \ystep) {};
			\node[] (S21) at (1*\xstep,0.15) {\normalsize $t_B$};
			\node[] (S22) at (1*\xstep,\numevents * \ystep) {};
			
			\draw[seqtrace] (S11) to (S12);
			\draw[seqtrace] (S21) to (S22);
			
			
			\node[event, draw=black, fill=white] (11) at (0*\xstep, 1*\ystep + 0*\ybias) {$\acq(\LockColorTwo{\lk_2})$};
			\node[event, draw=black, fill=white] (12) at (0*\xstep, 2*\ystep + 0*\ybias) {$\acq(\LockColorOne{\lk_1})$};
			\node[event, draw=black, fill=white, dotted] (13) at (0*\xstep, 3*\ystep + 0*\ybias) {$\cs(m_0,m_1)$};
			%\node[event, draw=black, fill=white, dotted] (14) at (0*\xstep, 4*\ystep + 0*\ybias) {$\acq(\lk')$};
			%\node[event, draw=black, fill=white, dotted] (15) at (0*\xstep, 5*\ystep + 0*\ybias) {$\rel(\lk')$};
			%\node[event, draw=black, fill=white, dotted] (16) at (0*\xstep, 6*\ystep + 0*\ybias) {$\rel(\lk)$};
			\node[event, draw=black, fill=white] (17) at (0*\xstep, 4*\ystep + 0*\ybias) {$\rel(\LockColorOne{\lk_1})$};
			\node[event, draw=black, fill=white] (18) at (0*\xstep, 5*\ystep + 0*\ybias) {$\rel(\LockColorTwo{\lk_2})$};
			
			\node[event, draw=black, fill=white] (19) at (0*\xstep, 6*\ystep + 1*\ybias) {$\acq(\LockColorOne{\lk_1})$};
			\node[event, draw=black, fill=white, dotted] (110) at (0*\xstep, 7*\ystep + 1*\ybias) {$\cs(m_0,m_1)$};
			%\node[event, draw=black, fill=white, dotted] (111) at (0*\xstep, 11*\ystep + 1*\ybias) {$\acq(\lk')$};
			%\node[event, draw=black, fill=white, dotted] (112) at (0*\xstep, 12*\ystep + 1*\ybias) {$\rel(\lk')$};
			%\node[event, draw=black, fill=white, dotted] (113) at (0*\xstep, 13*\ystep + 1*\ybias) {$\rel(\lk)$};
			\node[event, draw=black, fill=white] (114) at (0*\xstep, 8*\ystep + 1*\ybias) {$\rel(\LockColorOne{\lk_1})$};
			
			%%%%%%%%%%%%%%%%
			
			\node[event, draw=black, fill=white] (21) at (1*\xstep, 2*\ystep + 0*\ybias) {$\acq(\LockColorOne{\lk_1})$};
			\node[event, draw=black, fill=white, dotted] (22) at (1*\xstep, 3*\ystep + 0*\ybias) {$\cs(m_1,m_0)$};
			%\node[event, draw=black, fill=white, dotted] (23) at (1*\xstep, 4*\ystep + 0*\ybias) {$\acq(\lk)$};
			%\node[event, draw=black, fill=white, dotted] (24) at (1*\xstep, 5*\ystep + 0*\ybias) {$\rel(\lk)$};
			%\node[event, draw=black, fill=white, dotted] (25) at (1*\xstep, 6*\ystep + 0*\ybias) {$\rel(\lk')$};
			\node[event, draw=black, fill=white] (26) at (1*\xstep, 4*\ystep + 0*\ybias) {$\rel(\LockColorOne{\lk_1})$};
			
			\node[event, draw=black, fill=white] (27) at (1*\xstep, 5*\ystep + 1*\ybias) {$\acq_1(\LockColorTwo{\lk_2})$};
			\node[event, draw=black, fill=white, dotted] (28) at (1*\xstep, 6*\ystep + 1*\ybias) {$\cs(m_1,m)$};
			%\node[event, draw=black, fill=white, dotted] (29) at (1*\xstep, 11*\ystep + 1*\ybias) {$\acq(\lk)$};
			%\node[event, draw=black, fill=white, dotted] (210) at (1*\xstep, 12*\ystep + 1*\ybias) {$\rel(\lk)$};
			%\node[event, draw=black, fill=white, dotted] (211) at (1*\xstep, 13*\ystep + 1*\ybias) {$\rel(\lk')$};
			\node[event, draw=black, fill=white] (212) at (1*\xstep, 7*\ystep + 1*\ybias) {$\rel_1(\LockColorTwo{\lk_2})$};
			
			\begin{scope}[]
				\node[below left=of 212] (212b) {};
			\end{scope}
	
		\end{tikzpicture}
	}
	\caption{
		Reduction for OV-hardness proof from an instance of size $n = 2$ and $d = 2$.
		% We use the shortcut $\cs(\lk_i,\lk_j)$ to denote two nested critical sections on $\lk_i$ and $\lk_j$, as indicated in the figure.
	}
	\figlabel{ov-hardness}
\end{subfigure}
%\end{figure}
		\caption{
		Construction of $\W{1}$-hardness (\subref{fig:w1-hardness}) and OV-hardness (\subref{fig:ov-hardness}) results. 
		We use the shortcut $\cs(\lk,\lk')$ to denote two nested critical sections on $\lk$ and $\lk'$. That is,
	$\cs(\lk,\lk')=\acq(\lk) \cdot\acq(\lk') \cdot\rel(\lk') \cdot \rel(\lk)$.
		} 
		%\label{fig:main}
	\end{figure}


	\myparagraph{Construction}{
		%Let $G = (V, E)$ be the input graph and let $c$ be the parameter.
		Let $V = \set{v_1, v_2, \ldots, v_n}$.
		We assume a total ordering $<_E$ on the set of edges $E$.
		The trace $\tr$ we construct is a concatenation of $c$ sub-traces: 
		$
		\tr = \tr^{(1)} \cdot \tr^{(2)} \cdots \tr^{(c)}
		$
		and uses $c$ threads $\set{t_1, t_2, \ldots t_c}$ and
		$|E| + c$ locks $\set{\lk_{\set{u, v}}}_{\set{u, v} \in E} \uplus \set{\lk_0, \lk_1 \ldots, \lk_{c-1}}$.
		The $i^\text{th}$ sub-trace $\tr^{(i)}$ is a sequence of events performed by thread $t_i$, and
		is obtained by concatenation of $n = |V|$ sub-traces:
		$
		\tr^{(i)} = \tr^{(i)}_1 \cdot \tr^{(i)}_2 \cdots \tr^{(i)}_n
		$.
		Each sub-trace $\tr^{(i)}_j$ with $(i \leq c, j \leq n)$ 
		comprises of nested critical sections over locks of the 
		form $\lk_{\set{v_j, u}}$, where $u$ is a neighbor of $v_j$.
		Inside the nested block we have critical 
		sections on locks $\lk_{i \% c}$ and $\lk_{(i+1) \% c}$.
		Formally, let $\set{v_j, v_{k_1}}, \ldots, \set{v_j, v_{k_d}}$
		be the neighboring edges of $v_j$ (ordered according to $<_E$).
		Then, $\tr^{(i)}_j$ is the unique string generated by the
		grammar having $d+1$ non-terminals $S_0, S_1, \ldots, S_d$, start symbol $S_d$
		and the following production rules:
		\begin{itemize}
			\item $S_0 \to \ev{t_i, \acq(\lk_{i \% c})} \cdot \ev{t_i, \acq(\lk_{(i+1) \% c})} \cdot \ev{t_i, \rel(\lk_{(i+1) \% c})} \cdot \ev{t_i, \rel(\lk_{i \% c})}$.
			\item for each $1 \leq r \leq d$, $S_r \to \ev{t_i, \acq(\lk_{\set{v_j, v_{k_r}}})} \cdot S_{r-1} \cdot \ev{t_i, \rel(\lk_{\set{v_j, v_{k_r}}})}$.
		\end{itemize}
		\figref{w1-hardness} illustrates this construction for a graph with $3$ nodes
		and parameter $c = 3$.
		Finally, observe that the lock-nesting depth in $\tr$ is bounded by 2 + the degree of $G$.
	}
\end{proof}








\thmref{pattern-w1-hardness-pattern} implies that the problem is not only $\NP$-hard, 
but also unlikely to be \emph{fixed parameter tractable} in the size $k$ of the deadlock pattern.
In fact, under the well-believed Exponential Time Hypothesis (ETH), the parametrized
problem INDEPENDENT-SET(c) cannot be solved in time $f(c) \cdot n^{o(c)}$~\cite{chen2006strong}.
The above reduction preserves the parameter $k=c$, thus under ETH,
detecting deadlock patterns of size $k$ is unlikely to be solvable 
in time complexity $f(k) \cdot \NumEvents^{g(k)}$, where 
$g(k)$ is $o(k)$ (such as $g(k) = \sqrt{k}$ or even $g(k) = k/\log(k)$).
% It is known that INDEPENDENT-SET(c) cannot
% be solved in f (c) · n
% o(c)
% time under ETH [11]. As our reduction
% to rf-poset realizability and dynamic data-race prediction uses
% k = O(c) threads, each of these problems does not have a
% f (k) · n
% o(k)
% -time algorithm under ETH.
% In turn this means that it is unlikely to obtain an algorithm for detecting deadlock patterns that runs in time $O(\NumEvents^{c}\cdot f(k))$, where $c$ is a constant and $f$ is \emph{any} function independent of $n$.
% This is important as $k$ is typically constant; 
% $\W{1}$-hardness indicates that we can only expect algorithms with running time $O(\NumEvents^{g(k)})$, i.e., the degree of the polynomial depends on the size of the deadlock patterns which can be as high as the number of threads.
The problem of checking the existence of deadlock patterns 
is, intuitively, a precursor to the deadlock prediction problem. 
Thus, an approach for 
deadlock prediction that first identifies the existence of arbitrary deadlock patterns and 
then verifying their feasibility is unlikely to be tractable. 
In practice, the synchronization patterns corresponding to the hard instances are uncommon in executions, and our proposed algorithms (\secref{syncp} and \secref{otf}) can effectively expose predictable deadlocks (\secref{experiments}).


\myparagraph{Fine-grained hardness for deadlock pattern detection}{
\begin{comment}
	We now establish a fine-grained hardness for detecting
	deadlock patterns --- for each $k \geq 2$, we cannot check for the existence of 
	patterns of size $k$ in time $O(\NumEvents^{k-\epsilon})$ for any $\epsilon > 0$,
	under the popular Orthogonal Vectors hypothesis (OV).
	%This result is based on a reduction from the popular Orthogonal Vectors (OV) problem.
	For a fixed $k \geq 2$, the $k$-OV problem takes
	$k$ sets of $d$-dimensional vectors
	$A_1, A_2, \ldots, A_k \subseteq \set{0, 1}^d$, each of cardinality $|A_i|= n$
	($1 \leq i \leq k$) as input, 
	and asks if
	there are vectors $a_1 \in A_1, \ldots, a_k \in A_k$
	such that the extended dot product $a_1 \cdot a_2 \cdots a_k = \sum_{p =1}^d (a_1[p]\cdot a_2[p] \cdots a_k[p])= 0$.
	For a $k\geq 2$, the $k$-OV hypothesis states that for any 
	$\epsilon >0$, there is no $O(n^{k-\epsilon}\cdot \poly{d})$ algorithm for $k$-OV.
	The Strong Exponential Time Hypothesis (SETH) implies $k$-OV ~\cite{williams2005}.
	Our next theorem is based on showing that, for every $k \geq 2$, 
	detecting deadlock patterns of size $k$ is at least as hard as solving $k$-OV.
 Note the difference between \thmref{pattern-w1-hardness-pattern} and \thmref{pattern-ov-hardness}:~the former allows algorithms with running time of the form $\NumEvents^{k/2}$ (even under ETH), but the latter excludes them, requiring that $k$ fully appears in the exponent.
 The two results are based on different hypotheses and also incomparable since 
 it could turn out that $k$-OV is false but ETH is true.
 We thus establish both results, in order to develop a deeper understanding of the intricacies of the problem.
 \end{comment}
% %!TEX root = main.tex

\begin{proof}
	We show a fine-grained reduction from the Orthogonal Vectors Problem to
	the problem of checking for deadlock patterns of size $k$.
	For this, we start with two sets
	$A_1, A_2, \ldots, A_k \subseteq \set{0, 1}^d$
	of $d$-dimensional vectors with $|A_i| = n$ for every $1 \leq i \leq k$.
	We write the $j^{th}$ vector in $A_i$ as $A_{i, j}$.
	
	%%!TEX root = ../main.tex
%\begin{figure}
\begin{subfigure}[b]{0.475\textwidth}
	\newcommand{\xdisposition}{0}
	\newcommand{\ydisposition}{0}
	\newcommand{\xtstep}{0.75}
	\newcommand{\ytstep}{1}
	\newcommand{\ybias}{-0.3 }
	\newcommand{\xstep}{2.5}
	\newcommand{\ystep}{-0.475}
	\newcommand{\xtscale}{0.8}
	\def \numevents{9.5}
	\newcommand{\eventA}[4]{
		\node[event, draw=black, fill=white] (A#1) at (#1*\xstep, #2*\ystep) {\footnotesize $#2(x_{#3})$};
		%\node[] at (0*\xstep-\xtstep, {#1*\ystep}) {\small $#2(x_{#3})$};
	}
	\scalebox{0.9}{
		\begin{tikzpicture}[thick,
			pre/.style={<-,shorten >= 2pt, shorten <=2pt, very thick},
			post/.style={->,shorten >= 3pt, shorten <=3pt,   thick},
			seqtrace/.style={line width=2},
			und/.style={very thick, draw=gray},
			event/.style={rectangle, minimum height=0.8mm, minimum width=15mm,  line width=1pt, inner sep=0.5,},
			virt/.style={circle,draw=black!50,fill=black!20, opacity=0}]
			\footnotesize
			
			
%			\node[] at (0.5*\xstep, -1.5*\ystep){
%				\small
%				$\cs(\lk_i,\lk_j)=\acq(\lk_i) \cdot\acq(\lk_j) \cdot\rel(\lk_j) \cdot \rel(\lk_i)$
%			};
			

				
			\begin{scope}[shift={(0,-3*\ystep)}]
				
				\draw[dashed] (-0.9*\xstep,6.5*\ystep) rectangle (-0.5*\xstep,8.5*\ystep);
				\node (A1) at (-0.7*\xstep,7*\ystep) {\normalsize [1, 1]};
				\node (A2) at (-0.7*\xstep,8*\ystep) {\normalsize [1, 0]};
				\node (A) at (-0.7*\xstep,9*\ystep) {\large  $A$};
				
				\draw[dashed] (1.5*\xstep,6.5*\ystep) rectangle (1.9*\xstep,8.5*\ystep);
				\node (B1) at (1.7*\xstep,7*\ystep) {\normalsize [1, 0]};
				\node (B2) at (1.7*\xstep,8*\ystep) {\normalsize [0, 1]};
				\node (B) at (1.7*\xstep,9*\ystep) {\large  $B$};
				
				
			\end{scope}
			

			
			\node[] (S11) at (0*\xstep,0.15) {\normalsize $t_A$};
			\node[] (S12) at (0*\xstep,\numevents * \ystep) {};
			\node[] (S21) at (1*\xstep,0.15) {\normalsize $t_B$};
			\node[] (S22) at (1*\xstep,\numevents * \ystep) {};
			
			\draw[seqtrace] (S11) to (S12);
			\draw[seqtrace] (S21) to (S22);
			
			
			\node[event, draw=black, fill=white] (11) at (0*\xstep, 1*\ystep + 0*\ybias) {$\acq(\LockColorTwo{\lk_2})$};
			\node[event, draw=black, fill=white] (12) at (0*\xstep, 2*\ystep + 0*\ybias) {$\acq(\LockColorOne{\lk_1})$};
			\node[event, draw=black, fill=white, dotted] (13) at (0*\xstep, 3*\ystep + 0*\ybias) {$\cs(m_0,m_1)$};
			%\node[event, draw=black, fill=white, dotted] (14) at (0*\xstep, 4*\ystep + 0*\ybias) {$\acq(\lk')$};
			%\node[event, draw=black, fill=white, dotted] (15) at (0*\xstep, 5*\ystep + 0*\ybias) {$\rel(\lk')$};
			%\node[event, draw=black, fill=white, dotted] (16) at (0*\xstep, 6*\ystep + 0*\ybias) {$\rel(\lk)$};
			\node[event, draw=black, fill=white] (17) at (0*\xstep, 4*\ystep + 0*\ybias) {$\rel(\LockColorOne{\lk_1})$};
			\node[event, draw=black, fill=white] (18) at (0*\xstep, 5*\ystep + 0*\ybias) {$\rel(\LockColorTwo{\lk_2})$};
			
			\node[event, draw=black, fill=white] (19) at (0*\xstep, 6*\ystep + 1*\ybias) {$\acq(\LockColorOne{\lk_1})$};
			\node[event, draw=black, fill=white, dotted] (110) at (0*\xstep, 7*\ystep + 1*\ybias) {$\cs(m_0,m_1)$};
			%\node[event, draw=black, fill=white, dotted] (111) at (0*\xstep, 11*\ystep + 1*\ybias) {$\acq(\lk')$};
			%\node[event, draw=black, fill=white, dotted] (112) at (0*\xstep, 12*\ystep + 1*\ybias) {$\rel(\lk')$};
			%\node[event, draw=black, fill=white, dotted] (113) at (0*\xstep, 13*\ystep + 1*\ybias) {$\rel(\lk)$};
			\node[event, draw=black, fill=white] (114) at (0*\xstep, 8*\ystep + 1*\ybias) {$\rel(\LockColorOne{\lk_1})$};
			
			%%%%%%%%%%%%%%%%
			
			\node[event, draw=black, fill=white] (21) at (1*\xstep, 2*\ystep + 0*\ybias) {$\acq(\LockColorOne{\lk_1})$};
			\node[event, draw=black, fill=white, dotted] (22) at (1*\xstep, 3*\ystep + 0*\ybias) {$\cs(m_1,m_0)$};
			%\node[event, draw=black, fill=white, dotted] (23) at (1*\xstep, 4*\ystep + 0*\ybias) {$\acq(\lk)$};
			%\node[event, draw=black, fill=white, dotted] (24) at (1*\xstep, 5*\ystep + 0*\ybias) {$\rel(\lk)$};
			%\node[event, draw=black, fill=white, dotted] (25) at (1*\xstep, 6*\ystep + 0*\ybias) {$\rel(\lk')$};
			\node[event, draw=black, fill=white] (26) at (1*\xstep, 4*\ystep + 0*\ybias) {$\rel(\LockColorOne{\lk_1})$};
			
			\node[event, draw=black, fill=white] (27) at (1*\xstep, 5*\ystep + 1*\ybias) {$\acq_1(\LockColorTwo{\lk_2})$};
			\node[event, draw=black, fill=white, dotted] (28) at (1*\xstep, 6*\ystep + 1*\ybias) {$\cs(m_1,m)$};
			%\node[event, draw=black, fill=white, dotted] (29) at (1*\xstep, 11*\ystep + 1*\ybias) {$\acq(\lk)$};
			%\node[event, draw=black, fill=white, dotted] (210) at (1*\xstep, 12*\ystep + 1*\ybias) {$\rel(\lk)$};
			%\node[event, draw=black, fill=white, dotted] (211) at (1*\xstep, 13*\ystep + 1*\ybias) {$\rel(\lk')$};
			\node[event, draw=black, fill=white] (212) at (1*\xstep, 7*\ystep + 1*\ybias) {$\rel_1(\LockColorTwo{\lk_2})$};
			
			\begin{scope}[]
				\node[below left=of 212] (212b) {};
			\end{scope}
	
		\end{tikzpicture}
	}
	\caption{
		Reduction for OV-hardness proof from an instance of size $n = 2$ and $d = 2$.
		% We use the shortcut $\cs(\lk_i,\lk_j)$ to denote two nested critical sections on $\lk_i$ and $\lk_j$, as indicated in the figure.
	}
	\figlabel{ov-hardness}
\end{subfigure}
%\end{figure}
	
	\myparagraph{Construction}{
		We will construct a trace $\tr$ such that $\tr$ has a deadlock
		pattern of length $k$ iff $(A_1, A_2, \ldots, A_k)$ is a positive $k$-OV instance.
		The trace $\tr$ is of the form $\tr = \tr^{A_1} \cdot \tr^{A_2} \cdots \tr^{A_k}$
		and uses $k$ threads $\set{t_{A_1}, \ldots, t_{A_k}}$ and $d+k$ distinct 
		locks $\lk_1, \ldots, \lk_d, m_1, m_2, \ldots, m_k$. 
		Intuitively, the sub-trace $\tr^{A_i}$ encodes the given set of vectors $A_i$.
		The sub-traces $\tr^{A_i} = \tr^{A_i}_1 \cdot \tr^{A_i}_2 \cdots \tr^{A_i}_n$
		% and $\tr^B = \tr^B_1 \cdot \tr^B_2 \cdots \tr^B_n$ 
		are defined as follows.
		For each $j \in \set{1, 2, \ldots, n}$
		the sub-trace $\tr^{A_i}_j$ is the unique string generated by the
		grammar having $d+1$ non-terminals $S_0, S_1, \ldots, S_d$, start symbol $S_d$
		and the following production rules:
		\begin{itemize}
			\item $S_0 \to \ev{t_Z, \acq(m_i)} \cdot \ev{t_Z, \acq(m_{i \% k + 1})} \cdot \ev{t_Z, \rel(m_{i \% k + 1})} \cdot \ev{t_Z, \rel(m_i)}$.
			% where $(m, m') = (\lk, \lk')$ if $Z = A$, and $(m, m') = (\lk', \lk)$.
			\item for each $1 \leq p \leq d$, $S_p \to S_{p-1}$ if $A_{i, j}[p] = 0$.
			Otherwise (if $A_{i, j}[p] =1$), $S_p \to \ev{t_{A_i}, \acq(\lk_p)} \cdot S_{p-1} \cdot \ev{t_{A_i}, \rel(\lk_p)}$.
		\end{itemize}
		In words, all events of $\tr^{A_i}$ are performed by thread $t_{A_i}$.
		 % and those in $\tr^B$ are performed by $t_B$.
		Next, the $j^{th}$ sub-trace of $\tr^{A_i}$, denoted $\tr^{A_i}_j$ corresponds to the vector $A_{i, j}$
		as follows --- $\tr^{A_i}_j$ is a nested block of critical sections,
		with the innermost critical section being on lock $m_{i \% k + 1}$,
		which is immediately enclosed in a  critical section on lock $m_i$.
		Further, in the sub-trace $\tr^{A_i}_j$, 
		the lock $\lk_p$ occurs iff $A_{i, j}[p] = 1$.
		% The sub-traces $\tr^B_i$ is similarly constructed, except that the order
		% of the two innermost critical sections is inverted.
		\figref{ov-hardness} illustrates the construction for an OV-instance with $k=2$ $n=2$ and $d=2$.
	}
\end{proof}

% \myparagraph{The complexity of size-$2$ patterns}{
We next consider the problem of
detecting deadlock patterns of size $2$, as these form the most common case in practice~\cite{Lu08}.
Observe that \thmref{pattern-w1-hardness-pattern} has no implications on this case, as here $k$ is fixed.
The problem admits a folklore $O(\NumEvents^2)$ time algorithm, by iterating over all pairs of lock-acquisition events of the input trace, and checking whether any such pair forms a deadlock pattern.
Perhaps surprisingly, here we show that, despite its simplicity, this algorithm is optimal, i.e., we cannot hope to improve over this quadratic bound.
This result is based on a reduction from the popular Orthogonal Vectors (OV) problem.
Given two sets of $d$-dimensional vectors
$A, B \subseteq \set{0, 1}^d$
of cardinality $|A| = |B| = n$, the OV problem asks if there are $a \in A, b \in B$
such that $a \cdot b = \sum_i a[i]\cdot b[i]= 0$.
The OV hypothesis states that for any 
$\epsilon >0$, there is no $O(n^{2-\epsilon}\cdot \poly{d})$ algorithm for solving OV.
This is also a consequence of the famous Strong Exponential Time Hypothesis (SETH)~\cite{williams2005}.
We next show that detecting deadlock patterns of size $2$ is at least as hard as solving OV.
% }

\begin{restatable}{theorem}{patternovhardness}
\thmlabel{pattern-ov-hardness}
Given a trace $\tr$ of size $\NumEvents$ and $\NumLocks$ locks, 
for any $\epsilon>0$,
there is no algorithm that determines in $O(\NumEvents^{2-\epsilon}\cdot \poly{\NumLocks})$ time
whether $\tr$ has a deadlock pattern of size $2$, under the OV hypothesis.
\end{restatable}

% %!TEX root = main.tex

\begin{proof}
	We show a fine-grained reduction from the Orthogonal Vectors Problem to
	the problem of checking for deadlock patterns of size $k$.
	For this, we start with two sets
	$A_1, A_2, \ldots, A_k \subseteq \set{0, 1}^d$
	of $d$-dimensional vectors with $|A_i| = n$ for every $1 \leq i \leq k$.
	We write the $j^{th}$ vector in $A_i$ as $A_{i, j}$.
	
	%%!TEX root = ../main.tex
%\begin{figure}
\begin{subfigure}[b]{0.475\textwidth}
	\newcommand{\xdisposition}{0}
	\newcommand{\ydisposition}{0}
	\newcommand{\xtstep}{0.75}
	\newcommand{\ytstep}{1}
	\newcommand{\ybias}{-0.3 }
	\newcommand{\xstep}{2.5}
	\newcommand{\ystep}{-0.475}
	\newcommand{\xtscale}{0.8}
	\def \numevents{9.5}
	\newcommand{\eventA}[4]{
		\node[event, draw=black, fill=white] (A#1) at (#1*\xstep, #2*\ystep) {\footnotesize $#2(x_{#3})$};
		%\node[] at (0*\xstep-\xtstep, {#1*\ystep}) {\small $#2(x_{#3})$};
	}
	\scalebox{0.9}{
		\begin{tikzpicture}[thick,
			pre/.style={<-,shorten >= 2pt, shorten <=2pt, very thick},
			post/.style={->,shorten >= 3pt, shorten <=3pt,   thick},
			seqtrace/.style={line width=2},
			und/.style={very thick, draw=gray},
			event/.style={rectangle, minimum height=0.8mm, minimum width=15mm,  line width=1pt, inner sep=0.5,},
			virt/.style={circle,draw=black!50,fill=black!20, opacity=0}]
			\footnotesize
			
			
%			\node[] at (0.5*\xstep, -1.5*\ystep){
%				\small
%				$\cs(\lk_i,\lk_j)=\acq(\lk_i) \cdot\acq(\lk_j) \cdot\rel(\lk_j) \cdot \rel(\lk_i)$
%			};
			

				
			\begin{scope}[shift={(0,-3*\ystep)}]
				
				\draw[dashed] (-0.9*\xstep,6.5*\ystep) rectangle (-0.5*\xstep,8.5*\ystep);
				\node (A1) at (-0.7*\xstep,7*\ystep) {\normalsize [1, 1]};
				\node (A2) at (-0.7*\xstep,8*\ystep) {\normalsize [1, 0]};
				\node (A) at (-0.7*\xstep,9*\ystep) {\large  $A$};
				
				\draw[dashed] (1.5*\xstep,6.5*\ystep) rectangle (1.9*\xstep,8.5*\ystep);
				\node (B1) at (1.7*\xstep,7*\ystep) {\normalsize [1, 0]};
				\node (B2) at (1.7*\xstep,8*\ystep) {\normalsize [0, 1]};
				\node (B) at (1.7*\xstep,9*\ystep) {\large  $B$};
				
				
			\end{scope}
			

			
			\node[] (S11) at (0*\xstep,0.15) {\normalsize $t_A$};
			\node[] (S12) at (0*\xstep,\numevents * \ystep) {};
			\node[] (S21) at (1*\xstep,0.15) {\normalsize $t_B$};
			\node[] (S22) at (1*\xstep,\numevents * \ystep) {};
			
			\draw[seqtrace] (S11) to (S12);
			\draw[seqtrace] (S21) to (S22);
			
			
			\node[event, draw=black, fill=white] (11) at (0*\xstep, 1*\ystep + 0*\ybias) {$\acq(\LockColorTwo{\lk_2})$};
			\node[event, draw=black, fill=white] (12) at (0*\xstep, 2*\ystep + 0*\ybias) {$\acq(\LockColorOne{\lk_1})$};
			\node[event, draw=black, fill=white, dotted] (13) at (0*\xstep, 3*\ystep + 0*\ybias) {$\cs(m_0,m_1)$};
			%\node[event, draw=black, fill=white, dotted] (14) at (0*\xstep, 4*\ystep + 0*\ybias) {$\acq(\lk')$};
			%\node[event, draw=black, fill=white, dotted] (15) at (0*\xstep, 5*\ystep + 0*\ybias) {$\rel(\lk')$};
			%\node[event, draw=black, fill=white, dotted] (16) at (0*\xstep, 6*\ystep + 0*\ybias) {$\rel(\lk)$};
			\node[event, draw=black, fill=white] (17) at (0*\xstep, 4*\ystep + 0*\ybias) {$\rel(\LockColorOne{\lk_1})$};
			\node[event, draw=black, fill=white] (18) at (0*\xstep, 5*\ystep + 0*\ybias) {$\rel(\LockColorTwo{\lk_2})$};
			
			\node[event, draw=black, fill=white] (19) at (0*\xstep, 6*\ystep + 1*\ybias) {$\acq(\LockColorOne{\lk_1})$};
			\node[event, draw=black, fill=white, dotted] (110) at (0*\xstep, 7*\ystep + 1*\ybias) {$\cs(m_0,m_1)$};
			%\node[event, draw=black, fill=white, dotted] (111) at (0*\xstep, 11*\ystep + 1*\ybias) {$\acq(\lk')$};
			%\node[event, draw=black, fill=white, dotted] (112) at (0*\xstep, 12*\ystep + 1*\ybias) {$\rel(\lk')$};
			%\node[event, draw=black, fill=white, dotted] (113) at (0*\xstep, 13*\ystep + 1*\ybias) {$\rel(\lk)$};
			\node[event, draw=black, fill=white] (114) at (0*\xstep, 8*\ystep + 1*\ybias) {$\rel(\LockColorOne{\lk_1})$};
			
			%%%%%%%%%%%%%%%%
			
			\node[event, draw=black, fill=white] (21) at (1*\xstep, 2*\ystep + 0*\ybias) {$\acq(\LockColorOne{\lk_1})$};
			\node[event, draw=black, fill=white, dotted] (22) at (1*\xstep, 3*\ystep + 0*\ybias) {$\cs(m_1,m_0)$};
			%\node[event, draw=black, fill=white, dotted] (23) at (1*\xstep, 4*\ystep + 0*\ybias) {$\acq(\lk)$};
			%\node[event, draw=black, fill=white, dotted] (24) at (1*\xstep, 5*\ystep + 0*\ybias) {$\rel(\lk)$};
			%\node[event, draw=black, fill=white, dotted] (25) at (1*\xstep, 6*\ystep + 0*\ybias) {$\rel(\lk')$};
			\node[event, draw=black, fill=white] (26) at (1*\xstep, 4*\ystep + 0*\ybias) {$\rel(\LockColorOne{\lk_1})$};
			
			\node[event, draw=black, fill=white] (27) at (1*\xstep, 5*\ystep + 1*\ybias) {$\acq_1(\LockColorTwo{\lk_2})$};
			\node[event, draw=black, fill=white, dotted] (28) at (1*\xstep, 6*\ystep + 1*\ybias) {$\cs(m_1,m)$};
			%\node[event, draw=black, fill=white, dotted] (29) at (1*\xstep, 11*\ystep + 1*\ybias) {$\acq(\lk)$};
			%\node[event, draw=black, fill=white, dotted] (210) at (1*\xstep, 12*\ystep + 1*\ybias) {$\rel(\lk)$};
			%\node[event, draw=black, fill=white, dotted] (211) at (1*\xstep, 13*\ystep + 1*\ybias) {$\rel(\lk')$};
			\node[event, draw=black, fill=white] (212) at (1*\xstep, 7*\ystep + 1*\ybias) {$\rel_1(\LockColorTwo{\lk_2})$};
			
			\begin{scope}[]
				\node[below left=of 212] (212b) {};
			\end{scope}
	
		\end{tikzpicture}
	}
	\caption{
		Reduction for OV-hardness proof from an instance of size $n = 2$ and $d = 2$.
		% We use the shortcut $\cs(\lk_i,\lk_j)$ to denote two nested critical sections on $\lk_i$ and $\lk_j$, as indicated in the figure.
	}
	\figlabel{ov-hardness}
\end{subfigure}
%\end{figure}
	
	\myparagraph{Construction}{
		We will construct a trace $\tr$ such that $\tr$ has a deadlock
		pattern of length $k$ iff $(A_1, A_2, \ldots, A_k)$ is a positive $k$-OV instance.
		The trace $\tr$ is of the form $\tr = \tr^{A_1} \cdot \tr^{A_2} \cdots \tr^{A_k}$
		and uses $k$ threads $\set{t_{A_1}, \ldots, t_{A_k}}$ and $d+k$ distinct 
		locks $\lk_1, \ldots, \lk_d, m_1, m_2, \ldots, m_k$. 
		Intuitively, the sub-trace $\tr^{A_i}$ encodes the given set of vectors $A_i$.
		The sub-traces $\tr^{A_i} = \tr^{A_i}_1 \cdot \tr^{A_i}_2 \cdots \tr^{A_i}_n$
		% and $\tr^B = \tr^B_1 \cdot \tr^B_2 \cdots \tr^B_n$ 
		are defined as follows.
		For each $j \in \set{1, 2, \ldots, n}$
		the sub-trace $\tr^{A_i}_j$ is the unique string generated by the
		grammar having $d+1$ non-terminals $S_0, S_1, \ldots, S_d$, start symbol $S_d$
		and the following production rules:
		\begin{itemize}
			\item $S_0 \to \ev{t_Z, \acq(m_i)} \cdot \ev{t_Z, \acq(m_{i \% k + 1})} \cdot \ev{t_Z, \rel(m_{i \% k + 1})} \cdot \ev{t_Z, \rel(m_i)}$.
			% where $(m, m') = (\lk, \lk')$ if $Z = A$, and $(m, m') = (\lk', \lk)$.
			\item for each $1 \leq p \leq d$, $S_p \to S_{p-1}$ if $A_{i, j}[p] = 0$.
			Otherwise (if $A_{i, j}[p] =1$), $S_p \to \ev{t_{A_i}, \acq(\lk_p)} \cdot S_{p-1} \cdot \ev{t_{A_i}, \rel(\lk_p)}$.
		\end{itemize}
		In words, all events of $\tr^{A_i}$ are performed by thread $t_{A_i}$.
		 % and those in $\tr^B$ are performed by $t_B$.
		Next, the $j^{th}$ sub-trace of $\tr^{A_i}$, denoted $\tr^{A_i}_j$ corresponds to the vector $A_{i, j}$
		as follows --- $\tr^{A_i}_j$ is a nested block of critical sections,
		with the innermost critical section being on lock $m_{i \% k + 1}$,
		which is immediately enclosed in a  critical section on lock $m_i$.
		Further, in the sub-trace $\tr^{A_i}_j$, 
		the lock $\lk_p$ occurs iff $A_{i, j}[p] = 1$.
		% The sub-traces $\tr^B_i$ is similarly constructed, except that the order
		% of the two innermost critical sections is inverted.
		\figref{ov-hardness} illustrates the construction for an OV-instance with $k=2$ $n=2$ and $d=2$.
	}
\end{proof}

%!TEX root = main.tex


\begin{proof}
	We show a fine-grained reduction from the Orthogonal Vectors Problem to
	the problem of checking for deadlock patterns of size $2$.
	For this, we start with two sets
	$A, B \subseteq \set{0, 1}^d$
	of $d$-dimensional vectors with $|A| = |B| = n$.
	We write the $i^{th}$ vector in $A$ as $A_i$ and that in $B$ as $B_i$.
	
	%%!TEX root = ../main.tex
%\begin{figure}
\begin{subfigure}[b]{0.475\textwidth}
	\newcommand{\xdisposition}{0}
	\newcommand{\ydisposition}{0}
	\newcommand{\xtstep}{0.75}
	\newcommand{\ytstep}{1}
	\newcommand{\ybias}{-0.3 }
	\newcommand{\xstep}{2.5}
	\newcommand{\ystep}{-0.475}
	\newcommand{\xtscale}{0.8}
	\def \numevents{9.5}
	\newcommand{\eventA}[4]{
		\node[event, draw=black, fill=white] (A#1) at (#1*\xstep, #2*\ystep) {\footnotesize $#2(x_{#3})$};
		%\node[] at (0*\xstep-\xtstep, {#1*\ystep}) {\small $#2(x_{#3})$};
	}
	\scalebox{0.9}{
		\begin{tikzpicture}[thick,
			pre/.style={<-,shorten >= 2pt, shorten <=2pt, very thick},
			post/.style={->,shorten >= 3pt, shorten <=3pt,   thick},
			seqtrace/.style={line width=2},
			und/.style={very thick, draw=gray},
			event/.style={rectangle, minimum height=0.8mm, minimum width=15mm,  line width=1pt, inner sep=0.5,},
			virt/.style={circle,draw=black!50,fill=black!20, opacity=0}]
			\footnotesize
			
			
%			\node[] at (0.5*\xstep, -1.5*\ystep){
%				\small
%				$\cs(\lk_i,\lk_j)=\acq(\lk_i) \cdot\acq(\lk_j) \cdot\rel(\lk_j) \cdot \rel(\lk_i)$
%			};
			

				
			\begin{scope}[shift={(0,-3*\ystep)}]
				
				\draw[dashed] (-0.9*\xstep,6.5*\ystep) rectangle (-0.5*\xstep,8.5*\ystep);
				\node (A1) at (-0.7*\xstep,7*\ystep) {\normalsize [1, 1]};
				\node (A2) at (-0.7*\xstep,8*\ystep) {\normalsize [1, 0]};
				\node (A) at (-0.7*\xstep,9*\ystep) {\large  $A$};
				
				\draw[dashed] (1.5*\xstep,6.5*\ystep) rectangle (1.9*\xstep,8.5*\ystep);
				\node (B1) at (1.7*\xstep,7*\ystep) {\normalsize [1, 0]};
				\node (B2) at (1.7*\xstep,8*\ystep) {\normalsize [0, 1]};
				\node (B) at (1.7*\xstep,9*\ystep) {\large  $B$};
				
				
			\end{scope}
			

			
			\node[] (S11) at (0*\xstep,0.15) {\normalsize $t_A$};
			\node[] (S12) at (0*\xstep,\numevents * \ystep) {};
			\node[] (S21) at (1*\xstep,0.15) {\normalsize $t_B$};
			\node[] (S22) at (1*\xstep,\numevents * \ystep) {};
			
			\draw[seqtrace] (S11) to (S12);
			\draw[seqtrace] (S21) to (S22);
			
			
			\node[event, draw=black, fill=white] (11) at (0*\xstep, 1*\ystep + 0*\ybias) {$\acq(\LockColorTwo{\lk_2})$};
			\node[event, draw=black, fill=white] (12) at (0*\xstep, 2*\ystep + 0*\ybias) {$\acq(\LockColorOne{\lk_1})$};
			\node[event, draw=black, fill=white, dotted] (13) at (0*\xstep, 3*\ystep + 0*\ybias) {$\cs(m_0,m_1)$};
			%\node[event, draw=black, fill=white, dotted] (14) at (0*\xstep, 4*\ystep + 0*\ybias) {$\acq(\lk')$};
			%\node[event, draw=black, fill=white, dotted] (15) at (0*\xstep, 5*\ystep + 0*\ybias) {$\rel(\lk')$};
			%\node[event, draw=black, fill=white, dotted] (16) at (0*\xstep, 6*\ystep + 0*\ybias) {$\rel(\lk)$};
			\node[event, draw=black, fill=white] (17) at (0*\xstep, 4*\ystep + 0*\ybias) {$\rel(\LockColorOne{\lk_1})$};
			\node[event, draw=black, fill=white] (18) at (0*\xstep, 5*\ystep + 0*\ybias) {$\rel(\LockColorTwo{\lk_2})$};
			
			\node[event, draw=black, fill=white] (19) at (0*\xstep, 6*\ystep + 1*\ybias) {$\acq(\LockColorOne{\lk_1})$};
			\node[event, draw=black, fill=white, dotted] (110) at (0*\xstep, 7*\ystep + 1*\ybias) {$\cs(m_0,m_1)$};
			%\node[event, draw=black, fill=white, dotted] (111) at (0*\xstep, 11*\ystep + 1*\ybias) {$\acq(\lk')$};
			%\node[event, draw=black, fill=white, dotted] (112) at (0*\xstep, 12*\ystep + 1*\ybias) {$\rel(\lk')$};
			%\node[event, draw=black, fill=white, dotted] (113) at (0*\xstep, 13*\ystep + 1*\ybias) {$\rel(\lk)$};
			\node[event, draw=black, fill=white] (114) at (0*\xstep, 8*\ystep + 1*\ybias) {$\rel(\LockColorOne{\lk_1})$};
			
			%%%%%%%%%%%%%%%%
			
			\node[event, draw=black, fill=white] (21) at (1*\xstep, 2*\ystep + 0*\ybias) {$\acq(\LockColorOne{\lk_1})$};
			\node[event, draw=black, fill=white, dotted] (22) at (1*\xstep, 3*\ystep + 0*\ybias) {$\cs(m_1,m_0)$};
			%\node[event, draw=black, fill=white, dotted] (23) at (1*\xstep, 4*\ystep + 0*\ybias) {$\acq(\lk)$};
			%\node[event, draw=black, fill=white, dotted] (24) at (1*\xstep, 5*\ystep + 0*\ybias) {$\rel(\lk)$};
			%\node[event, draw=black, fill=white, dotted] (25) at (1*\xstep, 6*\ystep + 0*\ybias) {$\rel(\lk')$};
			\node[event, draw=black, fill=white] (26) at (1*\xstep, 4*\ystep + 0*\ybias) {$\rel(\LockColorOne{\lk_1})$};
			
			\node[event, draw=black, fill=white] (27) at (1*\xstep, 5*\ystep + 1*\ybias) {$\acq_1(\LockColorTwo{\lk_2})$};
			\node[event, draw=black, fill=white, dotted] (28) at (1*\xstep, 6*\ystep + 1*\ybias) {$\cs(m_1,m)$};
			%\node[event, draw=black, fill=white, dotted] (29) at (1*\xstep, 11*\ystep + 1*\ybias) {$\acq(\lk)$};
			%\node[event, draw=black, fill=white, dotted] (210) at (1*\xstep, 12*\ystep + 1*\ybias) {$\rel(\lk)$};
			%\node[event, draw=black, fill=white, dotted] (211) at (1*\xstep, 13*\ystep + 1*\ybias) {$\rel(\lk')$};
			\node[event, draw=black, fill=white] (212) at (1*\xstep, 7*\ystep + 1*\ybias) {$\rel_1(\LockColorTwo{\lk_2})$};
			
			\begin{scope}[]
				\node[below left=of 212] (212b) {};
			\end{scope}
	
		\end{tikzpicture}
	}
	\caption{
		Reduction for OV-hardness proof from an instance of size $n = 2$ and $d = 2$.
		% We use the shortcut $\cs(\lk_i,\lk_j)$ to denote two nested critical sections on $\lk_i$ and $\lk_j$, as indicated in the figure.
	}
	\figlabel{ov-hardness}
\end{subfigure}
%\end{figure}
	
	\myparagraph{Construction}{
		We will construct a trace $\tr$ such that $\tr$ has a deadlock
		pattern of length $2$ iff $(A, B)$ is a positive OV instance.
		The trace $\tr$ is of the form $\tr = \tr^A \cdot \tr^B$
		and uses $2$ threads $\set{t_A, t_B}$ and $d+2$ distinct locks $\lk_1, \ldots, \lk_d, m_0, m_1$. 
		Intuitively, $\tr^A$ and $\tr^B$  encode the given sets of vectors $A$ and $B$.
		The sub-traces $\tr^A = \tr^A_1 \cdot \tr^A_2 \cdots \tr^A_n$
		and $\tr^B = \tr^B_1 \cdot \tr^B_2 \cdots \tr^B_n$ are defined as follows.
		For each $i \in \set{1, 2, \ldots, n}$ and $Z \in \set{A, B}$, the
		sub-trace $\tr^Z_i$ is the unique string generated by the
		grammar having $d+1$ non-terminals $S_0, S_1, \ldots, S_d$, start symbol $S_d$
		and the following production rules:
		\begin{itemize}
			\item $S_0 \to \ev{t_Z, \acq(m)} \cdot \ev{t_Z, \acq(m')} \cdot \ev{t_Z, \rel(m')} \cdot \ev{t_Z, \rel(m)}$,
			where $(m, m') = (m_0, m_1)$ if $Z = A$, and $(m, m') = (m_1, m_0)$.
			\item for each $1 \leq j \leq d$, $S_j \to S_{j-1}$ if $Z_i[j] = 0$.
			Otherwise (if $Z_i[j] =1$), $S_j \to \ev{t_Z, \acq(\lk_j)} \cdot S_{j-1} \cdot \ev{t_Z, \rel(\lk_j)}$.
		\end{itemize}
		In words, all events of $\tr^A$ are performed by thread $t_A$ and those
		in $\tr^B$ are performed by $t_B$.
		Next, the $i^{th}$ sub-trace of $\tr^A$, denoted $\tr^A_i$ corresponds to the vector $A_i$
		as follows --- $\tr^A_i$ is a nested block of critical sections,
		with the innermost critical section being on lock $\lk'$,
		which is immediately enclosed in a  critical section on lock $\lk$.
		Further, in the sub-trace $\tr^A_i$, the lock $\lk_j$
		occurs iff $A_i[j] = 1$.
		The sub-traces $\tr^B_i$ is similarly constructed, except that the order
		of the two innermost critical sections is inverted.
		\figref{ov-hardness} illustrates the construction for an OV-instance with $n=2$ and $d=2$.}
\end{proof}



\myparagraph{The complexity of deadlock prediction}{
Finally, we settle the complexity of the prediction problem for deadlocks,
and show that, even for deadlock patterns of size $2$, the problem is $\W{1}$-hard
parameterized by the number of threads.
In contrast, recall that the $\W{1}$-hardness of 
\thmref{pattern-w1-hardness-pattern} concerns deadlock patterns of arbitrary size.
Our result is based on a similar hardness that was established recently for predicting data races~\cite{Mathur2020b}.


\begin{restatable}{theorem}{wonehardness}
\thmlabel{w1-hardness-pattern}
The problem of checking if a trace $\tr$ has a predictable deadlock of size $2$
is $\W{1}$-hard in the number of threads $\NumThreads$ appearing in $\tr$, and thus
is also $\NP$-hard.
\end{restatable}
}
%!TEX root = main.tex

\section{Synchronization-Preserving Deadlocks and their Prediction}
\seclabel{syncp}
Having established the intractability of general deadlock prediction in \secref{lower-bounds}, we now define the subclass of predictable deadlocks called synchronization-preserving (\emph{sync-preserving}, for short) in \secref{syncp-def}.
The key benefit of sync-preserving deadlocks is that, unlike arbitrary deadlocks, they can be detected efficiently; we develop our algorithm $\SyncPDOffline$ for this task in Sections 4.2-4.5.
Our experiments later indicate that most predictable deadlocks are actually sync-preserving, hence the benefit of fast detection comes at the cost of little-to-no precision loss in practice.
%We start with a brief overview.

\myparagraph{Overview of the algorithm}{
	There are several insights behind our algorithm.
	First, given a deadlock pattern, one can verify 
	if it is a sync-preserving deadlock in linear time (\secref{verify-patterns});
	this is based on our sound and complete characterization of sync-preserving deadlocks (\secref{characterize-patterns}).
%	This is in sharp contrast with prior works
%such as~\cite{Cai2021,Kalhauge2018} that resort to
%SMT-solving or expensive (i.e., of large polynomial) graph-based analyses. 
	Next, instead of verifying single deadlock patterns one-by-one, 
	we consider \emph{abstract deadlock patterns},
	which are essentially collections of deadlock patterns
	that share the same signature;
	the formal definition is given in \secref{verify-abstract-patterns}.
	We show that our basic algorithm can
	be extended to \emph{incrementally}
	verify \emph{all} the concretizations of an abstract deadlock pattern in linear time (\secref{verify-abstract-patterns}), in a single pass (\lemref{abstract-pattern-linear-time}).
	%Finally, we extend this algorithm to all deadlock patterns in the trace,
	%by first enumerating all the abstract deadlock patterns by constructing an \emph{abstract lock graph},
	%and then verifying each of them using the algorithm thus far (\secref{enumerate-patterns}).
	Finally, we feed this algorithm all the abstract deadlock patterns of the input trace,
    by constructing an \emph{abstract lock graph} and enumerating cycles in it (\secref{enumerate-patterns}).
	
	% We next discuss these in detail \ucomment{and conclude with a summary of our algorithm in section 4.6}.
}

\subsection{Synchronization-Preserving Deadlocks}
\seclabel{syncp-def}

Our notion of sync-preserving deadlocks builds on the recently introduced concept of sync-preserving correct reorderings~\cite{Mathur2021}.

\begin{definition}[Sync-preserving Correct Reordering]
\deflabel{syncp-correct-reordering}
A correct reordering $\rho$ of a trace $\tr$ is
\emph{sync-preserving} if for every lock $\lk \in \locks{\rho}$
and every two acquire events $e_1 \neq e_2 \in \events{\rho}$
with $\OpOf{e_1} = \OpOf{e_2} = \acq(\lk)$, the order of $e_1$ and
$e_2$ is the same in $\tr$ and $\rho$, i.e., 
$e_1 \trord{\rho} e_2$ iff $e_1 \trord{\tr} e_2$.
\end{definition} 
A sync-preserving correct reordering
preserves the order of those critical sections (on the same lock) that actually appear in 
the reordering, but allows 
intermediate critical sections to be dropped completely.
% 
% Since a correct reordering $\pi$ of $\tr$ may only contain
% a prefix of events in some thread $t$, some critical section of $\tr$
% in thread $t$ may not be present in $\pi$.
% In particular, sync-preserving correct reorderings allow intermediate 
% critical sections (of the same lock) to be dropped, 
% as long as the order on the remaining ones is preserved.
This style of reasoning is more permissive than
the space of reorderings %(and their prefixes) 
induced by the Happens-Before partial order~\cite{Lamport78},
that implicitly enforces that all intermediate critical sections on a lock
be present.
Sync-preserving deadlocks can now be defined naturally.
\begin{definition}[Sync-preserving Deadlocks]
\deflabel{syncp-deadlock}
Let $\tr$ be a trace and $D = \pattern{e_0, e_1, \ldots, e_{k-1}}$
be a deadlock pattern.
We say that $D$ is a sync-preserving deadlock of $\tr$
if there is a sync-preserving correct reordering $\rho$ of $\tr$
such that each of $e_0, \ldots, e_{k-1}$ is $\tr$-enabled in $\rho$.
\end{definition}

%!TEX root=../main.tex

\begin{figure}
	\centering
	\begin{subfigure}[t]{0.17\textwidth}
		\scalebox{0.9}{
			\execution{3}{
				\figev{1}{$\acq(\LockColorOne{\lk_1})$}
				\figev{1}{$\acq(\LockColorTwo{\lk_2})$}
				\figev{1}{$\rel(\LockColorTwo{\lk_2})$}
				\figev{1}{$\acq(\LockColorTwo{\lk_2})$}
				\figev{1}{$\wt(y)$}
				\figev{1}{$\rel(\LockColorTwo{\lk_2})$}
				\figev{1}{$\rel(\LockColorOne{\lk_1})$}
				\figev{2}{$\acq(\LockColorThree{\lk_3})$}
				\figev{2}{$\wt(x)$}
				\figev{2}{$\rd(y)$}
				\figev{2}{$\rel(\LockColorThree{\lk_3})$}
				\figev{3}{$\acq(\LockColorTwo{\lk_2})$}
				\figev{3}{$\acq(\LockColorThree{\lk_3})$}
				\figev{3}{$\rd(x)$}
				\figev{3}{$\rel(\LockColorThree{\lk_3})$}
				\figev{3}{$\acq(\LockColorOne{\lk_1})$}
			}
		}
	\end{subfigure}
	\hspace*{\fill}
	\begin{subfigure}[t]{0.17\textwidth}
		\scalebox{0.9}{
			\execution{3}{
				\figevoffset{16}{3}{$\wt(v)$}
				\figevoffset{16}{3}{$\rel(\LockColorOne{\lk_1})$}
				\figevoffset{16}{3}{$\acq(\LockColorOne{\lk_1})$}
				\figevoffset{16}{3}{$\rel(\LockColorOne{\lk_1})$}
				\figevoffset{16}{3}{$\rel(\LockColorTwo{\lk_2})$}
				\figevoffset{16}{2}{$\acq(\LockColorFour{\lk_4})$}
				\figevoffset{16}{2}{$\acq(\LockColorOne{\lk_1})$}
				\figevoffset{16}{2}{$\wt(z)$}
				\figevoffset{16}{2}{$\rd(v)$}
				\figevoffset{16}{2}{$\rel(\LockColorOne{\lk_1})$}
				\figevoffset{16}{2}{$\rel(\LockColorFour{\lk_4})$}
				\figevoffset{16}{1}{$\acq(\LockColorOne{\lk_1})$}
				\figevoffset{16}{1}{$\acq(\LockColorTwo{\lk_2})$}
				\figevoffset{16}{1}{$\rd(z)$}
				\figevoffset{16}{1}{$\rel(\LockColorTwo{\lk_2})$}
				\figevoffset{16}{1}{$\rel(\LockColorOne{\lk_1})$}
			}
		}
	\end{subfigure}
	\hspace*{\fill}
	\begin{subfigure}[t]{0.35\textwidth}
		\begin{tikzpicture}[baseline=-180pt]
			\tikzset{rectangle/.append style={draw=black}}
			\scalebox{0.84}{
				\def\xstep{2.2}
				\def\ystep{0.5}
				\node[align=left] (tt1) at (0*\xstep, -3.5){
					$
					\eta_1=
					\tuple{t_1, \LockColorTwo{\lk_2}, \{ \LockColorOne{\lk_1} \},
						\sequence{e_2, e_4, e_{29}}}
					$\\
					$
					\eta_2=
					\tuple{
						t_2, \LockColorOne{\lk_1}, \{ \LockColorFour{\lk_4} \},
						\sequence{e_{23}}}
					$\\
					$
					\eta_3=
					\tuple{
						t_3, \LockColorOne{\lk_1}, \{ \LockColorTwo{\lk_2} \},
						\sequence{e_{16}, e_{19}}}
					$\\
					$
					\eta_4=
					\tuple{
						t_3, \LockColorThree{\lk_3}, \{ \LockColorTwo{\lk_2} \},
						\sequence{e_{13}}}
					$
				};
				
				\node[align=left] (tt2) at (0*\xstep, -4.8){
					$\abst{D}=\pattern{\eta_1, \eta_3}$
				};
				
				\node[align=left] (tt3) at (-0.5*\xstep, -6){
					$D_1=\pattern{e_2, e_{16}}$\\
					$D_2=\pattern{e_2, e_{19}}$\\
					$D_3=\pattern{e_4, e_{16}}$\\
				};
				\node[align=left] (tt4) at (0.5*\xstep, -6){
					$D_4=\pattern{e_4, e_{19}}$\\
					$D_5=\pattern{e_{29}, e_{16}}$\\
					$D_6=\pattern{e_{29}, e_{19}}$\\
				};
			}
		\end{tikzpicture}
	\end{subfigure}
	% \vspace{-40pt}
	\caption{
		A trace $\tr_3$, its abstract acquires $\eta_i$, unique abstract deadlock pattern $\abst{D}$, concrete patterns $D_i\in \abst{D}$. 
	}
	% \vspace{-10pt}
	\figlabel{syncp_example}
\end{figure}


\begin{example}
\exlabel{sp-deadlocks}
Consider the trace $\tr_2$ in \figref{motivating-dl}.
The deadlock pattern $D = \pattern{e_4, e_{18}}$ is a sync-preserving deadlock, witnessed by the sync-preserving correct reordering
$\rho_3 = e_1e_2e_3 e_8 e_9 \,e_{12}..e_{15}\,e_{16} e_{17}$.
Now consider the trace $\tr_3$ from \figref{syncp_example}
and the deadlock pattern $D_5 = \pattern{e_{29}, e_{16}}$.
This is a predictable deadlock, witnessed by the correct reordering
$\rho_5 = e_1..e_7e_8..e_{11}e_{12}..e_{15}\,e_{28}$.
Observe that $\rho_5$ is a sync-preserving reordering, which makes $D_5$ a sync-preserving deadlock.
%\hunkar{
A key aspect in $\rho_5$ is that the events $e_{22}..e_{27}$ are dropped, as otherwise $e_{16}$ cannot be $\sigma_3$-enabled.
%}
%Another reordering that witnesses the deadlock $D_5$ is
%$\rho_5 = e_8\,e_1..e_3\,e_9\,e_4..e_7\,e_{12}e_{10}\,e_{11}e_{28}\,e_{13}..e_{15}$.
%This reordering $\rho_5$ inverts the order between some events,
%but is sync-preserving because
%it preserves the order between all pairs of critical sections
%on the same lock. % is the same as in $\tr_2$.
% Hence, $\rho_5$ is also a sync-preserving correct reordering.
A similar reasoning applies for the deadlock pattern $D_6$, and it is also a sync-preserving deadlock.
The other deadlock patterns ($D_1, D_2, D_3, D_4$)
are not predictable deadlocks.
Intuitively, the reason for this is that realizing these deadlock patterns require executing the read event $e_{14}$, which then enforces to execute the events $e_8..e_{11}$ and $e_1..e_6$.
This prevents the deadlocks from becoming realizable as the events 
$e_2$ or $e_4$ that appear in these deadlock patterns are no longer $\sigma_3$-enabled.
This point is detailed in \exref{syncp-ex}.
\end{example}






%!TEX root = main.tex

\subsection{Characterizing Sync-Preserving Deadlocks}
\seclabel{characterize-patterns}

There are two fundamental tasks in 
searching for a correct reordering that witnesses a deadlock --- 
(i)~determining the set of events in the correct reordering, and 
(ii)~identifying a total order on such events --- both of which are intractable~\cite{Mathur2020b}.
On the contrary, for sync-preserving deadlocks, we show that
(a)~the search for a correct reordering can be reduced
to the problem of checking if some well-defined set of events (\defref{spclosure})
does not contain the events appearing in the deadlock pattern (\lemref{spclosure-deadlock}),
and that (b)~this set can be constructed efficiently.
%These results directly extend the results of prior work
%on race detection~\cite{Mathur2021} to the more general
%case of checking concurrency of \emph{arbitrarily many} events.
%Let us define this set next.
\begin{definition}[Sync-Preserving Closure]
\deflabel{spclosure}
Let $\tr$ be a trace and $S \subseteq \events{\tr}$.
The sync-preserving closure of $S$, denoted
$\SPClosure{\tr}(S)$ is the smallest set $S'$ such that
\begin{enumerate*}[label=(\alph*)]
	\item \itmlabel{def-item-spclosure-a} $S \subseteq S'$,
	\item \itmlabel{def-item-spclosure-b} for every $e, e' \in \events{\tr}$ such that 
	$e \stricttho{\tr} e'$ or $e = \rf{\tr}(e')$, if $e' \in S'$, then $e \in S'$, and
	\item\itmlabel{acq-order-closure} for every lock $\lk$ and every two distinct events $e, e'\in S'$ with $\OpOf{e} = \OpOf{e'} = \acq(\lk)$,
	if $e \trord{\tr} e'$ then $\match{\tr}(e) \in S'$.
\end{enumerate*}
\end{definition}


\defref{spclosure} resembles 
the notion of correct reorderings (\defref{syncp-correct-reordering}). 
Indeed, \lemref{spclosure-necessary-sufficient} justifies using this set --- it is
both a necessary and a sufficient set for sync-preserving correct reorderings.

\begin{restatable}{lemma}{spclosureNecessarySufficient}
\lemlabel{spclosure-necessary-sufficient}
Let $\tr$ be a trace and let $S \subseteq \events{\tr}$.
For any sync-preserving correct reordering $\rho$ of $\tr$, 
	if $S \subseteq \events{\rho}$, then $\SPClosure{\tr}(S) \subseteq \events{\rho}$.
	Further, there is a sync-preserving correct reordering $\rho$ of 
	$\tr$ such that $\events{\rho} = \SPClosure{\tr}(S)$.
% \begin{enumerate}[label=(\arabic*)]
% 	\item For any sync-preserving correct reordering $\rho$ of $\tr$, 
% 	if $S \subseteq \events{\rho}$, then $\SPClosure{\tr}(S) \subseteq \events{\rho}$.
% 	\item There is a sync-preserving correct reordering $\rho$ of 
% 	$\tr$ such that $\events{\rho} = \SPClosure{\tr}(S)$.
% \end{enumerate}
\end{restatable}

%\hunkar{
For an intuition, consider again \cref{fig:syncp_example} and the sync-preserving correct reordering $\rho_5 = e_1..e_7e_8..e_{11}e_{12}..e_{15}\,e_{28}$ computed in \cref{ex:sp-deadlocks}.
According to \cref{lem:spclosure-necessary-sufficient}, $\SPClosure{\tr_3}(S) \subseteq \events{\rho_5}$ holds for all $S$ such that $S \subseteq \events{\rho_5}$.
For example, if we take $S=\set{e_1, e_{15}}$ then observe that $S \subseteq \events{\rho_5} $ and  $\SPClosure{\tr_3}(S)=\set{e_1, \ldots, e_6, \, e_8, \ldots e_{15}} \subseteq \events{\rho_5}$ holds.
%}

%\hunkar{
%Can we provide any intuition on the second part of the lemma?
%}
Based on~\lemref{spclosure-necessary-sufficient}, we 
present a sound and complete characterization of sync-preserving deadlocks
(\lemref{spclosure-deadlock}).
For a set $S \subseteq \events{\tr}$, 
we let $\prev{\tr}(S)$ denote the set of immediate thread
predecessors of events in $S$. 
That is,
$\prev{\tr}(S) = \setpred{e \in \events{\tr}}{\exists f \in S, e \stricttrord{\tr} f \text{ and } \forall e' \stricttrord{\tr} f, e' \tho{\tr} e}$.
% \begin{restatable}{lemma}{spclosure-iff-syncp-correct-reordering}
% \lemlabel{spclosure-iff-syncp-correct-reordering}
% Let $\tr$ be a trace and let $S \subseteq \events{\tr}$.
% There is a sync-preserving correct reordering $\rho$ of $\tr$
% in which each event $e \in S$ is $\tr$-enabled iff 
% $\SPClosure{\tr}(\prev{\tr}(S)) \cap S = \emptyset$.
% \end{restatable}

\begin{restatable}{lemma}{spclosureDeadlock}
\lemlabel{spclosure-deadlock}
Let $\tr$ be a trace and let $D = \pattern{e_0, \ldots, e_{k-1}}$
be a deadlock pattern of size $k$ in $\tr$.
$D$ is a sync-preserving deadlock of $\tr$ iff
$\SPClosure{\tr}(\prev{\tr}(S)) \cap S = \emptyset$, where
$S = \set{e_0, \ldots, e_{k-1}}$.
\end{restatable}

%!TEX root = main.tex

\begin{example}
\exlabel{syncp-ex}
%We illustrate sync-preserving closure on our example traces.
% We will use the notation 
% $\prev{\tr}^*(S) =
% \setpred{e \in \events{\tr}}{\exists f \in S, e \tho{\tr} f}$ 
% for a trace $\tr$ and a subset $S \subseteq \events{\tr}$.
% For ease of notation, given a trace $\tr$ and set of events $S\subseteq \events{\tr}$, we denote by $\LowerSet_{\tr}(S)=\{e\colon \exists e'. e\tho{\tr}e' \}$ these events and their local predecessors in $\tr$.\\
%
Consider the trace $\tr_2$ in \figref{motivating-dl}, 
and the deadlock pattern $D=\pattern{e_{4}, e_{18}}$.
We have 
$\SPClosure{\tr_2}(\prev{\tr_2}(\{ e_{4}, e_{18} \})) = 
\set{e_1, e_2, e_3, e_8, e_9, e_{12}, \ldots, e_{17}}$.
Since we have that $e_{4}, e_{18} \not \in$ $\SPClosure{\tr_2}( \prev{\tr_2}(\set{e_{4},  e_{18}}))$, $D$ is a sync-preserving deadlock. 
%
Now consider the trace $\tr_3$ in \figref{syncp_example},
and the deadlock patterns
$D_1=\allowbreak\pattern{e_2, e_{16}}$, 
$D_5=\allowbreak\pattern{e_{29}, e_{16}}$, and
$D_6=\allowbreak\pattern{e_{29}, e_{19}}$.
We have
$\SPClosure{\tr_3}(\prev{\tr_3}(\{ e_{2},e_{16}\}))\allowbreak=
\set{e_1, \ldots, e_6, e_8, \ldots, e_{15}}$, 
$\SPClosure{\tr_3}(\prev{\tr_3}(\allowbreak\{ e_{29},e_{16}\}))=
\set{
e_1, \ldots, e_{15}, e_{28}}$, and $\SPClosure{\tr_3}(\prev{\tr_3}(\{ e_{29},e_{19}\}))=
\set{
e_1, \ldots, e_{18}, e_{28}}$.
Since $e_2 \in \allowbreak \SPClosure{\tr_3}(\allowbreak\prev{\tr_3}(\{ e_{2},e_{16}\}))$,
$D_1$ is not a sync-preserving deadlock.
However, $e_{29}, e_{16} \not \in \SPClosure{\tr_3}(\prev{\tr_3}(\{ e_{29},e_{16}\}))$, and $e_{29}, e_{19} \not \in \SPClosure{\tr_3}(\prev{\tr_3}(\{ e_{29},e_{19}\}))$, thus $D_5$ and $D_6$ are sync-preserving deadlocks (as we also concluded in \exref{sp-deadlocks}).
\end{example}

%!TEX root = main.tex

\subsection{Verifying Deadlock Patterns}
\seclabel{verify-patterns}

Given a deadlock pattern, we check if it constitutes a sync-preserving deadlock
by constructing the sync-preserving closure (\lemref{spclosure-deadlock}) in linear time.
Based on \defref{spclosure}, this can be done in 
an iterative manner. 
%start with the set of strict predecessors
%of events in the deadlock pattern, and in each iteration, add
%$\tho{}$ and $\rf{}$ predecessors of the current set of events,
%and additionally identify the release events that must be added to the set.
We 
(i)~start with the set of $\tho{}$ predecessors of the events in the deadlock pattern, and
(ii)~iteratively  add
$\tho{}$ and $\rf{}$ predecessors of the current set of events.
Additionally, we identify and add the release events that must be included in the set.
% 
% \defref{spclosure} suggests that the closure can be computed in
% an iterative manner as follows --- starting with the set of strict predecessors
% of the events in the deadlock pattern,  until a fixpoint is reached,
% in each iteration, we can add events that are $\tho{}$
% or $\rf{}$ predecessors of the current set of events,
% and additionally, identify the release events that must be added to the set.
%In order to ensure that the \emph{entire} fixpoint computation is
%performed in linear time, we use \emph{timestamps}.
%\hunkar{
We utilize \emph{timestamps} to ensure that the \emph{entire} fixpoint computation is
performed in linear time.
%}
% to ensure that indeed each iteration is performed in time 
% proportional to only the number of events actually added in that iteration.

\newcommand{\TS}[1]{\mathsf{TS}_{\tr}}

\myparagraph{Thread-read-from timestamps}{
	Given a set $\threads{}$ of threads, a timestamp 
	is simply a mapping $T : \threads{} \to \nats$.
	Given timestamps $T_1, T_2$, we 
	use the notations 
	$T_1 \cle T_2$ and $T_1 \mx T_2$ for pointwise comparison
	and pointwise maximum, respectively.
	% $T_1 \cle T_2 \equiv \forall p \in P, T_1(p) \leq T_2(p)$
	% and
	% $T_1 \mx T_2 = \lambda p \in P, \max\set{T_1(p), T_2(p)}$.
	For a set $U$ of timestamps, we write $\bigsqcup U$ to denote the
	pointwise maximum over all elements of $U$.
	% Further, for a set $U$ of $\threads{}$-indexed timestamps, we write $\bigsqcup U$ to denote
	% $\lambda p \in P, \max\set{T(p)}_{T \in U}$.
	Let $\trf{\tr}$ be the reflexive transitive closure
	of the relation $(\tho{\tr} \cup \setpred{(\rf{\tr}(e), e)}{\exists x \in \vars{\tr}, \OpOf{e} = \rd(x)})$;
	observe that
	$\trf{\tr}$ is a partial order.
	We define the 
	% threads-read-from
	 timestamp $\TS{\tr}^e$ of an event $e$ in $\tr$
	to be a $\threads{\tr}$-indexed timestamp as follows: $\TS{\tr}^e(t) = |\setpred{f}{f \trf{\tr} e}|$.
	% such that for every $t \in \threads{\tr}$,
	% \[
	% 	\TS{\tr}^e(t) = |\setpred{f}{f \trf{\tr} e}|
	% \]
	This ensures that
	for two events $e, e' \in \events{\tr}$,
	$e \trf{\tr} e'$ iff $\TS{\tr}^e \cle \TS{\tr}^{e'}$.
	For a set $S \subseteq \events{\tr}$, we overload the notation
	and say the timestamp of $S$ is $\TS{\tr}^S = \bigsqcup \set{\TS{\tr}^e}_{e \in S}$.
	% Observe that, for a set $S$ that is $\trf{\tr}$-downward closed, we have $\TS{\tr}^e \cle \TS{\tr}^S$ is true iff $e \in S$.
	% In other words, timestamps allow for a succinct representation of 
	%$\trf{\tr}$-downward closed sets.
	Given a trace $\tr$ with $\NumEvents$ events and
	$\NumThreads$ threads we can compute these timestamps for all the events in $O(\NumEvents \cdot \NumThreads)$ time, using a simple vector clock algorithm~\cite{Mattern89,Fidge91}.
}


\myparagraph{Computing sync-preserving closures}{
	Recall the basic template of the fixpoint computation.
	In each iteration, we identify the set of release events
	that must be included in the set, together with their $\trf{\tr}$-closure.
	In order to identify such events efficiently, for every
	thread $t$ and lock $\lk$, we maintain
	a FIFO queue $\AcqLst_{t, \lk}$ (\emph{critical section history} of $t$ and $\lk$) 
	to store 
	the sequence of events that acquire $\lk$ in thread $t$.
	In each iteration, we traverse each list
	to determine the last acquire event that belongs to the current set.
	For a given lock, we need to add the matching release events of
	all thus identified events to the closure, 
	except possibly the matching release event of the latest acquire event (see~\defref{spclosure}).
	This computation is performed using timestamps, as shown in \algoref{compute-closure}.
	Starting with a set $S$, the algorithm
	runs in time $O(|S|\cdot\NumThreads + \NumThreads\cdot\NumAcquires)$,
	where $\NumThreads$ and $\NumAcquires$ are respectively
	the number of threads and acquire events in $\tr$.
}

\begin{minipage}{0.45\textwidth}
	%!TEX root = main.tex

% Input: trace tr, ad a deadlock pattern D = <e0, e1, .., ek>,

% SPClosure(tr, S):
% 	closure = emptyset
% 	while closure changes:

\small
%\setlength{\textfloatsep}{0pt}
\begin{algorithm}[H]
\Input{Trace $\tr$, Timestamp $T_0$}
\BlankLine
% \myfun{\fixpoint{$\sigma$, $T_0$, $\set{\AcqLst_{t,\lk}}_{\lk\in \locks{\tr}, t \in \threads{\tr}}$}}{
	\Let $\set{\AcqLst_{t,\lk}}_{\lk\in \locks{\tr}, t \in \threads{\tr}}$ be the lock-acquisition histories in $\tr$ \;
	$T \gets T_0$ \;
	\Repeat{$T$ does not change}{
		\For{$\lk \in \locks{}$}{
			\ForEach{$t \in \threads{}$}{
				\Let $e_t$ be the last event in $\AcqLst_{t,\lk}$ with $\TS{\tr}^{e_t} \cle T$ \;
				Remove all earlier events in $\AcqLst_{t, \lk}$ 
			}
			\Let $e_{t*}$ be the last event in $\set{e_t}_{t\in \threads{\tr}}$ according to $\trord{\tr}$ \;
			$T$ := $T \mx \bigsqcup \setpred{\TS{\tr}^{\match{\tr}(e_t)}}{e_t \neq e_{t*}}$	\;
		}
	}
	\Return $T$
% }
\caption{CompSPClosure:\\ Computing sync-preserving closure.}
\algolabel{compute-closure}
\end{algorithm}
\normalsize

\end{minipage}
\hfill ~ %\unskip\ \vrule\ \hfill
\begin{minipage}{0.48\textwidth}
	%!TEX root = main.tex

\small
\begin{algorithm}[H]
%\Input{Trace $\tr$, Abstract deadlock pattern $\abst{D}$ of length $k$}
\Input{Trace $\tr$, $\abst{D}$ of length $k$}
\BlankLine
	\Let $F_0, \ldots, F_{k-1}$ be the sequences of acquires in $\abst{D}$ \;
	\Let $n_0, \ldots, n_{k-1}$ be the lengths of $F_0, \ldots, F_{k-1}$ \;
	\lForEach{$j \in \set{0, \ldots, k-1}$}{
		$i_j \gets 1$
	}
	$T \gets \lambda t, 0$ %\tcp{Timestamp of sync-preserving closure}
	\While{$\bigwedge\limits_{j=0}^{k-1} i_j < n_j$}{
		\Let $e_0 = F_0[i_0], \ldots, e_{k-1} = F_{k-1}[i_{k-1}]$ \;
		$S \gets \prev{\tr}{\set{e_0, \ldots, e_{k-1}}}$\;
		$T \gets$ \fixpoint{$\tr, T \mx \TS{\tr}^S$}\\   \linelabel{call-comp-closure}
		%\tcp{\algoref{compute-closure}}
		\If{$\forall j < k, \TS{\tr}^{e_{j}} \cle T$}{
			\report pattern $D = e_0, \ldots, e_{k-1}$ and \exit
		}
		\ForEach{$j \in \set{0, \ldots, k-1}$}{
			$i_j = \min\setpred{l \leq n_j}{\TS{\tr}^{F_j[l]} \not\cle T}$ %\tcp{Update index in $F_j$}
		}
	}
\caption{CheckAbsDdlck:\\ Checking an abstract deadlock pattern.}
\algolabel{abstract-pattern}
\end{algorithm}
\normalsize

\end{minipage}



\myparagraph{Checking a deadlock pattern}{
	After computing the timestamp $T$ of
	the closure (output of \algoref{compute-closure}, starting with the set of events in the given deadlock pattern),
	determining whether a given deadlock pattern
	$D = e_0, \ldots, e_{k-1}$ is a sync-preserving deadlock 
	can be performed in time $O(k\cdot\NumThreads)$ --- simply check if  $\forall i, \TS{\tr}(e_i) \not\cle T$.
	This gives an algorithm for checking if a deadlock pattern
	of length $k$ is sync-preserving that runs in time 
	$O(\NumThreads\cdot \NumEvents + k\cdot\NumThreads + \NumThreads\cdot \NumAcquires) = O(\NumEvents\cdot \NumThreads)$.
}
%!TEX root = main.tex

\subsection{Verifying Abstract Deadlock Patterns}
\seclabel{verify-abstract-patterns}

\myparagraph{Abstract acquires and abstract deadlock patterns}{
Given a thread $t$, a lock $\lk$ and a set of locks 
$L\subseteq \locks{\tr}\neq \emptyset$ with $\lk\not \in L$,
we define the \emph{abstract acquire} $\eta =\tuple{t, \lk, L, F}$, 
where $F=\sequence{e_1,\dots, e_n}$ is the sequence of all events
$e_i\in \events{\tr}$ (in trace-order) such that for each $i$, we have 
(i)~$\ThreadOf{e_i}=t$,
(ii)~$\OpOf{e_i} = \acq(\lk)$, and
(iii)~$\lheld{\tr}(e_i) = L$.
In words, the abstract acquire $\eta$ contains the sequence of all acquire events 
%of the same \emph{signature}.
%\hunkar{
of a specific thread, that access a specific lock and hold the same set of locks when executed, ordered as per thread order.
%}
An \emph{abstract deadlock pattern} of size $k$ in a trace $\tr$ is a sequence
$
\abst{D} = \eta_0,\dots,\eta_{k-1}
$
of abstract acquires $\eta_i=\tuple{t_i, \lk_i, L_i, F_i}$ such that
$t_0, \ldots, t_{k-1}$ are distinct threads, 
$\lk_0,  \ldots, \lk_{k-1}$ are distinct locks,
and $L_0, L_1, \ldots, L_{k-1} \subseteq \locks{\tr}$ are
such that $\lk_i \not\in L_i$, $\lk_i \in L_{(i+1)\% k}$ for every $i$,
and $L_i \cap L_j= \emptyset$ for every $i\neq j$.
Thus, an abstract deadlock pattern $\abst{D}$ succinctly encodes all 
concrete deadlock patterns $F_0\times F_1\times \dots \times F_{k-1}$,
called \emph{instantiations} of $\abst{D}$.
We also write $D\in \abst{D}$ to denote that 
$D\in F_0\times F_1\times \dots \times F_{k-1}$.
We say that $\abst{D}$ contains a sync-preserving deadlock if there 
exists some instantiation $D\in \abst{D}$ that is a sync-preserving deadlock.
See \figref{syncp_example} for an example.
}
Our next result is stated below, followed by its proof idea.
\begin{restatable}{lemma}{abstractlineartime}
\lemlabel{abstract-pattern-linear-time}
Consider a trace $\tr$ with $\NumEvents$ events and $\NumThreads$ threads, 
and an abstract deadlock pattern $\abst{D}$ of $\tr$.
We can determine if $\abst{D}$ contains a sync-preserving deadlock in $O(\NumThreads\cdot \NumEvents)$ time.
\end{restatable}

%The algorithm for checking if a given abstract pattern
%$\abst{D}$ contains a sync-preserving deadlock is presented in 
%\algoref{abstract-pattern}.
%The algorithm relies on exploiting monotonicity properties
%enjoyed by sync-preserving deadlocks that allow us to
%avoid a naive enumeration of all $O(\NumEvents^k)$ instantiations
%of $\abst{D}$.
%A detailed description of the algorithm is presented in
%\appref{verify-abstract-patterns}.

%\hcomment{moved the text from the appendix -- start}


An abstract deadlock pattern of length $k \geq 2$ can have $\NumEvents^k$ instantiations,
giving a naive enumerate-and-check algorithm running 
in time $O(\NumThreads\cdot\NumEvents^{k+1})$, which is prohibitively large.
Instead, we exploit 
(i)~the monotonicity properties of 
the sync-preserving closure (\propref{spclosure-monotone}) and
(ii)~instantiations of an abstract pattern (\corref{pattern-monotone}) 
that allow for an \emph{incremental}
algorithm that iteratively checks successive instantiations
of a given abstract deadlock pattern, while spending total $O(\NumEvents\cdot \NumThreads)$ time.
% , i.e., the running time is linear in $\NumEvents$, which is the dominating parameter.
The first observation allows us to re-use a prior
computation when checking later deadlock patterns.

% We first observe that the sync-preserving closure is monotonic
% with respect to the thread-order, and thus, when computing closure
% of later events, one can avoid re-computation.
\begin{restatable}{proposition}{spclosureMonotone}
	\proplabel{spclosure-monotone}
	For a trace $\tr$ and sets $S, S' \subseteq \events{\tr}$.
	If for every event $e \in S$, there is an event $e' \in S'$
	such that $e \tho{\tr} e'$, then
	$\SPClosure{\tr}(S) \subseteq \SPClosure{\tr}(S')$.
\end{restatable}

%\begin{proof}
%Follows from the definition of $\SPClosure{\tr}(S)$.
%\end{proof}

%\hunkar{
Consider $\tr_3$ in~\cref{fig:syncp_example} and let $S=\prev{\tr_3}{(\set{e_{29}, e_{16}})}$, and $S'=\prev{\tr_3}{(\set{e_{29}, e_{19}})}$.
The sets $S, S'$ satisfy the conditions of~\cref{prop:spclosure-monotone}, hence $\SPClosure{\tr_3}(S) \subseteq \SPClosure{\tr_3}(S')$, as computed in~\cref{ex:syncp-ex}.
%}
Next, we extend \propref{spclosure-monotone} to avoid redundant
computations when a sync-preserving deadlock is not found and later deadlock patterns must be checked.
Given two deadlock patterns
$D_1 = e_0, \ldots, e_{k-1}$ and $D_2 = f_0, \ldots, f_{k-1}$
of the same length $k$, 
we say $D_1 \prec D_2$ if
they are instantiations of a common abstract pattern $\abst{D}$
(i.e., $D_1, D_2 \in \abst{D}$) and 
for every $i<k$, $e_i \tho{\tr} f_i$.
\begin{restatable}{corollary}{pattern-monotone}
	\corlabel{pattern-monotone}
	Let $\tr$ be a trace and let $D_1 = e_0, \ldots, e_{k-1}$ and 
	$D_2 = f_0, \ldots f_{k-1}$ be deadlock
	patterns of size $k$ in $\tr$ such that $D_1 \prec D_2$.
	Let $S_1 = \set{e_0, \ldots, e_{k-1}}$ and $S_2 = \set{f_0, \ldots, f_{k-1}}$.
	If $\SPClosure{\tr}(\prev{\tr}(S_1))\cap S_2 \neq \emptyset$,
	then $\SPClosure{\tr}(\prev{\tr}(S_2))\cap S_2 \neq \emptyset$.
\end{restatable}


We now describe how \cref{prop:spclosure-monotone} and \cref{cor:pattern-monotone} are used in our algorithms, 
and illustrate them later in  \cref{ex:ex4}. % in \cref{sec:enumerate-patterns} provides an example.
% \hunkar{
% , we illustrate the application of  \propref{spclosure-monotone} and \corref{pattern-monotone}.
% }
% \hunkar{Provide an intuition here too?}
\algoref{abstract-pattern} checks if an abstract deadlock 
pattern contains a sync-preserving deadlock.
The algorithm iterates over the sequences $F_0, \ldots, F_{k-1}$ of 
acquires (one for each abstract acquire) in trace order.
For this, it maintains
indices $i_0, \ldots, i_{k-1}$ that point to entries in $F_0, \ldots, F_{k-1}$.
%At each step, it determines whether the current
%deadlock pattern $D = e_0, \ldots, e_{k-1}$
%by computing the sync-preserving closure
%of the predecessors of the deadlock pattern, and checks
%if it is disjoint from the deadlock.
%If so, it terminates.
%\hunkar{
At each step, it determines whether the current deadlock pattern $D = e_0, \ldots, e_{k-1}$
constitutes a sync-preserving deadlock by computing the sync-preserving closure
of the thread-local predecessors of the events of the deadlock pattern.
%The algorithm reports a deadlock if the sync-preserving closure is disjoint from the timestamps of the events $e_0, \ldots, e_{k-1}$.
The algorithm reports a deadlock if the sync-preserving closure does not contain any of $e_0, \ldots, e_{k-1}$.
%}
Otherwise, it looks for the next eligible deadlock pattern,
which it determines based on \corref{pattern-monotone}.
In particular, it advances the pointer $i_j$
all the way until an entry which is outside of the
closure computed so far.
Observe that the timestamp $T$ of the closure
computed in an iteration is being used in later iterations;
this is a consequence of \propref{spclosure-monotone}.
%\hunkar{
Furthermore, in the call to the \algoref{compute-closure} at \lineref{call-comp-closure}, we ensure that the list of acquires $\AcqLst_{t, \lk}$,
used in the function \fixpoint
is reused across iterations, and not
re-assigned to the original list of all acquire events.
%}
The correctness of this optimization follows from~\propref{spclosure-monotone}.
Let us now calculate the running time of \algoref{abstract-pattern}.
Each of the $\AcqLst_{t, \lk}$ in \fixpoint is traversed
at most once. 
Next, each element of the sequences $F_0, \ldots, F_{k-1}$
is also traversed at most once.
For each of these acquires, the algorithm spends $O(\NumThreads)$
time for vector clock updates.
The total time required is thus $O(\NumEvents \cdot \NumThreads)$.
This concludes the proof of \lemref{abstract-pattern-linear-time}.

%\hcomment{moved the text from the appendix -- end}
%!TEX root = main.tex

\subsection{The Algorithm \SyncPDOffline}
\seclabel{enumerate-patterns}

\begin{comment}
As with prior work, we model the problem of finding deadlock patterns as a cycle-detection problem
over a graph. % $\lkevgraph{\tr}$ whose nodes are elements of the trace.
%However, instead of a naive lock graph (with locks as nodes),
%the nodes in our graph store additional metadata that helps us
%to rule out false positives, and eventually enumerate all deadlock
%patterns, for which we need information about concrete dynamic events in the trace.
However, instead of considering locks as vertices (as in prior works~\cite{Bensalem2006,Bensalem2005}),
we use an \emph{abstract lock graph} $\lkevgraph{\tr}$.
% of the graph, our 
% representation, called \emph{abstract lock graph} $\lkevgraph{\tr}$ 
% effectively exploits \lemref{abstract-pattern-linear-time}.
% The key novelty compared to such graphs in the literature
% comes from combining it with the notion of sync-preserving deadlocks.
Every node of $\lkevgraph{\tr}$ is an abstract acquire of $\tr$, and thus
every deadlock pattern in $\tr$ appears in some abstract deadlock pattern defined by a cycle in $\lkevgraph{\tr}$.
Hence, we can detect \emph{all} sync-preserving deadlocks in $\tr$ by enumerating each of the cycles of $\lkevgraph{\tr}$, and if it constitutes an abstract deadlock pattern, use \lemref{abstract-pattern-linear-time} to check if it contains a sync-preserving deadlock in $\Otilde(\NumEvents)$ time.
% Hence, in linear time, our algorithm checks a whole abstract deadlock pattern, that typically summarizes polynomially many concrete deadlock patterns. 
\end{comment}

We now present the final ingredients of \SyncPDOffline.
We construct the \emph{abstract lock graph},
enumerate cycles in it, check whether
any cycle is an abstract deadlock pattern,
and if so, whether it contains sync-preserving deadlocks.

%!TEX root=../main.tex

\setlength{\textfloatsep}{10pt}
\begin{figure}[t]
%\begin{wrapfigure}{r}{0.35\textwidth}
\begin{subfigure}[t]{0.2\textwidth}
\scalebox{0.84}{
\begin{tikzpicture}[yscale=0.9]
	\tikzset{rectangle/.append style={draw=black}}
		\def\xstep{2.2}
		\def\ystep{0.5}
		\node[rectangle, align=center] (t1) at (-1, -7){
			$t_1, \LockColorTwo{\lk_2}, \{ \LockColorOne{\lk_1}\}, \sequence{e_{2}}$
		};
		\node[rectangle, align=center] (t31) at (-1, -8.7){
			$t_2, \LockColorOne{\lk_1}, \{ \LockColorTwo{\lk_2} \}, \sequence{e_{8}}$
		};
		\draw[->, thick, -Latex, bend right=10] ([xshift=-3mm]t1.south) to ([xshift=-3mm]t31.north);% (t31.north);
		\draw[->, thick, -Latex, bend right=10] (t31) to (t1);
\end{tikzpicture}
}
\end{subfigure}
~~~~
\begin{subfigure}[t]{0.2\textwidth}
\scalebox{0.84}{
\begin{tikzpicture}[yscale=0.9]
	\tikzset{rectangle/.append style={draw=black}}
		\def\xstep{2.2}
		\def\ystep{0.5}
		\node[rectangle, align=center] (t1) at (-1, -7){
			$t_2, \LockColorThree{\lk_3}, \{ \LockColorTwo{\lk_2}\}, \sequence{e_{4}}$
		};
		\node[rectangle, align=center] (t31) at (-1, -8.7){
			$t_3, \LockColorTwo{\lk_2}, \{ \LockColorThree{\lk_3} \}, \sequence{e_{18}}$
		};
		\draw[->, thick, -Latex, bend right=10] ([xshift=-3mm]t1.south) to ([xshift=-3mm]t31.north);% (t31.north);
		\draw[->, thick, -Latex, bend right=10] (t31) to (t1);
\end{tikzpicture}
}
\end{subfigure}
~~~~
% \hspace*{\fill}
\begin{subfigure}[t]{0.4\textwidth}
\scalebox{0.84}{
\begin{tikzpicture}[yscale=0.9]
	\tikzset{rectangle/.append style={draw=black}}
		\node[rectangle, align=center] (t1) at (-1, -7){
					$t_1, \LockColorTwo{\lk_2}, \{ \LockColorOne{\lk_1} \}$,
					$\sequence{e_2, e_4, e_{29}}$
				};
				\node[rectangle,align=center] (t2) at (2.5, -7){
					$t_2, \LockColorOne{\lk_1}, \{ \LockColorFour{\lk_4} \}$,
					$\sequence{e_{23}}$
				};
				\node[rectangle, align=center] (t31) at (-1, -8.7){
					$t_3, \LockColorOne{\lk_1}, \{ \LockColorTwo{\lk_2} \}$,
					$\sequence{e_{16}, e_{19}}$
				};
				\node[rectangle,align=center] (t32) at (2.5, -8.7){
					$t_3, \LockColorThree{\lk_3}, \{ \LockColorTwo{\lk_2} \}$,
					$\sequence{e_{13}}$
				};
				\draw[->, thick, -Latex] (t2) -- (t1);
				\draw[->, thick, -Latex] ([xshift=6mm]t1.south) to (t32.west);
				\draw[->, thick, -Latex, bend right=10] ([xshift=-3mm]t1.south) to ([xshift=-3mm]t31.north);
				\draw[->, thick, -Latex, bend right=10] (t31) to (t1);
\end{tikzpicture}
}
\end{subfigure}
% 
% \caption{
% A  trace $\tr$ with a sync-preserving deadlock, stalling $t_2$ on $e_4$ and $t_3$ on $e_{18}$.
% Underlined events indicate the slice of $\tr$ that forms a sync-preserving lower set of $\tr$,
% and serves as a valid witness for the deadlock.
% \figlabel{motivating}
% }
\caption{
Abstract lock graphs of the traces from~\figref{motivating-no-dl} (left), \figref{motivating-dl} (middle) and \figref{syncp_example} (right).
% \Andreas{@Umang: You can use this space for more figures, examples}
\figlabel{motivating-graph}
}
 \end{figure}
%\end{wrapfigure}
\myparagraph{Abstract lock graph}{
The abstract lock graph of $\tr$ is a directed graph
$
\lkevgraph{\tr} = (V_\tr, E_\tr)
$,
where
\begin{itemize}
	\item $V_\tr=\{\tuple{t_1, \lk_1, L_1, F_1},\dots, \tuple{t_k, \lk_k, L_k, F_k}\}$ is the set of abstract acquires of $\tr$, and
	\item for every $\eta_1 {=} \tuple{t_1, \lk_1, L_1, F_1}, 
	\eta_2 {=} \tuple{t_2, \lk_2, L_2, F_2} \in V_\tr$, we have  $(\eta_1, \eta_2) \in E_\tr$ iff 
	$t_1 \neq t_2$, $\lk_1 \in L_2$, and $L_1 \cap L_2 = \emptyset$.
\end{itemize}
A node $\tuple{t_1, \lk_1, L_1, F_1}$ signifies that there is an event $\acq_1(\lk_1)$ performed by thread $t_1$ while holding the locks in $L_1$.
The last component $F_1$ is a list which contains all such events $\acq_1$ in order of appearance in $\tr$.
An edge $(\eta_1, \eta_2)$ signifies that the lock $\lk_1$ acquired by each of the
events $\acq_1 \in F_1$ was held by $t_2$ when it executed each of $\acq_2\in F_2$ while not holding a common lock.
%See \figref{lock_graph} for the lock graph to the trace $\tr_2$ of \figref{syncp_example}.
The abstract lock graph can be constructed incrementally
as new events appear in $\tr$.
For $\NumEvents$ events, $\NumLocks$ locks and nesting depth $\NestingDepth$,
the graph has
$|V_\tr| = O\big(\NumThreads \cdot \NumLocks^{\NestingDepth}\big)$ vertices,
$|E_\tr| = O(|V_\tr|\cdot \NumLocks^{\NestingDepth - 1})$ edges and can be
constructed in $O(\NumEvents \cdot \NestingDepth)$ time.
See \figref{motivating-graph} for examples.
In the left graph, the cycle marks an abstract deadlock pattern and its single concrete deadlock pattern
$\abst{D}=\{ e_2 \} \times \{ e_{8} \}$, and similarly for the middle graph where $\abst{D}=\{ e_4 \} \times \{ e_{18} \}$.
%In the graph on the right, we have the abstract lock graph of the trace $\tr_3$ of  \figref{syncp_example}.
In the right graph, there is a unique cycle which marks an abstract deadlock pattern of $6$ concrete deadlock patterns
$\abst{D}=\{ e_2, e_4, e_{29} \} \times \{ e_{16}, e_{19} \}$.


% \begin{example}
% \figref{lock_graph} illustrates the lock graph of the trace $\tr_2$ in \figref{syncp_example}.
% We have a unique cycle, of length $2$, that marks a corresponding abstract deadlock pattern of $6$ concrete deadlock patterns
% $\abst{D}=\{ e_2, e_4, e_{29} \} \times \{ e_{16}, e_{19} \}$.
% \end{example}
}
%!TEX root = main.tex
\setlength{\textfloatsep}{10pt}
\begin{algorithm}[h]
\small
\DontPrintSemicolon
\SetInd{0.4em}{0.4em}
\KwIn{
A trace $\tr$.
}
\KwOut{All abstract deadlock patterns of $\tr$ that contain a sync-preserving deadlock.}
\BlankLine
Construct the abstract lock graph $\lkevgraph{\tr}$\\
\ForEach{cycle $C=\tuple{\eta_0,\dots, \eta_{k-1}}$ in $G$}{
Let $\eta_i=\tuple{t_i, \lk_i, L_i, F_i}$\\
\uIf(\tcp*[h]{$C$ is an abstract deadlock pattern}){$\forall i\neq j$ we have $t_i\neq t_j$ and $\lk_i\neq \lk_j$ and $L_i\cap L_j=\emptyset$ }{
%
\lIf{$\checkAbsDeadP(C)$}{
%Report the deadlock pattern in $C$ that constitutes a sync-preserving deadlock
Report that $C$ contains a sync-preserving deadlock
}
}
}
\caption{
Algorithm $\SyncPDOffline$.
\algolabel{offline}
}
\end{algorithm}
\normalsize

% \vspace{-0.2in}
\myparagraph{Algorithm $\SyncPDOffline$}{
	It is straightforward to verify that every abstract deadlock pattern of $\tr$ appears as a (simple) cycle in $\lkevgraph{\tr}$.
	However, the opposite is not true.
	A cycle $C = \eta_0, \eta_1, \ldots, \eta_{k-1}$ of $\lkevgraph{\tr}$,
	where $\eta_i = \tuple{t_i, \lk_i, L_i, F_i}$ defines an abstract deadlock pattern
	if additionally every thread $t_i$ is distinct, all every lock $\lk_i$ is distinct, and all sets $L_i$ are pairwise disjoint.
	This gives us a simple recipe for enumerating all abstract deadlock patterns, by using
	Johnson's algorithm~\cite{Johnson1975} to enumerate
	every simple cycle $C$ in $\lkevgraph{\tr}$, and check whether $C$ is an abstract deadlock pattern.
	We thus arrived at our offline algorithm $\SyncPDOffline$ (\algoref{offline}).
	The running time depends linearly on the length of $\tr$
	and the number of cycles in $\lkevgraph{\tr}$.% (\thmref{syncp_offline}).
}




\begin{restatable}{theorem}{syncpoffline}
\thmlabel{syncp_offline}
Consider a trace $\tr$ of $\NumEvents$ events, $\NumThreads$ threads and $\textsf{Cyc}_\tr$ cycles in $\lkevgraph{\tr}$.
The algorithm $\SyncPDOffline$ reports all sync-preserving deadlocks of $\tr$ in time $O(\NumEvents \cdot \NumThreads\cdot \textsf{Cyc}_\tr)$.
% When the number of threads and lock-nesting depth is small, $\SyncPDOffline$ runs in $\Otilde(\NumEvents \cdot \poly{\NumLocks})$ time.
\end{restatable}

Although, in principle, we can have exponentially many cycles in $\lkevgraph{\tr}$,
because the nodes of $\lkevgraph{\tr}$ are \emph{abstract} acquire events (as opposed to \emph{concrete}), we expect that the number of cycles (and thus abstract deadlock patterns) in $\lkevgraph{\tr}$ remains small, even though the number of \emph{concrete} deadlock patterns can grow exponentially.
Since $\SyncPDOffline$ spends linear time per abstract deadlock pattern, we  have an efficient procedure overall 
for constant $\NumThreads$ and $\NumLocks$.
We evaluate $\textsf{Cyc}_\tr$ experimentally in \secref{experiments}, and confirm that it is very small compared to the number of concrete deadlock patterns in $\tr$.
%\hunkar{
Nevertheless, $\textsf{Cyc}_\tr$ can become exponential when $\NumThreads$ and $\NumLocks$ are large, making~\cref{algo:offline} run in exponential time.
Note that this barrier is unavoidable in general, as proven in \thmref{pattern-w1-hardness-pattern}.
%}

%!TEX root = main.tex

\begin{example}
    \exlabel{ex4}
    We illustrate how the lock graph is integrated inside $\SyncPDOffline$.	
    Consider the trace $\tr_3$ in \figref{syncp_example}.
    It contains $6$ concrete deadlock patterns $D_1 \ldots D_6$.
    A naive algorithm would enumerate each pattern explicitly until it finds a deadlock. 
    However, the tight interplay between the abstract lock graph and sync-preservation enables a more efficient procedure.
    $\SyncPDOffline$ starts by computing the sync-preserving closure of $D_1$,
    $\SPClosure{\tr_3}(\prev{\tr_3}(\{ e_{2},e_{16}\}))=\set{e_1, \ldots, e_6, \, e_8, \ldots, e_{15}}$.
    As $e_2 \in \SPClosure{\tr_3}(\prev{\tr_3}(\{ e_{2},e_{16}\}))$, we conclude that $D_1$ is not a sync-preserving deadlock.
    The algorithm further deduces that the deadlock patterns $D_2$, $D_3$ and $D_4$ are also not sync-preserving deadlocks, as follows.
    $D_2=\pattern{e_2, e_{19}}$ shares a common event $e_2$ with $D_1$ but contains the event $e_{19}$ instead of $e_{16}$, while $e_5 \in \SPClosure{\tr_3}(\prev{\tr_3}(\{ e_{2},e_{16}\}))$.
    Since $e_{16} \tho{\tr_3} e_{19}$, and the sync-preserving closure grows monotonically (\propref{spclosure-monotone}), the sync-preserving closure of $e_2$ and $e_{19}$ will also contain $e_5$ (and thus $e_2$).
    Therefore, $D_2$ cannot be a sync-preserving deadlock.
    This reasoning is formalized in \corref{pattern-monotone}, and also applies to $D_3$ and $D_4$.
    Next, the algorithm proceeds with $D_5$.
    The above reasoning does not hold for $D_5$ as $\SPClosure{\tr_3}(\prev{\tr_3}(\{ e_{2},e_{16}\})) \cap S_5 = \emptyset$ 
    where $S_5=\set{e_{29}, e_{16}}$.
    The algorithm then computes the sync-preserving closure of $D_5$, reports a deadlock (\exref{syncp-ex}) and stops analyzing this abstract deadlock pattern.
    %Hence, $D_6$ is also not considered.
    %As the remaining pattern $D_6$ originates from the same abstract deadlock pattern, it need not be considered.
    In the end, we have only explicitly enumerated the deadlock patterns $D_1$ and $D_5$.
\end{example}

\begin{remark}
Although the concept of lock graphs exists in the literature~\cite{Havelund2000,Bensalem2005,Cai2020,cai14magiclock},
our notion of \emph{abstract} lock graphs is novel and tailored to sync-preserving deadlocks.
The closest concept to abstract lock graphs is that of equivalent cycles~\cite{cai14magiclock}. 
However, equivalent cycles unify all the concrete patterns of a given abstract pattern and lead to unsound deadlock detection, which was indeed their use. 
\end{remark}
%%!TEX root = main.tex
\subsection{Overall Algorithm}
\seclabel{overall-offline}





%!TEX root = main.tex

\section{On-the-fly Deadlock Prediction}
\seclabel{otf}

Although $\SyncPDOffline$ is efficient, both theoretically (\thmref{syncp_offline})
and in practice (Table~\ref{tab:expr-results}),
it runs in two passes, akin to other predictive deadlock-detection methods~\cite{Kalhauge2018,Cai2021}.
%In other words, it is \emph{offline} and requires access to the input trace $\tr$ upfront.
In a runtime monitoring setting, it is desirable to operate in an \emph{online} fashion.
Recall that $\checkAbsDeadP(\cdot)$ indeed operates online (\secref{verify-abstract-patterns}),
while the offline nature of $\SyncPDOffline$
is tied to the offline construction of the abstract lock graph $\lkevgraph{\tr}$.
% , and performs cycle detection in order to find abstract deadlock patterns.
%We note that $\lkevgraph{\tr}$ may also be maintained online, and be periodically checked for cycles, as, e.g., in~\cite{Havelund2000,Agarwal2005,Cai2020}.
%However, this operation becomes costly when there are many cycles.
%as there can be many cycles of short length, while long cycles take long to discover.
To achieve the golden standard of online, linear-time, 
sound deadlock prediction, we focus on deadlocks of size $2$.
This focus is barely restrictive as most deadlocks in the wild have size $2$~\cite{Lu08}.
Further, deadlocks of size $2$ enjoy the following computational benefits:
\begin{enumerate*}[label=(\alph*)]
\item cycles of length $2$ can be detected instantaneously without performing graph traversals, and
\item \emph{every} cycle of length $2$ is an abstract deadlock pattern.
\end{enumerate*}


\myparagraph{Algorithm $\SyncPDOnline$}{
% We are now ready to outline our algorithm $\SyncPDOnline$ for the online detection of all sync-preserving deadlocks of size 2. 
%For completeness we refer to \appref{app_online} for the pseudocode.
The algorithm $\SyncPDOnline$ maintains all abstract acquires of the form $\eta=\tuple{t, \lk_1, \{\lk_2\}, F}$,
i.e., we only focus on one lock $\lk_2$ that is protecting each such acquire.
When a new acquire event $e=\ev{t,\acq(\lk_1)}$ is encountered, 
the algorithm iterates over all the locks $\lk_2\in\lheld{\tr}(e)$ that are held in $e$,
and append the event $e$ to the sequence $F$ of the corresponding abstract acquire $\eta=\tuple{t, \lk_1, \{ \lk_2 \}, F}$; $F$ is maintained as FIFO queue.
Recall that we use timestamps on the acquire events in $F$ to determine membership in our closure computation.
Our online algorithm computes these timestamps on-the-fly and stores them in these queues together with the events.
%together
%, instead of storing the actual event.
% \Andreas{@Umang: I am missing some logic about vector clocks here, please update accordingly}
Then the algorithm calls $\checkAbsDeadP(C)$ on the abstract deadlock pattern $C$ formed between $\eta$ and $\eta'=\tuple{t'\neq t, \lk_2, \{ \lk_1 \}, F'}$, in order to check for sync-preserving deadlocks between the deadlock patterns in $F\times F'$.
If $\checkAbsDeadP(C)$ reports no deadlock, the contents of $F'$ are emptied
%as the monotonicity of \corref{pattern-monotone} 
as we are guaranteed that $F'$ will not cause a sync-preserving deadlock with any further acquire of thread $t$ on lock $\lk_1$.
% 
For a trace with $\NumEvents$ events, $\NumThreads$ threads and $\NumLocks$ locks.
The algorithm calls $\checkAbsDeadP$ for each of the $O(\NumThreads^2\cdot \NumLocks^2)$ abstract
deadlock patterns of size $2$, each call taking $O(\NumEvents\cdot \NumThreads)$ time.

%!TEX root = main.tex

\small
\begin{algorithm*}[t!]
% \BlankLine
 \vspace*{-0.35cm}
\begin{multicols}{2}
\iffalse
\myfun{\maxLB{$U$, $Lst$}}{
	$(g_{\max}, C_{\max}, C'_{\max})$ := $(0, \bot, \bot)$ \; 
	\While{$\NOT\, \isEmpty{Lst}$}{
		$(g, C, C')$ := $\first{Lst}$ \;
		\If{$C \cle U$}{
			$(g_{\max}, C_{\max}, C'_{\max})$ := $(g, C, C')$
		}
		\lElse{
			\Break
		}
		$\removeFirst{Lst}$
	}
	\Return $(g_{\max}, C_{\max}, C'_{\max})$
}

\BlankLine

\myfun{\fixpoint{$I$, $\tuple{t_1, \lk_1, t_2, \lk_2}$}}{
	\Repeat{$I$ does not change}{
		\ForEach{$\lk \in \locks{}$}{
			\ForEach{$t \in \threads{}$}{
				($g_{\lk, t}, C_{\lk, t}, C'_{\lk, t}$) := 
				\maxLB{$I$, $\view{\AcqLst}{t, \lk}{\tuple{t_1, \lk_1, t_2, \lk_2}}$} \;
			}
			$t_{\max}$ := argmax$_{t \in \threads{}}\set{g_{\lk, t}}$ \;
			$I$ := $I \mx \bigsqcup_{t \neq t_{\max} \in \threads{}} C'_{\lk, t}$	
		}
	}
	\Return $I$
}
\fi

\myfun{\checkDeadlock{$Lst$, $I$, $\tuple{t_1, \lk_1, t_2, \lk_2}$}}{
	\While{$\NOT\, \isEmpty{Lst}$}{
		($C_{\prev{}}, C$) := $\first{Lst}$ \;
		$I$ := \fixpoint{$I \mx C_{\prev{}}$, $\tuple{t_1, \lk_1, t_2, \lk_2}$} \;
		\lIf{$C \not\cle I$}{ \linelabel{check-containment}
			\declare `Deadlock' and \Break
		}
		$\removeFirst{Lst}$
	}
	\Return $I$
}

\myhandler{\wthandler{$t$, $x$}}{
	$\LW_x$ := $\Cc_t$ \linelabel{clkwt}\\
	$\Cc_t$ := $\Cc_t[t \mapsto \Cc_t(t)+1]$ \linelabel{clkctw}
}

\myhandler{\relhandler{$t$, $\lk$}}{
	$\Cc_t$ := $\Cc_t[t \mapsto \Cc_t(t)+1]$ \linelabel{clkrel} \;
	\ForEach{$t_1, t_2 \in \threads{}, \lk_1, \lk_2 \in  \locks{}$}{
		$\updateRel{\lst{\view{\AcqLst}{t, \lk}{\tuple{t_1, \lk_1, t_2, \lk_2}}}}{$\Cc_t$}$
	}
}

\myhandler{\acqhandler{$t$, $\lk$}}{
	$C_{\prev{}}$ := $\Cc_t$ \;
	$\Cc_t$ := $\Cc_t[t \mapsto \Cc_t(t)+1]$\linelabel{clkacq}\\ $\gId_\lk$ := $\gId_\lk + 1$ \linelabel{gidlk}\;
	\ForEach{$t_1, t_2 \in \threads{}, \lk_1, \lk_2 \in \locks{}$}{
		$\addLst{\view{\AcqLst}{t, \lk}{\tuple{t_1, \lk_1, t_2, \lk_2}}}{$(\gId_\lk, \Cc_t, \bot)$}$
	}
	\ForEach{$u\in \threads{}$, $\lk' \in \lheld{}$}{
		$\addLst{\view{\DPLst}{t, \lk, \lk'}{\tuple{u}}}{($C_{\prev{}}$, $\Cc_t$)}$ \linelabel{acqhist}
	}
	\ForEach{$u \neq t \in \threads{}$, $\lk' \in \lheld{}$}{
		$I$ := $\view{\Ii}{}{\tuple{u, \lk' t, \lk}} \mx C_{\prev{}}$ \;
		$\view{\Ii}{}{\tuple{u, \lk' t, \lk}}$ := \checkDeadlock{$\view{\DPLst}{u, \lk', \lk\,}{\tuple{t}}$, $I$, $\tuple{u, \lk' t, \lk}$} \linelabel{checkdlk}
	}

}

\myhandler{\rdhandler{$t$, $x$}}{
	$\Cc_t$ := $\Cc_t \mx \LW_x$ \linelabel{clkrd}
}

\BlankLine

\end{multicols}
 \vspace*{-0.25cm}
\normalsize
\caption{$\SyncPDOnline$.}
\algolabel{online}
\end{algorithm*}
\normalsize

$\SyncPDOnline$ is shown in detail in \algoref{online}.
The pseudocode contains handlers for processing the different events of $\tr$ in a streaming fashion, as well as a helper function for checking deadlocks.
The main data structures of the algorithm are
(i)~vector clocks $\Cc_t$, $\LW_x$, and $\view{\Ii}{}{\tuple{t_1, \lk_1, t_2, \lk_2}}$,
(ii)~scalar $\gId_\lk$, and
(iii)~FIFO queues of vector clocks $\view{\AcqLst}{t, \lk}{\tuple{t_1, \lk_1, t_2, \lk_2}}$ and $\view{\DPLst}{t, \lk_1, \lk_2}{\tuple{u}}$,
where $t, t_1, t_2, u$ range over threads, $x$ ranges over variables, and $\lk, \lk_1, \lk_2$ range over locks.
$\Cc_t$ stores the timestamp $\TS{\tr}^e$ where $e$ is the last event in thread $t$.
%(Lines \ref{line:clkctw},  \ref{line:clkrel}, \ref{line:clkacq}, \ref{line:clkrd}).
$\LW_x$ keeps track of the $\TS{\tr}^e$ where $e$ is the last event such that $\OpOf{e} = \wt(x)$. %(\cref{line:clkwt}).
$\view{\Ii}{}{\tuple{t_1, \lk_1, t_2, \lk_2}}$ stores the computed sync-preserving closures for every tuple $\tuple{t_1, \lk_1, t_2, \lk_2}$.
%(\cref{line:checkdlk}).
The scalar variable $\gId_\lk$ keeps track of the index of the last acquire event on lock $\lk$. %(\cref{line:gidlk}).
Similar to \cref{algo:compute-closure}, the FIFO queue $\view{\AcqLst}{t, \lk}{\tuple{t_1, \lk_1, t_2, \lk_2}}$ is maintained to keep track of the critical section history of thread $t$ and lock $\lk$.
Lastly, for an acquire event $e$, $\view{\DPLst}{t, \lk_1, \lk_2}{\tuple{u}}$ maintains a queue of tuples of the form $\tuple{C_{\prev{}}, \Cc_t}$ where $C_{\prev{}}$ and $\Cc_t$ are the timestamps of $\prev{\sigma}(e)$ and $e$, respectively
%(\cref{line:acqhist}).
These tuples are utilized when checking for deadlocks (\cref{line:checkdlk}).
}

\begin{restatable}{theorem}{syncponline}
\thmlabel{syncp_online}
Consider a trace $\tr$ of $\NumEvents$ events, $\NumThreads$ threads and $\NumLocks$ locks.
The online $\SyncPDOnline$ algorithm reports all sync-preserving deadlocks of size $2$ of $\tr$ in $O(\NumEvents\cdot  \NumThreads^3\cdot \NumLocks^2)$ time.
%$O(\NumEvents\cdot \NumThreads+ \NumThreads^3\cdot \NumLocks^2\cdot \NumAcquires)$ time.
\end{restatable}


%!TEX root = main.tex

\section{Experimental Evaluation}
\seclabel{experiments}

We first evaluated our algorithms
in an offline setting~(\secref{offline-expr}), where we record execution traces and evaluate different approaches on the \emph{same} input.
This eliminates biases due to non-deterministic thread scheduling.
Next, we consider an online setting~(\secref{online-expr}),
where we instrument programs and perform the analyses during runtime.
We conducted all our experiments on a standard laptop with \SI{1.8}{GHz} Intel Core i7 processor and \SI{16}{GB} RAM.

% We evaluated our algorithms in two experimental settings. 
% The first setting is offline experiments~(\secref{offline-expr}),
% in which we record execution traces and evaluate different approaches.
% This has the benefit that different approaches can be compared on the \emph{same} input,
% thereby eliminating biases due to non-deterministic thread scheduling.
% The second setting is online experiments~(\secref{online-expr}),
% in which we instrument programs and perform the analyses during runtime.
% We conducted all our experiments on a standard laptop with 1.8GHz Intel Core i7 processor and 16GB RAM.

%We compare our predictive algorithm with the \dlfuzzer tool. 
%This setting enables us to assess the applicability of our technique in a runtime monitoring system.
%The state-of-the-art deadlock predictors \dirk and \seqc work offline by design.
%Hence, they are not applicable for a comparison in this setting.
%As discussed in~\secref{otf}, offline methods are not directly applicable in runtime monitoring systems.

%Talk about implementation - name of tool + prog lang + traces are logged + filtering phase + cycle detection phase + online vs offline + trace conversion for seqc
%
%Setup - benchmarks + log traces using so-and-so-tools + 1-trace-per-benchmark  + cluster details + how many runs per trace + Timeout (if any)
%
%benchmark details: number of shared locks
%\subsection{Evaluation}
%
%Some suggested experiments:
%
%\begin{itemize}
%	\item Comparison with other tool(s) - \dirk and \seqc
%	\item Comparison of online vs online algorithms
%	\item Scalability with number of events
%	\item Scalability with number of threads in the trace
%	\item Scalability with number of threads in the deadlock pattern
%	\item Time spent in each kinds of events (this will be useful to guide future research)
%\end{itemize}
%
%\Andreas{For matching reports, we slightly adapt the algorithm to report all sets of program locations that contain a deadlock pattern (hence we might have more than one reports per abstract deadlock pattern, corresponding to different program locations)}


% Here we report on an implementation and experimental evaluation of our algorithms.

% \subsection{Experimental Setup}



\subsection{Offline Experiments}
\seclabel{offline-expr}


\Paragraph{Experimental setup}
The goal of the first set of experiments is to evaluate 
$\SyncPDOffline$, and compare it
against prior algorithms for dynamic deadlock prediction.
In order for our evaluation to be precise we evaluate all algorithms on the \emph{same} execution trace.
We implemented $\SyncPDOffline$ in Java inside the \toolname analysis tool~\cite{rapid}, 
following closely the pseudocode in \algoref{offline}.
\toolname takes as input execution traces, as defined in \secref{prelim}.
These also include fork, join, and lock-request events.
We compare $\SyncPDOffline$ with two state-of-the-art, 
theoretically-sound albeit computationally more expensive, deadlock predictors,
\seqc~\cite{Cai2021} and \dirk~\cite{Kalhauge2018}, both of which also work on execution traces.




On the theoretical side, the complexity of \seqc is $\Otilde(\NumEvents^4)$, 
as opposed to the $\Otilde(\NumEvents)$ complexity of $\SyncPDOffline$. 
Moreover, \seqc only predicts deadlocks of size $2$, and though it could be extended to handle deadlocks of any size, this would degrade performance further.
\seqc may miss sync-preserving deadlocks even of size $2$, 
but can  detect deadlocks that are not sync-preserving.
Thus \seqc and $\SyncPDOffline$ are theoretically incomparable in their detection capability.
We refer to\begin{pldi}~\cite{arxiv}\end{pldi}\begin{arxiv}~\appref{incomp}\end{arxiv} for examples.
We noticed that \seqc fails on traces with non-well-nested locks --- we encountered one such case in our dataset.
\dirk's algorithm is theoretically complete, i.e., it can find all predictable deadlocks in a trace.
In addition, it can find deadlocks beyond the predictable ones, by reasoning about event values.
However, \dirk relies on heavyweight SMT-solving and
employs windowing techniques to scale to large traces. 
Due to windowing, it can miss deadlocks between events that are outside the given window. 
%\hunkar{
As with previous works~\cite{Cai2021, Kalhauge2018}, we set a window size of $10$K for \dirk.
%}

Our dataset consists of several benchmarks 
from standard benchmark suites --- IBM Contest suite~\cite{Farchi03}, Java Grande suite~\cite{Smith01},
DaCapo~\cite{Blackburn06}, and
SIR~\cite{doESE05} ---
and recent literature~\cite{Kalhauge2018, Cai2021, jula2008deadlock, Joshi2009}.
Each benchmark was instrumented with RV-Predict~\cite{rvpredict} or Wiretap~\cite{Kalhauge2018} and
executed in order to log a single execution trace.

%!TEX root = ../main.tex

\begin{table}[h!]
\caption{
%\hcomment{Should we remove some not very important rows from this table?}
Trace characteristics, abstract lock graph statistics and performance comparison.
%Column 1 denotes the name of the benchmark.
Columns 2-6 show the number of events, threads, variables, locks
and total number of lock acquire and request events.
Columns 7-9 show the number of cycles, abstract and concrete deadlock patterns in the abstract lock graph.
Columns 10 - 15 show the number of deadlocks reported and the times (in seconds) taken. 
by \dirk, \seqc, and \SyncPDOffline.
%Statistics on the lock graph $\lkevgraph{\tr}$ of each trace $\tr$.
% We denote by $\mathcal{N}$, $\mathcal{T}$, $\mathcal{M}$, $\mathcal{L}$ and $\NumAcquires + \mathcal{R}$ the total number of events, number of threads, number of memory locations, number of locks and number of acquire and request events in the benchmarks, respectively. 
%All times are in seconds. 
Time out (T.O) was set to $3$h.
F stands for technical failure.
\label{tab:expr-results}
}
\vspace{-0.17cm}
\setlength\tabcolsep{3pt}
\renewcommand{\arraystretch}{0.91}
\centering
\scalebox{0.86}{
\begin{tabular}{|r|c|c|c|c|c||c|c|c||c|c||c|c||c|c|}
\hline
1 & 2 & 3 & 4 & 5 & 6 & 7 & 8 & 9 & 10 & 11 & 12 & 13 & 14 & 15 \\
\hline
\multirow{2}{*}{\textbf{Benchmark}}& 
\multirow{2}{*}{$\mathcal{N}$} & \multirow{2}{*}{$\mathcal{T}$} & \multirow{2}{*}{$\mathcal{V}$} & \multirow{2}{*}{$\mathcal{L}$} & \multirow{2}{*}{$\NumAcquires/\mathcal{R}$} 
& \multicolumn{3}{c||}{ \textsf{A. Lock Graph }}
& \multicolumn{2}{c||}{ {\dirk} } 
& \multicolumn{2}{c||}{ {\seqc}} 
& \multicolumn{2}{c|}{ {\textsf{$\SyncPDOffline$}}} \\
\cline{7-15}
& & & & & 
& \textsf{|$\textsf{Cyc}$|}
& \textsf{\textsf{A. P.}}
& \textsf{\textsf{C. P.}}
& \textbf{Dlk} 
& {\textbf{Time}} 
&{\textbf{Dlk}} 
& {\textbf{Time}}
&{\textbf{Dlk}} 
&{\textbf{Time}} \\
\hline
Deadlock & 39 & 3 & 4 & 3 & 8 & 1 & 1 & 1 & 1 & 0.02 & 0 & 0.09 & 0 & 0.16\\
NotADeadlock & 60 & 3 & 4 & 5 & 16 & 1 & 1 & 1 & 0 & 0.02 & 0 & 0.09 & 0 & 0.16\\
Picklock & 66 & 3 & 6 & 6 & 20 & 2 & 2 & 2 & 1 & 0.02 & 1 & 0.10 & 1 & 0.18\\
Bensalem & 68 & 4 & 5 & 5 & 22 & 2 & 2 & 2 & 1 & 0.02 & 1 & 0.12 & 1 & 0.16\\
Transfer & 72 & 3 & 11 & 4 & 12 & 1 & 1 & 1 & 1 & 0.02 & 0 & 0.09 & 0 & 0.15\\
Test-Dimmunix & 73 & 3 & 9 & 7 & 26 & 2 & 2 & 2 & 2 & 0.02 & 2 & 0.10 & 2 & 0.17\\
StringBuffer & 74 & 3 & 14 & 4 & 16 & 1 & 3 & 6 & 2 & 0.02 & 2 & 0.12 & 2 & 0.19\\
Test-Calfuzzer & 168 & 5 & 16 & 6 & 48 & 2 & 1 & 1 & 1 & 0.02 & 1 & 0.12 & 1 & 0.17\\
DiningPhil & 277 & 6 & 21 & 6 & 100 & 1 & 1 & 3K & 1 & 1.60 & 0 & 0.09 & 1 & 0.17\\
HashTable & 318 & 3 & 5 & 3 & 174 & 1 & 2 & 43 & 2 & 0.19 & 2 & 0.12 & 2 & 0.19\\
Account & 706 & 6 & 47 & 7 & 134 & 3 & 1 & 12 & 0 & 0.19 & 0 & 0.09 & 0 & 0.18\\
Log4j2 & 1K & 4 & 334 & 11 & 43 & 1 & 1 & 1 & 1 & 0.65 & 1 & 0.11 & 1 & 0.20\\
Dbcp1 & 2K & 3 & 768 & 5 & 56 & 2 & 2 & 3 & - & F & 2 & 0.11 & 2 & 0.19\\
Dbcp2 & 2K & 3 & 592 & 10 & 76 & 1 & 2 & 4 & - & F & 0 & 0.10 & 0 & 0.18\\
Derby2 & 3K & 3 & 1K & 4 & 16 & 1 & 1 & 1 & 1 & 0.23 & 1 & 0.10 & 1 & 0.17\\
RayTracer & 31K & 5 & 5K & 15 & 976 & 0 & 0 & 0 & - & F & 0 & 0.15 & 0 & 0.19\\
jigsaw & 143K & 21 & 8K & 2K & 67K & 172 & 12 & 70 & - & F & 2 & 0.36 & 1 & 1.55\\
elevator & 246K & 5 & 727 & 52 & 48K & 0 & 0 & 0 & 0 & 1.65 & 0 & 0.33 & 0 & 0.27\\
hedc & 410K & 7 & 109K & 8 & 32 & 0 & 0 & 0 & 0 & 2.09 & 0 & 0.50 & 0 & 0.24\\
JDBCMySQL-1 & 442K & 3 & 73K & 11 & 13K & 2 & 4 & 6 & 2 & 28.45 & 2 & 0.24 & 2 & 0.48\\
JDBCMySQL-2 & 442K & 3 & 73K & 11 & 13K & 4 & 4 & 9 & 1 & 3.37 & 1 & 0.22 & 1 & 0.33\\
JDBCMySQL-3 & 443K & 3 & 73K & 13 & 13K & 5 & 8 & 16 & 1 & 31.23 & 1 & 0.25 & 1 & 0.45\\
JDBCMySQL-4 & 443K & 3 & 73K & 14 & 13K & 5 & 10 & 18 & 2 & 5.51 & 2 & 0.28 & 2 & 0.49\\
cache4j & 775K & 2 & 46K & 20 & 35K & 0 & 0 & 0 & 0 & 5.86 & 0 & 0.46 & 0 & 0.39\\
ArrayList & 3M & 801 & 121K & 802 & 176K & 9 & 3 & 672 & 3 & 8.7K & 3 & 21.98 & 3 & 1.68\\
IdentityHashMap & 3M & 801 & 496K & 802 & 162K & 1 & 3 & 4 & 1 & 443.93 & 1 & 8.51 & 1 & 1.45\\
Stack & 3M & 801 & 118K & 2K & 405K & 9 & 3 & 481 & 1 & T.O & 3 & 25.34 & 3 & 2.94\\
Sor & 3M & 301 & 2K & 3 & 719K & 0 & 0 & 0 & 0 & 15.89 & 0 & 44.12 & 0 & 0.61\\
LinkedList & 3M & 801 & 290K & 802 & 176K & 9 & 3 & 10K & 3 & 4.7K & 3 & 48.02 & 3 & 2.06\\
HashMap & 3M & 801 & 555K & 802 & 169K & 1 & 3 & 10K & 3 & 4.4K & 2 & 504.36 & 2 & 1.65\\
WeakHashMap & 3M & 801 & 540K & 802 & 169K & 1 & 3 & 10K & - & T.O & 2 & 499.68 & 2 & 1.70\\
Swing & 4M & 8 & 31K & 739 & 2M & 0 & 0 & 0 & - & F & 0 & 0.72 & 0 & 0.88\\
Vector & 4M & 3 & 15 & 4 & 800K & 1 & 1 & 1B & - & T.O & 1 & 1.52 & 1 & 1.90\\
LinkedHashMap & 4M & 801 & 617K & 802 & 169K & 1 & 3 & 10K & 2 & 40.74 & 2 & 492.87 & 2 & 1.69\\
montecarlo & 8M & 3 & 850K & 3 & 26 & 0 & 0 & 0 & 0 & 2.6K & 0 & 1.81 & 0 & 0.79\\
TreeMap & 9M & 801 & 493K & 802 & 169K & 1 & 3 & 10K & 2 & 105.45 & 2 & 480.11 & 2 & 1.92\\
hsqldb & 20M & 46 & 945K & 403 & 419K & 0 & 0 & 0 & - & F & - & - & 0 & 2.38\\
sunflow & 21M & 16 & 2M & 12 & 1K & 0 & 0 & 0 & - & F & 0 & 8.35 & 0 & 1.62\\
jspider & 22M & 11 & 5M & 15 & 10K & 0 & 0 & 0 & - & F & 0 & 8.49 & 0 & 1.95\\
tradesoap & 42M & 236 & 3M & 6K & 245K & 2 & 1 & 4 & - & F & 0 & 108.16 & 0 & 7.06\\
tradebeans & 42M & 236 & 3M & 6K & 245K & 2 & 1 & 4 & - & F & 0 & 116.23 & 0 & 7.26\\
eclipse & 64M & 15 & 10M & 5K & 377K & 9 & 5 & 280 & - & F & 0 & 26.67 & 0 & 9.90\\
TestPerf & 80M & 50 & 599 & 9 & 197K & 0 & 0 & 0 & 0 & 795.04 & 0 & 47.56 & 0 & 4.30\\
Groovy2 & 120M & 13 & 13M & 10K & 69K & 0 & 0 & 0 & 0 & 1.7K & 0 & 38.06 & 0 & 8.92\\
Tsp & 200M & 6 & 24K & 3 & 882 & 0 & 0 & 0 & 0 & 7.6K & 0 & 72.62 & 0 & 12.70\\
lusearch & 203M & 7 & 3M & 98 & 273K & 0 & 0 & 0 & 0 & 1.3K & 0 & 75.88 & 0 & 14.44\\
biojava & 221M & 6 & 121K & 79 & 16K & 0 & 0 & 0 & - & F & 0 & 63.79 & 0 & 12.65\\
graphchi & 241M & 20 & 25M & 61 & 1K & 0 & 0 & 0 & - & F & 0 & 102.05 & 0 & 25.25\\
\hline\hline\textbf{Totals} & \textbf{1B} & \textbf{7K} & \textbf{70M} & \textbf{37K} & \textbf{8M} & \textbf{256} & \textbf{93} & \textbf{1B} & \textbf{35} & \textbf{>18h} & \textbf{40} & \textbf{2801} & \textbf{40} & \textbf{135}\\\hline
\end{tabular}
}
\end{table}

\Paragraph{Evaluation}
\cref{tab:expr-results} presents our results.
A bug identifies a unique tuple of source
code locations corresponding to events participating in the deadlock.
%Columns 2-6 present the characteristics of our benchmark traces.
Trace lengths vary vastly from $39$ to about $241$M, while the number of threads ranges from $3$ to about $800$,
which are representative features of real-world settings.
\texttt{Hsqldb} contains critical sections that are not well nested, 
and \seqc was not able to handle this benchmark;
our algorithm does not have such a restriction.
%We were unable to run \dirk on certain benchmarks due to \dirk's technical issues, which are marked with F.


\vspace{-0.1cm}
\SubParagraph{\underline{Abstract vs Concrete Patterns}}
Columns 7-9 present statistics on the 
abstract lock graph $\lkevgraph{\tr}$ of each trace $\tr$.
Many traces have a large number of concrete deadlock patterns 
but much fewer abstract deadlock patterns;
a single abstract deadlock pattern can 
comprise up to an order of $10^4$ more concrete patterns (Column $8$ v/s Column $9$).
Unlike all prior sound techniques, 
our algorithms analyze 
abstract deadlock patterns, instead of concrete ones. 
We thus expect our algorithms to be much more scalable in practice.


\SubParagraph{\underline{Deadlock-detection capability}}
In total, both \seqc and $\SyncPDOffline$ reported 40 deadlocks.
\seqc misses a deadlock of size $5$ in \texttt{DiningPhil},
which is detected by $\SyncPDOffline$,
and $\SyncPDOffline$ misses a deadlock in \texttt{jigsaw} which is detected by \seqc.
As $\SyncPDOffline$ is complete for sync-preserving deadlocks, we conclude that there are no more such deadlocks in our dataset.
Overall, $\SyncPDOffline$ and \seqc miss only three deadlocks reported by \dirk. 
On closer inspection, we found that these deadlocks are not witnessed by correct reorderings, and require reasoning about event values.
On the other hand, \dirk struggles to analyze even moderately-sized benchmarks and times out in $3$ of them. %(timeout is 3h).
This results in \dirk failing to report 5 deadlocks after $9$ hours, all of which are reported by $\SyncPDOffline$ in under a minute.
Similar conclusions were recently made in~\cite{Cai2021}.  
Overall, our results strongly indicate that the notion of sync-preservation characterizes most deadlocks that other tools are able to predict.


\SubParagraph{\underline{Unsoundness of \dirk}}
In our evaluation, we discovered that the soundness guarantee 
underlying \dirk~\cite{Kalhauge2018} is broken, resulting in it reporting false positives.
% Although \dirk is portrayed as sound, we have found two independent sources of unsoundness 
% resulting in it reporting false positives.
First, its constraint formulation~\cite{Kalhauge2018} 
does not rule out deadlock patterns when acquire events in the pattern hold common locks, 
in which case mutual exclusion disallows such a pattern to be a real predictable deadlock.
Second, \dirk also models conditional statements, allowing it to reason about witnesses beyond correct reorderings.
While this relaxation allows \dirk to predict additional deadlocks in \texttt{Transfer}, \texttt{Deadlock} and \texttt{HashMap}, 
its formalization is not precise and its implementation is erroneous.
We elaborate these aspects further in\begin{pldi}~\cite{arxiv}. \end{pldi}\begin{arxiv}~\appref{unsound-dirk}.\end{arxiv}


%We noticed that \dirk reports false positives in certain cases, which contradicts its theoretical soundness guarantee.
%We provide two such cases in~\appref{unsound-dirk}.
%The first case is a modified version of \texttt{Transfer}.
%Here, \dirk falsely reports a deadlock because it inadequately models conditional statements, which are used to allow more trace reorderings that can expose a deadlock.
%This relaxation enables \dirk to predict deadlocks in the benchmarks \texttt{Transfer} and \texttt{Deadlock}, which are missed by $\SyncPDOffline$ and \seqc. 
%However, this relaxation is not formalized and its implementation is erroneous.
%In the second case, \dirk falsely reports a deadlock due to missing that a common lock protects an otherwise deadlock pattern.

\SubParagraph{\underline{Running time}}
Our experimental results indicate that \dirk, backed by SMT solving, 
is the least efficient technique in terms of running time ---
it takes considerably longer or times out on large benchmark instances.
$\SyncPDOffline$ analyzed the entire set of traces $\sim\!\!\!21\times$ faster than \seqc.
On the most demanding benchmarks, such as 
HashMap and TreeMap, $\SyncPDOffline$ is more than $200\times$ faster than \seqc.
Although \seqc employs a polynomial-time algorithm for deadlock prediction, 
and thus significantly faster than the SMT-based \dirk,
the large polynomial complexity in its running time hinders scalability on 
execution traces coming from benchmarks that are more representative of realistic workloads.
In contrast, the linear time guarantees of $\SyncPDOffline$ are realized in practice, 
allowing it to scale on even the most challenging inputs.
More importantly, the improved performance comes while preserving essentially the same precision.


%\begin{figure}[!h]
\def\scatterscale{0.192}
\def\scatterwidth{0.23\textwidth}

\centering
\begin{subfigure}[b]{\scatterwidth}
\includegraphics[scale=\scatterscale]{figures/LinkedListDeadlockTest.pdf}
\label{subfig:sca-one-lock-clique}
\caption{LinkedList}
\end{subfigure}
\begin{subfigure}[b]{\scatterwidth}
\includegraphics[scale=\scatterscale]{figures/StackDeadlockTest.pdf}
\caption{Stack}
\label{subfig:sca-skewed}
\end{subfigure}


\caption{
Scalability experiments.
}
\label{fig:scalability}
\end{figure}


%\Andreas{We might revisit/remove the following}

%\textcolor{red}{check the numbers again.}
\SubParagraph{\underline{False negatives}}
Our benchmark set contains $93$ abstract deadlock patterns, $40$ of which are confirmed sync-preserving deadlocks.
We inspected the remaining $53$ abstract patterns to see if any of them are predictable deadlocks
missed by our sync-preserving criterion, independently of the compared tools.
$48$ of these $53$ patterns are in fact not predictable deadlocks ---
for every such pattern $D$, 
the set $S_D$ of events in the downward-closure of $\prev{}(D)$ with respect to $\tho{}$ and $\rf{}$,
already contains an event from $D$, disallowing any correct reordering
(sync-preserving or not) in which $D$ can be enabled.
Of the remaining, $4$ deadlock patterns obey the following scheme:
there are two acquire events $\acq_1, \acq_2$ participating in the deadlock pattern, 
each $\acq_i$ is preceded by a critical section on a lock that appears in 
$\lheld{}(\acq_{3-i})$, again disallowing a correct reordering that witnesses the pattern.
Thus, \emph{only one} predictable deadlock is not sync-preserving in our whole dataset.
This analysis supports that the notion of sync-preservation is not overly conservative in practice.
%\hunkar{I think here we should clarify that this analysis is not overly conservative as far as predictable deadlocks go.}

%\hunkar{
The above analysis concerns false negatives wrt. predictable deadlocks.
Some deadlocks are beyond the common notion of predictability we have adopted here, as they can only be exposed by reasoning about event values and control-flow dependencies, a problem that is $\NP$-hard even for 3 threads~\cite{Gibbons1997}.
We noticed $3$ such deadlocks in our dataset, found by \dirk,
though, as mentioned above, \dirk's reasoning for capturing such deadlocks is unsound in practice.
%}

%\hunkar{Maybe stress again that this reasoning about event values is non-trivial as it relies on heavyweight SMT solving. Also, maybe have the section "Unsoundness of Dirk" after this one and say that next we show that it is also tricky to implement in practice.}
%sync-preserving deadlocks are likely to be permissive and not lead to a high false negative rate.
%We conduct an analyses based on our benchmark set and characterize potential false negatives of our technique.
%Recall the conditions imposed on correct reorderings (see \secref{prelim}).
%One of the main conditions imposed by the definition of a correct reordering is that
%every read event reads from the same write as in the original trace. 
%We investigated the effect of this condition on the potential false negatives.
%Our benchmark set contains $93$ abstract deadlock patterns and our technique is able to find $42$ deadlocks.
%The remaining $51$ deadlock patterns constitute potential false negatives.
%If we only impose the above condition, then $46$ of the deadlock patterns cannot be realized to a real deadlock.
%The additional conditions that are specific to sync-preserving deadlocks (e.g., the order of critical sections on the same lock cannot be reversed) prevents the remaining $5$ deadlock patterns from being realizable.
%We remark that the definition of correct reorderings adopted in this work originate from the standard model that is widely used in this domain~\cite{Smaragdakis12,serbanuta2013,Koushik05}.
%Hence, this analyses further supports that given this standard program model, the additional restrictions introduced with the
%sync-preserving deadlocks are likely to be permissive and not lead to a high false negative rate.


\subsection{Online Experiments}
\seclabel{online-expr}

\Paragraph{Experimental setup}
%Dynamic analyses can incur high runtime overheads.
The objective of our second set of experiments is to evaluate 
the performance 
of our proposed algorithms in an \emph{online} setting.
For this, we implemented our $\SyncPDOnline$ algorithm inside the
framework of \dlfuzzer~\cite{Joshi2009} following closely the pseudocode in \algoref{online}. 
This framework instruments a concurrent program so that it can
perform analysis on-the-fly while executing it.
If a deadlock occurs during execution, it is reported and the execution halts.
However, if a deadlock is predicted in an alternate interleaving, 
then this deadlock is reported and the execution continues to search further deadlocks.
We used the same dataset as in \secref{offline-expr}, 
after discarding some benchmarks that could not be instrumented by \dlfuzzer.

To the best of our knowledge, all prior deadlock prediction techniques work offline.
For this reason, we only compared our online tool with the randomized 
scheduling technique of~\cite{Joshi2009} already
implemented inside the same \dlfuzzer framework.	
At a high level, this random scheduling technique works as follows.
Initially, it
(i)~executes the input program with a random scheduler, 
(ii)~constructs a \emph{lock dependency relation}, and 
(iii)~runs a cycle detection algorithm to discover deadlock patterns. 
For each deadlock pattern thus found, it spawns new executions that attempt 
to realize it as an actual deadlock.
To increase the likelihood of hitting the deadlock,
\dlfuzzer biases the random scheduler by pausing threads at specific locations.
%(e.g., before acquiring a certain lock).

The second, confirmation phase of~\cite{Joshi2009}
acts as a best-effort proxy for sound deadlock prediction.
On the other hand, $\SyncPDOnline$ is already sound and predictive, and thus does not require
additional confirmation runs, making it more efficient.
Towards effective prediction, we also implemented a simple bias to the scheduler.
If a thread $t$ attempts to write on a shared variable $x$ while holding a lock, then 
our procedure randomly decides to pause this operation for a short duration.
This effectively explores race conditions in different orders.
Overall, implementing $\SyncPDOnline$ inside \dlfuzzer provided the added advantage of supplementing a powerful prediction technique with a biased randomized scheduler.
%As an added advantage of implementing our algorithms in the
% framework, we are able to achieve
%the complementary objective of evaluating how well prediction supplements
%the effectiveness of an otherwise naive
%controlled concurrency testing
%technique like randomized scheduling.
To our knowledge, our work is the first to effectively 
combine these two orthogonal techniques.
We also remark that such a bias is of no benefit to \dlfuzzer itself
since it does not employ any predictive reasoning.

% a randomized testing procedure, even though the latter is rather simple.

For this experiment, we run \dlfuzzer on each benchmark, and for each deadlock pattern found in the initial execution, 
we let it spawn $3$ new executions trying to realize the deadlock, 
as per standard (\href{https://github.com/ksen007/calfuzzer}{https://github.com/ksen007/calfuzzer}).
We repeated this process $50$ times and recorded the total time taken.
Then, we allocated the same time for $\SyncPDOnline$ to repeatedly execute the same program and perform deadlock prediction.
We measured all deadlocks found by each technique.
%We also noticed that \dlfuzzer fails to work properly if the given input program deadlocks in the initial run.
%This results in \dlfuzzer itself to go into deadlock without producing any deadlock reports.  
%Hence, we made a modest modification on \dlfuzzer allowing it to work properly and report deadlocks in such cases.
%We used this modified version of \dlfuzzer in our evaluation.





%!TEX root = ../main.tex

\begin{table*}
\caption{
Performance comparison of $\SyncPDOnline$ ($\syncpdshort$) and $\dlfuzzer$ ($\dlfshort$).
%%Column 1 denotes the name of the benchmark.
Columns 2-3 show the total number of bug reports. 
Columns 4-6 show the total number of unique bugs found by each tool, and their union. 
Columns 7-12 show the hit rate on each bug.
Columns 13-16 show the runtime overhead of the tools.
%We refer to $\dlfuzzer$ as $\dlfshort$, to $\SyncPDOnline$ as $\syncpdshort$, and the instrumentation phases using the suffix $\mathsf{\tt -I}$.
%Performance comparison of $\SyncPDOnline$ ($\syncpdshort$) and $\dlfuzzer$ ($\dlfshort$).
%Column 1 denotes the name of the benchmark.
%Columns 2-3 show the total number of bug reports. 
%Columns 4-6 show the total number of unique bugs found by each tool, and their union. 
%Columns 7-12 show the hit rate on each bug.
%Columns 13-17 show the runtime overhead of the tools.
%We use the suffix $\mathsf{\tt -I}$ for the instrumentation phases.
\label{tab:expr-dlf-results}
}
\vspace{-0.15cm}
\setlength\tabcolsep{3pt}
\renewcommand{\arraystretch}{0.9}
\small
\centering
\scalebox{1}{
\begin{tabular}{|r|c|c|c|c|c|c|c|c|c|c|c|c|c|c|c|c|c|c|c|c|c|c|c|c|c|c|c|}
\hline
1 & 2 & 3 & 4 & 5 & 6 & 7 & 8 & 9 & 10 & 11 & 12 & 13 & 14 & 15 & 16 \\
\hline
\multirow{2}{*}{\textbf{Benchmark}}& 
\multicolumn{2}{c|}{ {\textsf{Bug Hits}}} & 
\multicolumn{3}{c|}{ {\textsf{Unique Bugs}}} 
& \multicolumn{2}{c|}{ \textsf{Bug 1}} &
\multicolumn{2}{c|}{ \textsf{Bug 2}} &
\multicolumn{2}{c|}{ \textsf{Bug 3}} & \multicolumn{4}{c|}{ {\textsf{Runtime Overhead}}}  \\
\cline{7-16}
\cline{2-6}
& \textsf{$\syncpdshort$} & \textsf{$\dlfshort$} & \textsf{$\syncpdshort$} & \textsf{$\dlfshort$}  & All 
& \textsf{\textsf{$\syncpdshort$}} 
& \textbf{$\dlfshort$} 
& {\textbf{$\syncpdshort$}} 
&{\textbf{$\dlfshort$}} 
& {\textbf{$\syncpdshort$}}
&{\textbf{$\dlfshort$}} & {\textbf{$\syncpdshortinstr$}}  & {\textbf{$\syncpdshort$}}  & 
{\textbf{$\dlfshortinstr$}}  &
 {\textbf{$\dlfshort$}} \\
\hline
%\multirow{2}{*}{\textbf{Benchmark}}& \multirow{2}{*}{$\mathcal{N}$} & \multirow{2}{*}{$\mathcal{T}$} & \multirow{2}{*}{$\mathcal{M}$}
%& \multirow{2}{*}{$\mathcal{L}$} &
%\multirow{2}{*}{$\NumAcquires+\mathcal{R}$} &\multicolumn{2}{c||}{ {\seqc}} & \multicolumn{2}{c||}{ {\textsf{$\SyncPDOffline$}}} & \multicolumn{2}{c|}{ {\textsf{$\SyncPDOnline$}}} \\
%\cline{7-12}
%& & & & & &  \textbf{$\#$Deadlocks} & { \textbf{Time (s)}} &{ \textbf{$\#$Deadlocks}} & { \textbf{Time (s)}} &{ \textbf{$\#$Deadlocks}} &{ \textbf{Time (s)}} \\
\hline
Deadlock & 50 & 50 & 1 & 1 & 1 & 50 & 50 & - & - & - & - & 2$\times$ & 3$\times$ & 2$\times$ & 4$\times$\\
Picklock & 227 & 97 & 2 & 1 & 2 & 226 & 97 & 1 & 0 & - & - & 2$\times$ & 2$\times$ & 2$\times$ & 5$\times$\\
Bensalem & 355 & 32 & 2 & 1 & 2 & 8 & 0 & 347 & 32 & - & - & 2$\times$ & 2$\times$ & 2$\times$ & 6$\times$\\
Transfer & 54 & 50 & 1 & 1 & 1 & 54 & 50 & - & - & - & - & 2$\times$ & 2$\times$ & 1$\times$ & 4$\times$\\
Test-Dimmunix & 702 & 0 & 2 & 0 & 2 & 351 & 0 & 351 & 0 & - & - & 2$\times$ & 2$\times$ & 2$\times$ & 4$\times$\\
StringBuffer & 153 & 131 & 2 & 2 & 2 & 128 & 118 & 25 & 13 & - & - & - & - & - & -\\
Test-Calfuzzer & 177 & 44 & 1 & 1 & 1 & 177 & 44 & - & - & - & - & 2$\times$ & 2$\times$ & 2$\times$ & 4$\times$\\
DiningPhil & 162 & 100 & 1 & 1 & 1 & 162 & 100 & - & - & - & - & - & - & - & -\\
HashTable & 169 & 120 & 2 & 2 & 2 & 82 & 21 & 87 & 99 & - & - & - & - & - & -\\
Account & 19 & 188 & 1 & 1 & 1 & 19 & 188 & - & - & - & - & 2$\times$ & 8$\times$ & 2$\times$ & 16$\times$\\
Log4j2 & 290 & 100 & 2 & 1 & 2 & 145 & 100 & 145 & 0 & - & - & - & - & - & -\\
Dbcp1 & 265 & 138 & 2 & 2 & 2 & 264 & 61 & 1 & 77 & - & - & - & - & - & -\\
Dbcp2 & 129 & 126 & 2 & 2 & 2 & 86 & 99 & 43 & 27 & - & - & - & - & - & -\\
RayTracer & 0 & 0 & 0 & 0 & 0 & - & - & - & - & - & - & 122$\times$ & 124$\times$ & 109$\times$ & 111$\times$\\
Tsp & 0 & 0 & 0 & 0 & 0 & - & - & - & - & - & - & 47$\times$ & 60$\times$ & 37$\times$ & 40$\times$\\
jigsaw & 1189 & 1 & 1 & 1 & 2 & 1189 & 0 & 0 & 1 & - & - & - & - & - & -\\
elevator & 0 & 0 & 0 & 0 & 0 & - & - & - & - & - & - & 2$\times$ & 2$\times$ & 2$\times$ & 2$\times$\\
JDBCMySQL-1 & 349 & 117 & 2 & 3 & 3 & 1 & 21 & 0 & 4 & 348 & 92 & 3$\times$ & 4$\times$ & 2$\times$ & 13$\times$\\
JDBCMySQL-2 & 559 & 73 & 1 & 1 & 1 & 559 & 73 & - & - & - & - & 2$\times$ & 4$\times$ & 2$\times$ & 18$\times$\\
JDBCMySQL-3 & 560 & 224 & 1 & 1 & 1 & 560 & 224 & - & - & - & - & 2$\times$ & 5$\times$ & 2$\times$ & 24$\times$\\
JDBCMySQL-4 & 1717 & 101 & 3 & 1 & 3 & 95 & 0 & 834 & 0 & 788 & 101 & 3$\times$ & 5$\times$ & 2$\times$ & 31$\times$\\
hedc & 0 & 0 & 0 & 0 & 0 & - & - & - & - & - & - & 2$\times$ & 2$\times$ & 1$\times$ & 2$\times$\\
cache4j & 0 & 0 & 0 & 0 & 0 & - & - & - & - & - & - & 2$\times$ & 2$\times$ & 2$\times$ & 2$\times$\\
lusearch & 0 & 0 & 0 & 0 & 0 & - & - & - & - & - & - & 16$\times$ & 17$\times$ & 13$\times$ & 16$\times$\\
ArrayList & 47 & 45 & 3 & 3 & 3 & 20 & 22 & 3 & 5 & 24 & 18 & 50$\times$ & 69$\times$ & 18$\times$ & 79$\times$\\
Stack & 44 & 27 & 3 & 3 & 3 & 18 & 13 & 8 & 4 & 18 & 10 & 69$\times$ & 91$\times$ & 64$\times$ & 86$\times$\\
IdentityHashMap & 68 & 62 & 2 & 2 & 2 & 13 & 47 & 55 & 15 & - & - & 4$\times$ & 8$\times$ & 3$\times$ & 10$\times$\\
LinkedList & 48 & 26 & 3 & 2 & 3 & 21 & 17 & 7 & 0 & 20 & 9 & 16$\times$ & 28$\times$ & 14$\times$ & 32$\times$\\
Swing & 0 & 0 & 0 & 0 & 0 & - & - & - & - & - & - & 5$\times$ & 6$\times$ & 4$\times$ & 6$\times$\\
Sor & 0 & 0 & 0 & 0 & 0 & - & - & - & - & - & - & 2$\times$ & 7$\times$ & 2$\times$ & 2$\times$\\
HashMap & 46 & 44 & 2 & 2 & 2 & 18 & 11 & 28 & 33 & - & - & 7$\times$ & 11$\times$ & 4$\times$ & 8$\times$\\
Vector & 126 & 50 & 1 & 1 & 1 & 126 & 50 & - & - & - & - & 2$\times$ & 2$\times$ & 2$\times$ & 3$\times$\\
LinkedHashMap & 57 & 43 & 2 & 2 & 2 & 22 & 10 & 35 & 33 & - & - & 10$\times$ & 10$\times$ & 4$\times$ & 8$\times$\\
WeakHashMap & 29 & 40 & 2 & 2 & 2 & 6 & 11 & 23 & 29 & - & - & 7$\times$ & 12$\times$ & 4$\times$ & 8$\times$\\
montecarlo & 0 & 0 & 0 & 0 & 0 & - & - & - & - & - & - & 16$\times$ & 100$\times$ & 13$\times$ & 126$\times$\\
TreeMap & 42 & 47 & 2 & 2 & 2 & 16 & 15 & 26 & 32 & - & - & 9$\times$ & 12$\times$ & 5$\times$ & 9$\times$\\
eclipse & 0 & 0 & 0 & 0 & 0 & - & - & - & - & - & - & 2$\times$ & 2$\times$ & 2$\times$ & 2$\times$\\
TestPerf & 0 & 0 & 0 & 0 & 0 & - & - & - & - & - & - & 2$\times$ & 2$\times$ & 2$\times$ & 2$\times$\\
\hline\hline\textbf{Total} & \textbf{7633} & \textbf{2076} & \textbf{49} & \textbf{42} & \textbf{51} & - & - & - & - & - & - & - & - & - & - \\ 
\hline
\end{tabular}
}
\end{table*}


\Paragraph{Evaluation}
\cref{tab:expr-dlf-results} presents our experimental results.
%A bug identifies a unique tuple of source code locations corresponding to events
%participating in the deadlock.
Columns $2$-$3$ of the table display the total number of bug hits,
which is the total number of times a bug was predicted by $\SyncPDOnline$ in the entire duration,
or was confirmed in any trial of \dlfuzzer.
Columns $4$-$6$ display the unique bugs (i.e., unique tuples of source code locations leading to a deadlock) 
found by the techniques.
The employed techniques are able to find a maximum of $3$ unique bugs for each benchmark
in our benchmark set. 
Respectively, columns $7$-$12$ display the detailed information on the number 
of times a particular bug was found by each technique.
Runtime overheads are displayed in the columns $13$-$16$, with $\mathsf{\tt -I}$ denoting the instrumentation phase only.


\SubParagraph{\underline{Deadlock-detection capability}}
\dlfuzzer had $2076$ bug reports in total, and it found $42$ unique bugs.
In contrast, $\SyncPDOnline$ flagged $7633$ bug reports, corresponding to $49$ unique bugs.
In more detail, \dlfuzzer missed $9$ bugs reported by \SyncPDOnline whereas 
$\SyncPDOnline$ missed $2$ bugs reported by \dlfuzzer.
Also, \SyncPDOnline significantly outperformed \dlfuzzer in total number of bugs hits.
Our experiments again support that the notion of sync-preservation  captures most deadlocks that occur in practice, to the extent that other state-of-the-art techniques can capture.
%\hunkar{
A further observation is that in the offline experiments, \SyncPDOffline  was not able to find deadlocks in \texttt{Transfer} and \texttt{Deadlock}. 
However, the random scheduling procedure allowed \SyncPDOnline 
to navigate to executions from which deadlocks can be predicted.
This demonstrates the potential of combining predictive dynamic
techniques with controlled concurrency testing.
%; a direction we find promising to be pursued further by the community.


\SubParagraph{\underline{Runtime overhead}}
We have also measured the runtime overhead of both \SyncPDOnline and \dlfuzzer,
both as incurred by instrumentation, as well as by the deadlock analysis.
The latter is the time taken by \algoref{online} for the case of $\SyncPDOnline$,
and the overhead introduced due to the new executions in the second confirmation phase for the case of \dlfuzzer.
Our results show that the instrumentation overhead of \SyncPDOnline is, in fact, comparable to that of \dlfuzzer, though somewhat larger. 
This is expected, as \SyncPDOnline needs to also instrument memory access events, while \dlfuzzer only instruments lock events, but at the same time surprising because the number of memory access events
is typically much larger than the number of lock events.
On the other hand, the analysis overhead is often larger for \dlfuzzer, 
even though it reports fewer bugs.
It was not possible to measure the runtime overhead in certain benchmarks as 
either they were always deadlocking or the computation was running indefinitely.



Distributed systems have been maintaining their importance for the last several decades due to the increase in the need for scalable and reliable distributed applications while preserving high performance. 
To analyze distributed systems comprehensively and compare them in terms of features and services, various surveys and evaluations have been published in the past. Surveys on cloud providers, data warehouses, distributed file systems, or metadata services can be counted among them. 

Cloud providers are analyzed and evaluated in terms of elasticity \cite{CMART}, computing power \cite{comperative-benchmarking}, and cost to performance efficiency \cite{fair-benchmarking} in previous efforts. Widely used distributed services are also analyzed in many works, such as a survey on stream processing \cite{stream-benchmarking} or performance and dependability evaluation of MapReduce systems \cite{MapReduce-benchmarking}. Similarly, different aspects of distributed systems are studied in several surveys, like reliability analysis on distributed systems \cite{reliability-survey} and load balancing characteristics of known systems \cite{load-balancing-survey}.

As a big part of distributed systems, data warehouses and file systems are studied for many specifications. Evaluation of distributed data warehouses for the cost-effectiveness of different hardware configurations \cite{ALOJA} and query performance of distinct design choices\cite{benchmarking-data-warehouse} are among the known efforts in these works. Distributed file systems are examined in many past works for general concepts \cite{file-systems-concepts,file-systems-gen1} or specific applications such as distributed access control \cite{access-control-file-systems}. Due to the differences in optimization, design techniques, and the complex interactions between the file systems and other system components like the kernel or operating system, benchmarking distributed file systems is not trivial. To identify the important metrics for the evaluation of distributed file systems, researchers also studied benchmarking file systems \cite{File-system-benchmarking,benchmarking-file-rocket}. 

Analysis of distributed coordination services in terms of general characteristics and importance of coordination \cite{importance-of-coordination} and the comparison of existing algorithms \cite{paxos-made-simple} are among the published works. However, to the best of our knowledge, there is no published work on the evaluation of distributed coordination systems. As mentioned in the Introduction, due to the lack of standard benchmarking tools for distributed coordination services, developers widely use their ad-hoc benchmarks, which are prone to unfair comparisons or limited results for the evaluation of the systems. This study is unique in identifying the metrics and parameters for the evaluation of distributed coordination systems, discussing how each system uses these metrics and parameters for its evaluation, pinpointing the deficiencies of well-known benchmarking suites in evaluating distributed computing systems, and finally discussing the features of an ideal distributed coordination benchmark. 
% ;; -*- coding: iso-latin-1; TeX-PDF-mode: t; TeX-master: "main" -*-%

We have proposed a framework for the automatic generation of initial
contexts for deadlock-guided symbolic execution. Such initial contexts
are composed of the interfering tasks which, according to a static
deadlock analyzer, might lead to deadlock. Given the initial contexts,
we can drive symbolic execution towards paths that are more likely to
manifest a deadlock, discarding safe contexts.
%
There is a large body of work on deadlock detection including both
dynamic and static approaches.  Much of the existing work, both for
asynchronous programs \cite{FloresAG13-short,GGLLW13} and
thread-based programs
\cite{DBLP:conf/pdd/MasticolaR91,DBLP:journals/tocs/SavageBNSA97}, is
based on static analysis techniques. Although we have used the static
analysis of \cite{FloresAG13-short}, the information provided by other
deadlock analyzers could be used in an analogous way.
Deadlock detection has been also studied in the context of dynamic
testing and model checking
\cite{DBLP:conf/icst/ChristakisGS13,DBLP:conf/pldi/JoshiPSN09,lockout},
where sometimes has been combined with static information
\cite{DBLP:conf/hvc/AgarwalWS05,DBLP:conf/sigsoft/JoshiNSG10}. 
The initial contexts generated by our framework are of interest also
in these approaches. Deadlock detection is even more challenging in
the context of thread-based concurrency model. As future work, we plan to
investigate how our framework could be adapted to this model.

% Static analysis can ensure the
% absence of errors, however it works on approximations (especially for
% pointer aliasing) which might lead to a ``don't
% know'' answer. Our work complements static analysis techniques and can
% be used to look for deadlock paths when static analysis is not able to
% prove deadlock freedom. Using our method, we try to find a deadlock by
% generating initial contexts which include the conflicting tasks given
% by our deadlock detection algorithm that relies on the static
% information.




 %  The
%  approach in \cite{lockout} consists in generating binaries that are
%  more likely to manifest a deadlock such that by relying on standard
%  testing it is easier to capture deadlock derivations. 
% As regards combined approaches, the approach in
% \cite{DBLP:conf/sigsoft/JoshiNSG10} first performs a transformation of
% the program into a trace program that only keeps the instructions that
% are relevant for deadlock and then dynamic testing is performed on
% such program.


% The approach is fundamentally different from ours: in
% their case, since model checking is performed on the trace program
% (that over-approximates the deadlock behaviour), the method can detect
% deadlocks that do not exist in the program, while in our case this is
% not possible since the testing is performed on the original program
% and the analysis information is only used to drive the execution. 




\begin{acks}
Andreas Pavlogiannis was partially supported by a research grant (VIL42117) from VILLUM FONDEN.
Umang Mathur was partially supported by the Simons Institute for the Theory of Computing, and by a Singapore Ministry of Education (MoE) Academic Research Fund (AcRF) Tier 1 grant.
Mahesh Viswanathan was partially supported by NSF SHF 1901069 and NSF CCF 2007428.
\end{acks}

\section*{Data and Software Availability Statement}

The artifact developed for this work is available~\cite{artifact}, which contains all source codes and experimental data necessary to
reproduce our evaluation in~\cref{sec:experiments}, excluding the results of \seqc. 




%We provide an artifact~\cite{artifact} which contains all source codes and experimental data necessary to
%reproduce our evaluation in~\cref{sec:experiments}, excluding the results of \seqc. 
%In particular, source codes of the tools \SyncPDOffline and \SyncPDOnline as well as all the benchmarks used in our evaluations are provided. 
%The compared techniques \dirk~\cite{Kalhauge2018}, and DeadlockFuzzer~\cite{Joshi2009} are also included. 
%The artifact contains scripts that fully automate the process of reproducing~\cref{tab:expr-results} and~\cref{tab:expr-dlf-results}.
%Instructions on the usage of these scripts and customization of the evaluation are also provided with the artifact.

%\pagebreak
\bibliography{references}



%% Acknowledgments
%\begin{acks}                            %% acks environment is optional
                                        %% contents suppressed with 'anonymous'
  %% Commands \grantsponsor{<sponsorID>}{<name>}{<url>} and
  %% \grantnum[<url>]{<sponsorID>}{<number>} should be used to
  %% acknowledge financial support and will be used by metadata
  %% extraction tools.
  %This material is based upon work supported by the
  %\grantsponsor{GS100000001}{National Science
  %  Foundation}{http://dx.doi.org/10.13039/100000001} under Grant
  %No.~\grantnum{GS100000001}{nnnnnnn} and Grant
  %No.~\grantnum{GS100000001}{mmmmmmm}.  Any opinions, findings, and
  %conclusions or recommendations expressed in this material are those
  %of the author and do not necessarily reflect the views of the
  %National Science Foundation.
%\end{acks}


%% Bibliography
%\bibliography{bibfile}


%% Appendix
\begin{arxiv}
\pagebreak
\appendix
%!TEX root = main.tex

\section{Proofs of \secref{lower-bounds}}\seclabel{sec:app_proofs_lower_bounds}

In this section we present rigorous proofs of our hardness results in \thmref{pattern-w1-hardness-pattern}, \thmref{pattern-ov-hardness}, and \thmref{w1-hardness-pattern}.


%!TEX root = main.tex

%\subsection{$\W{1}$-hardness for Detecting Deadlock Patterns}
\subsection{Proof of \thmref{pattern-w1-hardness-pattern}}
\seclabel{w1-hardness-pattern}

Checking whether a trace of length
$\NumEvents$ has a deadlock pattern of size $k$
can be solved using a simple enumeration based algorithm
in time $O(\poly{\NumThreads} \cdot \NumEvents^\NumThreads)$.
In this section, we show that this problem is, in fact,
$\W{1}$-hard
% \footnote{We remark that $\W{1}$-hardness in a parameter also implies $\NP$-hardness in general} in the parameter $k$,
and thus is unlikely to be FPT (fixed parameter tractable).
In other words, an algorithm with running time $O(f(\NumThreads) \cdot \poly{\NumEvents})$ 
is unlikely to exist for this problem, for any computable function $f$.
%We now state our result.



\patternwonehardness*
\begin{proof}
	We show that there is a polynomial-time FPT-reduction from
	INDEPENDENT-SET(c) to the problem of checking the existence of deadlock-patterns
	of size $c$.
	Our reduction takes as input an undirected graph $G$ and outputs a trace $\tr$
	such that $G$ has an independent set of size $c$ iff $\tr$
	has a deadlock pattern of size $c$.
	%Here present the construction and refer to \appref{proofs} for the proof of correctness.
	
% 	\begin{figure}[]
% 		\input{figures/construction-w1-hardness}
% 		\hfill ~ \unskip\ \vrule\ \hfill %!TEX root = ../main.tex
%\begin{figure}
\begin{subfigure}[b]{0.475\textwidth}
	\newcommand{\xdisposition}{0}
	\newcommand{\ydisposition}{0}
	\newcommand{\xtstep}{0.75}
	\newcommand{\ytstep}{1}
	\newcommand{\ybias}{-0.3 }
	\newcommand{\xstep}{2.5}
	\newcommand{\ystep}{-0.475}
	\newcommand{\xtscale}{0.8}
	\def \numevents{9.5}
	\newcommand{\eventA}[4]{
		\node[event, draw=black, fill=white] (A#1) at (#1*\xstep, #2*\ystep) {\footnotesize $#2(x_{#3})$};
		%\node[] at (0*\xstep-\xtstep, {#1*\ystep}) {\small $#2(x_{#3})$};
	}
	\scalebox{0.9}{
		\begin{tikzpicture}[thick,
			pre/.style={<-,shorten >= 2pt, shorten <=2pt, very thick},
			post/.style={->,shorten >= 3pt, shorten <=3pt,   thick},
			seqtrace/.style={line width=2},
			und/.style={very thick, draw=gray},
			event/.style={rectangle, minimum height=0.8mm, minimum width=15mm,  line width=1pt, inner sep=0.5,},
			virt/.style={circle,draw=black!50,fill=black!20, opacity=0}]
			\footnotesize
			
			
%			\node[] at (0.5*\xstep, -1.5*\ystep){
%				\small
%				$\cs(\lk_i,\lk_j)=\acq(\lk_i) \cdot\acq(\lk_j) \cdot\rel(\lk_j) \cdot \rel(\lk_i)$
%			};
			

				
			\begin{scope}[shift={(0,-3*\ystep)}]
				
				\draw[dashed] (-0.9*\xstep,6.5*\ystep) rectangle (-0.5*\xstep,8.5*\ystep);
				\node (A1) at (-0.7*\xstep,7*\ystep) {\normalsize [1, 1]};
				\node (A2) at (-0.7*\xstep,8*\ystep) {\normalsize [1, 0]};
				\node (A) at (-0.7*\xstep,9*\ystep) {\large  $A$};
				
				\draw[dashed] (1.5*\xstep,6.5*\ystep) rectangle (1.9*\xstep,8.5*\ystep);
				\node (B1) at (1.7*\xstep,7*\ystep) {\normalsize [1, 0]};
				\node (B2) at (1.7*\xstep,8*\ystep) {\normalsize [0, 1]};
				\node (B) at (1.7*\xstep,9*\ystep) {\large  $B$};
				
				
			\end{scope}
			

			
			\node[] (S11) at (0*\xstep,0.15) {\normalsize $t_A$};
			\node[] (S12) at (0*\xstep,\numevents * \ystep) {};
			\node[] (S21) at (1*\xstep,0.15) {\normalsize $t_B$};
			\node[] (S22) at (1*\xstep,\numevents * \ystep) {};
			
			\draw[seqtrace] (S11) to (S12);
			\draw[seqtrace] (S21) to (S22);
			
			
			\node[event, draw=black, fill=white] (11) at (0*\xstep, 1*\ystep + 0*\ybias) {$\acq(\LockColorTwo{\lk_2})$};
			\node[event, draw=black, fill=white] (12) at (0*\xstep, 2*\ystep + 0*\ybias) {$\acq(\LockColorOne{\lk_1})$};
			\node[event, draw=black, fill=white, dotted] (13) at (0*\xstep, 3*\ystep + 0*\ybias) {$\cs(m_0,m_1)$};
			%\node[event, draw=black, fill=white, dotted] (14) at (0*\xstep, 4*\ystep + 0*\ybias) {$\acq(\lk')$};
			%\node[event, draw=black, fill=white, dotted] (15) at (0*\xstep, 5*\ystep + 0*\ybias) {$\rel(\lk')$};
			%\node[event, draw=black, fill=white, dotted] (16) at (0*\xstep, 6*\ystep + 0*\ybias) {$\rel(\lk)$};
			\node[event, draw=black, fill=white] (17) at (0*\xstep, 4*\ystep + 0*\ybias) {$\rel(\LockColorOne{\lk_1})$};
			\node[event, draw=black, fill=white] (18) at (0*\xstep, 5*\ystep + 0*\ybias) {$\rel(\LockColorTwo{\lk_2})$};
			
			\node[event, draw=black, fill=white] (19) at (0*\xstep, 6*\ystep + 1*\ybias) {$\acq(\LockColorOne{\lk_1})$};
			\node[event, draw=black, fill=white, dotted] (110) at (0*\xstep, 7*\ystep + 1*\ybias) {$\cs(m_0,m_1)$};
			%\node[event, draw=black, fill=white, dotted] (111) at (0*\xstep, 11*\ystep + 1*\ybias) {$\acq(\lk')$};
			%\node[event, draw=black, fill=white, dotted] (112) at (0*\xstep, 12*\ystep + 1*\ybias) {$\rel(\lk')$};
			%\node[event, draw=black, fill=white, dotted] (113) at (0*\xstep, 13*\ystep + 1*\ybias) {$\rel(\lk)$};
			\node[event, draw=black, fill=white] (114) at (0*\xstep, 8*\ystep + 1*\ybias) {$\rel(\LockColorOne{\lk_1})$};
			
			%%%%%%%%%%%%%%%%
			
			\node[event, draw=black, fill=white] (21) at (1*\xstep, 2*\ystep + 0*\ybias) {$\acq(\LockColorOne{\lk_1})$};
			\node[event, draw=black, fill=white, dotted] (22) at (1*\xstep, 3*\ystep + 0*\ybias) {$\cs(m_1,m_0)$};
			%\node[event, draw=black, fill=white, dotted] (23) at (1*\xstep, 4*\ystep + 0*\ybias) {$\acq(\lk)$};
			%\node[event, draw=black, fill=white, dotted] (24) at (1*\xstep, 5*\ystep + 0*\ybias) {$\rel(\lk)$};
			%\node[event, draw=black, fill=white, dotted] (25) at (1*\xstep, 6*\ystep + 0*\ybias) {$\rel(\lk')$};
			\node[event, draw=black, fill=white] (26) at (1*\xstep, 4*\ystep + 0*\ybias) {$\rel(\LockColorOne{\lk_1})$};
			
			\node[event, draw=black, fill=white] (27) at (1*\xstep, 5*\ystep + 1*\ybias) {$\acq_1(\LockColorTwo{\lk_2})$};
			\node[event, draw=black, fill=white, dotted] (28) at (1*\xstep, 6*\ystep + 1*\ybias) {$\cs(m_1,m)$};
			%\node[event, draw=black, fill=white, dotted] (29) at (1*\xstep, 11*\ystep + 1*\ybias) {$\acq(\lk)$};
			%\node[event, draw=black, fill=white, dotted] (210) at (1*\xstep, 12*\ystep + 1*\ybias) {$\rel(\lk)$};
			%\node[event, draw=black, fill=white, dotted] (211) at (1*\xstep, 13*\ystep + 1*\ybias) {$\rel(\lk')$};
			\node[event, draw=black, fill=white] (212) at (1*\xstep, 7*\ystep + 1*\ybias) {$\rel_1(\LockColorTwo{\lk_2})$};
			
			\begin{scope}[]
				\node[below left=of 212] (212b) {};
			\end{scope}
	
		\end{tikzpicture}
	}
	\caption{
		Reduction for OV-hardness proof from an instance of size $n = 2$ and $d = 2$.
		% We use the shortcut $\cs(\lk_i,\lk_j)$ to denote two nested critical sections on $\lk_i$ and $\lk_j$, as indicated in the figure.
	}
	\figlabel{ov-hardness}
\end{subfigure}
%\end{figure}
% 		\caption{
% 		Construction of $\W{1}$-hardness (\subref{fig:w1-hardness}) and OV-hardness (\subref{fig:ov-hardness}) results. 
% 		We use the shortcut $\cs(\lk_i,\lk_j)$ to denote two nested critical sections on $\lk_i$ and $\lk_j$. That is,
% 	$\cs(\lk_i,\lk_j)=\acq(\lk_i) \cdot\acq(\lk_j) \cdot\rel(\lk_j) \cdot \rel(\lk_i)$.
% 		} 
% 		%\label{fig:main}
% 	\end{figure}


	\myparagraph{Construction}{
		%Let $G = (V, E)$ be the input graph and let $c$ be the parameter.
		We assume  that the vertices are indexed from $1$ through $n = |V|$: $V = \set{v_1, v_2, \ldots, v_n}$.
		We also assume a total ordering $<_E$ on the set of edges $E$.
		The trace $\tr$ we construct is a concatenation of $c$ sub-traces: 
		$
		\tr = \tr^{(1)} \cdot \tr^{(2)} \cdots \tr^{(c)}
		$
		and uses $c$ threads $\set{t_1, t_2, \ldots t_c}$ and
		$|E| + c$ locks $\set{\lk_{\set{u, v}}}_{\set{u, v} \in E} \uplus \set{\lk_0, \lk_1 \ldots, \lk_{c-1}}$.
		The $i^\text{th}$ sub-trace $\tr^{(i)}$ is a sequence of events performed by thread $t_i$, and
		is obtained by concatenation of $n = |V|$ sub-traces:
		$
		\tr^{(i)} = \tr^{(i)}_1 \cdot \tr^{(i)}_2 \cdots \tr^{(i)}_n
		$.
		Each sub-trace $\tr^{(i)}_j$ with $(i \leq c, j \leq n)$ 
		comprises of nested critical sections over locks of the 
		form $\lk_{\set{v_j, u}}$, where $u$ is a neighbor of $v_j$.
		Inside the nested block we have critical 
		sections on locks $\lk_{i \% c}$ and $\lk_{(i+1) \% c}$.
		Formally, let $\set{v_j, v_{k_1}}, \ldots, \set{v_j, v_{k_d}}$
		be the neighboring edges of $v_j$ (ordered according to $<_E$).
		Then, $\tr^{(i)}_j$ is the unique string generated by the
		grammar having $d+1$ non-terminals $S_0, S_1, \ldots, S_d$, start symbol $S_d$
		and the following production rules:
		\begin{itemize}
			\item $S_0 \to \ev{t_i, \acq(\lk_{i \% c})} \cdot \ev{t_i, \acq(\lk_{(i+1) \% c})} \cdot \ev{t_i, \rel(\lk_{(i+1) \% c})} \cdot \ev{t_i, \rel(\lk_{i \% c})}$.
			\item for each $1 \leq r \leq d$, $S_r \to \ev{t_i, \acq(\lk_{\set{v_j, v_{k_r}}})} \cdot S_{r-1} \cdot \ev{t_i, \rel(\lk_{\set{v_j, v_{k_r}}})}$.
		\end{itemize}
		\figref{w1-hardness} illustrates this construction for a graph with $3$ nodes
		and parameter $c = 3$.
		Finally, observe that the lock-nesting depth in $\tr$ is bounded by 2 + the degree of $G$.
	}
	
\myparagraph{Correctness}{
	We now argue for the correctness of the construction.
	First, consider the case when $G$ has an independent set $I = \set{v_{j_1}, \ldots v_{j_c}}$ of size $\geq c$.
	Let  be $c$ distinct vertices from the independent set.
	Let $e^{(i)}_{j}$ be the innermost acquire (on lock $\lk_{(i+1) \% c}$) in the sub-trace 
	$\tr^{(i)}_{j}$.
	We show that the sequence $D = \pattern{e^{(1)}_{j_1}, e^{(2)}_{j_2} \ldots, e^{(c)}_{j_c}}$ 
	is a deadlock pattern.
	Observe that $\ThreadOf{e^{(i)}_{j_i}} = t_i \neq t_{i'} = \ThreadOf{e^{(i')}_{j_{i'}}}$ for every $i \neq i'$.
	Similarly, the locks acquired in $e^{(i)}_{j_i}$ and $e^{(i')}_{j_{i'}}$ are also distinct
	when $i \neq i'$.
	Finally, we want to show that $\lheld{\tr}(e^{(i)}_{j_i}) \cap \lheld{\tr}(e^{(i')}_{j_{i'}}) = \emptyset$ for every $i \neq i'$.
	Assume on the contrary that there is a lock $\lk \in \lheld{\tr}(e^{(i)}_{j_i}) \cap \lheld{\tr}(e^{(i')}_{j_{i'}})$.
	Clearly, $\lk$ cannot be of the form $\lk_m$ for some $0 \leq m < c$ as the only such lock held
	at $e^{(i)}_{j_i}$ is $\lk_{i \% c}$ and the only such lock held at $e^{(i')}_{j_{i'}}$ is $\lk_{i' \% c}$
	which are different.
	Thus, it must be of the form $\lk_{\set{u, v}}$ for some $\set{u, v} \in E$.
	Since it is acquired in both sub-traces  $\tr^{(i)}_{j_i}$ and $\tr^{(i')}_{j_{i'}}$,
	we have $\lk = \lk_{\set{v_i, v_{i'}}}$.
	This means that $\set{v_i, v_{i'}} \in E$ contradicting that $I$ is an independent set.

	Now, consider the case when there is a deadlock pattern of size $c$:
	$D = \pattern{e_0, e_1, \ldots, e_{c-1}}$.
	Since there are exactly $c$ threads in $\tr$, there must be one event in each thread.
	We first argue that if there is one $e_i$ that acquires a lock of the form $\lk_{\set{u, v}}$ then all events in the deadlock pattern do so.
	This follows from a simple inductive argument --- $e_{i+1}$ must acquire a lock held at $e_i$
	and thus must be of the form $\lk_{\set{u', v'}}$, and so on.
	Next, all locks of the form $\lk_{\set{u, v}}$ are acquired in the order consistent with
	the total order $<_E$.
	Hence, they cannot form a cyclic dependency.
	Thus, the locks acquired at $e_0, \ldots e_{c-1}$ are $\lk_0, \ldots, \lk_{c-1}$.
	Symmetry ensures that, you can assume, w.l.o.g that $\OpOf{e_i} = \acq(\lk_{i \% c})$.
	Let $f(i)$ be the index such that $e_i$ is in the subtrace $e^{(i)}_{f(i)}$.
	Then, we show that the set $A = \set{v_{f(1)}, v_{f(2)}, \ldots, v_{f(c)}}$ is an independent set of size $c$.
	Assume on the contrary that there are $i \neq i'$ such that $\set{v_{f(i)}, v_{f(i')}} \in E$.
	Then, our construction ensures that
	the lock $\lk_\set{v_{f(i)}, v_{f(i')}}$ will be held at both $e_i$ and $e_{i'}$
	contradicting that $D$ is a deadlock pattern.
}	

\myparagraph{Time complexity}{
		%Let us assume that the graph $G$ is presented as an adjacency list.
		The size of the trace is $O\big(c \cdot (2 + \sum_{v \in V} \mathsf{degree}(v)) \big)$
		$ = O(c \cdot (|V| + |E|))$
		and the time taken is also the same.
		%Also, the parameter $c$ is the same in both problems.
	}
\end{proof}







% %!TEX root = main.tex

\subsection{Proof of \thmref{pattern-ov-hardness}}
\seclabel{ov-hardness}




\patternovhardness*
\begin{proof}
	
	We show a fine-grained reduction from the Orthogonal Vectors Problem to
	the problem of checking for deadlock patterns of size $k$.
	For this, we start with two sets
	$A_1, A_2, \ldots, A_k \subseteq \set{0, 1}^d$
	of $d$-dimensional vectors with $|A_i| = n$ for every $1 \leq i \leq k$.
	We write the $j^{th}$ vector in $A_i$ as $A_{i, j}$.
	
	%%!TEX root = ../main.tex
%\begin{figure}
\begin{subfigure}[b]{0.475\textwidth}
	\newcommand{\xdisposition}{0}
	\newcommand{\ydisposition}{0}
	\newcommand{\xtstep}{0.75}
	\newcommand{\ytstep}{1}
	\newcommand{\ybias}{-0.3 }
	\newcommand{\xstep}{2.5}
	\newcommand{\ystep}{-0.475}
	\newcommand{\xtscale}{0.8}
	\def \numevents{9.5}
	\newcommand{\eventA}[4]{
		\node[event, draw=black, fill=white] (A#1) at (#1*\xstep, #2*\ystep) {\footnotesize $#2(x_{#3})$};
		%\node[] at (0*\xstep-\xtstep, {#1*\ystep}) {\small $#2(x_{#3})$};
	}
	\scalebox{0.9}{
		\begin{tikzpicture}[thick,
			pre/.style={<-,shorten >= 2pt, shorten <=2pt, very thick},
			post/.style={->,shorten >= 3pt, shorten <=3pt,   thick},
			seqtrace/.style={line width=2},
			und/.style={very thick, draw=gray},
			event/.style={rectangle, minimum height=0.8mm, minimum width=15mm,  line width=1pt, inner sep=0.5,},
			virt/.style={circle,draw=black!50,fill=black!20, opacity=0}]
			\footnotesize
			
			
%			\node[] at (0.5*\xstep, -1.5*\ystep){
%				\small
%				$\cs(\lk_i,\lk_j)=\acq(\lk_i) \cdot\acq(\lk_j) \cdot\rel(\lk_j) \cdot \rel(\lk_i)$
%			};
			

				
			\begin{scope}[shift={(0,-3*\ystep)}]
				
				\draw[dashed] (-0.9*\xstep,6.5*\ystep) rectangle (-0.5*\xstep,8.5*\ystep);
				\node (A1) at (-0.7*\xstep,7*\ystep) {\normalsize [1, 1]};
				\node (A2) at (-0.7*\xstep,8*\ystep) {\normalsize [1, 0]};
				\node (A) at (-0.7*\xstep,9*\ystep) {\large  $A$};
				
				\draw[dashed] (1.5*\xstep,6.5*\ystep) rectangle (1.9*\xstep,8.5*\ystep);
				\node (B1) at (1.7*\xstep,7*\ystep) {\normalsize [1, 0]};
				\node (B2) at (1.7*\xstep,8*\ystep) {\normalsize [0, 1]};
				\node (B) at (1.7*\xstep,9*\ystep) {\large  $B$};
				
				
			\end{scope}
			

			
			\node[] (S11) at (0*\xstep,0.15) {\normalsize $t_A$};
			\node[] (S12) at (0*\xstep,\numevents * \ystep) {};
			\node[] (S21) at (1*\xstep,0.15) {\normalsize $t_B$};
			\node[] (S22) at (1*\xstep,\numevents * \ystep) {};
			
			\draw[seqtrace] (S11) to (S12);
			\draw[seqtrace] (S21) to (S22);
			
			
			\node[event, draw=black, fill=white] (11) at (0*\xstep, 1*\ystep + 0*\ybias) {$\acq(\LockColorTwo{\lk_2})$};
			\node[event, draw=black, fill=white] (12) at (0*\xstep, 2*\ystep + 0*\ybias) {$\acq(\LockColorOne{\lk_1})$};
			\node[event, draw=black, fill=white, dotted] (13) at (0*\xstep, 3*\ystep + 0*\ybias) {$\cs(m_0,m_1)$};
			%\node[event, draw=black, fill=white, dotted] (14) at (0*\xstep, 4*\ystep + 0*\ybias) {$\acq(\lk')$};
			%\node[event, draw=black, fill=white, dotted] (15) at (0*\xstep, 5*\ystep + 0*\ybias) {$\rel(\lk')$};
			%\node[event, draw=black, fill=white, dotted] (16) at (0*\xstep, 6*\ystep + 0*\ybias) {$\rel(\lk)$};
			\node[event, draw=black, fill=white] (17) at (0*\xstep, 4*\ystep + 0*\ybias) {$\rel(\LockColorOne{\lk_1})$};
			\node[event, draw=black, fill=white] (18) at (0*\xstep, 5*\ystep + 0*\ybias) {$\rel(\LockColorTwo{\lk_2})$};
			
			\node[event, draw=black, fill=white] (19) at (0*\xstep, 6*\ystep + 1*\ybias) {$\acq(\LockColorOne{\lk_1})$};
			\node[event, draw=black, fill=white, dotted] (110) at (0*\xstep, 7*\ystep + 1*\ybias) {$\cs(m_0,m_1)$};
			%\node[event, draw=black, fill=white, dotted] (111) at (0*\xstep, 11*\ystep + 1*\ybias) {$\acq(\lk')$};
			%\node[event, draw=black, fill=white, dotted] (112) at (0*\xstep, 12*\ystep + 1*\ybias) {$\rel(\lk')$};
			%\node[event, draw=black, fill=white, dotted] (113) at (0*\xstep, 13*\ystep + 1*\ybias) {$\rel(\lk)$};
			\node[event, draw=black, fill=white] (114) at (0*\xstep, 8*\ystep + 1*\ybias) {$\rel(\LockColorOne{\lk_1})$};
			
			%%%%%%%%%%%%%%%%
			
			\node[event, draw=black, fill=white] (21) at (1*\xstep, 2*\ystep + 0*\ybias) {$\acq(\LockColorOne{\lk_1})$};
			\node[event, draw=black, fill=white, dotted] (22) at (1*\xstep, 3*\ystep + 0*\ybias) {$\cs(m_1,m_0)$};
			%\node[event, draw=black, fill=white, dotted] (23) at (1*\xstep, 4*\ystep + 0*\ybias) {$\acq(\lk)$};
			%\node[event, draw=black, fill=white, dotted] (24) at (1*\xstep, 5*\ystep + 0*\ybias) {$\rel(\lk)$};
			%\node[event, draw=black, fill=white, dotted] (25) at (1*\xstep, 6*\ystep + 0*\ybias) {$\rel(\lk')$};
			\node[event, draw=black, fill=white] (26) at (1*\xstep, 4*\ystep + 0*\ybias) {$\rel(\LockColorOne{\lk_1})$};
			
			\node[event, draw=black, fill=white] (27) at (1*\xstep, 5*\ystep + 1*\ybias) {$\acq_1(\LockColorTwo{\lk_2})$};
			\node[event, draw=black, fill=white, dotted] (28) at (1*\xstep, 6*\ystep + 1*\ybias) {$\cs(m_1,m)$};
			%\node[event, draw=black, fill=white, dotted] (29) at (1*\xstep, 11*\ystep + 1*\ybias) {$\acq(\lk)$};
			%\node[event, draw=black, fill=white, dotted] (210) at (1*\xstep, 12*\ystep + 1*\ybias) {$\rel(\lk)$};
			%\node[event, draw=black, fill=white, dotted] (211) at (1*\xstep, 13*\ystep + 1*\ybias) {$\rel(\lk')$};
			\node[event, draw=black, fill=white] (212) at (1*\xstep, 7*\ystep + 1*\ybias) {$\rel_1(\LockColorTwo{\lk_2})$};
			
			\begin{scope}[]
				\node[below left=of 212] (212b) {};
			\end{scope}
	
		\end{tikzpicture}
	}
	\caption{
		Reduction for OV-hardness proof from an instance of size $n = 2$ and $d = 2$.
		% We use the shortcut $\cs(\lk_i,\lk_j)$ to denote two nested critical sections on $\lk_i$ and $\lk_j$, as indicated in the figure.
	}
	\figlabel{ov-hardness}
\end{subfigure}
%\end{figure}
	
	\myparagraph{Construction}{
		We will construct a trace $\tr$ such that $\tr$ has a deadlock
		pattern of length $k$ iff $(A_1, A_2, \ldots, A_k)$ is a positive $k$-OV instance.
		The trace $\tr$ is of the form $\tr = \tr^{A_1} \cdot \tr^{A_2} \cdots \tr^{A_k}$
		and uses $k$ threads $\set{t_{A_1}, \ldots, t_{A_k}}$ and $d+k$ distinct 
		locks $\lk_1, \ldots, \lk_d, m_1, m_2, \ldots, m_k$. 
		Intuitively, the sub-trace $\tr^{A_i}$ encodes the given set of vectors $A_i$.
		The sub-traces $\tr^{A_i} = \tr^{A_i}_1 \cdot \tr^{A_i}_2 \cdots \tr^{A_i}_n$
		% and $\tr^B = \tr^B_1 \cdot \tr^B_2 \cdots \tr^B_n$ 
		are defined as follows.
		For each $j \in \set{1, 2, \ldots, n}$
		the sub-trace $\tr^{A_i}_j$ is the unique string generated by the
		grammar having $d+1$ non-terminals $S_0, S_1, \ldots, S_d$, start symbol $S_d$
		and the following production rules:
		\begin{itemize}
			\item $S_0 \to \ev{t_Z, \acq(m_i)} \cdot \ev{t_Z, \acq(m_{i \% k + 1})} \cdot \ev{t_Z, \rel(m_{i \% k + 1})} \cdot \ev{t_Z, \rel(m_i)}$.
			% where $(m, m') = (\lk, \lk')$ if $Z = A$, and $(m, m') = (\lk', \lk)$.
			\item for each $1 \leq p \leq d$, $S_p \to S_{p-1}$ if $A_{i, j}[p] = 0$.
			Otherwise (if $A_{i, j}[p] =1$), $S_p \to \ev{t_{A_i}, \acq(\lk_p)} \cdot S_{p-1} \cdot \ev{t_{A_i}, \rel(\lk_p)}$.
		\end{itemize}
		In words, all events of $\tr^{A_i}$ are performed by thread $t_{A_i}$.
		 % and those in $\tr^B$ are performed by $t_B$.
		Next, the $j^{th}$ sub-trace of $\tr^{A_i}$, denoted $\tr^{A_i}_j$ corresponds to the vector $A_{i, j}$
		as follows --- $\tr^{A_i}_j$ is a nested block of critical sections,
		with the innermost critical section being on lock $m_{i \% k + 1}$,
		which is immediately enclosed in a  critical section on lock $m_i$.
		Further, in the sub-trace $\tr^{A_i}_j$, 
		the lock $\lk_p$ occurs iff $A_{i, j}[p] = 1$.
		% The sub-traces $\tr^B_i$ is similarly constructed, except that the order
		% of the two innermost critical sections is inverted.
		\figref{ov-hardness} illustrates the construction for an OV-instance with $k=2$ $n=2$ and $d=2$.
	}

\myparagraph{Correctness}{
We now argue for the correctness of the construction.
First, if $(A_1, A_2, \ldots, A_k)$ is a positive $OV$ instance, then
there are $k$ indices $\alpha_1, \alpha_2, \ldots, \alpha_k \in \set{1 ,\ldots, n}$
such that $\sum_{p=1}^d \prod_{i=1}^k A_{i, \alpha_i}[p] = 0$.
 % and $B_\beta$ are orthogonal.
Thus, for every $1 \leq p \leq d$, there is an $i$ such that $A_{i, \alpha_i}[p] = 0$.
% or $B_\beta[j] = 0$.
Based on the construction, this implies that for every lock $\lk_p$
(with $1 \leq p \leq d$),
$\lk_p$ occurs in at most one of the sub-traces $\tr^{A_1}_{\alpha_1}, \ldots, \tr^{A_k}_{\alpha_k}$.
 % and $\tr^B_\beta$.
This gives a deadlock pattern $D = \pattern{e_1, e_2, \ldots, e_k}$,
where $e_i$ is the unique event in the sub-trace $\tr^{A_i}_{\alpha_i}$ with $\OpOf{e} = \acq(m_{i\%k +1})$.
% and $e'$ is the unique event in $\tr^B_\beta$ with $\OpOf{e'} = \acq(\lk)$.
This is because $m_i \in \lheld{\tr}(e_i)$,
and each lock $\lk_p$ exists in at most one of $\set{\lheld{\tr}(e_i)}_{i\leq n}$
% or $\lheld{\tr}(e')$ 
(and thus $\lheld{\tr}(e_i) \cap \lheld{\tr}(e_{i'}) = \emptyset$ for every $i \neq i' \leq n$).

Now, consider the case when there is a deadlock pattern $D = \pattern{e_1, e_2, \ldots, e_k}$ in $\tr$.
W.l.o.g, we can assume that
$\ThreadOf{e_i} = t_{A_i}$.
This means there are $\alpha_1, \alpha_2, \ldots, \alpha_k$
 % and $\beta$ such that
 such that
$e_i$ occurs in $\tr^{A_i}_{\alpha_i}$.
 % and $e'$ occurs in $\tr^B_\beta$.
Next, observe that the events $\set{e_i}_{i \leq k}$
 % and $e'$ 
 can only be acquire events on locks $\set{m_1, m_2, \ldots, m_k}$
and not on the locks $\set{\lk_1, \ldots, \lk_d}$.
This is because the order of acquisition amongst $\lk_1, \ldots, \lk_d$
is always fixed and further, each of these are never acquired inside of critical sections
of $m_1, \ldots, m_k$.
 % or $\lk'$.
By the choice of threads, $\OpOf{e_i} = \acq(m_{i\%k + 1})$ for every $i \leq k$.
% and $\OpOf{e'} = \acq(\lk)$.
Since $\lheld{\tr}(e_i) \cap \lheld{\tr}(e_{i'}) = \emptyset$ for every $i \neq i' \leq k$, 
we also have that
for every $p \leq d$, $\lk_p \not\in \lheld{\tr}(e) \cap \lheld{\tr}(e')$.
This means that at, for each $p \leq d$, there is atmost one $i \leq k$
for which $A_{i, \alpha_i}[p] = 1$.
This means that the vectors $A_{1, \alpha_1}, A_{2, \alpha_2} \ldots, A_{k, \alpha_k}$
witness that the given OV instance is a positive instance.
}


	\myparagraph{Time complexity}{
		For any fixed $k\geq 2$, the total size of the trace is $\NumEvents = O(n\cdot d)$
		and it has $\NumLocks = d + k $ locks.
		The construction can be performed in time $O(n\cdot d)$.
		Thus, if for some $\epsilon > 0$, 
		there is a $O(\poly{\NumLocks} \cdot \NumEvents^{k-\epsilon}) = O(\poly{d} \cdot n^{k-\epsilon})$ 
		algorithm for the problem of detecting deadlock patterns of length $k$,
		we would get a $O(\poly{d} \cdot n^{k-\epsilon} + n\cdot d) = O(\poly{d} \cdot n^{k-\epsilon})$
		time algorithm for $k$-OV, falsifying the $k$-OV hypothesis.
	}
\end{proof}

\input{2-ov-hardness-pattern-old}

%!TEX root = main.tex

\subsection{$\W{1}$-hardness for Deadlock Prediction}
\seclabel{w1-hardness}

Finally, we prove that predicting deadlocks is a $\W{1}$-hard problem parameterized by the number of threads $\NumThreads$.

\wonehardness*
\begin{proof}
    We establish a simple reduction from the problem of data-race prediction to the problem of predicting deadlocks of size $2$.
	Given a trace $\tr$ and a deadlock pattern $D=\pattern{\acq_1, \acq_2}$ of size $2$, we choose a fresh variable $x\not \in \vars{\tr}$, and construct two events $\wt_1=\ev{\ThreadOf{\acq_1}, \wt(x)}$ and $\wt_2=\ev{\ThreadOf{\acq_2}, \wt(x)}$, where $x$.
	We construct a trace $\tr'$ by replacing $\acq_1$ and $\acq_2$ in $\tr$ with $\wt_1$ and $\wt_2$, respectively.
	It is straightforward to verify that any correct $\rho$ reordering of $\tr'$ that witnesses $(\wt_1, \wt_2)$ as a data race is also a correct reordering of $\tr$ that witnesses the deadlock $(\acq_1, \acq_2)$.
	Since predicting data races is $\W{1}$-hard on $\NumThreads$\cite{Mathur2020b}, the desired result follows.
\end{proof}






%!TEX root = main.tex

\section{Proofs of \secref{syncp}}
\seclabel{sec:app_proofs_characterize_patterns}

We first prove \lemref{spclosure-necessary-sufficient}.

\spclosureNecessarySufficient*

% \begin{restatable}{lemma}{spclosureNecessarySufficient}
% \lemlabel{spclosure-necessary-sufficient}
% Let $\tr$ be a trace and let $S \subseteq \events{\tr}$.
% For any sync-preserving correct reordering $\rho$ of $\tr$, 
% 	if $S \subseteq \events{\rho}$, then $\SPClosure{\tr}(S) \subseteq \events{\rho}$.
% 	Further, there is a sync-preserving correct reordering $\rho$ of 
% 	$\tr$ such that $\events{\rho} = \SPClosure{\tr}(S)$.
% \end{restatable}

\begin{proof}
Let $\rho$ be a sync-preserving correct reordering of $\tr$ (and thus $\events{\rho} \subseteq \events{\tr}$)
such that $S \subseteq  \events{\rho}$.
Since, $\rho$ is a correct-reordering of $\tr$, for every event $e \in \events{\rho}$,
we must have $f \in \events{\rho}$ for every $f \tho{\tr} e$.
Likewise, for every read event $e \in \events{\rho}$, we must also have $\rf{\tr}(e) \in \events{\rho}$,
again because $\rho$ is a correct reordering of $\tr$.
Finally, let $e \in \events{\rho}$ be an acquire event (with $\OpOf{e} = \acq(\lk)$), and let
$f \neq e$ be another acquire event (with $\OpOf{e} = \acq(\lk)$)
with $f \trord{\tr} e$.
Since $\rho$ is a sync-preserving correct reordering of $\rho$, we must also have
$\match{\tr}(f) \in \events{\rho}$.
Thus, $\rho$ must contain $\SPClosure{\tr}(S)$.

Let us now argue that we can construct a sync-preserving correct reordering whose events are precisely $\SPClosure{\tr}(S)$.
Consider the sequence of events $\rho$ obtained by projecting $\tr$
to $\SPClosure{\tr}(S)$.
We will argue that (a) $\rho$ is a well-formed trace, 
(b) $\rho$ is a correct reordering of $\tr$, and 
(c) $\rho$ preserves the order of critical section on the same lock as in $\tr$.
For (a), we need to argue that critical sections on the same lock do not overlap in $\rho$.
This follows because if they did, then we have two acquire events $e_1 \neq e_2 \in \SPClosure{\tr}(S)$
such that $\OpOf{e_1} = \OpOf{e_2} = \acq(\lk)$ (for some $\lk$),
$e_1 \trord{\tr} e_2$ (and thus $e_1$ also appears before $e_2$ in $\rho$ by construction),
but $\match{\tr}(e_1)$ either does not appear in $\rho$ or appears later than $e_2$.
In the first case, we arrive at a contradiction to the definition of $\SPClosure{\tr}(S)$,
and in the second case, we arrive at a contradiction that $\rho$ is obtained without changing the relative order of the residual events.
Arguing (b) is straightforward: $\SPClosure{\tr}(S)$ is closed under both $\tho{\tr}$ as well as $\rf{\tr}$,
and further $\rho$ preserves the order of events as in $\tr$.
Likewise,  (c) is easy to argue because $\events{\rho} = \SPClosure{\tr}(S)$ are closed
under sync-preservation and the order of events in $\rho$ does not change with respect to $\tr$. 
\end{proof}

Let us now turn our attention to \lemref{spclosure-deadlock}.

\spclosureDeadlock*

% \begin{restatable}{lemma}{spclosureDeadlock}
% \lemlabel{spclosure-deadlock}
% Let $\tr$ be a trace and let $D = \pattern{e_0, \ldots, e_{k-1}}$
% be a deadlock pattern of size $k$ in $\tr$.
% $D$ is a sync-preserving deadlock of $\tr$ iff
% $\SPClosure{\tr}(\prev{\tr}(S)) \cap S = \emptyset$, where
% $S = \set{e_0, \ldots, e_{k-1}}$.
% \end{restatable}

\begin{proof}
($\Rightarrow$). Assume $D$ is a sync-preserving deadlock of $\tr$.
Let $\rho$ be the sync-preserving correct reordering of $\tr$ that witnesses $D$.
Thus, each of $e_0, \ldots, e_{k-1}$ are enabled in $\rho$.
This means that $\prev{\tr}(S) \subseteq \events{\rho}$, and also $S \cap \events{\rho} = \emptyset$.
By \lemref{spclosure-necessary-sufficient}, we have
$\SPClosure{\tr}(\prev{\tr}(S)) \subseteq \events{\rho}$, giving us
$\SPClosure{\tr}(\prev{\tr}(S)) \cap S = \emptyset$.

($\Leftarrow$). Assume $\SPClosure{\tr}(\prev{\tr}(S)) \cap S = \emptyset$.
From \lemref{spclosure-necessary-sufficient}, there exists a sync-preserving correct reordering
$\rho$ with $\prev{\tr}(S) \subseteq \SPClosure{\tr}(\prev{\tr}(S)) = \events{\rho}$.
Since $S \cap \events{\rho} = \emptyset$, the entire deadlock pattern $D$ is also enabled in $\rho$.
Thus, $D$ is a sync-preserving deadlock of $\tr$.
\end{proof}

The proof of the following proposition follows directly from the definition of sync-preserving closure.

\spclosureMonotone*

\begin{proof}
Follows from the definition of $\SPClosure{\tr}(S)$.
\end{proof}
%%!TEX root = main.tex

\section{Details of Algorithms}

%\subsection{Verifying Abstract Deadlock Patterns}
%\applabel{verify-abstract-patterns}

%Given an abstract deadlock pattern of length $k \geq 2$,
%there can be $O(\NumEvents^k)$ instantiations of it,
%giving a naive enumerate-and-check algorithm running 
%in time $O(\NumThreads\cdot\NumEvents^{k+1})$.
%Instead we exploit monotonicity properties of 
%the sync-preserving closure (\propref{spclosure-monotone})
%and instantiations of an abstract pattern (\corref{pattern-monotone}) 
%that allow for an \emph{incremental}
%algorithm that iteratively checks successive instantiations
%of a given abstract deadlock pattern, while still spending total $O(\NumEvents\cdot \NumThreads)$ time.
%% , i.e., the running time is linear in $\NumEvents$, which is the dominating parameter.
%The first observation allows us to re-use a prior
%computation when checking later deadlock patterns.
%
%% We first observe that the sync-preserving closure is monotonic
%% with respect to the thread-order, and thus, when computing closure
%% of later events, one can avoid re-computation.
%\begin{restatable}{proposition}{spclosure-monotone}
%\proplabel{spclosure-monotone}
%For a trace $\tr$ and sets $S, S' \subseteq \events{\tr}$.
%If for every event $e \in S$, there is an event $e' \in S$
%such that $e \tho{\tr} e'$, then
%$\SPClosure{\tr}(S) \subseteq \SPClosure{\tr}(S')$.
%\end{restatable}
%
%%\begin{proof}
%%Follows from the definition of $\SPClosure{\tr}(S)$.
%%\end{proof}
%
%Next, we extend \propref{spclosure-monotone} to avoid redundant
%computations when a sync-preserving deadlock is not found and later deadlock patterns must be checked.
%Given two deadlock patterns
%$D_1 = e_0, \ldots, e_{k-1}$ and $D_2 = f_0, \ldots, f_{k-1}$
%of the same length $k$, 
%we say $D_1 \prec D_2$ if
%they are instantiations of a common abstract pattern $\abst{D}$
%(i.e., $D_1, D_2 \in \abst{D}$) and 
%for every $i<k$, $e_i \tho{\tr} f_i$.
%\begin{restatable}{corollary}{pattern-monotone}
%\corlabel{pattern-monotone}
%Let $\tr$ be a trace and let $D_1 = e_0, \ldots, e_{k-1}$ and 
%$D_2 = f_0, \ldots f_{k-1}$ be deadlock
%patterns of size $k$ in $\tr$ such that $D_1 \prec D_2$.
%Let $S_1 = \set{e_0, \ldots, e_{k-1}}$ and $S_2 = \set{f_0, \ldots, f_{k-1}}$.
%If $\SPClosure{\tr}(\prev{\tr}(S_1))\cap S_2 \neq \emptyset$,
%then $\SPClosure{\tr}(\prev{\tr}(S_2))\cap S_2 \neq \emptyset$.
%\end{restatable}
%
%The algorithm for checking if an abstract deadlock 
%pattern contains a sync-preserving instantiation
%is presented in~\algoref{abstract-pattern}.
%At a high level, it iterates over the sequences $F_0, \ldots, F_{k-1}$ of 
%acquires (one for each abstract acquire) in trace order.
%For this, the algorithm maintains
% indices $i_0, \ldots, i_{k-1}$ that point to entries in $F_0, \ldots, F_{k-1}$.
%At each step, it determines whether the current
%deadlock pattern $D = e_0, \ldots, e_{k-1}$
%by computing the sync-preserving closure
%of the predecessors of the deadlock pattern, and checks
%if it is disjoint from the deadlock.
%If so, it terminates.
%Otherwise, it looks for the next eligible deadlock,
%which it determines based on \corref{pattern-monotone}.
%In particular, it advances the pointer $i_j$
%all the way until an entry which is outside of the
%closure computed so far.
%Also observe that the timestamp $T$ of the closure
%computed in an iteration is being used in later iterations;
%this is a consequence of \propref{spclosure-monotone}.
%Further, we ensure that the list of acquires $\AcqLst_{t, \lk}$
%used in the function \fixpoint (\algoref{compute-closure})
%be reused across iterations, and not need be
%re-assigned to the original list of all acquire events;
%correctness of this optimization follows from~\propref{spclosure-monotone}.
%Let us now calculate the running time for this algorithm.
%Each of the $\AcqLst_{t, \lk}$ in \fixpoint is traversed
%at most once. 
%Next, each element of the sequences $F_0, \ldots, F_{k-1}$
%is also traversed atmost once.
%For each of these acquires, the algorithm spends $O(\NumThreads)$
%time for vector clock updates.
%The total time required is thus $O(\NumEvents \cdot \NumThreads)$.
%This concludes the proof of \lemref{abstract-pattern-linear-time}.

\subsection{The Online $\SyncPDOffline$ Algorithm}
\applabel{app_online}
%\seclabel{app_online}
\Andreas{We can consider bringing this in the main paper. We should have space now. Some reviewers requested this, and it does make the paper self contained.}
In this section we give a detailed pseudocode description of the online $\SyncPDOffline$ algorithm (\algoref{online}).
%!TEX root = main.tex

\small
\begin{algorithm*}[H]
% \BlankLine
\begin{multicols}{2}
\myfun{\init{}}{
	\ForEach{$t \in \threads{}$ $\cdot$}{
		$\Cc_t$ := $\bot$
	}
	\ForEach{$x \in \vars{}$ $\cdot$}{
		$\LW_x$ := $\bot$
	}
	\ForEach{$\lk \in \locks{}$ $\cdot$}{
		$\gId_\lk$ := $0$
	}
	\ForEach{$t_1 {\neq}t_2 \in \threads{}, \lk_1 {\neq} \lk_2 \in \locks{}$}{
		$\view{\Ii}{}{\tuple{t_1, \lk_1, t_2, \lk_2}}$ := $\bot$ \;
		\ForEach{$t \in \threads{}, \lk \in \locks{}$}{
			$\view{\AcqLst}{t, \lk}{\tuple{t_1, \lk_1, t_2, \lk_2}}$ := $\epsilon$
		}
	}
	\ForEach{$u \in \threads{}$}{
		\ForEach{$t\in \threads{}, \lk_1, \lk_2 \in \locks{}$}{
			$\view{\DPLst}{t, \lk_1, \lk_2}{\tuple{u}}$ := $\epsilon$
		}
	}
}
\BlankLine
% {\tiny \tcp{$Lst$ is a list of triplets $(g, C, C')$ in increasing order of $C$}}
% {\tiny \tcp{Returns the last element $(g_{\max}, C_{\max}, C'_{\max})$ s.t. $C_{\max} \cle U$}}
% {\tiny \tcp{Removes all other elements $(g,C, C')$ from $Lst$ with $C \cle U$}}
\myfun{\maxLB{$U$, $Lst$}}{
	$(g_{\max}, C_{\max}, C'_{\max})$ := $(0, \bot, \bot)$ \; 
	\While{$\NOT\, \isEmpty{Lst}$}{
		$(g, C, C')$ := $\first{Lst}$ \;
		\If{$C \cle U$}{
			$(g_{\max}, C_{\max}, C'_{\max})$ := $(g, C, C')$
		}
		\Else{
			\Break
		}
		$\removeFirst{Lst}$
	}
	\Return $(g_{\max}, C_{\max}, C'_{\max})$
}

\BlankLine

\myfun{\fixpoint{$I$, $\tuple{t_1, \lk_1, t_2, \lk_2}$}}{
	\Repeat{$I$ does not change}{
		\ForEach{$\lk \in \locks{}$}{
			\ForEach{$t \in \threads{}$}{
				($g_{\lk, t}, C_{\lk, t}, C'_{\lk, t}$) := 
				\maxLB{$I$, $\view{\AcqLst}{t, \lk}{\tuple{t_1, \lk_1, t_2, \lk_2}}$} \;
			}
			$t_{\max}$ := argmax$_{t \in \threads{}}\set{g_{\lk, t}}$ \;
			$I$ := $I \mx \bigsqcup_{t \neq t_{\max} \in \threads{}} C'_{\lk, t}$	
		}
	}
	\Return $I$
}

\myfun{\checkDeadlock{$Lst$, $I$, $\tuple{t_1, \lk_1, t_2, \lk_2}$}}{
	\While{$\NOT\, \isEmpty{Lst}$}{
		($C_{\prev{}}, C$) := $\first{Lst}$ \;
		$I$ := \fixpoint{$I \mx C_{\prev{}}$, $\tuple{t_1, \lk_1, t_2, \lk_2}$} \;
		\If{$C \not\cle I$}{ \linelabel{check-containment}
			\declare `Deadlock' \;
			\Break
		}
		$\removeFirst{Lst}$
	}
	\Return $I$
}

\myhandler{\rdhandler{$t$, $x$}}{
	$\Cc_t$ := $\Cc_t \mx \LW_x$ \linelabel{clkrd}
}

\myhandler{\wthandler{$t$, $x$}}{
	$\LW_x$ := $\Cc_t$ \linelabel{clkwt};\\
	$\Cc_t$ := $\Cc_t[t \mapsto \Cc_t(t)+1]$
}

\myhandler{\acqhandler{$t$, $\lk$}}{
	$C_{\prev{}}$ := $\Cc_t$ \;
	$\Cc_t$ := $\Cc_t[t \mapsto \Cc_t(t)+1]$\\ $\gId_\lk$ := $\gId_\lk + 1$ \linelabel{clkacq}\;
	\ForEach{$t_1, t_2 \in \threads{}, \lk_1, \lk_2 \in \locks{}$}{
		$\addLst{\view{\AcqLst}{t, \lk}{\tuple{t_1, \lk_1, t_2, \lk_2}}}{$(\gId_\lk, \Cc_t, \bot)$}$
	}
	\ForEach{$u\in \threads{}$, $\lk' \in \lheld{}$}{
		$\addLst{\view{\DPLst}{t, \lk, \lk'}{\tuple{u}}}{($C_{\prev{}}$, $\Cc_t$)}$
	}
	\ForEach{$u \neq t \in \threads{}$, $\lk' \in \lheld{}$}{
		$I$ := $\view{\Ii}{}{\tuple{u, \lk' t, \lk}} \mx C_{\prev{}}$ \;
		$\view{\Ii}{}{\tuple{u, \lk' t, \lk}}$ := \checkDeadlock{$\view{\DPLst}{u, \lk', \lk\,}{\tuple{t}}$, $I$, $\tuple{u, \lk' t, \lk}$}
	}

}

\myhandler{\relhandler{$t$, $\lk$}}{
	$\Cc_t$ := $\Cc_t[t \mapsto \Cc_t(t)+1]$ \linelabel{clkrel} \;
	\ForEach{$t_1, t_2 \in \threads{}, \lk_1, \lk_2 \in  \locks{}$}{
		$\updateRel{\lst{\view{\AcqLst}{t, \lk}{\tuple{t_1, \lk_1, t_2, \lk_2}}}}{$\Cc_t$}$
	}
}


\end{multicols}
% \vspace*{0.25cm}
\normalsize
\caption{$\SyncPDOnline$.}
\algolabel{online}
\end{algorithm*}
\normalsize
%%!TEX root = ./main.tex

\section{Further Experiments}
%\seclabel{app_experiments}
\applabel{app_experiments}

For the sake of completeness, here we report on a small experimental comparison with the SMT-based tool Dirk~\cite{Kalhauge2018}.
As the complexity of Dirk is polynomial (due to SMT solving), it uses windowing to examine segments of a large trace.
In our experiments we set a window size of 1.5K and timeout of $1$ hour.
Finally, our experiments were run on a subset of our dataset, as we faced technical challenges in executing Dirk on some benchmarks.
We nevertheless ensured that all tools were executed on identical inputs.
\cref{tab:expr-results-dirk} shows the obtained results.
We confirm that Dirk takes extremely long to complete even this subset and fails to report deadlocks on larger instances due to timeout, 
while it detects two deadlocks missed by our algorithms due to the fact that it also detects deadlocks not witnessed by correct reorderings.
In particular, Dirk reports a deadlock in the benchmarks Deadlock and Transfer.
The reordering that exposes these deadlocks changes the reads-from function so that reads observe different writes in the witness trace (which, nevertheless, write the same value as the old writes, and thus the witness is sound).
We remark that reasoning about reorderings based on values is a very intractable problem (well-known to be $\NP$-hard even for 3 threads), and thus methods based on such reasoning are not scalable.
Indeed, Dirk timed out after 1h on each benchmark that required at least 1 second in $\SyncPDOnline$ and $\SyncPDOffline$.
%!TEX root = ../main.tex

\begin{table*}[H]
\caption{
Trace characteristics, Lock graph statistics and performance comparison.
Column 1 denotes the name of the benchmark.
Columns 2, 3, 4, 5 and 6 respectively denote the number of events, threads, memory locations, locks
and the total number of acquire and request events.
Columns 7, 8 and 9 denote the number of nodes, edges and cycles in the abstract lock graph,
while Columns 10 and 11 denote the number of abstract and concrete deadlock patterns.
Columns 12 - 17 denote the times taken (in seconds) and the number of deadlocks reported
by each of Dirk, \SyncPDOnline and \SyncPDOffline.
%Statistics on the lock graph $\lkevgraph{\tr}$ of each trace $\tr$.
% We denote by $\mathcal{N}$, $\mathcal{T}$, $\mathcal{M}$, $\mathcal{L}$ and $\NumAcquires + \mathcal{R}$ the total number of events, number of threads, number of memory locations, number of locks and number of acquire and request events in the benchmarks, respectively. 
\label{tab:expr-results-dirk}
}
\setlength\tabcolsep{3pt}
\renewcommand{\arraystretch}{0.95}
\small
\centering
\scalebox{0.96}{
\begin{tabular}{|c|c|c|c|c|c||c|c|c|c|c||c|c||c|c||c|c|}
\hline
1 & 2 & 3 & 4 & 5 & 6 & 7 & 8 & 9 & 10 & 11 & 12 & 13 & 14 & 15 & 16 & 17 \\
\hline
\multirow{2}{*}{\textbf{Benchmark}}& 
\multirow{2}{*}{$\mathcal{N}$} & \multirow{2}{*}{$\mathcal{T}$} & \multirow{2}{*}{$\mathcal{M}$} & \multirow{2}{*}{$\mathcal{L}$} & \multirow{2}{*}{$\NumAcquires+\mathcal{R}$} 
& \multicolumn{5}{c||}{ \textsf{Abstract Lock Graph } $\lkevgraph{\tr}$}
& \multicolumn{2}{c||}{ {\textsf{Dirk}} } 
& \multicolumn{2}{c||}{ {\textsf{$\SyncPDOffline$}}} 
& \multicolumn{2}{c|}{ {\textsf{$\SyncPDOnline$}}} \\
\cline{7-17}
& & & & & 
& \textsf{|$V_\tr$|}
& \textsf{|$E_\tr$|}
& \textsf{|$\textsf{Cyc}_\tr$|}
& \textsf{\textsf{Abs. Pat.}}
& \textsf{\textsf{Conc. Pat.}}
& \textbf{\#Dlk} 
& {\textbf{Time (s)}} 
&{\textbf{\#Dlk}} 
& {\textbf{Time (s)}}
&{\textbf{\#Dlk}} 
&{\textbf{Time (s)}} \\
\hline
Deadlock & 50 & 3 & 4 & 3 & 4 & 2 & 2 & 1 & 1 & 1 & 1 & 0.02 & 0 & 0.06 & 0 & 0.03 
\\
NotADeadlock & 74 & 3 & 3 & 4 & 10 & 2 & 2 & 1 & 1 & 1 & 0 & 0.02 & 0 & 0.06 & 0 & 0.03 
\\
Picklock & 80 & 3 & 5 & 5 & 12 & 6 & 8 & 2 & 2 & 2 & 1 & 0.02 & 1 & 0.06 & 1 & 0.02 
\\
Bensalem & 86 & 4 & 4 & 4 & 12 & 6 & 6 & 2 & 2 & 2 & 1 & 0.03 & 1 & 0.06 & 1 & 0.03 
\\
Transfer & 93 & 3 & 10 & 3 & 8 & 2 & 2 & 1 & 1 & 1 & 1 & 0.03 & 0 & 0.06 & 0 & 0.03 
\\
Test & 95 & 3 & 8 & 6 & 14 & 7 & 7 & 2 & 2 & 2 & 2 & 0.06 & 2 & 0.06 & 2 & 0.03 
\\
StringBuffer & 116 & 3 & 14 & 4 & 6 & 2 & 2 & 1 & 1 & 2 & 2 & 0.03 & 2 & 0.06 & 2 & 0.03 
\\
DiningPhil & 398 & 6 & 21 & 6 & 50 & 5 & 5 & 1 & 1 & 3.1K & 1 & 2.0 & 1 & 0.07 & 0 & 0.03 
\\
HashTable & 522 & 3 & 5 & 3 & 65 & 2 & 0 & 0 & 0 & 0 & 2 & 0.23 & 2 & 0.08 & 2 & 0.03 
\\
DirkAccount & 897 & 6 & 46 & 6 & 72 & 8 & 12 & 3 & 3 & 12 & 0 & 0.27 & 0 & 0.08 & 0 & 0.03 
\\
Log4j2 & 2.5K & 4 & 334 & 11 & 15 & 1 & 0 & 0 & 0 & 0 & 1 & 0.78 & 1 & 0.07 & 1 & 0.03 
\\
Derby2 & 3.0K & 3 & 1.0K & 4 & 2 & 1 & 0 & 0 & 0 & 0 & 1 & 0.33 & 1 & 0.08 & 1 & 0.04 
\\
JDBCMySQL-1 & 674.5K & 3 & 73.1K & 11 & 4.1K & 15 & 10 & 1 & 1 & 4 & 2 & 34.5 & 2 & 0.33 & 2 & 0.26 
\\
JDBCMySQL-2 & 674.7K & 3 & 73.2K & 11 & 1.3K & 11 & 2 & 0 & 0 & 0 & 1 & 37.1 & 1 & 0.24 & 1 & 0.29 
\\
JDBCMySQL-3 & 675.9K & 3 & 73.2K & 12 & 4.1K & 22 & 23 & 5 & 5 & 16 & 1 & 115.6 & 1 & 0.34 & 1 & 0.33 
\\
JDBCMySQL-4 & 676.0K & 3 & 73.3K & 13 & 4.1K & 21 & 21 & 5 & 4 & 18 & 2 & 189.2 & 2 & 0.51 & 2 & 0.34 
\\
ArrayList & 4.1M & 801 & 120.6K & 802 & 81.4K & 115 & 8 & 4 & 4 & 301 & - & - & 3 & 1.9 & 3 & 1.6 
\\
Stack & 4.4M & 801 & 118.3K & 1.6K & 165.9K & 1.9K & 246 & 3 & 3 & 201 & - & - & 3 & 4.2 & 3 & 3.1 
\\
IdentityHashMap & 4.5M & 801 & 496.1K & 802 & 80.9K & 40 & 2 & 1 & 1 & 4 & - & - & 1 & 1.7 & 1 & 1.7 
\\
LinkedList & 4.8M & 801 & 289.9K & 802 & 81.3K & 114 & 8 & 4 & 4 & 10.2K & - & - & 3 & 2.5 & 3 & 2.1 
\\
HashMap & 5.3M & 801 & 554.7K & 802 & 84.6K & 40 & 2 & 1 & 1 & 10.2K & - & - & 2 & 2.0 & 2 & 2.0 
\\
Vector & 5.4M & 3 & 14 & 3 & 400.0K & 2 & 2 & 1 & 1 & 1.4B & - & - & 1 & 2.2 & 1 & 1.2 
\\
LinkedHashMap & 6.2M & 801 & 617.0K & 802 & 84.6K & 40 & 2 & 1 & 1 & 10.2K & - & - & 2 & 1.9 & 2 & 2.3 
\\
WeakHashMap & 6.3M & 801 & 540.4K & 802 & 84.5K & 40 & 2 & 1 & 1 & 10.2K & - & - & 2 & 1.9 & 2 & 2.1 
\\
TreeMap & 19.5M & 801 & 493.4K & 802 & 84.6K & 40 & 2 & 1 & 1 & 10.2K & - & - & 2 & 3.0 & 2 & 2.5 
\\
TestPerf & 200.1M & 50 & 599 & 9 & 98.3K & 0 & 0 & 0 & 0 & 0 & - & - & 0 & 16.0 & 0 & 7.9 
\\
Groovy2 & 212.8M & 13 & 12.7M & 9.9K & 10.8K & 13 & 0 & 0 & 0 & 0 & - & - & 0 & 15.5 & 0 & 13.3 
\\
\hline
\hline
\textbf{Total} & \textbf{476M} & & \textbf{16M} & \textbf{17k} & \textbf{1.2M} 
& \textbf{2.4K} & \textbf{376} & \textbf{42} & \textbf{41} & \textbf{1.4B}
& \textbf{19} & \textbf{>380} & \textbf{36} & \textbf{55} & \textbf{35} & \textbf{41}
\\
\hline
\end{tabular}
}
\end{table*}


\newpage
%!TEX root = ./main.tex

\section{Incomparability of $\SyncPDOffline$ and \seqc}
\applabel{incomp}
In this section we provide two example traces and discuss the deadlocks found and missed on them by \seqc and \SyncPDOffline.
\exref{incomp-ex1} discusses a trace which contains a deadlock found by \SyncPDOffline but missed by \seqc.
\exref{incomp-ex2} discusses a trace which contains a deadlock found by \seqc but missed by \SyncPDOffline.
These examples showcase the incomparability of the techniques \SyncPDOffline and \seqc.

\begin{example}
\exlabel{incomp-ex1}
\figref{incomp1} displays a trace $\sigma$ in which \SyncPDOffline reports a deadlock whereas \seqc misses the deadlock.
In this trace, there is a single deadlock pattern $\pattern{e_4, e_{14}}$.
This deadlock pattern can be realized by the sync-preserving correct reordering $\rho = e_3, e_{12}, e_{13}, e_8, e_9$.
However, this deadlock is missed by \seqc.
The reason for this can be explained as follows.
First, observe that a correct reordering of $\sigma$ would need to include the event $e_{13}$ and consequently $e_9 = \rf{\sigma}(e_{13})$.
A strategy that \seqc employs when it is producing correct reorderings is that it aims to close all the critical sections that appear in the reordering.
In order to fulfill this constraint, \seqc closes the critical section on lock $\LockColorOne{\lk_1}$ in thread $t_1$.
This requires \seqc to execute the read event $e_{10}$ as well as the event $e_7 = \rf{\sigma}(e_{10})$.
This in turn causes \seqc to execute the events $e_3\ldots e_6$.
Observe that the first event that appears in the deadlock pattern, $e_4$, can no longer be scheduled concurrently with $e_{14}$.
Hence, \seqc is unable to report a deadlock in this trace.
\end{example}

\begin{example}
\exlabel{incomp-ex2}
\figref{incomp2} display a trace $\sigma$ in which \seqc reports one more deadlock than \SyncPDOffline.
In this trace there are two deadlock patterns $\pattern{e_2, e_{6}}$ and $\pattern{e_2, e_{8}}$.
The first deadlock pattern $\pattern{e_2, e_{6}}$ can be realized to a real deadlock by both \seqc and \SyncPDOffline. 
The second deadlock pattern  $\pattern{e_2, e_{8}}$ cannot be realized to a real deadlock by \SyncPDOffline.
Whereas, \seqc is able to report this as a real deadlock.
The reason why \SyncPDOffline misses the second deadlock can be explained as follows.
The computed sync-preserving closure for the second deadlock pattern is $\SPClosure{\tr}(\prev{\tr}(\{ e_{2},e_{8}\})) = \set{e_1,\ldots, e_4,e_5,\ldots,e_7}$.
Since the \SyncPDOffline algorithm cannot reverse the order of critical sections on a given lock, 
it always enforces to close the critical section on lock $\LockColorOne{\lk_1}$ in thread $t_2$ after the one in $t_1$.
This brings the event $e_2$ into the sync-preserving closure of $\pattern{e_2, e_{8}}$
and disables \SyncPDOffline from realizing this deadlock pattern.
\end{example}


%!TEX root=../main.tex

\begin{figure}[t]
\centering
\scalebox{1}{
\execution{4}{
\figev{4}{$\acq(\LockColorOne{\lk_1})$}
\figev{4}{$\rel(\LockColorOne{\lk_1})$}
\figev{3}{$\acq(\LockColorTwo{\lk_2})$}
\figev{3}{$\mathbf{\Bacq(\LockColorThree{\lk_3})}$}
\figev{3}{$\rel(\LockColorThree{\lk_3})$}
\figev{3}{$\rel(\LockColorTwo{\lk_2})$}
\figev{3}{$\wt(y)$}
\figev{1}{$\acq(\LockColorOne{\lk_1})$}
\figev{1}{$\wt(x)$}
\figev{1}{$\rd(y)$}
\figev{1}{$\rel(\LockColorOne{\lk_1})$}
\figev{2}{${\acq(\LockColorThree{\lk_3})}$}
\figev{2}{$\rd(x)$}
\figev{2}{$\mathbf{\Bacq(\LockColorTwo{\lk_2})}$}
\figev{2}{$\rel(\LockColorTwo{\lk_2})$}
\figev{2}{$\rel(\LockColorThree{\lk_3})$}
}
}
% \caption{
% A  trace $\tr$ with a sync-preserving deadlock, stalling $t_2$ on $e_4$ and $t_3$ on $e_{18}$.
% Underlined events indicate the slice of $\tr$ that forms a sync-preserving lower set of $\tr$,
% and serves as a valid witness for the deadlock.
% \figlabel{motivating}
% }
\caption{
A trace $\tr$ with a sync-preserving deadlock, stalling $t_2$ on $e_{14}$ and $t_3$ on $e_{4}$.
\figlabel{incomp1}
}
\end{figure}
%!TEX root=../main.tex

\begin{figure}[t]
\centering
\scalebox{1}{
\execution{2}{
\figev{1}{$\acq(\LockColorOne{\lk_1})$}
\figev{1}{$\mathbf{\Bacq(\LockColorTwo{\lk_2})}$}
\figev{1}{$\rel(\LockColorTwo{\lk_2})$}
\figev{1}{$\rel(\LockColorOne{\lk_1})$}
\figev{2}{${\acq(\LockColorTwo{\lk_2})}$}
\figev{2}{$\mathbf{\Bacq(\LockColorOne{\lk_1})}$}
\figev{2}{$\rel(\LockColorOne{\lk_1})$}
\figev{2}{$\mathbf{\Bacq(\LockColorOne{\lk_1})}$}
\figev{2}{$\rel(\LockColorOne{\lk_1})$}
\figev{2}{$\rel(\LockColorTwo{\lk_2})$}
}
}
% \caption{
% A  trace $\tr$ with a sync-preserving deadlock, stalling $t_2$ on $e_4$ and $t_3$ on $e_{18}$.
% Underlined events indicate the slice of $\tr$ that forms a sync-preserving lower set of $\tr$,
% and serves as a valid witness for the deadlock.
% \figlabel{motivating}
% }
\caption{
A  trace $\tr$ with a sync-preserving deadlock, stalling $t_1$ on $e_{2}$ and $t_2$ on $e_{6}$ and a predictable but not sync-preserving deadlock stalling $t_1$ on $e_{2}$ and $t_2$ on $e_{8}$.
\figlabel{incomp2}
}
\end{figure}
%!TEX root = ./main.tex

\section{Unsoundness Of \dirk}
\applabel{unsound-dirk}
In this section we provide two example programs that showcase the unsoundness of \dirk.
\exref{dirk-unsound1} discusses a program which is a modified version of the Transfer benchmark.
\dirk reports a false deadlock in this example due to that it is aiming to relax the standard definition of correct
reorderings and allow a read event to read from a different write than in the original trace.
Note that this relaxation enables \dirk to predict deadlocks in the benchmarks Transfer and Deadlock, which are missed by $\SyncPDOffline$ and \seqc. 
However, this relaxation is not fully formalized and the implementation of it contains errors.
\exref{dirk-unsound2} discusses a crafted program in which \dirk is reporting a false deadlock due to missing that a common lock protects an otherwise deadlock pattern.

\begin{example}
\exlabel{dirk-unsound2}
	\figref{dirk-unsound-2} displays a program in which \dirk makes a false deadlock report. 
	The program spawns three threads \texttt{T1}, \texttt{T2}, and \texttt{T3} and  
	there exist three locks \texttt{L1}, \texttt{L2}, and \texttt{L3} which are used in these threads.
	In the program, \texttt{T1} acquires the lock \texttt{L1} and then forks \texttt{T2}, and it holds the lock \texttt{L1} until \texttt{T2} is joined.
	Even though the locks \texttt{L2} and \texttt{L3} are acquired in a cyclic manner 
	in the threads  \texttt{T2} and \texttt{T3}, 
	this program does not exhibit a deadlock due to the common lock \texttt{L1}.
	Therefore, this program does not exhibit a deadlock.
	However, \dirk fails to recognize the common lock \texttt{L1} and reports a false deadlock.
	This observed behavior is not an implementation error, but it is in alignment with the underlying theory of \dirk~\cite{Kalhauge2018}.
	In particular, \dirk makes SMT queries to decide the existence of deadlocks, but its encoding unsoundly neglects the condition that deadlocking events must have disjoint lock sets.  
\end{example}

\begin{example}
\exlabel{dirk-unsound1}
	\figref{dirk-unsound-1} displays a program in which \dirk makes a false deadlock report. 
	This program is a modified version of the Transfer benchmark, which was used as a running example to explain \dirk~\cite{Kalhauge2018}.
	The example contains the classes \texttt{Transfer1}, \texttt{Transfer2}, \texttt{Store} and \texttt{UserObject}. 
	The \texttt{Store} class contains two synchronized methods, \texttt{transferTo} and \texttt{add}.
	The \texttt{transferTo} takes another \texttt{Store} object as argument and calls the \texttt{add} method by using this object.
	Notice that if two different \texttt{Store} objects simultaneously call the \texttt{transferTo} 
	method by passing each other as the argument, then this can potentially lead to a deadlock.
	Observe that in the given program there exists such a scenario where \texttt{transferTo} is called 
	with opposite \texttt{Store} objects in the classes \texttt{Transfer1} and \texttt{Transfer2}.
	However, in this case the program can never actually deadlock.
	The reason is that the \texttt{transferTo} method in \texttt{Transfer2} class can only be executed if the \texttt{UserObject.data} can be cast to an integer value.
	The \texttt{UserObject.data} is a volatile variable that is shared by the classes \texttt{Transfer1} and \texttt{Transfer2}. 
	This variable is initialized to a string value and an integer value is assigned to this variable only after the
	\texttt{transferTo} method has been executed in \texttt{Transfer1}. 
	Hence, this program does not exhibit a deadlock.
	However, \dirk reports a deadlock in this program.
	The reason is that \dirk unsoundly concludes that reading the value of \texttt{UserObject.data} at Line $54$ does not affect the control flow of the program.
	Hence, in case the value of \texttt{UserObject.data} object in Line $54$ is read from the write event on Line $39$, 
	\dirk ignores this dependency and concludes that Lines $38$ and $54$ can execute concurrently and reports a false deadlock. 
\end{example}



%!TEX root = ./main.tex


\begin{figure}[h!]
	\hspace{1cm}
	\begin{subfigure}[t]{.43\linewidth}
		\lstinputlisting[language = Java,firstline=0,lastline=32]{figures/FalsePositive1.java}
	\end{subfigure}
	\hfill
	\begin{subfigure}[t]{.45\linewidth}
		\lstinputlisting[language = Java,firstline=32,lastline=65,firstnumber=33]{figures/FalsePositive1.java}
	\end{subfigure}
	\caption{FalseDeadlock1.java}
	\figlabel{dirk-unsound-2}
\end{figure}

\begin{figure}[h!]
	\hspace{1cm}
	\begin{subfigure}[t]{.43\linewidth}
		\lstinputlisting[language = Java,firstline=0,lastline=41]{figures/FalsePositive2.java}
	\end{subfigure}
\hfill
	\begin{subfigure}[t]{.43\linewidth}
		\lstinputlisting[language = Java,firstline=42,lastline=82,firstnumber=42]{figures/FalsePositive2.java}
	\end{subfigure}
	\caption{FalseDeadlock2.java}
	\figlabel{dirk-unsound-1}
\end{figure}

\end{arxiv}

\end{document}
